\documentclass[10pt,openright,10pt]{book}

\usepackage{iftex}
\ifPDFTeX
  \usepackage[T1]{fontenc}
  \usepackage[utf8]{inputenc}
  \usepackage{textcomp}
  \usepackage{lmodern}
\else
  \usepackage{fontspec}
  \setmainfont{lmroman10-regular.otf}[
    BoldFont=lmroman10-bold.otf,
    ItalicFont=lmroman10-italic.otf,
    BoldItalicFont=lmroman10-bolditalic.otf
  ]
  \setsansfont{lmsans10-regular.otf}[
    BoldFont=lmsans10-bold.otf,
    ItalicFont=lmsans10-oblique.otf
  ]
  \setmonofont{lmmono10-regular.otf}[Scale=MatchLowercase]
\fi
\usepackage{microtype}
\usepackage[paperwidth=7in,paperheight=10in,inner=0.7in,outer=0.55in,top=0.65in,bottom=0.75in]{geometry}
\usepackage{setspace}
\usepackage{titlesec}
\usepackage{fancyhdr}
\usepackage{graphicx}
\usepackage[export]{adjustbox}
\providecommand{\pandocbounded}[1]{%
  \begin{adjustbox}{max width=\linewidth,max totalheight=0.62\textheight,center}
  #1%
  \end{adjustbox}%
}
\usepackage[table]{xcolor}
\usepackage{hyperref}
\usepackage{bookmark}
\usepackage{longtable}
\usepackage{booktabs}
\usepackage{array}
\usepackage{calc}
\usepackage{etoolbox}
\usepackage{enumitem}
\usepackage{amssymb}
\usepackage{fvextra}
\usepackage{mdframed}

\providecommand{\tightlist}{%
  \setlength{\itemsep}{0pt}\setlength{\parskip}{0pt}}
\newcounter{none}

\hypersetup{
  colorlinks=true,
  linkcolor=blue!50!black,
  urlcolor=blue!50!black,
  citecolor=black,
  pdftitle={Principles of Data Engineering},
  pdfauthor={Matthias Wong}
}

\setstretch{1.08}
\setlength{\parindent}{0pt}
\setlength{\parskip}{0.7em}
\setlist{itemsep=0.25em, topsep=0.4em}

\definecolor{pbiHeader}{HTML}{EAF2F8}
\definecolor{pbiStripe}{HTML}{F7FAFC}
\definecolor{pbiGrid}{HTML}{C9D3DD}
\definecolor{codeBack}{HTML}{F6F8FA}
\definecolor{codeFrame}{HTML}{C9D3DD}
\definecolor{codeText}{HTML}{17202A}
\arrayrulecolor{pbiGrid}
\setlength{\arrayrulewidth}{0.35pt}
\renewcommand{\arraystretch}{1.18}
\setlength{\tabcolsep}{5pt}
\renewcommand{\toprule}{\arrayrulecolor{pbiGrid}\hline\rowcolor{pbiHeader}}
\renewcommand{\midrule}{\hline}
\renewcommand{\bottomrule}{\hline}
\AtBeginEnvironment{longtable}{%
  \small
  \setlength{\LTpre}{0.45em}
  \setlength{\LTpost}{0.9em}
  \rowcolors{2}{white}{pbiStripe}
}
\AtEndEnvironment{longtable}{\rowcolors{2}{}{}}

\RecustomVerbatimEnvironment{verbatim}{Verbatim}{%
  breaklines,
  breakanywhere,
  fontsize=\small,
  commandchars=\\\{\}
}
\BeforeBeginEnvironment{verbatim}{%
  \begin{mdframed}[
    backgroundcolor=codeBack,
    linecolor=codeFrame,
    linewidth=0.35pt,
    roundcorner=2pt,
    innerleftmargin=7pt,
    innerrightmargin=7pt,
    innertopmargin=6pt,
    innerbottommargin=6pt,
    skipabove=0.8em,
    skipbelow=0.9em
  ]\color{codeText}%
}
\AfterEndEnvironment{verbatim}{\end{mdframed}}

\pagestyle{fancy}
\fancyhf{}
\fancyfoot[C]{\thepage}
\renewcommand{\headrulewidth}{0pt}

\titleformat{\chapter}[display]
  {\normalfont\bfseries\Huge}
  {}
  {0pt}
  {\Huge}
\titlespacing*{\chapter}{0pt}{0pt}{24pt}


\begin{document}

\frontmatter

\thispagestyle{empty}

\vspace*{0.25\textheight}

\begin{center}

{\Huge\bfseries Principles of Data Engineering\\}

\vspace{1.5cm}

{\Large Matthias Wong}

\end{center}

\cleardoublepage

\thispagestyle{empty}

\vspace*{0.35\textheight}

\noindent The book is available on
\href{https://principlesofdataengineering.org}{principlesofdataengineering.org}.

\vspace{1em}

\noindent Downloadable PDF and EPUB editions are available for offline
reading. The online version is canonical and may be updated over time.

\cleardoublepage

\tableofcontents

\cleardoublepage

\mainmatter

\chapter{The author}\label{docs-about-the-author}

\section{About the author}\label{about-the-author}

Matthias Wong\textquotesingle s work examines how organisations see,
understand, and act in complex environments. His writing spans
organisational theory, data, AI, business, and people leadership.

As part of this work, he has led and trained advanced data engineering
teams working on high-stakes systems, with a focus on transforming
fragmented operational data into usable, decision-ready information.

\emph{Principles of Data Engineering} forms the technical segment of
this broader work: an account of how organisations make reality visible
through data, and how engineers can build the representations
organisations rely on to exercise sound judgement.

He is especially passionate about uplifting staff from junior to
advanced capability, and providing senior advisory to organisational
leaders on the effective use of data to steer complex institutions.

Matthias holds a PhD in Computational Mathematics.

For contact information, see \href{https://matthiaswong.me}{the
author\textquotesingle s website}.

\cleardoublepage

\chapter{The book}\label{docs-about-the-book}

\section{Dedication and
acknowledgements}\label{dedication-and-acknowledgements}

This book is dedicated to N.H. and G.H.

It was born from years of ongoing dialogue with developers and business
stakeholders at the Department of Agriculture, Fisheries and Forestry
between 2017 and 2025.

It would not have been possible without the many data engineers and
analysts who contributed their thinking, challenged assumptions, and
worked through complex problems in practice. Their experiences,
questions, and solutions form the foundation of much of what is written
here.

Special thanks to E.D. and M.T. for their contribution and support
throughout this journey, and to M.C., Z.B., D.S., S.S., S.M., and R.D.
for their review and feedback.

\section{Origin and release}\label{origin-and-release}

\emph{Principles of Data Engineering} was written in the Department of
Agriculture, Fisheries and Forestry, and released under Freedom of
Information.

This website is a web edition of that manuscript. It has been
transcribed, edited, restructured, and supplemented for online
publication by Matthias Wong.

\section{Attribution}\label{attribution}

Department of Agriculture, Fisheries and Forestry. 2026.
\emph{Principles of Data Engineering manual}
\href{https://www.agriculture.gov.au/sites/default/files/documents/35670.pdf}{FOI
Disclosure Log Reference 35670}.

This web edition includes adapted and additional material by Matthias
Wong.

© Department of Agriculture, Fisheries and Forestry 2026 Web adaptation
and additional material © Matthias Wong 2026

\cleardoublepage

\thispagestyle{empty}
\phantomsection
\hypertarget{docs-foundations}{}
\addcontentsline{toc}{part}{Foundations}
\vspace*{0.34\textheight}
\begin{center}
{\Huge\bfseries Foundations}
\end{center}
\cleardoublepage

\chapter{What is data
engineering?}\label{docs-foundations-what-is-data-engineering}

\section{The aim is insight}\label{the-aim-is-insight}

The aim of data engineering is insight.

But what is insight? Insight is information analysed in the light of
intent.

\begin{quote}
Data engineering, therefore, is an activity in which business intent is
applied to data to create first information, then insight.
\end{quote}

\section{Why is it engineering}\label{why-is-it-engineering}

There are other activities besides data engineering that seek insight.
These include staff consultation, customer feedback, and market
research. In the case of government agencies, these activities include
overt and covert means such as diplomacy and espionage.

Data engineering distinguishes itself from other means of gaining
insight through the nature of its methodology.

Like other engineering disciplines, it focuses on systematic rigour,
efficiency, and reproducibility---achieving these through the
application of scientific and mathematical principles. As with other
engineering disciplines, data engineering arose in
today\textquotesingle s technological society as an effective tool to
tackle complexity.

\section{Where data engineering is
different}\label{where-data-engineering-is-different}

Despite similarities in method, data engineering differs from other
engineering disciplines in crucial ways.

\subsection{Easy substitutes}\label{easy-substitutes}

There are easy substitutes to data engineering.

The business engaging a data engineer is seeking insight, and in many
instances, there are seemingly simple alternatives to serious data
engineering. For example, business clients may choose to rely on their
own judgement or use one of the many available self-service data tools
to rapidly create a visually stunning report with little effort.

This is not the case in other disciplines. The client who engages a
bridge engineer needs a bridge. The client who engages a software
engineer needs software. There are no easy alternatives to achieving
these aims, and therefore no easy alternatives to the engineering
effort.

\subsection{Open-ended questions}\label{open-ended-questions}

The search for business insight is open-ended.

If a client asks for a bridge, the bridge is the end of the engagement.
A bridge is not likely to extend into a longer bridge, nor is the client
likely to settle for half a bridge. This is not the case with business
insight: a question provokes another question, and a question may be
half answered.

This interminability of aim means that a data engineering project is
likewise open-ended. In no other engineering discipline is the engineer
working with such a degree of open-endedness.

\subsection{Ambiguous goals}\label{ambiguous-goals}

The search for business insight is ambiguous.

Quite often, when stakeholders begin a data engineering project, they do
not know what business insights they are after. Partly this is because
business areas first need to discover the information before they know
what questions to ask. Partly it is due to the wide-ranging and
sometimes conflicting business intents---are we concerned about profit,
cost, market share, or negative externalities?

And partly it is because of the wide range of stakeholders who need to
be engaged---senior executives and operational staff have valid but
different concerns---yet the data engineer is expected to meet all
interests.

This ambiguity of aim is akin to the challenge of planning a new city
suburb to satisfy the needs of all residents. Again, this is a unique
challenge for a data engineer.

\section{The fluid nature of data
engineering}\label{the-fluid-nature-of-data-engineering}

The search for business insights is shaped by three intrinsic
characteristics:

\begin{itemize}
\tightlist
\item
  the availability of easy substitutes
\item
  the open-endedness of the search
\item
  the ambiguity of aim
\end{itemize}

We summarise these as the \textbf{fluid nature} of data engineering.

Certainly, the data engineering team can shut its eyes to these problems
and insist on hard-and-fast requirements from stakeholders. This
approach is unlikely to lead to quality outcomes. Instead, the data
engineer\textquotesingle s approach must adapt to the fluid nature of
business insights.

At the end of the day, technical approach alone is insufficient.
Instead, the data engineering team has an unparalleled need to combine
flexibility, endurance of vision, and the ability to negotiate.

Moreover, compared to other engineering contexts, the non-technical
capabilities sit in a small team---sometimes within a sole data
engineer. A data engineering team must confidently wield these
non-technical capabilities if it is to achieve excellence in an
organisation.

The search for business insights is both demanding in complexity and
fluid in nature. This dual aspect makes data engineering one of the most
engaging and rewarding disciplines.

\begin{quote}
\textbf{Key ideas.}

Data engineering is an activity in which business intent is applied to
data to create first information, then insight.

The availability of easy substitutes, the open-endedness of the search,
and the ambiguity of aim define the fluid nature of data engineering.
\end{quote}

\cleardoublepage

\chapter{What is data?}\label{docs-foundations-what-is-data}

\section{The central premise}\label{the-central-premise}

Insight is information analysed in the light of intent. However, the
data engineer starts one step further back than information. Data
engineering starts with data.

Data is what is given before interpretation. It has no fixed meaning by
itself, and is not yet information. More precisely:

\begin{quote}
Data is a fragment of reality captured by process.
\end{quote}

Consider a customer whose seat is cancelled because a flight was
overbooked. The experience is real, but it is internal to the customer.
The decision-maker does not know about it unless the experience is
captured as data and made available to someone who was not there.

Without data, a decision-maker attempting to understand overbooked
flights must rely on guesswork, intuition, anecdote, or high-profile
complaints. These can all be important, but none is systematic.

The same applies to shopping experiences, cargo delays, transactions,
legal decisions, medical observations, and countless other events.
Nothing becomes systematically available to an organisation unless it is
recorded as data.

This definition means data is neither reality itself nor mere
record-keeping. It sits between the world and the organisation's
understanding of the world.

This becomes especially important at scale. In a small organisation,
leaders may retain a direct sense of what is happening. In a large
organisation, that direct visibility breaks down. Data becomes the means
by which reality remains visible across distance, hierarchy, and
complexity.

\section{How is data captured?}\label{how-is-data-captured}

Data is captured by process.

This creates two sources of imperfection. The first is the process. The
second is the capture. Between both sits the digital system: the medium
through which business activity is encoded into data.

A process may be imperfect because it was designed for operational work
rather than analytical understanding. Most workflows exist to get
something done. The digital systems that support those workflows are
usually built around that operational goal. Therefore, the process may
not try to record everything a future analyst may need for business
intent. This is especially the case when the source of data is not
controlled by the business seeking to analyse it.

The capture itself may also be imperfect. It can be technically
difficult, costly, or impractical to record business events accurately.
This challenge is intensified by a rapidly evolving technological
landscape.

Data is therefore never a simple copy of the world. It is an imperfect
projection of the real world onto the data world. The work of the data
engineer begins by understanding these imperfections.

\section{How is data used?}\label{how-is-data-used}

Data is captured by business processes. Once collected, it can be used
to return understanding to the business.

In this perspective, data engineering is the task of taking data
projected by business processes and reshaping it into a form required by
business intent. This can be summarised in Figure 1.

\pandocbounded{\includegraphics[keepaspectratio]{book/epub-assets/diagram-001.png}}

\emph{Figure 1. Business reality is projected into the data world, then
shaped back into business understanding.}

We will return frequently to the concepts illustrated in Figure 1
throughout this book.

\begin{quote}
\textbf{Key ideas.}

Data is a fragment of reality captured by process. At scale, data is the
means by which reality remains visible.

The data world can be seen as an imperfect projection of the business
world.

Data engineering reshapes the data world into forms the business world
can use.
\end{quote}

\cleardoublepage

\chapter{Data and
organisations}\label{docs-foundations-data-and-organisations}

\section{Data and organisational
self-understanding}\label{data-and-organisational-self-understanding}

Data is a fragment of reality. This includes the reality of the
environment in which they operate, but also the reality of the
organisation itself.

This has implications for how data engineering work is done in practice.
The veteran data engineer instinctively works within the entanglement
between data and organisation.

To understand this, it is useful to visualise an organisation as a
feedback loop.

\section{The organisation agency
cycle}\label{the-organisation-agency-cycle}

The central premise "Data is a fragment of reality captured by process"
can be interpreted as a feedback loop around three components: business
processes, digital systems, and captured data. This relationship is
illustrated in Figure 1.

\pandocbounded{\includegraphics[keepaspectratio]{book/epub-assets/diagram-002.png}}

\emph{Figure 1. Business intent defines business process, business
process defines digital systems, digital systems capture data, and data
feeds back into business understanding.}

Like many cycles, an organisation\textquotesingle s feedback loop can
become vicious or virtuous. Business intent is the determinant. An
organisation moving toward its intent is in a virtuous cycle; an
organisation drifting away from its intent is in a vicious cycle.

A key point of the cycle is that data is not passive ``facts and
figures'' from the world. Data consists of records of the world the
organisation is interacting with, and records of that interaction. The
feedback loop is therefore a model of how an organisation exercises its
goals in the world.

The implication is that what the organisation does is shaped by what it
sees, and what it sees is shaped by what it does. Data is the medium of
this double movement.

Consequently, data is central to an organisation\textquotesingle s
self-understanding. Data engineering is not one capability among others.
It is the capability through which other capabilities become visible.

\section{Implications for data
engineering}\label{implications-for-data-engineering}

There are two important implications for data engineering work.

First, recorded data is not simply ``what is in the system.'' It is
downstream of business processes and digital capture. The data engineer
is always working inside this context, and this context shapes the scope
and options for action within a project.

Second, the data engineer's work itself is not neutral. The output of
data engineering is often the first time an organisation encounters the
outcome of its own work in a stable form. For this reason, the output
may be contested, especially when multiple business areas are involved.
Interpretation becomes inseparable from the work: what is good, what is
bad, what is reportable, what is not, which record is authoritative, and
which definition should be canonical?

Organisations do not simply use data. Data is also where meaning is
negotiated.

The data engineer cannot stand outside this negotiation. The engineer
can help the organisation navigate meaning with expertise, defer to
consensus without judgement, or allow the work to fail because meaning
was never clear.

The agency cycle therefore shows why no part of data work is neutral.
The business process that defines the digital system is not neutral. The
digital capture is not neutral. Most importantly, the data engineer's
work to create feedback is not neutral. All are evaluated by intent.
This is why data engineering reshapes data in light of intent.

\begin{quote}
\textbf{Key ideas.}

What the organisation does is shaped by what it sees; what it sees is
shaped by what it does. Data is the medium of this loop.

Organisations do not simply use data; data is also where meaning is
negotiated.
\end{quote}

\cleardoublepage

\thispagestyle{empty}
\phantomsection
\hypertarget{docs-creating-information}{}
\addcontentsline{toc}{part}{Creating information}
\vspace*{0.34\textheight}
\begin{center}
{\Huge\bfseries Creating information}
\end{center}
\cleardoublepage

\chapter{Expressiveness and fragment
modelling}\label{docs-creating-information-expressiveness-and-fragment-modelling}

\section{From data to reusable
information}\label{from-data-to-reusable-information}

Data engineering reshapes
\hyperref[docs-foundations-what-is-data]{fragments of reality} into
forms the business can use.

This work must be done efficiently, reproducibly, and over sustained
periods of time. The goal is not merely to produce isolated reports or
one-off datasets, but to create a cost-effective and robust pipeline of
valuable, reusable blocks of information that are conducive to business
insight.

How does the data engineer do this effectively?

This book introduces five principles of data engineering. The first two
are \textbf{expressiveness} and \textbf{fragment modelling}. The two
principles are:

\begin{quote}
Instead of accepting garbage-in, garbage-out with raw data, add value
through expressive entities.

Instead of building giant tables, create meaningful fragments.
\end{quote}

\section{Expressiveness}\label{expressiveness}

Insight is
\hyperref[docs-foundations-what-is-data-engineering]{information
analysed in the light of intent}. A data engineer who takes this
seriously becomes invested in the business world, asks how data relates
to it, and is driven by the need to organise data in a way that makes
sense for business decisions.

This is the task of applying business intent to reshape data. When done
successfully, the output is expressive of business intent.

The first and most important principle of data engineering is therefore
the search for \textbf{expressiveness}.

The data engineer's work is expressive when its outputs do not merely
reflect data as it was found, but correspond strongly to real business
entities and processes. It is expressive because the consumer or
reviewer of the model can clearly recognise the world the model is
attempting to approximate.

If a reasonably competent layperson cannot easily relate the model to
the real world, then the model is not expressive. Correspondence to the
world by the organisation trying to influence it is therefore the final
arbitrator of whether a data model is successful.

Expressiveness has three qualities.

It is \textbf{faithful}: it preserves meaningful correspondence with the
business reality it represents.

It is \textbf{legible}: a competent business user can recognise what the
model is attempting to show.

It is \textbf{operable}: the model can be used in workflows or decisions
with minimal rework.

Expressiveness encompasses all aspects of a good data engineer's work.
It lies in everything: the way the engineer designs tables, chooses
names for tables and columns, connects entities through primary and
foreign keys, and describes the intent of code to others.

\section{Fragment modelling}\label{fragment-modelling}

The \hyperref[docs-foundations-what-is-data-engineering]{fluid nature of
business insight} means that the data engineer cannot know in advance
exactly what needs to be built.

At a micro level, this applies to a single data engineering project.
Requirements shift as information becomes visible and stakeholders
encounter their own business through data. At a macro level, discovering
business insight can be a journey of years over multiple teams. In a
large organisation, different teams will also want to see the same
information in different ways to reflect specific business interests.

This is why the goal of a data engineer is not only to build complete
data products, but to produce reusable blocks of information.

When the data engineer shifts focus from building complete products to
creating blocks that can be reused for multiple products, the warehouse
is no longer seen merely as a repository of data models. It becomes a
fragment store of valuable information.

A data engineering approach focused on creating such a store is
\textbf{fragment modelling}.

Modularity of code is a cornerstone in software engineering. Fragment
modelling applies this concept to the modularity of information creation
in a data pipeline. The result, as in software engineering, is clarity,
flexibility, and reuse.

\begin{quote}
The data engineer therefore sees the pipeline primarily as software
code, and tables as functions.
\end{quote}

A good table in a pipeline performs a defined transformation---it
accepts upstream data, applies a specific piece of business meaning, and
produces an output that can be reused downstream.

In this sense, much of the wisdom in the Zen of Python carries to
fragment modelling.

Fragment modelling is concerned not only with the informational value of
a fragment, but also with its intended use. The data engineer does not
rest until the fragment becomes plug-and-play for the end user who can
make use of the information with minimal additional engineering.

In most cases, this means designing fragments that integrate cleanly
into traditional dimensional models for self-service business
intelligence. However, it also means ensuring that fragments can be
adapted for other analytical scenarios, such as data science feature
engineering.

This combined focus on both information and usability makes fragment
modelling especially suited to the demands of the data age, where data
analytics is employed in a growing range of contexts.

Fragment modelling is also an important step towards expressiveness.

A wide table with a large number of miscellaneous columns buries
multiple business concepts into one undifferentiated "thing" that does
not recognisably map to any real-world entity. On the other hand,
tightly defined fragments bring clarity to business attributes by giving
them sharp relief.

It is impossible to achieve a truly expressive data model without
deconstructing source data into meaningful fragments and reorganising
them to better reflect real-world entities.

This chapter introduces fragment modelling as a principle. The later
chapter
\hyperref[docs-creating-information-meaningful-fragments]{Meaningful
fragments} explores the practice in detail.

Fragment modelling can be disorienting for an engineer used to
traditional warehousing approaches. Once familiar, it becomes a powerful
approach to creating sustainable and flexible pipelines. The following
chapters explain how this is done.

\begin{quote}
\textbf{Key ideas.}

Expressiveness is the principle that data models should recognisably
correspond to the business world they represent.

Expressive data is faithful, legible, and operable.

Fragment modelling is the principle of creating reusable blocks of
information rather than giant tables.

In a well-designed pipeline, tables can be understood as functions that
transform fragments of data into reusable information.
\end{quote}

\cleardoublepage

\chapter{Mapping the data
world}\label{docs-creating-information-mapping-the-data-world}

\section{From records to entities}\label{from-records-to-entities}

Business is not interested in data for data's sake. It is interested in
the business reality behind the data.

From this perspective, data is not an end in itself. Its function is to
bridge the business processes that collect the data and the business
intent that seeks insight.

Consequently, the data engineer sees each data record as an intermediary
of the business entity that generated it.

A business entity can be a concrete object, such as a customer or a
product. It can also be an abstract event, such as a business
transaction.

This shift from data record to business entity is the first practical
step in expressiveness. For many data engineers, it is also the first
threshold: learning to see not the record itself, but the entity behind
it.

For a data engineer who starts with the data and reaches back to the
business world, the tool of choice is the primary key.

\section{Primary key}\label{primary-key}

Primary keys play a special role in establishing the correspondence
between business entities in the real world and data records in the data
world.

A primary key is the handle by which business processes use to
accurately retrieve a business entity's data record from a data store.

For example, an order number identifies a sales transaction, and a
customer number identifies a client record. These keys are visible to
end users who use them to relate back to business processes, such as
calling a company to follow up on an order.

Systems may also employ internal keys, invisible to users, to support
subsidiary information retrieval. In every case, the function remains
the same: reliable lookup of records.

The quality of a primary key is underpinned by its business process.
This process may be:

\begin{itemize}
\tightlist
\item
  \textbf{Rigorous:} A strict identity-proofing protocol, such as
  issuing social security numbers, ensures unique assignment of a key to
  a personal identity. Mistakes can occur. However, the overall system
  remains robust across a large population.
\item
  \textbf{Fuzzy:} Smaller organisations may register customers by phone
  number or email. When the database is small, this may return accurate
  results. However, increased numbers of customers may lead to duplicate
  records or accidental access if contact details overlap.
\item
  \textbf{Log-like:} Modern online shops often implement guest checkout
  to cater to buyer preference. Guest checkout creates a new record for
  every transaction. In this case, each row is effectively logged as a
  separate event.
\end{itemize}

The final example is a case of having no true primary key.

This absence may be by design, as in the case of guest checkout. The
data engineer can treat each row as its own entity and assign a random,
unique number for a primary key. If the business needs to know the
underlying business entity driving each record, such as the identity of
the guest shopper, the data engineer or data scientist may apply entity
resolution techniques to deduplicate the data.

The absence of a primary key can also be accidental. For example, staff
may record information each day on a spreadsheet without a clear key for
each row, leaving the data engineer to infer how those rows relate to
business entities.

This way of looking at the primary key departs from treating it only as
a technical constraint, or even as a pure link between a record and a
real-world entity.

The emphasis on lookup is significant. The ability of a primary key to
serve as a reliable link to the real world is directly proportional to
the quality of the lookup process that supports it. It helps the data
engineer see not only that primary keys can fail, but where the failure
is likely to sit within the business process.

Implementing and maintaining a strong primary-key process carries cost.
Quite often, this cost becomes a barrier to maintaining accurate
correspondence between data records and business entities. The result is
duplication or loss of information:

\begin{itemize}
\tightlist
\item
  When record creation is easier than updates, duplicate entries
  proliferate.
\item
  When creation or updates become burdensome, entities may go
  unrecorded, overwritten, or repurposed.
\end{itemize}

Any primary-key issue inevitably traces back to the underlying business
process.

If there are issues with the primary key, the data engineer must do more
than accept the issue at face value. The engineer must work to find the
root cause in the business process.

Given its foundational role, the primary key is the natural starting
point for a data engineer examining a new system.

Often the database does not come with keys implemented as constraints,
for example when uniqueness is enforced in the application layer. In
these cases, a data engineer investigating a new system would do well to
infer the primary keys for each table and then ask how those keys are
created, retrieved, and deleted in business processes.

A direct corollary of the primary key's significance is that relational
algebra becomes indispensable.

Experienced engineers think in keys. They rely on algebraic operations
to track and communicate data rather than memorise every detail of
source processes. Discussion of relational algebra lies beyond this
text's scope.

\section{Immutable and mutable
entities}\label{immutable-and-mutable-entities}

It is helpful to divide business entities into two broad types:
immutable and mutable.

Immutable entities are those that do not change. If they do change, the
change is treated as the creation of a different entity, and the latest
version is considered the correct one. In other words, neither the
entity nor its attributes evolve over time.

Examples include a completed bank transfer or an importer lodging an
import declaration into a country. If the import declaration is altered,
it is regarded as a new entity. Products can also be treated as
immutable. A new model of phone is considered a different product from
the previous version.

Mutable entities are those that may change over time while still being
considered the same entity.

A customer may change purchasing habits. An employee in an HR system may
update personal details or gain new qualifications. These changes are
important but, according to business intent, still part of the same
entity.

This distinction is helpful because most business processes are built
with one of these perspectives in mind.

A digital system may have sub-systems that manage both. For example, a
sales system may record purchases as immutable entities, treat products
themselves as immutable, but also register information about customers
whose attributes and behaviour are mutable.

One way to distinguish mutable from immutable entities is the presence
of a datetime component in the primary key to indicate the validity
period of the record.

In warehousing terminology, such a record is a Type II record. The
datetime signals that it is important to track changes to the entity's
attributes over time. This tracking may be implemented in the source
system and maintained by business processes, or it may be absent and
later added by a data engineer to meet business intent.

The two major types of business entity, immutable and mutable,
correspond to two major types of data engineering methods:
\textbf{entity processing} and \textbf{entity tracking}.

The distinction between immutable and mutable entities shapes the
overall approach for the downstream pipeline.

{\def\LTcaptype{none} % do not increment counter
\begin{longtable}[]{@{}
  >{\raggedright\arraybackslash}p{(\linewidth - 4\tabcolsep) * \real{0.3333}}
  >{\raggedright\arraybackslash}p{(\linewidth - 4\tabcolsep) * \real{0.3333}}
  >{\raggedright\arraybackslash}p{(\linewidth - 4\tabcolsep) * \real{0.3333}}@{}}
\toprule\noalign{}
\begin{minipage}[b]{\linewidth}\raggedright
\end{minipage} & \begin{minipage}[b]{\linewidth}\raggedright
Immutable entities
\end{minipage} & \begin{minipage}[b]{\linewidth}\raggedright
Mutable entities
\end{minipage} \\
\midrule\noalign{}
\endhead
\bottomrule\noalign{}
\endlastfoot
Basic question & What happened? & What changed? \\
Entity behaviour & Treated as fixed once created; changes are new
versions & Changes over time while remaining the same entity \\
Typical examples & Sales orders, bank transfers, import declarations &
Customers, employees, suppliers, accounts \\
Data engineering approach & Entity processing & Entity tracking \\
Pipeline focus & Process the entity through business logic & Track
attributes, behaviour, and state over time \\
\end{longtable}
}

These are the subjects of the next two chapters:
\hyperref[docs-creating-information-entity-processing]{Entity
processing} and
\hyperref[docs-creating-information-entity-tracking]{Entity tracking}.

\begin{quote}
\textbf{Key ideas.}

A data record is an intermediary of the business entity behind it.

Primary keys establish correspondence between data records and business
entities through reliable lookup.

The quality of a primary key depends on the quality of the business
process that creates, retrieves, and maintains it.

Immutable entities lead to entity processing.

Mutable entities lead to entity tracking.
\end{quote}

\cleardoublepage

\chapter{Entity
processing}\label{docs-creating-information-entity-processing}

\section{Processing immutable
entities}\label{processing-immutable-entities}

Entity processing is the main pattern for
\hyperref[docs-creating-information-mapping-the-data-world]{immutable
business entities}: entities that are treated as fixed once created.

Examples include a sales order, a bank transfer, or an import
declaration. If such an entity changes, the change is treated as though
it were a different entity. For immutable entities, the business
objective is to process the entity according to defined business logic.
Hence the term \textbf{entity processing}.

Immutable entities share common characteristics:

\begin{itemize}
\tightlist
\item
  \textbf{Versions}---each version can be considered a new entity, with
  the latest version replacing all others as the correct version.
\item
  \textbf{Header and details structure}---the header represents the main
  entity, while the details represent its components. A sales order and
  its line items are a typical example.
\item
  \textbf{Subprocesses}---multiple business processes are needed to
  complete the business activity on the entity.
\item
  \textbf{Good or bad}---the entity is evaluated by outcome, such as
  profitable or unprofitable, compliant or non-compliant, fraudulent or
  legitimate.
\item
  \textbf{Actors}---actors persist across multiple entities, such as a
  customer making multiple purchases over time.
\end{itemize}

These characteristics explain the goals of entity processing:

\begin{itemize}
\tightlist
\item
  Surface insight from the underlying business process.
\item
  Identify macro-level trends, such as sales by region.
\item
  Measure time between process milestones.
\item
  Provide drill-through to specific events for troubleshooting and
  root-cause analysis.
\item
  Provide a foundation for predictive analytics by engineering
  attributes that can be used to predict good or bad outcomes.
\item
  Understand actors whose behaviour affects business outcomes.
\end{itemize}

The first task is to identify the \textbf{entity of interest}.

Usually this is clear. Sometimes it is not. A sales transaction may have
multiple line items. Depending on the business question, the transaction
may be the entity of interest, or each line item may be the entity of
interest.

The correct lens is the one that supports business action.

Once the entity of interest is identified, the data engineer tells the
high-level story of that entity while preserving the necessary detail.
\hyperref[docs-creating-information-expressiveness-and-fragment-modelling]{Fragment
modelling} is suited to this task because different aspects of the
entity can be maintained in different fragments.

\section{Building the pipeline}\label{building-the-pipeline}

A pipeline for entity processing follows three broad passes:

\begin{itemize}
\tightlist
\item
  \textbf{Filter}---remove irrelevant data and establish the minimum
  meaningful structure.
\item
  \textbf{Map}---compute reusable blocks of business meaning.
\item
  \textbf{Reduce}---aggregate to the grain required for analysis.
\end{itemize}

\subsection{First pass---Filter}\label{first-passfilter}

The first pass works predominantly with incoming raw data. Its aim is to
establish structure and meaning, beginning with keys.

It has three stages: identifying keys, defining reference tables, and
extracting transactions.

\subsubsection{Identifying keys}\label{identifying-keys}

The first task is to identify the business keys that serve as primary
keys to the raw data. These may be defined as database constraints in
the source data, or inferred through exploratory analysis.

Defining primary and foreign keys is a simple but powerful way to
\hyperref[docs-creating-information-mapping-the-data-world]{bridge the
data world and the business world}.

In some cases, key columns need to be created. For immutable entities,
the two common scenarios are versioning and sequence numbers.

In versioning, a typical application might use \texttt{{[}Sales\ ID{]}}
and \texttt{{[}Previous\ sales\ ID{]}} to indicate that one record
replaces another. A clearer approach is to introduce
\texttt{{[}Original\ sales\ ID{]}} and
\texttt{{[}Sales\ version\ number{]}}, using partitioning and ordering
to establish lineage. The original \texttt{{[}Sales\ ID{]}} can still be
retained for joins, but should be renamed
\texttt{{[}Sales\ version\ ID{]}}.

In header-detail structures, detail rows are often stored as simple
lists. For example, a \texttt{{[}Sales\ ID{]}} may be associated with
multiple \texttt{{[}Sales\ item\ ID{]}} entries. Introducing
\texttt{{[}Sales\ item\ sequence\ number{]}} can clarify the
relationship between the order and its line items.

Keys do not need to be enforced as physical constraints. They can be
stored in metadata tables. What matters is that they communicate meaning
and can be looked up by consumers of the data.

\subsubsection{Defining reference
tables}\label{defining-reference-tables}

Reference tables are small, slow-moving, descriptive tables that
describe the content of larger, fast-moving transaction tables.

Examples include \texttt{RefProduct}, which describes the products being
sold, and \texttt{RefLocation}, which describes the store where the sale
occurred.

Reference tables are important for two reasons:

\begin{itemize}
\tightlist
\item
  \textbf{Expressiveness}---they provide a central place to describe the
  attributes of a business entity, rather than embedding all information
  in the transaction table.
\item
  \textbf{Efficiency}---updates can be made to a small reference table
  rather than to large volumes of transaction data.
\end{itemize}

The data engineer can use source-system reference tables for
inspiration, but should not be limited by them. Application reference
tables are usually designed to support system functionality, not
business insight.

The data engineer can make reference data more expressive by:

\begin{itemize}
\tightlist
\item
  Renaming incoming tables and columns to reflect business intent.
\item
  Combining fragmented application reference tables when this clarifies
  the business concept.
\item
  Creating new reference tables for analytical use.
\item
  Adding default rows such as \texttt{Unknown\ product} where analytical
  use requires them.
\item
  Using reference data to identify conformance across the warehouse.
\item
  Adding metadata such as business definitions and descriptions.
\end{itemize}

\subsubsection{Extracting transaction
tables}\label{extracting-transaction-tables}

Transaction tables are fast-moving tables that record business events.
They are the primary source of informational value for the business.

When extracting transaction tables, use names that pair with the
reference tables describing them. For example, \texttt{Sales} should be
matched with \texttt{RefSales}, \texttt{Inspection} with
\texttt{RefInspection}, and \texttt{AuditResult} with
\texttt{RefAuditResult}.

Filter out noisy data, including unnecessary columns and, where
justifiable and computationally simple, unnecessary rows.

Removing rows requires judgement. In large tables, reversing such a
decision later can be costly for the pipeline. A prudent approach is to
remove only rows that are clearly irrelevant, such as stray test
records, and handle more ambiguous filtering scenarios in later passes
through mapping or classification tables.

Filtering is not only for efficiency. It removes noise and improves
clarity for both the data engineer and the business.

Transactional inconsistency is an important case. During batch
extraction, inconsistency can arise within the batch. Depending on
business use, early records may need to be filtered out until the next
batch, provided the next extraction picks up the skipped records.

The transaction table should be as narrow as possible. Pass as much
descriptive information as possible to reference tables.

\textbf{Example extracted tables}

\texttt{SalesItem}

{\def\LTcaptype{none} % do not increment counter
\begin{longtable}[]{@{}llrr@{}}
\toprule\noalign{}
Sales ID & Item ID & Quantity & Sale value \\
\midrule\noalign{}
\endhead
\bottomrule\noalign{}
\endlastfoot
1001 & A & 2 & 40 \\
1001 & B & 1 & 20 \\
1002 & A & 3 & 60 \\
\end{longtable}
}

\texttt{ItemSupplierCost}

{\def\LTcaptype{none} % do not increment counter
\begin{longtable}[]{@{}lr@{}}
\toprule\noalign{}
Item ID & Supplier cost \\
\midrule\noalign{}
\endhead
\bottomrule\noalign{}
\endlastfoot
A & 12 \\
B & 15 \\
\end{longtable}
}

When these steps are applied, the data engineer will have:

\begin{itemize}
\tightlist
\item
  Extracted the minimum information needed to remove noise and preserve
  clarity.
\item
  Built a basic map of information through primary and foreign key
  relationships.
\item
  Created expressive meaning through clear naming and reference tables.
\item
  Applied change detection to incoming records to support downstream
  change tracking.
\item
  Set up a foundation for future optimisation, including incremental
  extraction.
\end{itemize}

\subsection{Second pass---Map}\label{second-passmap}

The first pass extracts raw data, clarifies basic meaning, and
establishes the foundation for deeper analytical work.

The second pass builds on this foundation. It creates useful
information---derived, calculated, and shaped for business insight.

This pass should avoid working directly from raw data. It should use the
outputs from the first pass.

The mindset here is storytelling. The data engineer should not be
limited by source-system concepts, or even by the business terms
initially provided by stakeholders. The work requires appropriate
creativity: reshaping language and logic so the data makes sense for
business decision-making.

The guiding question is:

\begin{quote}
What valuable pieces of information would be useful for business
insight?
\end{quote}

Often, this information is latent in the data and must be explicitly
calculated. For example, the concept of a non-compliant transaction may
not be stated in the source data. It must be inferred.

The result should be stored in standalone tables rather than added as
extra columns to extracted tables. This follows fragment modelling and
provides a clean container for complex logic.

As in the first pass, information should be offloaded to reference
tables and transaction tables kept narrow. This may require new
reference tables because the concept does not exist in the source data.

For example, if the first pass creates \texttt{Sales},
\texttt{SalesItem}, \texttt{ItemSupplierCost}, and
\texttt{SalesItemRefund}, the second pass may create
\texttt{SalesItemMargin} to compute profitability on each sales item.

The full set of tables would look like this.

\textbf{First pass from the source data}

\begin{itemize}
\tightlist
\item
  \texttt{Sales}---header for the sales event.
\item
  \texttt{SalesItem}---items sold in the sale and their prices.
\item
  \texttt{ItemSupplierCost}---cost to supply each item, stored as a Type
  II table because supply cost can change.
\item
  \texttt{SalesItemRefund}---refunds recorded against sales items.
\end{itemize}

\textbf{Second pass of derived computation}

\begin{itemize}
\tightlist
\item
  \texttt{SalesItemMargin}---profitability of each sales item after
  accounting for sale value, refunds, and costs.
\end{itemize}

\textbf{Example structure of \texttt{SalesItemMargin}}

{\def\LTcaptype{none} % do not increment counter
\begin{longtable}[]{@{}llrrr@{}}
\toprule\noalign{}
Sales ID & Item ID & Sale value & Supplier cost & Margin \\
\midrule\noalign{}
\endhead
\bottomrule\noalign{}
\endlastfoot
1001 & A & 40 & 24 & 16 \\
1001 & B & 20 & 15 & 5 \\
1002 & A & 60 & 36 & 24 \\
\end{longtable}
}

\subsection{Third pass---Reduce}\label{third-passreduce}

The third pass focuses on aggregation.

At this stage, the data engineer takes the detailed fragments from the
first and second passes and rolls them up to the grain of the entity of
interest. Source systems often contain fine-grained detail, while
business attention is centred on the entity itself.

This stage is a major opportunity for the data engineer to add value.
Aggregations can be technically complex, especially when they involve
multiple sources or complex business rules. By building them into the
pipeline, the data engineer delivers ready-made, trusted information
blocks that can be used without repeated manual work.

Continuing the earlier example, suppose the first pass produced
\texttt{Sales}, \texttt{SalesItem}, \texttt{ItemSupplierCost}, and
\texttt{SalesItemRefund}, and the second pass created
\texttt{SalesItemMargin}.

In the third pass, this information can be aggregated to the sales level
to produce \texttt{SalesProfit}. This table combines the total margin
from all items in the sale, applies any additional business rules, and
arrives at a final profit figure for the sale. A reference table,
\texttt{RefSalesProfit}, can include a binary flag indicating whether
the sale was profitable.

This pass is also where the data engineer computes when an entity
reaches specific milestones. A sales order may have milestones for
placement, confirmation, shipment, and completion. These dates can be
recorded in \texttt{SalesProcessMilestone}.

\textbf{Example structure of \texttt{SalesProfit}}

{\def\LTcaptype{none} % do not increment counter
\begin{longtable}[]{@{}lrl@{}}
\toprule\noalign{}
Sales ID & Total margin & Is profitable \\
\midrule\noalign{}
\endhead
\bottomrule\noalign{}
\endlastfoot
1001 & 21 & true \\
1002 & 24 & true \\
\end{longtable}
}

\textbf{Example structure of \texttt{SalesProcessMilestone}}

{\def\LTcaptype{none} % do not increment counter
\begin{longtable}[]{@{}
  >{\raggedright\arraybackslash}p{(\linewidth - 10\tabcolsep) * \real{0.1667}}
  >{\raggedright\arraybackslash}p{(\linewidth - 10\tabcolsep) * \real{0.1667}}
  >{\raggedright\arraybackslash}p{(\linewidth - 10\tabcolsep) * \real{0.1667}}
  >{\raggedright\arraybackslash}p{(\linewidth - 10\tabcolsep) * \real{0.1667}}
  >{\raggedright\arraybackslash}p{(\linewidth - 10\tabcolsep) * \real{0.1667}}
  >{\raggedright\arraybackslash}p{(\linewidth - 10\tabcolsep) * \real{0.1667}}@{}}
\toprule\noalign{}
\begin{minipage}[b]{\linewidth}\raggedright
Sales ID
\end{minipage} & \begin{minipage}[b]{\linewidth}\raggedright
Placed date
\end{minipage} & \begin{minipage}[b]{\linewidth}\raggedright
Confirmed date
\end{minipage} & \begin{minipage}[b]{\linewidth}\raggedright
Shipped date
\end{minipage} & \begin{minipage}[b]{\linewidth}\raggedright
Completed date
\end{minipage} & \begin{minipage}[b]{\linewidth}\raggedright
Is meeting SLA
\end{minipage} \\
\midrule\noalign{}
\endhead
\bottomrule\noalign{}
\endlastfoot
1001 & 2023-01-03 & 2023-01-04 & 2023-01-06 & 2023-01-09 & true \\
1002 & 2023-01-05 & 2023-01-06 & 2023-01-08 & 2023-01-15 & false \\
\end{longtable}
}

The outputs of this pass are valuable building blocks.

In business intelligence scenarios, they provide ready-to-use measures
for dashboards and reports. In data science scenarios, they are equally
important: much of feature engineering involves aggregating detailed
attributes back to the correct grain.

\subsection{Summary of the three
passes}\label{summary-of-the-three-passes}

When complete, the pipeline produces reusable information blocks:
expressive of business intent, supported by metadata, and delivered
through an efficient and robust process.

These blocks are ready for downstream use, including self-service
dimensional models and feature engineering.

An overview of the three passes and examples is summarised in Figure 1.

\pandocbounded{\includegraphics[keepaspectratio]{book/epub-assets/diagram-003.png}}

\emph{Figure 1. The three passes of entity processing, showing both the
purpose of each pass and example artefacts produced in a sales
pipeline.}

\section{Common problems}\label{common-problems}

New data engineers should watch for common mistakes in entity
processing.

\textbf{Starting with the details.} Work should follow the shape of the
business process. Start with the main entity, such as \texttt{Sales},
before moving to associated detail tables.

\textbf{Neglecting reference tables.} A common mistake is to store too
much information in the transaction table instead of keeping it narrow
and letting reference tables carry the descriptive load.

\textbf{Accepting ``garbage in, garbage out.''} Source systems often
have poor or cryptic names because presentation is handled in the user
interface. Renaming and reshaping data to reflect business meaning is
one of the most direct ways to make the model recognisable and useful.

\textbf{Skipping aggregation.} The third pass aggregates information to
the entity level. New engineers often stop at the first or second pass
without completing this step.

\textbf{Compounding logic.} Placing too much computation into a single
table makes the pipeline difficult to understand and maintain. Fragment
modelling suggests breaking these computations into standalone tables.

New engineers often mix the first and second passes: extracting raw data
while performing complex transformation. This compounds computational
complexity and weakens future options for incremental extraction.

The same temptation applies to all three passes. Convenience at the
beginning creates fragility later.

\begin{quote}
\textbf{Key ideas.}

Entity processing applies to immutable business entities.

The first task is to identify the entity of interest.

The three passes are filter, map, and reduce.

The reduce pass returns detailed fragments to the grain required by
business intent.
\end{quote}

\cleardoublepage

\chapter{Entity
tracking}\label{docs-creating-information-entity-tracking}

\section{Tracking mutable entities}\label{tracking-mutable-entities}

Entity tracking is the main pattern for
\hyperref[docs-creating-information-mapping-the-data-world]{mutable
business entities}: those that change over time while still being
considered the same entity.

Examples include a customer whose contact details, preferences, or
purchasing behaviour evolve; an employee whose qualifications, role, or
performance record changes; or a supplier whose compliance status shifts
over time.

\hyperref[docs-creating-information-entity-processing]{Entity
processing} asks what happened to a fixed entity. Entity tracking asks
how an entity changed while remaining itself.

Mutable entities share common characteristics:

\begin{itemize}
\tightlist
\item
  \textbf{Status}---Instead of "versions", the entity has statuses that
  change over time.
\item
  \textbf{Changing attributes}---the entity has attributes whose changes
  define its mutation.
\item
  \textbf{Actions}---the entity acts within business processes, and its
  performance or behaviour is of ongoing concern.
\item
  \textbf{Audit}---the entity may be audited periodically or in response
  to events.
\item
  \textbf{Registration}---the entity is usually underpinned by a
  registration process that establishes its initial presence in the
  system.
\end{itemize}

These characteristics explain the goals of entity tracking:

\begin{itemize}
\tightlist
\item
  Monitor macro-level trends in the number of entities over time,
  including changes in the composition of the population.
\item
  Track management activities associated with these entities, such as
  registration and audit.
\item
  Monitor when entities reach performance milestones defined by business
  thresholds.
\item
  Detect anomalous entities for targeted action and early intervention.
\item
  Provide drill-down information about the entity and the changes it has
  undergone.
\item
  Understand the effect of entities on other business outcomes.
\end{itemize}

\section{Two ways to track}\label{two-ways-to-track}

There are two common ways to track a mutable entity. The first is by its
attributes. The second is by its behaviour.

Tracking by attributes records what the entity is over time. An entity
is registered into a database, and tables record attributes such as
name, address, status, account type, account level, and so on. These
tables are usually Type II tables, where each row is labelled with a
validity period using columns such as \texttt{{[}Start\ date{]}} and
\texttt{{[}End\ date{]}}.

Tracking by behaviour records what the entity does over time. Behaviour
may be observed through audits, transactions, events, or actions taken
in business processes.

Many systems have both. A bank, for example, may track accounts in two
ways:

\begin{itemize}
\tightlist
\item
  \textbf{Attribute}---bank accounts have evolving attributes such as
  account holder, account type, and account level.
\item
  \textbf{Behaviour}---deposits, withdrawals, transfers, and interest
  accruals are actions associated with the account.
\end{itemize}

The same three-pass structure applies to both modes, but with different
emphasis. Attribute tracking introduces distinctive temporal machinery
early. Behaviour tracking enters in the third pass when actions are
reduced into entity-level signals over time.

\section{Building the pipeline}\label{building-the-pipeline-1}

Much of the approach from entity processing still applies. The data
engineer must identify the entity of interest, craft
\hyperref[docs-creating-information-reference-data]{expressive reference
tables}, preserve necessary detail, and tell a coherent story.

The difference is that mutable entities introduce time into the entity
itself. In entity processing, time usually appears around the entity:
orders, confirmation, shipment, completion. In entity tracking, time
defines the state of the entity.

\subsection{First pass---Filter and
compress}\label{first-passfilter-and-compress}

The first pass removes noise, establishes keys, and builds the basic
semantic map.

For attribute tracking, this pass also introduces \textbf{temporal
compression}. Temporal compression collapses consecutive rows that carry
identical information across adjacent time periods. It is akin to
run-length encoding.

Consider a bank account whose attributes are recorded over time.

\textbf{Example before compression}

{\def\LTcaptype{none} % do not increment counter
\begin{longtable}[]{@{}
  >{\raggedright\arraybackslash}p{(\linewidth - 10\tabcolsep) * \real{0.1667}}
  >{\raggedright\arraybackslash}p{(\linewidth - 10\tabcolsep) * \real{0.1667}}
  >{\raggedright\arraybackslash}p{(\linewidth - 10\tabcolsep) * \real{0.1667}}
  >{\raggedright\arraybackslash}p{(\linewidth - 10\tabcolsep) * \real{0.1667}}
  >{\raggedright\arraybackslash}p{(\linewidth - 10\tabcolsep) * \real{0.1667}}
  >{\raggedright\arraybackslash}p{(\linewidth - 10\tabcolsep) * \real{0.1667}}@{}}
\toprule\noalign{}
\begin{minipage}[b]{\linewidth}\raggedright
Account ID
\end{minipage} & \begin{minipage}[b]{\linewidth}\raggedright
Start date
\end{minipage} & \begin{minipage}[b]{\linewidth}\raggedright
End date
\end{minipage} & \begin{minipage}[b]{\linewidth}\raggedright
Account type
\end{minipage} & \begin{minipage}[b]{\linewidth}\raggedright
Account level
\end{minipage} & \begin{minipage}[b]{\linewidth}\raggedright
Account status
\end{minipage} \\
\midrule\noalign{}
\endhead
\bottomrule\noalign{}
\endlastfoot
12345 & 2022-01-01 & 2022-06-30 & Savings & Silver & Active \\
12345 & 2022-07-01 & 2023-06-30 & Savings & Silver & Active \\
\end{longtable}
}

\textbf{After compression}

{\def\LTcaptype{none} % do not increment counter
\begin{longtable}[]{@{}
  >{\raggedright\arraybackslash}p{(\linewidth - 10\tabcolsep) * \real{0.1667}}
  >{\raggedright\arraybackslash}p{(\linewidth - 10\tabcolsep) * \real{0.1667}}
  >{\raggedright\arraybackslash}p{(\linewidth - 10\tabcolsep) * \real{0.1667}}
  >{\raggedright\arraybackslash}p{(\linewidth - 10\tabcolsep) * \real{0.1667}}
  >{\raggedright\arraybackslash}p{(\linewidth - 10\tabcolsep) * \real{0.1667}}
  >{\raggedright\arraybackslash}p{(\linewidth - 10\tabcolsep) * \real{0.1667}}@{}}
\toprule\noalign{}
\begin{minipage}[b]{\linewidth}\raggedright
Account ID
\end{minipage} & \begin{minipage}[b]{\linewidth}\raggedright
Start date
\end{minipage} & \begin{minipage}[b]{\linewidth}\raggedright
End date
\end{minipage} & \begin{minipage}[b]{\linewidth}\raggedright
Account type
\end{minipage} & \begin{minipage}[b]{\linewidth}\raggedright
Account level
\end{minipage} & \begin{minipage}[b]{\linewidth}\raggedright
Account status
\end{minipage} \\
\midrule\noalign{}
\endhead
\bottomrule\noalign{}
\endlastfoot
12345 & 2022-01-01 & 2023-06-30 & Savings & Silver & Active \\
\end{longtable}
}

In addition to compression, the first pass should assign a surrogate key
to each validity period. For example, \texttt{{[}Account\ SK{]}} may
represent the combination of \texttt{{[}Account\ ID{]}} and
\texttt{{[}Start\ date{]}}.

The first pass should also normalise validity periods.

Open-ended histories should be given explicit boundaries. A missing
start date can be replaced with a minimum sentinel date such as
\texttt{1900-01-01}; a missing end date can be replaced with a maximum
sentinel date such as \texttt{9999-12-31}. This gives every row a
bounded interval and makes later timeline construction simpler.

The data engineer should also decide how to handle gaps. In some systems
a record may be deleted and later re-created even though the business
regards the entity as continuous. A row may need to be inserted to
indicate non-existence, with \texttt{lag} functions to fill the dates.

For behaviour tracking, the first pass usually resembles entity
processing. The data engineer filters and structures behaviour records,
such as transactions, audit results, observations, or actions. The
important temporal requirement is to preserve the event time of each
behaviour.

\subsection{Second pass---Build the
timeline}\label{second-passbuild-the-timeline}

An entity usually has attributes stored in multiple tables.

For immutable entities, retrieving attributes from multiple tables is
usually a matter of joining on the appropriate keys. This is not the
case for mutable entities. In a mutable entity, each attribute table may
be Type II, meaning each row is defined by a validity period.

Joining attributes in this scenario means solving the overlapping-window
problem.

The timeline table is the central artefact of attribute tracking.

A source system may track each changing attribute separately. Account
type, account level, and account status may each have their own history
table, with their own validity periods. The difficulty is that the
business does not usually ask for these histories separately. It asks
what the account was at a point in time.

The timeline table answers this question. It converts multiple
overlapping histories into one coherent sequence of entity states.

An account may have attributes stored across multiple tables:

\begin{itemize}
\tightlist
\item
  \texttt{Bank.AccountTypeHistory}
\item
  \texttt{Bank.AccountLevelHistory}
\item
  \texttt{Bank.AccountStatusHistory}
\end{itemize}

The goal is to create \texttt{Bank.AccountTimeline}, which expresses the
valid combination of attributes for each period.

\textbf{Example structure of \texttt{Bank.AccountTimeline}}

{\def\LTcaptype{none} % do not increment counter
\begin{longtable}[]{@{}
  >{\raggedright\arraybackslash}p{(\linewidth - 10\tabcolsep) * \real{0.1667}}
  >{\raggedright\arraybackslash}p{(\linewidth - 10\tabcolsep) * \real{0.1667}}
  >{\raggedright\arraybackslash}p{(\linewidth - 10\tabcolsep) * \real{0.1667}}
  >{\raggedright\arraybackslash}p{(\linewidth - 10\tabcolsep) * \real{0.1667}}
  >{\raggedright\arraybackslash}p{(\linewidth - 10\tabcolsep) * \real{0.1667}}
  >{\raggedright\arraybackslash}p{(\linewidth - 10\tabcolsep) * \real{0.1667}}@{}}
\toprule\noalign{}
\begin{minipage}[b]{\linewidth}\raggedright
Account ID
\end{minipage} & \begin{minipage}[b]{\linewidth}\raggedright
Account type SK
\end{minipage} & \begin{minipage}[b]{\linewidth}\raggedright
Account level SK
\end{minipage} & \begin{minipage}[b]{\linewidth}\raggedright
Account status SK
\end{minipage} & \begin{minipage}[b]{\linewidth}\raggedright
Start date
\end{minipage} & \begin{minipage}[b]{\linewidth}\raggedright
End date
\end{minipage} \\
\midrule\noalign{}
\endhead
\bottomrule\noalign{}
\endlastfoot
12345 & 1 & 1 & 1 & 2022-01-01 & 2023-06-30 \\
12345 & 1 & 2 & 1 & 2023-07-01 & 2024-12-31 \\
\end{longtable}
}

The core SQL pattern is an interval overlap join.

\begin{verbatim}
select
      th.[Account ID]
    , th.[Account type SK]
    , lh.[Account level SK]
    , sh.[Account status SK]
    , greatest(th.[Start date], lh.[Start date], sh.[Start date])   as [Start date]
    , least(   th.[End date],   lh.[End date],   sh.[End date]  )   as [End date]
from Bank.AccountTypeHistory   th
inner join Bank.AccountLevelHistory  lh  on   lh.[Account ID] = th.[Account ID]
                                         and greatest(th.[Start date], lh.[Start date])                   < least(th.[End date], lh.[End date])
inner join Bank.AccountStatusHistory sh  on   sh.[Account ID] = th.[Account ID]
                                         and greatest( th.[Start date], lh.[Start date], sh.[Start date]) < least(th.[End date], lh.[End date], sh.[End date] );
\end{verbatim}

The timeline table is produced through a temporal join. Notice that,
without the time component, this is an ordinary join on the entity key.
This is worth remembering. The timeline table is simply what the join
would have looked like "at that point in time", but pre-computed for all
time intervals in the most computationally compact form.

In this example, we use inner join for simplicity. Real systems also
require careful handling of gaps in history.

\begin{quote}
\textbf{Advanced note: two kinds of time.}

Some tracking problems involve more than one temporal dimension.

For example, a user may report an account status for a
business-effective period, then later change that reported period. The
status has business dates describing when it was effective, but the
system also has architectural dates describing when that report was
recorded or changed.

This creates questions such as:

\begin{quote}
How many valid account holders were there at time X, according to what
was known at time Z?
\end{quote}

These problems can become extremely messy because the data engineer is
no longer tracking only what was true in business time, but also what
was known in system time.

We do not develop the full pattern here. The important point is that the
timeline approach generalises. A timeline can be extended across an
additional time dimension, giving the user a disciplined way to answer
these questions while maintaining the clarity on the multiple time
dimensions.
\end{quote}

\subsection{Third pass---Create snapshots and infer
behaviour}\label{third-passcreate-snapshots-and-infer-behaviour}

The third pass produces entity-level outputs over time.

For attribute tracking, the main output is an end-of-period table: a
snapshot of the entity at regular reporting intervals. This is created
by joining the timeline to a calendar table and selecting the row that
represents the entity's state at the end of each period.

For example, joining \texttt{Bank.AccountTimeline} on the validity
period to a distinct list of end-of-month dates returns
\texttt{Bank.AccountEndOfMonth}.

\textbf{Example structure of \texttt{Bank.AccountEndOfMonth}}

{\def\LTcaptype{none} % do not increment counter
\begin{longtable}[]{@{}
  >{\raggedright\arraybackslash}p{(\linewidth - 8\tabcolsep) * \real{0.2000}}
  >{\raggedright\arraybackslash}p{(\linewidth - 8\tabcolsep) * \real{0.2000}}
  >{\raggedright\arraybackslash}p{(\linewidth - 8\tabcolsep) * \real{0.2000}}
  >{\raggedright\arraybackslash}p{(\linewidth - 8\tabcolsep) * \real{0.2000}}
  >{\raggedright\arraybackslash}p{(\linewidth - 8\tabcolsep) * \real{0.2000}}@{}}
\toprule\noalign{}
\begin{minipage}[b]{\linewidth}\raggedright
Account ID
\end{minipage} & \begin{minipage}[b]{\linewidth}\raggedright
Period end date
\end{minipage} & \begin{minipage}[b]{\linewidth}\raggedright
Account type SK
\end{minipage} & \begin{minipage}[b]{\linewidth}\raggedright
Account level SK
\end{minipage} & \begin{minipage}[b]{\linewidth}\raggedright
Account status SK
\end{minipage} \\
\midrule\noalign{}
\endhead
\bottomrule\noalign{}
\endlastfoot
12345 & 2023-06-30 & 1 & 1 & 1 \\
12345 & 2023-07-31 & 1 & 2 & 1 \\
12345 & 2023-08-31 & 1 & 2 & 1 \\
\end{longtable}
}

This table is a building block. It gives downstream consumers regular
access to the entity state without requiring them to solve overlapping
validity windows.

For behaviour tracking, the third pass aggregates behaviour back to the
entity over time.

Behaviour records are usually at the grain of the action. A bank
transaction table, for example, may be at the grain of
\texttt{{[}Transaction\ ID{]}}. The tracking question is usually at the
grain of \texttt{{[}Account\ ID{]}} over time.

The data engineer therefore aggregates behaviour into signals about the
entity: balances, counts, failures, first occurrences, repeated events,
thresholds, or changes in pattern.

Following the previous example, \texttt{Bank.Transaction} holds records
of an entity's actions, the withdrawals and deposits of an account. It
is at the grain of \texttt{{[}Transaction\ ID{]}}, which is much finer
than \texttt{{[}Account\ ID{]}}. To track behaviour of the bank account,
the data engineer can aggregate by \texttt{{[}Account\ ID{]}} to produce
\texttt{Bank.AccountValue}, described by \texttt{Bank.RefAccountValue}.
Example columns in \texttt{Bank.RefAccountValue} would be
\texttt{{[}Is\ a\ millionaire{]}} or
\texttt{{[}Is\ a\ new\ millionaire\ in\ the\ last\ 3\ months{]}}.

\textbf{Example structure of \texttt{Bank.AccountValue} with
\texttt{Bank.RefAccountvalue}}

{\def\LTcaptype{none} % do not increment counter
\begin{longtable}[]{@{}
  >{\raggedright\arraybackslash}p{(\linewidth - 8\tabcolsep) * \real{0.1875}}
  >{\raggedright\arraybackslash}p{(\linewidth - 8\tabcolsep) * \real{0.1875}}
  >{\raggedleft\arraybackslash}p{(\linewidth - 8\tabcolsep) * \real{0.2500}}
  >{\raggedright\arraybackslash}p{(\linewidth - 8\tabcolsep) * \real{0.1875}}
  >{\raggedright\arraybackslash}p{(\linewidth - 8\tabcolsep) * \real{0.1875}}@{}}
\toprule\noalign{}
\begin{minipage}[b]{\linewidth}\raggedright
Account ID
\end{minipage} & \begin{minipage}[b]{\linewidth}\raggedright
Period end date
\end{minipage} & \begin{minipage}[b]{\linewidth}\raggedleft
Closing balance
\end{minipage} & \begin{minipage}[b]{\linewidth}\raggedright
Is millionaire
\end{minipage} & \begin{minipage}[b]{\linewidth}\raggedright
Is new millionaire in last 3 months
\end{minipage} \\
\midrule\noalign{}
\endhead
\bottomrule\noalign{}
\endlastfoot
12345 & 2023-06-30 & 8700.00 & false & false \\
12345 & 2023-07-31 & 9200.00 & false & false \\
12345 & 2023-08-31 & 13700.00 & false & false \\
\end{longtable}
}

With this table, a consumer should not need to recalculate balances or
scan transaction history to find millionaire accounts. The pipeline
should produce a fragment where that business state is already visible.

At an advanced level, behavioural fragments can themselves become Type
II attributes. For example, once \texttt{Bank.AccountValue} identifies
whether an account is a millionaire at each period, the data engineer
can convert that repeated period-level result into a history of
millionaire status over time.

\subsection{Summary of the three
passes}\label{summary-of-the-three-passes-1}

When complete, the pipeline produces reusable tracking blocks:
expressive business entities, clear management of time, in
computationally compact form.

These blocks are ready for downstream use where users can access
information without struggling with error-prone process of writing
time-based queries.

An overview of the three passes and examples is summarised in Figure 1.

\pandocbounded{\includegraphics[keepaspectratio]{book/epub-assets/diagram-004.png}}

\emph{Figure 1. The three passes of entity tracking, showing both the
purpose of each pass and example artefacts produced in a bank account
pipeline.}

\section{Common problems}\label{common-problems-1}

Compared with entity processing, entity tracking has additional
difficulties.

\textbf{Hidden tracking problems.} It is not always obvious that the
business is dealing with an entity tracking problem. Sometimes tracking
is hidden within what appears to be a processing task.

\textbf{Unregistered entity.} A lack of entity registration undermines
the clarity of identity. If the entity is not clearly registered,
changes cannot be reliably attributed to the same thing over time.

\textbf{Undefined memory.} Behaviour informs performance. Memory of
behaviour informs the view of performance. But it may be unclear how
much behaviour history is relevant for analysis.

\textbf{Gaps in change history.} Tracking mutable entities requires
reliable records of historical change. This is difficult when the source
system does not track changes natively.

\textbf{Multiple timelines.} Confusion arises when there is a lag
between when a change occurs, when it is entered, and when it is
recorded by the system.

There are three common types of datetime columns:

\begin{itemize}
\tightlist
\item
  \textbf{Architectural}---system-generated timestamps, such as those
  from database triggers.
\item
  \textbf{Application}---timestamps defined by application logic, such
  as \texttt{{[}Update\ datetime{]}}.
\item
  \textbf{Business}---timestamps tied to the actual business event, such
  as \texttt{{[}Sales\ order\ date{]}}.
\end{itemize}

The data engineer must engage early with stakeholders to understand the
strengths and weaknesses of each type. Clear naming conventions and
consistent usage are essential to avoid confusion and promote
expressiveness.

Temporal representation introduces another important decision. Since
attributes can change, the data engineer must decide how to represent
the entity at a given point in time:

\begin{itemize}
\tightlist
\item
  \textbf{Type I}---show the latest available attribute value.
\item
  \textbf{Type II}---show the attribute value as it was at a specific
  point in time.
\item
  \textbf{Dynamic Type I}---show the latest value as defined by the
  user's context, such as a time filter.
\end{itemize}

The choice depends on the analytical scenario. The same attribute may
require different representations depending on the business question.
Fragment modelling and timeline tables are well suited to supporting
this flexibility.

Computational performance can also become a bottleneck, especially in
end-of-period reporting. Many of these challenges require detailed
handling rather than a ready-made approach, and they often require
direct engagement with the business.

\begin{quote}
\textbf{Key ideas.}

Entity tracking applies to mutable business entities.

Entity processing asks what happened to a fixed entity. Entity tracking
asks how an entity changed while remaining itself.

Mutable entities can be tracked through attributes, behaviour, or both.

Attribute tracking requires timelines and end-of-period snapshots.

Behaviour tracking aggregates actions into information about the entity
over time.

The central difficulty of entity tracking is representing change without
losing identity.
\end{quote}

\cleardoublepage

\chapter{Reference data}\label{docs-creating-information-reference-data}

\section{A stable place for meaning}\label{a-stable-place-for-meaning}

In the chapters
\hyperref[docs-creating-information-entity-processing]{Entity
processing} and
\hyperref[docs-creating-information-entity-tracking]{Entity tracking},
reference tables appeared as a way to make business meaning expressive.
They gave names, descriptions, and analytical attributes a stable place
to live.

This chapter takes the next step. The goal is to introduce reference
data into the warehouse and connect individual systems to it. Local
systems may continue using their own terminology, codes, and
classifications, while the warehouse provides shared points of reference
for business meaning.

Two constructs are important for this purpose: \textbf{reference tables}
and \textbf{mapping tables}. Reference tables define the meanings the
business wants to use. Mapping tables connect local system values to
those shared meanings.

This chapter focuses on applying these ideas within a single system. The
next chapter extends the same pattern across
\hyperref[docs-creating-information-conforming-systems]{multiple
systems}.

\section{A point of reference}\label{a-point-of-reference}

A reference table that applies to multiple business processes becomes a
shared point of reference.

The calendar is a simple example. Most business processes can map
directly to a calendar table through a date column. A country table is
another common example. Product categories, organisational structures,
geographic hierarchies, and customer classifications often play the same
role.

Such tables are commonly called \textbf{conformed reference tables}
because multiple business processes conform to the same definition. This
makes reference data one of the main ways a warehouse gives stable form
to \hyperref[docs-foundations-data-and-organisations]{negotiated
organisational meaning}.

The term \textbf{golden record} is also used, particularly when a single
reference table represents the authoritative version of a concept.

In simple cases, systems can map directly to the shared reference. In
more complex cases, each system requires its own mapping table.

The goal is not to force every source system to use identical codes. The
goal is to provide a common reference that allows information from
different sources to be understood consistently.

\section{Building and applying shared
references}\label{building-and-applying-shared-references}

It can be tempting to select a source-system table and declare it to be
the golden reference, asking all other systems to conform to it.

This should be avoided.

The golden reference should be treated as physically and conceptually
separate from any individual source system. This preserves flexibility
and reduces coupling. If a source system changes, the warehouse
reference remains stable.

A practical approach consists of three steps:

\begin{itemize}
\tightlist
\item
  Model the local system.
\item
  Build shared references.
\item
  Apply the references during presentation.
\end{itemize}

The following example uses a sales process from a fictional system
called \texttt{Cake}.

\subsection{Step 1---Model the local
system}\label{step-1model-the-local-system}

Begin by modelling the source system in isolation. The objective at this
stage is not integration. The objective is to understand and represent
the system faithfully.

The resulting tables may look like:

\begin{itemize}
\tightlist
\item
  \texttt{Cake.Sales}
\item
  \texttt{Cake.RefSales}
\item
  \texttt{Cake.RefCountry}
\end{itemize}

\textbf{Example structure of \texttt{Cake.RefCountry}}

{\def\LTcaptype{none} % do not increment counter
\begin{longtable}[]{@{}ll@{}}
\toprule\noalign{}
Cake country code & Cake country name \\
\midrule\noalign{}
\endhead
\bottomrule\noalign{}
\endlastfoot
AUS & Aus \\
NZ & Aotearoa \\
UK & Britain \\
\end{longtable}
}

The local model should stand on its own before any reference-data work
begins.

This separation allows the engineering work to progress while mapping
decisions are still being refined.

\subsection{Step 2---Build shared
references}\label{step-2build-shared-references}

Next, construct the warehouse reference table and the mapping table.

The warehouse reference captures the business concept independently of
any individual source system.

\textbf{Example structure of \texttt{Location.RefCountry}}

{\def\LTcaptype{none} % do not increment counter
\begin{longtable}[]{@{}rll@{}}
\toprule\noalign{}
Country ID & Country name & Region \\
\midrule\noalign{}
\endhead
\bottomrule\noalign{}
\endlastfoot
1 & Australia & Oceania \\
2 & New Zealand & Oceania \\
3 & United Kingdom & Europe \\
\end{longtable}
}

\textbf{Example structure of \texttt{Cake.CountryMap}}

{\def\LTcaptype{none} % do not increment counter
\begin{longtable}[]{@{}lr@{}}
\toprule\noalign{}
Cake country code & Country ID \\
\midrule\noalign{}
\endhead
\bottomrule\noalign{}
\endlastfoot
AUS & 1 \\
NZ & 2 \\
UK & 3 \\
\end{longtable}
}

The reference table can now hold information that may never have existed
in the source system. Additional classifications, descriptions,
metadata, and analytical attributes can be added without modifying
operational systems.

If multiple systems require mapping to the same country list, the
mappings may be stored separately:

\begin{itemize}
\tightlist
\item
  \texttt{Cake.CountryMap}
\item
  \texttt{Finance.CountryMap}
\item
  \texttt{Logistics.CountryMap}
\end{itemize}

An alternative is to combine them into a single table.

\textbf{Example structure of a consolidated mapping table}

{\def\LTcaptype{none} % do not increment counter
\begin{longtable}[]{@{}llr@{}}
\toprule\noalign{}
System & Source country code & Country ID \\
\midrule\noalign{}
\endhead
\bottomrule\noalign{}
\endlastfoot
Cake & AUS & 1 \\
Cake & NZ & 2 \\
Finance & AU & 1 \\
Finance & NZL & 2 \\
\end{longtable}
}

While convenient, this introduces a polymorphic key. The meaning of the
source code depends on the value of another column. Such designs are
practical but no longer correspond cleanly to traditional
entity-relationship semantics.

\subsection{Step 3---Apply during
presentation}\label{step-3apply-during-presentation}

The final step is to apply the mapping during presentation.

This usually occurs through a dimensional model, reporting view,
semantic model, or presentation layer.

The process is straightforward:

\begin{enumerate}
\def\labelenumi{\arabic{enumi}.}
\tightlist
\item
  Retrieve the local system value.
\item
  Look up the mapping table.
\item
  Retrieve the shared reference.
\end{enumerate}

The result is that reports and analytical models operate on shared
business concepts rather than system-specific codes.

For example, a report consuming \texttt{Cake.Sales} no longer sees
country code \texttt{AUS}. It sees the warehouse country reference and
any associated classifications such as region.

\textbf{Example application}

Suppose \texttt{Cake.Sales} contains the local country code:

{\def\LTcaptype{none} % do not increment counter
\begin{longtable}[]{@{}llr@{}}
\toprule\noalign{}
Sales ID & Country code & Sales value \\
\midrule\noalign{}
\endhead
\bottomrule\noalign{}
\endlastfoot
1001 & AUS & 125.00 \\
1002 & UK & 320.00 \\
\end{longtable}
}

The presentation layer can translate the local code into the warehouse
reference by joining through the mapping table:

\begin{verbatim}
select
      s.[Sales ID]
    , s.[Sales value]
    , rc.[Country name]
    , rc.[Region]
from      Cake.Sales s
left join Cake.CountryMap     cm on cm.[Country code] = s.[Country code]
left join Location.RefCountry rc on rc.[Country ID]   = cm.[Country ID];
\end{verbatim}

The result is:

{\def\LTcaptype{none} % do not increment counter
\begin{longtable}[]{@{}lrll@{}}
\toprule\noalign{}
Sales ID & Sales value & Country name & Region \\
\midrule\noalign{}
\endhead
\bottomrule\noalign{}
\endlastfoot
1001 & 125.00 & Australia & Oceania \\
1002 & 320.00 & United Kingdom & Europe \\
\end{longtable}
}

The report consumes the shared reference rather than the local system
code. This allows multiple systems to present countries consistently
even when their underlying codes differ.

In summary, the complete pattern introduces four tables, each table
having a distinct responsibility:

{\def\LTcaptype{none} % do not increment counter
\begin{longtable}[]{@{}
  >{\raggedright\arraybackslash}p{(\linewidth - 2\tabcolsep) * \real{0.5000}}
  >{\raggedright\arraybackslash}p{(\linewidth - 2\tabcolsep) * \real{0.5000}}@{}}
\toprule\noalign{}
\begin{minipage}[b]{\linewidth}\raggedright
Table
\end{minipage} & \begin{minipage}[b]{\linewidth}\raggedright
Responsibility
\end{minipage} \\
\midrule\noalign{}
\endhead
\bottomrule\noalign{}
\endlastfoot
\texttt{Cake.Sales} & Local business transaction \\
\texttt{Cake.RefCountry} & Local system country reference \\
\texttt{Cake.CountryMap} & Translation from local country code to shared
country reference \\
\texttt{Location.RefCountry} & Shared warehouse country reference \\
\end{longtable}
}

This structure provides flexibility, stability, and reuse while keeping
responsibilities clear.

\section{Why these separations}\label{why-these-separations}

The pattern depends on several separations:

\begin{itemize}
\tightlist
\item
  the local system is separated from the warehouse reference;
\item
  the reference table is separated from the mapping table;
\item
  the building of reference data is separated from its application.
\end{itemize}

These separations guard against a common mistake: introducing shared
reference data too early and too deeply into the local system model.

The local system should first be modelled on its own terms.
\texttt{Cake.Sales} and \texttt{Cake.RefCountry} describe how the Cake
system records sales and countries. \texttt{Location.RefCountry}
describes the warehouse's shared country reference.
\texttt{Cake.CountryMap} connects the two.

Keeping these responsibilities separate has practical advantages.

First, the local system can be developed while reference mapping is
still being reviewed. Mapping countries, product categories, and
organisational structures often requires detailed business attention and
cannot always be completed immediately. This matters in projects that
require tight iteration loops.

Second, the design remains modular. Reference tables and mapping tables
can be changed independently, which aligns with fragment modelling.

Third, reference data has an amplification effect. A change to a single
reference row may affect large portions of the warehouse. If shared
references are introduced too early into local transaction tables,
changes to the warehouse reference can ripple through the local model.
Delaying the dependency until presentation reduces unnecessary coupling.

\section{When to simplify}\label{when-to-simplify}

Smaller systems may not require the full pattern.

One option is to skip \texttt{Cake.RefCountry} if the mapping to the
golden reference can be created directly.

If a mapping can be expressed directly and maintained safely,
\texttt{Cake.CountryMap} may be replaced by code. Where possible, this
should be implemented through a temporary table or CTE rather than a
large inline conditional, so that the pattern is still visible in code.

The full pattern exists because it scales well. Simpler systems can
often adopt lighter variations without losing the underlying principle.

\section{The discipline of
separation}\label{the-discipline-of-separation}

The pattern is simple, but easy to damage in practice.

Most mistakes come from collapsing distinctions that should remain
separate: treating a source-system table as the warehouse reference,
hard-coding mappings into reports, applying shared references too early,
or merging mapping logic into the reference table itself.

The pattern in this chapter has been tested in complex settings. It is
the path of least regret.

Above all, it is a discipline of separation, not a mechanical
implementation recipe.

\begin{quote}
\textbf{Key ideas.}

Reference data gives business meaning a stable place to live.

Reference tables define shared business concepts.

Mapping tables connect local system values to those shared concepts.

Build the local system first, then construct shared references, then
apply them during presentation.

Shared references allow a warehouse to remain coherent even when source
systems use different codes and classifications.

Changes to shared references can amplify across the warehouse; delay the
dependency until presentation when possible.
\end{quote}

\cleardoublepage

\chapter{Conforming
systems}\label{docs-creating-information-conforming-systems}

\section{Integration without
distortion}\label{integration-without-distortion}

Integration is common, tempting, and dangerous because it invites the
data engineer to connect tables before clarifying what kind of sameness
is involved.

Business often needs to see multiple systems together. A legacy system
may be replaced by a new platform, but the business still needs to
analyse across both systems. Different business processes may also need
to be compared through common references. A company may want to compare
production and sales by region and month. A government agency may want
to compare applications, inspections, and enforcement actions by
location or program.

New engineers often fall into two traps: forcing a union of tables that
do not naturally fit, or performing large joins that create ambiguous
grain and duplicated data. Both approaches can produce outputs that run
but no longer mean what they appear to mean.

Conforming systems is the discipline of integrating information without
distorting the
\hyperref[docs-creating-information-mapping-the-data-world]{business
entities} being represented. Fragment modelling handles this by refusing
to collapse meaning into convenience.

There are two approaches:

\begin{itemize}
\tightlist
\item
  \textbf{Vertical integration}---the same kind of entity is recorded
  across multiple systems.
\item
  \textbf{Horizontal integration}---different entities can be compared
  through shared references.
\end{itemize}

\section{Vertical integration}\label{vertical-integration}

Vertical integration applies when the same kind of entity is recorded
across multiple systems.

This commonly occurs when a legacy system is replaced by a new platform.
The systems differ, but the business entity continues. The data
engineer's task is to preserve that continuity without pretending the
two systems are identical.

Take a cake company using two systems, \texttt{CakeV1} and
\texttt{CakeV2}, to record sales. The newer system adds fields and
improves logic, but both systems record sales.

Integration proceeds in three steps:

\begin{itemize}
\tightlist
\item
  Model the systems individually.
\item
  Build conformed reference tables.
\item
  Integrate the transaction tables.
\end{itemize}

\subsection{Step 1---Model systems
individually}\label{step-1model-systems-individually}

A common mistake is to integrate too early.

Each system should first be modelled on its own terms. This preserves
local meaning and makes the later integration safer.

Suppose both systems have sales and status, but \texttt{CakeV2} also has
a marketing campaign.

The local tables might be:

\begin{itemize}
\tightlist
\item
  \texttt{CakeV1.Sales}
\item
  \texttt{CakeV1.RefStatus}
\item
  \texttt{CakeV2.Sales}
\item
  \texttt{CakeV2.RefStatus}
\item
  \texttt{CakeV2.RefCampaign}
\end{itemize}

At this stage, no union is required. The goal is to represent each
incoming system clearly.

\subsection{Step 2---Build conformed
references}\label{step-2build-conformed-references}

The second step is to build
\hyperref[docs-creating-information-reference-data]{reference data} that
expresses shared business meaning across both systems.

For example, suppose the two systems record sales status differently.

\textbf{\texttt{CakeV1.RefStatus}}

{\def\LTcaptype{none} % do not increment counter
\begin{longtable}[]{@{}ll@{}}
\toprule\noalign{}
V1 status code & V1 status name \\
\midrule\noalign{}
\endhead
\bottomrule\noalign{}
\endlastfoot
O & Open \\
C & Completed \\
A & Abandoned \\
\end{longtable}
}

\textbf{\texttt{CakeV2.RefStatus}}

{\def\LTcaptype{none} % do not increment counter
\begin{longtable}[]{@{}ll@{}}
\toprule\noalign{}
V2 status code & V2 status name \\
\midrule\noalign{}
\endhead
\bottomrule\noalign{}
\endlastfoot
OP & Open \\
CO & Completed \\
WD & Withdrawn \\
RF & Refunded \\
\end{longtable}
}

A proposed conformed reference might be:

\textbf{\texttt{Cake.RefStatus}}

{\def\LTcaptype{none} % do not increment counter
\begin{longtable}[]{@{}rll@{}}
\toprule\noalign{}
Status ID & Status name & Is finished sale \\
\midrule\noalign{}
\endhead
\bottomrule\noalign{}
\endlastfoot
1 & Open & false \\
2 & Completed & true \\
3 & Withdrawn & true \\
4 & Refunded & true \\
\end{longtable}
}

The mapping table then connects each system-specific status to the
conformed reference.

\textbf{\texttt{Cake.StatusMap}}

{\def\LTcaptype{none} % do not increment counter
\begin{longtable}[]{@{}llr@{}}
\toprule\noalign{}
System & Source status code & Status ID \\
\midrule\noalign{}
\endhead
\bottomrule\noalign{}
\endlastfoot
CakeV1 & O & 1 \\
CakeV1 & C & 2 \\
CakeV1 & A & 3 \\
CakeV2 & OP & 1 \\
CakeV2 & CO & 2 \\
CakeV2 & WD & 3 \\
CakeV2 & RF & 4 \\
\end{longtable}
}

But this is only a hypothesis. It must be validated with the business.
Even subtle differences in meaning can undermine the integration.

In addition to aligning codes, the golden reference tables should:

\begin{itemize}
\tightlist
\item
  Include default rows for unknown values
\item
  Add analytical columns such as \texttt{{[}Is\ finished\ sale{]}} to
  support downstream use
\item
  Be documented with metadata to support clarity and reuse
\end{itemize}

\subsection{Step 3---Integrate the
transactions}\label{step-3integrate-the-transactions}

Only after the local systems have been modelled and the references
conformed should the transaction tables be integrated.

The integrated table might be:

\begin{itemize}
\tightlist
\item
  \texttt{Cake.Sales}
\end{itemize}

This table is created as a union of \texttt{CakeV1.Sales} and
\texttt{CakeV2.Sales}.

If a column exists only in one system, the other system should be
populated with an appropriate default value. This includes default
foreign keys to reference tables where required.

During the union, the mapping tables translate system-specific codes to
conformed references. This allows \texttt{Cake.Sales} to use shared
meanings rather than system-specific values.

Example SQL would be:

\begin{verbatim}
select
      v1.[Sales ID]
    , 'CakeV1'              as [Source system]
    , v1.[Sales date]
    , sm.[Status ID]        as [Status ID]
    , -1                    as [Campaign ID]
    , v1.[Sales value]
from      CakeV1.Sales      v1
left join Cake.StatusMap    sm on  sm.[System]             = 'CakeV1'
                              and sm.[Source status code] = v1.[Status code]

union all

select
      v2.[Sales ID]
    , 'CakeV2'              as [Source system]
    , v2.[Sales date]
    , sm.[Status ID]        as [Status ID]
    , cm.[Campaign ID]      as [Campaign ID]
    , v2.[Sales value]
from      CakeV2.Sales      v2
left join Cake.StatusMap    sm on  sm.[System]             = 'CakeV2'
                              and sm.[Source status code] = v2.[Status code]
left join Cake.CampaignMap  cm on  cm.[System]             = 'CakeV2'
                              and cm.[Source campaign code] = v2.[Campaign code];
\end{verbatim}

A sample result is:

{\def\LTcaptype{none} % do not increment counter
\begin{longtable}[]{@{}
  >{\raggedright\arraybackslash}p{(\linewidth - 10\tabcolsep) * \real{0.1429}}
  >{\raggedright\arraybackslash}p{(\linewidth - 10\tabcolsep) * \real{0.1429}}
  >{\raggedright\arraybackslash}p{(\linewidth - 10\tabcolsep) * \real{0.1429}}
  >{\raggedleft\arraybackslash}p{(\linewidth - 10\tabcolsep) * \real{0.1905}}
  >{\raggedleft\arraybackslash}p{(\linewidth - 10\tabcolsep) * \real{0.1905}}
  >{\raggedleft\arraybackslash}p{(\linewidth - 10\tabcolsep) * \real{0.1905}}@{}}
\toprule\noalign{}
\begin{minipage}[b]{\linewidth}\raggedright
Sales ID
\end{minipage} & \begin{minipage}[b]{\linewidth}\raggedright
Source system
\end{minipage} & \begin{minipage}[b]{\linewidth}\raggedright
Sales date
\end{minipage} & \begin{minipage}[b]{\linewidth}\raggedleft
Status ID
\end{minipage} & \begin{minipage}[b]{\linewidth}\raggedleft
Campaign ID
\end{minipage} & \begin{minipage}[b]{\linewidth}\raggedleft
Sales value
\end{minipage} \\
\midrule\noalign{}
\endhead
\bottomrule\noalign{}
\endlastfoot
10001 & CakeV1 & 2025-06-01 & 1 & -1 & 120.00 \\
10002 & CakeV1 & 2025-06-03 & 2 & -1 & 340.00 \\
10001 & CakeV2 & 2025-06-02 & 1 & 5 & 180.00 \\
10002 & CakeV2 & 2025-06-05 & 4 & 7 & 95.00 \\
\end{longtable}
}

\texttt{CakeV1} and \texttt{CakeV2} may use overlapping {[}Sales ID{]}
values, the primary key for the union is
\texttt{{[}Sales\ ID{]},\ {[}Source\ system{]}}. A new
\texttt{{[}Sales\ SK{]}} surrogate key may be necessary if a
single-column primary key is needed.

The union does not preserve the system-specific status codes. It
replaces them with the conformed \texttt{{[}Status\ ID{]}}.
\texttt{CakeV1} has no campaign, so it receives the default campaign
key. The resulting \texttt{Cake.Sales} table is therefore not merely a
stack of two source tables. It is a conformed sales table.

The final structure includes:

\begin{itemize}
\tightlist
\item
  \texttt{Cake.Sales}---the unified transaction table
\item
  \texttt{Cake.RefStatus}---the conformed status reference
\item
  \texttt{Cake.RefCampaign}---the conformed campaign reference
\end{itemize}

The separation into three steps---modelling, reference building, and
integration---allows each system to be developed independently, supports
incremental delivery, and makes it easier to refactor or extend the
model later. In practice, simplifications may be appropriate. The
decision to simplify should be left to the data engineer, guided by
expressiveness, fragment modelling, and the need to build a stable
pipeline.

The overall workflow would look like:

\pandocbounded{\includegraphics[keepaspectratio]{book/epub-assets/diagram-005.png}}

\emph{Figure 1. Vertical integration. Local systems are modelled
separately, mapped to conformed references, and then integrated into a
unified transaction table.}

\section{Horizontal integration}\label{horizontal-integration}

Horizontal integration applies when the entities differ conceptually but
share enough commonality to make comparison valuable. This pattern is
closely related to Kimball\textquotesingle s notion of conformed
dimensions and the bus matrix. Different business processes remain
separate while becoming comparable through shared references.

A cake company may both produce and sell cakes. The core tables are:

\begin{itemize}
\tightlist
\item
  \texttt{Cake.Production}
\item
  \texttt{Cake.Sales}
\end{itemize}

Production records manufacturing. Sales records transactions. However,
they share common attributes such as region, time, cost, and staff,
making comparison meaningful.

These should not be collapsed into a single abstract table such as
\texttt{Cake.Event}. This creates a model that is difficult to
understand and prone to error.

They can still be compared through shared references.

For example, both production and sales may use:

\begin{itemize}
\tightlist
\item
  \texttt{Cake.RefCalendar}
\item
  \texttt{Cake.RefRegion}
\end{itemize}

This allows comparisons such as:

\begin{itemize}
\tightlist
\item
  production volume and sales volume by region and month
\item
  production cost and sales value by region and month
\item
  staff count and sales performance by site
\end{itemize}

An example SQL would be

\begin{verbatim}
with production as (
    select
          [Production date]
        , [Region ID]
        , sum([Production volume])      as [Production volume]
        , sum([Production cost])        as [Production cost]
        , count(distinct [Staff ID])    as [Production staff count]
    from Cake.Production
    group by
          [Production date]
        , [Region ID]
),
sales as (
    select
          [Sales date]
        , [Region ID]
        , sum([Sales volume])           as [Sales volume]
        , sum([Sales value])            as [Sales value]
        , count(distinct [Staff ID])    as [Sales staff count]
    from Cake.Sales
    group by
          [Sales date]
        , [Region ID]
)
select
      c.[Month start date]
    , r.[Region name]
    , sum(p.[Production volume])        as [Production volume]
    , sum(s.[Sales volume])             as [Sales volume]
    , sum(p.[Production cost])          as [Production cost]
    , sum(s.[Sales value])              as [Sales value]
    , sum(p.[Production staff count])   as [Production staff count]
    , sum(s.[Sales staff count])        as [Sales staff count]
from       Cake.RefCalendar c
cross join Cake.RefRegion   r
left join  production       p on p.[Production date] = c.[Calendar date]
                              and p.[Region ID]       = r.[Region ID]
left join  sales            s on s.[Sales date]      = c.[Calendar date]
                              and s.[Region ID]      = r.[Region ID]
group by
      c.[Month start date]
    , r.[Region name];
\end{verbatim}

The result is:

{\def\LTcaptype{none} % do not increment counter
\begin{longtable}[]{@{}
  >{\raggedright\arraybackslash}p{(\linewidth - 14\tabcolsep) * \real{0.1000}}
  >{\raggedright\arraybackslash}p{(\linewidth - 14\tabcolsep) * \real{0.1000}}
  >{\raggedleft\arraybackslash}p{(\linewidth - 14\tabcolsep) * \real{0.1333}}
  >{\raggedleft\arraybackslash}p{(\linewidth - 14\tabcolsep) * \real{0.1333}}
  >{\raggedleft\arraybackslash}p{(\linewidth - 14\tabcolsep) * \real{0.1333}}
  >{\raggedleft\arraybackslash}p{(\linewidth - 14\tabcolsep) * \real{0.1333}}
  >{\raggedleft\arraybackslash}p{(\linewidth - 14\tabcolsep) * \real{0.1333}}
  >{\raggedleft\arraybackslash}p{(\linewidth - 14\tabcolsep) * \real{0.1333}}@{}}
\toprule\noalign{}
\begin{minipage}[b]{\linewidth}\raggedright
Month start date
\end{minipage} & \begin{minipage}[b]{\linewidth}\raggedright
Region name
\end{minipage} & \begin{minipage}[b]{\linewidth}\raggedleft
Production volume
\end{minipage} & \begin{minipage}[b]{\linewidth}\raggedleft
Sales volume
\end{minipage} & \begin{minipage}[b]{\linewidth}\raggedleft
Production cost
\end{minipage} & \begin{minipage}[b]{\linewidth}\raggedleft
Sales value
\end{minipage} & \begin{minipage}[b]{\linewidth}\raggedleft
Production staff count
\end{minipage} & \begin{minipage}[b]{\linewidth}\raggedleft
Sales staff count
\end{minipage} \\
\midrule\noalign{}
\endhead
\bottomrule\noalign{}
\endlastfoot
2025-01-01 & North & 12,500 & 11,900 & 82,000 & 145,000 & 18 & 7 \\
2025-01-01 & South & 9,800 & 10,200 & 63,000 & 126,000 & 15 & 6 \\
2025-02-01 & North & 13,200 & 12,600 & 86,500 & 152,000 & 19 & 8 \\
2025-02-01 & South & 10,400 & 10,100 & 66,000 & 128,500 & 15 & 6 \\
\end{longtable}
}

The comparison occurs through the shared references of calendar and
region. Neither \texttt{Cake.Production} nor \texttt{Cake.Sales} has
been altered or merged into a common transaction table. This is
integration without collapse of meaning.

\subsection{Horizontal integration through mapping
fragments}\label{horizontal-integration-through-mapping-fragments}

Sometimes shared references are not enough.

The business may want to understand how specific production batches
relate to specific sales. However, the source data may not record that
relationship directly.

Suppose the business proposes a rule: Cakes produced in the same region
and month are sold in that region and month, in production order.

This relationship is inferred. It is not native to either table.

New engineers may try to enforce the relationship by adding
\texttt{{[}Production\ ID{]}} to \texttt{Cake.Sales}, or
\texttt{{[}Sales\ ID{]}} to \texttt{Cake.Production}.

Both approaches damage the grain of the table.

A better approach is to create a mapping fragment.

\textbf{\texttt{Cake.SalesProductionMap}}

{\def\LTcaptype{none} % do not increment counter
\begin{longtable}[]{@{}llll@{}}
\toprule\noalign{}
Sales ID & Production ID & Sales month & Production batch \\
\midrule\noalign{}
\endhead
\bottomrule\noalign{}
\endlastfoot
S1001 & P5501 & June & June, Batch 1 \\
S1002 & P5501 & June & June, Batch 1 \\
S1003 & P5502 & June & June, Batch 2 \\
S1004 & P5502 & June & June, Batch 2 \\
S1005 & P5601 & July & July, Batch 1 \\
S1006 & P5601 & July & July, Batch 1 \\
\end{longtable}
}

Sales in the same region and month are allocated to production batches
in production order. The relationship is stored as a fragment rather
than forced into either \texttt{Cake.Sales} or \texttt{Cake.Production}.

\texttt{Cake.Sales} remains a sales table. \texttt{Cake.Production}
remains a production table. The inferred relationship lives in its own
fragment, where its logic can be tested, reviewed, and refined.

This is especially important when the relationship is fuzzy, inferred,
many-to-many, or likely to change.

\section{Choosing between vertical and horizontal
integration}\label{choosing-between-vertical-and-horizontal-integration}

The distinction between vertical and horizontal integration is not
always obvious. A useful starting point is to ask two questions.

\textbf{First, would a union produce relatively few null columns?}

If the two systems record largely the same attributes, a union is often
a good sign. If most columns would be null for large portions of the
data, the systems may not represent the same business entity.

\textbf{Second, can the resulting union be given a meaningful name?}

For example, it is natural to combine \texttt{CakeV1.Sales} and
\texttt{CakeV2.Sales} into \texttt{Cake.Sales}. Likewise,
\texttt{LegacyCustomer} and \texttt{ModernCustomer} may become
\texttt{Customer}.

In contrast, combining \texttt{Production} and \texttt{Sales} into a
table called \texttt{Event} is technically possible but conceptually
weak. The name exists to support the union rather than describe a
genuine business entity.

When both questions can be answered positively, vertical integration is
often the stronger option. It preserves continuity across systems and
usually produces a simpler model for downstream use.

When the answer to either question is no, then vertical integration is
likely to distort meaning. Horizontal integration would be more
appropriate.

\section{Key ideas}\label{key-ideas}

\begin{quote}
\textbf{Key ideas.}

Integration should begin by asking what kind of sameness is involved.

Vertical integration applies when the same kind of entity is recorded
across multiple systems.

Horizontal integration applies when different entities can be compared
through shared references.

Mapping fragments preserve grain when relationships between entities are
inferred, fuzzy, or many-to-many.

A quick test for vertical integration: would the union produce few null
columns, and can the resulting table be given a meaningful business
name?

The central danger of integration is collapsing meaning for technical
convenience.
\end{quote}

\cleardoublepage

\chapter{Storytelling}\label{docs-creating-information-storytelling}

\section{Meaning beyond correctness}\label{meaning-beyond-correctness}

Source systems record operational events. They do not necessarily expose
the outcomes, patterns, and journeys that matter for decision making.

A production system records inspections. A help desk system records
escalations. A sales system records transactions. None automatically
tell the business whether things are going well or how entities move
through business processes.

Data engineering is not finished when the data is correct. It is
finished when the business can see.

Storytelling is the process of transforming operational detail into
forms that allow the business to recognise outcomes, patterns, and
journeys.

This chapter explores three approaches:

\begin{itemize}
\tightlist
\item
  \textbf{Good and bad entities}---making business judgement explicit.
\item
  \textbf{Trading detail for insight}---reducing noise into categories,
  combinations, and special cases.
\item
  \textbf{Storytelling dimensions}---combining these interpretations
  into reusable reference data.
\end{itemize}

\section{Good and bad entities}\label{good-and-bad-entities}

Business often sees the world through the lens of ``good'' and ``bad.''
A trade makes a profit or a loss. A customer is retained or lost. A
production run succeeds or fails.

However, operational systems are often designed to support workflow
rather than business judgement. The outcome that matters to the business
may not exist explicitly in the source system.

One of the most valuable contributions of a data engineer is therefore
to define outcomes that help the business recognise success and failure
consistently. Sometimes the definition is obvious to the business user.
Other times, it is latent in the way stakeholders talk about the data.
The data engineer can proactively surface this latent definition by
listening carefully during discussions and asking the right questions.

As a simplified example, a \texttt{Cake.Production} process may contain
inspection results and operational notes, but not a clear concept of
production failure. The data engineer can introduce that concept through
a reference table and fragment.

\textbf{Example structure of \texttt{Cake.RefProductionFailure}}

{\def\LTcaptype{none} % do not increment counter
\begin{longtable}[]{@{}rll@{}}
\toprule\noalign{}
Production failure ID & Production failure outcome & Failure reason \\
\midrule\noalign{}
\endhead
\bottomrule\noalign{}
\endlastfoot
1 & Production failure & Spoiled cake \\
2 & Production failure & Undercooked cake \\
3 & Production failure & Damaged packaging \\
4 & Production failure & Missing allergen label \\
5 & No production failure & No failure detected \\
\end{longtable}
}

\textbf{Example structure of \texttt{Cake.ProductionFailure}}

{\def\LTcaptype{none} % do not increment counter
\begin{longtable}[]{@{}lr@{}}
\toprule\noalign{}
Production ID & Production failure ID \\
\midrule\noalign{}
\endhead
\bottomrule\noalign{}
\endlastfoot
P1001 & 5 \\
P1002 & 1 \\
P1003 & 5 \\
P1004 & 3 \\
P1005 & 4 \\
\end{longtable}
}

The source system may record production events and inspection notes, but
the business needs to see whether a production run failed and why. The
data engineer creates that interpretation as a reusable fragment.

The definition of good and bad can be a significant point of debate in a
complex organisation with different business viewpoints. In these
situations, it can be extremely challenging to arrive at a consensus.
Many organisations fail to do so, crippling their ability to see issues
consistently. In this case, it is often
\hyperref[docs-foundations-data-and-organisations]{the creativity and
technical expertise of the data engineer} that can broker between
parties by showing the way forward.

\section{Trading detail for insight}\label{trading-detail-for-insight}

Not every detail contributes equally to understanding. Often the fastest
path to insight is to reduce detail rather than increase it.

A sales process may appear noisy when analysed by individual city but
reveal a clear trend when grouped by region. A daily pattern may become
obvious only after aggregating into weeks or sales seasons. In both
cases, detail is traded for visibility.

During development, business stakeholders will often ask for the nth
detail. This is understandable, but it can obscure an essential
perspective. The problem becomes acute in fast-paced projects where
teams tick off requirements and move on. The data engineer must
proactively compensate for this tendency. Aggregation should be built
into the work plan from the outset.

The data engineer has three common ways of trading details for insights:

\begin{itemize}
\tightlist
\item
  Creating categories
\item
  Pivoting combinations
\item
  Highlighting special cases
\end{itemize}

\subsection{Creating categories}\label{creating-categories}

Creating categories is the simplest way to surface insight that is
hidden by noise. For example, using age bands rather than individual
ages, or grouping days into seasons, can reveal patterns that are
otherwise obscured.

A special case of this is the use of binary columns. Suppose a help desk
workflow has cases in \texttt{Helpdesk.Case} and another
\texttt{Helpdesk.Escalation} table to record escalation events.

\textbf{Example structure of \texttt{Helpdesk.RefEscalationTier}}

{\def\LTcaptype{none} % do not increment counter
\begin{longtable}[]{@{}ll@{}}
\toprule\noalign{}
Escalation tier & Is escalated \\
\midrule\noalign{}
\endhead
\bottomrule\noalign{}
\endlastfoot
Tier 1 & false \\
Tier 2 & true \\
Tier 3 & true \\
Tier 4 & true \\
\end{longtable}
}

Creating a column such as \texttt{{[}Is\ escalated{]}}, defined as false
for \texttt{Tier\ 1} and true for \texttt{Tier\ 2} and above, can be
immediately useful for business users seeking to understand escalation
behaviour.

A further special case of binary columns is to link categories to good
and bad. For example, an inspection is a common business process. The
system records inspection results selected from a finite list.

Binary columns also play an important role in filtering and creating
measures, as will be discussed in following chapters.

\textbf{Example structure of \texttt{Cake.RefInspectionResult}}

{\def\LTcaptype{none} % do not increment counter
\begin{longtable}[]{@{}ll@{}}
\toprule\noalign{}
Inspection result & Is bad cake \\
\midrule\noalign{}
\endhead
\bottomrule\noalign{}
\endlastfoot
Excellent cake & false \\
Good cake & false \\
Spoiled cake & true \\
Not tasty & true \\
\end{longtable}
}

The grouping looks straightforward, but it has a large influence
downstream in designing reports, exposing filters, and designing
measures.

\subsection{Pivoting combinations}\label{pivoting-combinations}

Another approach is to pivot options in the entity's detail rows back to
the entity grain.

Continuing with the help desk workflow example, escalation tiers can be
pivoted into columns such as \texttt{{[}Tier\ 1{]}},
\texttt{{[}Tier\ 2{]}}, \texttt{{[}Tier\ 3{]}}, and
\texttt{{[}Tier\ 4{]}} to represent whether each tier was reached.

\textbf{Source detail rows in \texttt{Helpdesk.Escalation}}

{\def\LTcaptype{none} % do not increment counter
\begin{longtable}[]{@{}ll@{}}
\toprule\noalign{}
Case ID & Escalation tier \\
\midrule\noalign{}
\endhead
\bottomrule\noalign{}
\endlastfoot
H1001 & Tier 1 \\
H1001 & Tier 2 \\
H1001 & Tier 3 \\
H1002 & Tier 1 \\
H1003 & Tier 4 \\
H1004 & Tier 1 \\
H1004 & Tier 2 \\
\end{longtable}
}

Would be converted to a table of the entity grain.

\textbf{Pivoted structure at the case grain}

{\def\LTcaptype{none} % do not increment counter
\begin{longtable}[]{@{}lllll@{}}
\toprule\noalign{}
Case ID & Tier 1 & Tier 2 & Tier 3 & Tier 4 \\
\midrule\noalign{}
\endhead
\bottomrule\noalign{}
\endlastfoot
H1001 & true & true & true & false \\
H1002 & true & false & false & false \\
H1003 & false & false & false & true \\
H1004 & true & true & false & false \\
\end{longtable}
}

The detail rows, including information such as the timing of the
escalation, have been traded for a compact representation of the path.
The model can now describe the case at the case grain without
double-counting escalation events.

Once pivoted, the data engineer can summarise the escalation path in a
variety of ways.

\textbf{Example escalation path}

{\def\LTcaptype{none} % do not increment counter
\begin{longtable}[]{@{}llllll@{}}
\toprule\noalign{}
Case ID & Tier 1 & Tier 2 & Tier 3 & Tier 4 & Escalation path \\
\midrule\noalign{}
\endhead
\bottomrule\noalign{}
\endlastfoot
H1001 & true & true & true & false & Progressively escalated \\
H1002 & true & false & false & false & Not escalated \\
H1003 & false & false & false & true & Direct escalation \\
H1004 & true & true & false & false & Progressively escalated \\
\end{longtable}
}

The column \texttt{{[}Escalation\ path{]}} communicates an operational
pattern that is difficult to see from detail rows alone. A business user
does not need to inspect a deluge of tier values to understand whether
the case stayed at the first tier, progressed through escalation, or
jumped directly to a higher tier.

\subsection{Highlighting special
cases}\label{highlighting-special-cases}

Sometimes a single result tells most of the story. For example, if there
are multiple results, the most important one may be the worst result,
first result, final result, or best result.

In the escalation example, the highest tier reached is often more
important than every individual escalation event.

\textbf{Extended structure}

{\def\LTcaptype{none} % do not increment counter
\begin{longtable}[]{@{}llllll@{}}
\toprule\noalign{}
Case ID & Tier 1 & Tier 2 & Tier 3 & Tier 4 & Highest escalation \\
\midrule\noalign{}
\endhead
\bottomrule\noalign{}
\endlastfoot
H1001 & true & true & true & false & Tier 3 \\
H1002 & true & false & false & false & Tier 1 \\
H1003 & false & false & false & true & Tier 4 \\
H1004 & true & true & false & false & Tier 2 \\
\end{longtable}
}

The \texttt{{[}Highest\ escalation{]}} column highlights the most
significant outcome of the escalation journey. Users do not need to
inspect every event individually because the most important result has
already been surfaced.

A common scenario is the header-detail structure where the header is the
entity and the detail is the result.

For example, in \texttt{Cake.InspectionResult}, a single cake may have
multiple inspection criteria and therefore multiple inspection results.
A simple yet powerful way of storytelling is to compute the worst of the
inspection result.

\textbf{Source detail rows in \texttt{Cake.InspectionResult}}

{\def\LTcaptype{none} % do not increment counter
\begin{longtable}[]{@{}lll@{}}
\toprule\noalign{}
Cake ID & Inspection criterion & Inspection result \\
\midrule\noalign{}
\endhead
\bottomrule\noalign{}
\endlastfoot
C1001 & Appearance & Excellent cake \\
C1001 & Taste & Good cake \\
C1001 & Packaging & Good cake \\
C1002 & Appearance & Excellent cake \\
C1002 & Taste & Not tasty \\
C1002 & Packaging & Good cake \\
C1003 & Appearance & Good cake \\
C1003 & Taste & Good cake \\
C1003 & Packaging & Spoiled cake \\
\end{longtable}
}

Suppose the reference table includes a binary classification:

\textbf{\texttt{Cake.RefInspectionResult}}

{\def\LTcaptype{none} % do not increment counter
\begin{longtable}[]{@{}ll@{}}
\toprule\noalign{}
Inspection result & Is bad cake \\
\midrule\noalign{}
\endhead
\bottomrule\noalign{}
\endlastfoot
Excellent cake & false \\
Good cake & false \\
Not tasty & true \\
Spoiled cake & true \\
\end{longtable}
}

The worst inspection result for each cake can then be calculated as the
maximum of \texttt{{[}Is\ bad\ cake{]}}.

\textbf{Result at the cake grain}

{\def\LTcaptype{none} % do not increment counter
\begin{longtable}[]{@{}ll@{}}
\toprule\noalign{}
Cake ID & Is bad cake \\
\midrule\noalign{}
\endhead
\bottomrule\noalign{}
\endlastfoot
C1001 & false \\
C1002 & true \\
C1003 & true \\
\end{longtable}
}

The business no longer needs to inspect every individual inspection
result. The most significant outcome has been surfaced at the grain of
the cake itself.

\section{Storytelling dimensions}\label{storytelling-dimensions}

A storytelling dimension brings multiple acts of interpretation together
into a single reusable business view.

Categories, combinations, special cases, and journeys can each be useful
individually. A storytelling dimension combines them into a coherent way
of seeing the entity. The aim is to tell the overall story of the
entity. This means stepping back from what is written in the database
and asking what kind of journey the entity undertook through business
processes. The emphasis is on the business view of the situation. The
story is then expressed through
\hyperref[docs-creating-information-reference-data]{reference data}
rather than in the transaction table itself.

Continuing the help desk example:

\textbf{Example structure of \texttt{Helpdesk.RefCaseEscalation}}

{\def\LTcaptype{none} % do not increment counter
\begin{longtable}[]{@{}
  >{\raggedleft\arraybackslash}p{(\linewidth - 18\tabcolsep) * \real{0.1250}}
  >{\raggedright\arraybackslash}p{(\linewidth - 18\tabcolsep) * \real{0.0938}}
  >{\raggedright\arraybackslash}p{(\linewidth - 18\tabcolsep) * \real{0.0938}}
  >{\raggedright\arraybackslash}p{(\linewidth - 18\tabcolsep) * \real{0.0938}}
  >{\raggedright\arraybackslash}p{(\linewidth - 18\tabcolsep) * \real{0.0938}}
  >{\raggedright\arraybackslash}p{(\linewidth - 18\tabcolsep) * \real{0.0938}}
  >{\raggedright\arraybackslash}p{(\linewidth - 18\tabcolsep) * \real{0.0938}}
  >{\raggedright\arraybackslash}p{(\linewidth - 18\tabcolsep) * \real{0.0938}}
  >{\raggedright\arraybackslash}p{(\linewidth - 18\tabcolsep) * \real{0.0938}}
  >{\raggedleft\arraybackslash}p{(\linewidth - 18\tabcolsep) * \real{0.1250}}@{}}
\toprule\noalign{}
\begin{minipage}[b]{\linewidth}\raggedleft
Case escalation ID
\end{minipage} & \begin{minipage}[b]{\linewidth}\raggedright
Tier 1
\end{minipage} & \begin{minipage}[b]{\linewidth}\raggedright
Tier 2
\end{minipage} & \begin{minipage}[b]{\linewidth}\raggedright
Tier 3
\end{minipage} & \begin{minipage}[b]{\linewidth}\raggedright
Tier 4
\end{minipage} & \begin{minipage}[b]{\linewidth}\raggedright
Is escalated
\end{minipage} & \begin{minipage}[b]{\linewidth}\raggedright
Highest escalation
\end{minipage} & \begin{minipage}[b]{\linewidth}\raggedright
Escalation path
\end{minipage} & \begin{minipage}[b]{\linewidth}\raggedright
Escalation summary
\end{minipage} & \begin{minipage}[b]{\linewidth}\raggedleft
Display order
\end{minipage} \\
\midrule\noalign{}
\endhead
\bottomrule\noalign{}
\endlastfoot
1 & true & false & false & false & false & Tier 1 & Not escalated &
Closed without escalation & 1 \\
2 & true & true & false & false & true & Tier 2 & Progressive escalation
& Closed after escalation & 2 \\
3 & true & true & true & false & true & Tier 3 & Progressive escalation
& Closed after escalation & 2 \\
4 & true & true & true & true & true & Tier 4 & Progressive escalation &
Closed after escalation & 2 \\
5 & false & true & false & false & true & Tier 2 & Direct escalation &
Closed, started with escalation & 3 \\
6 & false & false & true & false & true & Tier 3 & Direct escalation &
Closed, started with escalation & 3 \\
7 & false & false & false & true & true & Tier 4 & Direct escalation &
Closed, started with escalation & 3 \\
8 & true & false & true & false & true & Tier 3 & Escalated with skipped
tier & Closed after escalation & 2 \\
... & ... & ... & ... & ... & ... & ... & ... & ... & ... \\
16 & false & false & false & false & false & Not applicable & Not
applicable & Case still open & 99 \\
\end{longtable}
}

This single reference table contains several layers of interpretation:

\begin{itemize}
\tightlist
\item
  Tier combinations
\item
  Special cases
\item
  Journey descriptions
\item
  Business summaries
\end{itemize}

Together they tell the story of the case.

Creating and applying a storytelling dimension follows three basic
steps:

\begin{itemize}
\tightlist
\item
  Build the reference table
\item
  Map to the facts
\item
  Visual check
\end{itemize}

\subsection{Step 1---Build the reference
table}\label{step-1build-the-reference-table}

The storytelling dimension begins with a reference table that describes
the journeys that matter to the business.

Returning to the help desk example, the starting point is the escalation
pattern. The data engineer can represent whether each tier was reached
using columns such as \texttt{{[}Tier\ 1{]}}, \texttt{{[}Tier\ 2{]}},
\texttt{{[}Tier\ 3{]}}, and \texttt{{[}Tier\ 4{]}}.

Those tier columns can then support \texttt{{[}Highest\ escalation{]}},
which identifies the most severe escalation reached by the case.

The same combination can also support \texttt{{[}Escalation\ path{]}},
with values such as:

\begin{itemize}
\tightlist
\item
  Not escalated
\item
  Progressive escalation
\item
  Direct escalation
\item
  Escalated with skipped tiers
\end{itemize}

These values describe how the case moved through the escalation process,
not merely where it ended.

The story can then be enriched with \texttt{{[}Escalation\ summary{]}},
which combines escalation path, highest escalation, and case status into
a small set of business-facing categories.

Example summaries include:

\begin{itemize}
\tightlist
\item
  Closed without escalation
\item
  Closed after escalation
\item
  Closed, started with escalation
\item
  Case still open
\end{itemize}

When creating a summary, resist the temptation to include too much
detail. A summary is useful because it compresses many possibilities
into a handful of categories the business can hold mentally. Adding
detail creates more combinations and defeats the point of the summary.

These summaries often imply a ranking. A \texttt{{[}Display\ order{]}}
column should be included to support this ranking in visualisations and
filters.

The reference table should be heavily annotated with useful business
descriptions. Each column should describe the business rule it
represents in a way that is close to pseudo-code for business
understanding.

Such tables are often composed from underlying binary flags. The exact
SQL depends on the platform, but the pattern is usually:

\begin{itemize}
\tightlist
\item
  generate all combinations of the binary flags;
\item
  assign each combination a stable key;
\item
  compute mechanical columns such as \texttt{{[}Highest\ escalation{]}};
\item
  add business interpretation columns manually or through controlled
  rules.
\end{itemize}

\begin{verbatim}

select
          1
        + cast([Tier 1].value as int) * 1
        + cast([Tier 2].value as int) * 2
        + cast([Tier 3].value as int) * 4
        + cast([Tier 4].value as int) * 8        as [Case escalation ID]

    ,     cast([Tier 1].value as bit)            as [Tier 1]
    ,     cast([Tier 2].value as bit)            as [Tier 2]
    ,     cast([Tier 3].value as bit)            as [Tier 3]
    ,     cast([Tier 4].value as bit)            as [Tier 4]

    ,     ...                                    as [Is escalated]
    ,     ...                                    as [Highest escalation]
    ,     ...                                    as [Escalation path]
    ,     ...                                    as [Escalation summary]
    ,     ...                                    as [Display order]

from        string_split('0,1', ',') [Tier 1]
cross apply string_split('0,1', ',') [Tier 2]
cross apply string_split('0,1', ',') [Tier 3]
cross apply string_split('0,1', ',') [Tier 4];
\end{verbatim}

In dimensional modelling, this type of reference table is sometimes
called a junk dimension or transaction profile dimension.

\subsection{Step 2---Map to facts}\label{step-2map-to-facts}

Once the storytelling dimension exists, the entity of interest can be
mapped to it.

\textbf{Example structure of \texttt{Helpdesk.CaseEscalation}}

{\def\LTcaptype{none} % do not increment counter
\begin{longtable}[]{@{}lr@{}}
\toprule\noalign{}
Case ID & Case escalation ID \\
\midrule\noalign{}
\endhead
\bottomrule\noalign{}
\endlastfoot
H1001 & 3 \\
H1002 & 1 \\
H1003 & 7 \\
H1004 & 2 \\
\end{longtable}
}

The calculation should be done at the grain of the entity of interest
during the reduce stage of the pipeline, as described in the chapter on
\hyperref[docs-creating-information-entity-processing]{Entity
Processing}.

The calculation can be added to an existing aggregation table if
appropriate or implemented as a new fragment. Continuing the help desk
example, this would result in a \texttt{Helpdesk.CaseEscalation} table
paired with \texttt{Helpdesk.RefCaseEscalation}. \textbf{Example result
after joining to \texttt{Helpdesk.RefCaseEscalation}}

{\def\LTcaptype{none} % do not increment counter
\begin{longtable}[]{@{}
  >{\raggedright\arraybackslash}p{(\linewidth - 8\tabcolsep) * \real{0.2000}}
  >{\raggedright\arraybackslash}p{(\linewidth - 8\tabcolsep) * \real{0.2000}}
  >{\raggedright\arraybackslash}p{(\linewidth - 8\tabcolsep) * \real{0.2000}}
  >{\raggedright\arraybackslash}p{(\linewidth - 8\tabcolsep) * \real{0.2000}}
  >{\raggedright\arraybackslash}p{(\linewidth - 8\tabcolsep) * \real{0.2000}}@{}}
\toprule\noalign{}
\begin{minipage}[b]{\linewidth}\raggedright
Case ID
\end{minipage} & \begin{minipage}[b]{\linewidth}\raggedright
Is escalated
\end{minipage} & \begin{minipage}[b]{\linewidth}\raggedright
Highest escalation
\end{minipage} & \begin{minipage}[b]{\linewidth}\raggedright
Escalation path
\end{minipage} & \begin{minipage}[b]{\linewidth}\raggedright
Escalation summary
\end{minipage} \\
\midrule\noalign{}
\endhead
\bottomrule\noalign{}
\endlastfoot
H1001 & true & Tier 3 & Progressive escalation & Closed after
escalation \\
H1002 & false & Tier 1 & Not escalated & Closed without escalation \\
H1003 & true & Tier 4 & Direct escalation & Closed, started with
escalation \\
H1004 & true & Tier 2 & Progressive escalation & Closed after
escalation \\
\end{longtable}
}

A table like this tells the business far more about the case than a
simple list of escalation events is able to.

\subsection{Step 3---Visual check}\label{step-3visual-check}

Storytelling dimensions should always be reviewed visually.

A quick chart or summary table often reveals whether the intended story
is actually visible.

For example:

{\def\LTcaptype{none} % do not increment counter
\begin{longtable}[]{@{}lr@{}}
\toprule\noalign{}
Escalation summary & Cases \\
\midrule\noalign{}
\endhead
\bottomrule\noalign{}
\endlastfoot
Closed without escalation & 8,450 \\
Closed after escalation & 2,130 \\
Closed, started with escalation & 18 \\
Case still open & 742 \\
\end{longtable}
}

The distribution may reveal unexpected patterns, data-quality issues, or
opportunities to refine the storytelling dimension itself. A quick bar
chart may reveal that there are almost no directly escalated records,
and these cases occur only due to data entry errors. Alternatively, the
chart may show a high number of directly escalated cases, suggesting a
deeper operational issue or a more nuanced story worth exploring.

Visual checks are therefore not merely validation. They are part of the
iterative design process.

\subsection{Multiple storytelling
dimensions}\label{multiple-storytelling-dimensions}

Complex business processes may require multiple storytelling dimensions.

A help desk workflow may need one dimension describing escalation
behaviour and another describing service-level agreement outcomes. A
production process may need one dimension describing inspection outcomes
and another describing manufacturing quality. \#\# Storytelling and the
organisation

Storytelling is important because correctness is not enough. A table can
be accurate and still fail to show the business what happened.

The data engineer's task is to create forms of information that make the
business visible: its outcomes, patterns, and journeys.

This involves decisions and compromise: details to discard, or journeys
to emphasise. Consequently, storytelling is a central place where the
\hyperref[docs-foundations-data-and-organisations]{organisation
negotiates meaning}, and where the data engineer plays an influential
role in how organisation understands its own business.

\section{Key ideas}\label{key-ideas-1}

\begin{quote}
\textbf{Key ideas.}

Data engineering is not finished when the data is correct. It is
finished when the business can see.

Storytelling creates business meaning that is not explicitly present in
the source system.

Defining outcomes is one of the most powerful ways to make business
performance visible.

Trading detail for insight reveals patterns that are often hidden at the
lowest level of granularity.

Storytelling dimensions combine categories, combinations, special cases,
and journeys into reusable business interpretations.

Storytelling dimensions should be implemented as reference data and
mapped back to the entity of interest.

Visual checks are an essential part of developing storytelling
dimensions.
\end{quote}

\cleardoublepage

\chapter{Meaningful
fragments}\label{docs-creating-information-meaningful-fragments}

\section{Fragment modelling}\label{fragment-modelling-1}

The earlier chapters introduced
\hyperref[docs-creating-information-entity-processing]{entity
processing}, \hyperref[docs-creating-information-entity-tracking]{entity
tracking}, \hyperref[docs-creating-information-reference-data]{reference
data},
\hyperref[docs-creating-information-conforming-systems]{conforming
systems}, and
\hyperref[docs-creating-information-storytelling]{storytelling}. Each
depended on the same discipline: break complexity into individual
fragments and connect them deliberately.

This approach is \textbf{fragment modelling}.

This chapter closes the section on creating information by making
fragment modelling explicit. It steps back from those specific patterns
and names the broader practice. It then introduces additional fragment
patterns that commonly appear in real business pipelines: milestones,
current-version selectors, timelines, mappings, distribution weights,
and hubs.

The key idea is summarised by the
\hyperref[docs-creating-information-expressiveness-and-fragment-modelling]{second
principle of data engineering}:

\begin{quote}
Instead of building giant tables, create meaningful fragments.
\end{quote}

\section{Why giant tables fail}\label{why-giant-tables-fail}

New engineers often build giant tables to answer queries. This feels
convenient because everything is in one place, and there are fewer
artefacts to create and deploy.

However, the convenience is temporary.

Wide tables accumulate unrelated meanings. Business logic becomes buried
and entangled. Change detection becomes harder because every column is
treated as part of the same object. Being wide, they also become a
magnet for dependencies for downstream use, creating additional risks.

The result is a pipeline that looks simple from far away and becomes
fragile up close.

\section{Meaningful fragments}\label{meaningful-fragments}

Fragment modelling takes the opposite approach.

A \textbf{meaningful fragment} is a narrow table that packages one
reusable piece of business meaning. Each fragment should answer a clear
question:

\begin{quote}
What information does this table create?
\end{quote}

The fragment then isolates the calculation of this information to its
own table, and gives the table an expressive name.

By packaging up complex calculation into a meaningful fragment, it
becomes operable.

Consistent use of meaningful fragments makes the pipeline easier to
understand, test, change, and reuse.

\subsection{Restaurant analogy}\label{restaurant-analogy}

Fragment modelling is like \emph{mise en place} in a restaurant kitchen.

The raw data are the whole potatoes and unwashed greens. Fragments are
the prepared ingredients: chopped vegetables, cured sauces, and
marinated proteins. They are not the final dish, but they make service
possible.

A good pipeline works the same way. Once the fragments are well
prepared, insights can be assembled quickly, consistently, and without
repeating the same preparation work.

\subsection{The impact on legibility}\label{the-impact-on-legibility}

A further benefit of meaningful fragments is legibility.

Creating a fragment and giving it a name elevates a concept into a
first-class object within the warehouse. Users become aware that the
concept exists and can discover it through the schema itself.

For example, milestone datetimes could be implemented as additional
columns within \texttt{Cake.Sales}. While technically correct, users
browsing the warehouse may not realise that milestone information has
been prepared and curated. By creating a dedicated
\texttt{Cake.SalesMilestone} fragment, the concept becomes visible. The
table name itself communicates that milestones are important and
available for use.

In this sense, fragments do more than package information for easy
operability. They help communicate the structure of the business domain.
A well-designed warehouse is not merely a collection of tables, but a
map of the concepts that the organisation has chosen to recognise.

\section{Symptoms of poor
fragmentation}\label{symptoms-of-poor-fragmentation}

When meaningful fragments are missing, the symptoms are usually visible:

\begin{itemize}
\tightlist
\item
  the same logic appears repeatedly in \texttt{where} clauses, joins,
  reports, or measures;
\item
  tables become wider as constructed columns are bolted onto fact
  tables;
\item
  reference data is duplicated into transaction tables;
\item
  users need excessive joins, filters, and deduplication to answer basic
  questions;
\item
  complex transformations are hidden inside one large table;
\item
  small changes create large downstream effects.
\end{itemize}

These symptoms indicate that the pipeline has not been decomposed into
its minimal informational components, and made easy for access.

\section{Common fragment patterns}\label{common-fragment-patterns}

The following examples illustrate common fragment patterns. They are not
exhaustive. Any self-contained piece of reusable business meaning can
become a fragment.

\subsection{Summaries of details and
storytelling}\label{summaries-of-details-and-storytelling}

Complex business processes generate detail rows that must be summarised
to become useful. The computation for these summaries can be complex and
should be placed in a separate fragment rather than compounded onto
existing tables.

For example, a cake may undergo multiple inspections.

\textbf{\texttt{Cake.InspectionResult}}

{\def\LTcaptype{none} % do not increment counter
\begin{longtable}[]{@{}lll@{}}
\toprule\noalign{}
Inspection ID & Cake ID & Inspection result \\
\midrule\noalign{}
\endhead
\bottomrule\noalign{}
\endlastfoot
I1001 & C1001 & Excellent cake \\
I1002 & C1001 & Good cake \\
I1003 & C1001 & Good cake \\
I1004 & C1002 & Good cake \\
I1005 & C1002 & Spoiled cake \\
\end{longtable}
}

The business may only need to know whether the cake had a bad result.

\textbf{\texttt{Cake.CakeStory}}

{\def\LTcaptype{none} % do not increment counter
\begin{longtable}[]{@{}ll@{}}
\toprule\noalign{}
Cake ID & Is bad cake \\
\midrule\noalign{}
\endhead
\bottomrule\noalign{}
\endlastfoot
C1001 & false \\
C1002 & true \\
\end{longtable}
}

The fragment turns scattered detail into entity-level meaning.

This is a special case of the the storytelling dimension studied in the
last chapter.

Once computed, the storytelling dimension and the fragment often becomes
the primary window through which the business sees the data.

It may become the default filter in reports and queries because it lets
users stay at the grain of the entity of interest. They do not need to
go lower into subprocess tables, joining across grain, and risk
double-counting or inconsistent logic.

Thus, the storytelling fragment absorbs that complexity once, so
ordinary queries can remain simple. Across a warehouse with many users,
the option to skip over the detailed grain can vastly reduce errors and
improve maintenance.

\subsection{Milestone datetimes}\label{milestone-datetimes}

Business is always interested in measuring the time between milestones.
This can be difficult if steps repeat or spread across multiple tables.
The data engineer can add value by precomputing these milestones into a
dedicated fragment.

A milestone datetimes fragment is simply a table at the grain of the
entity of interest, with one datetime column per milestone.

A help desk may want to know when a case was received, escalated, and
closed.

\textbf{\texttt{Helpdesk.CaseEvent}}

{\def\LTcaptype{none} % do not increment counter
\begin{longtable}[]{@{}lll@{}}
\toprule\noalign{}
Case ID & Event & Datetime \\
\midrule\noalign{}
\endhead
\bottomrule\noalign{}
\endlastfoot
H1001 & Received & 2025-01-10 09:05 \\
H1001 & Escalated & 2025-01-11 11:20 \\
H1001 & Closed & 2025-01-12 15:40 \\
\end{longtable}
}

This can be simplified into a milestone fragment.

\textbf{\texttt{Helpdesk.CaseMilestone}}

{\def\LTcaptype{none} % do not increment counter
\begin{longtable}[]{@{}llll@{}}
\toprule\noalign{}
Case ID & Received datetime & Escalated datetime & Closed datetime \\
\midrule\noalign{}
\endhead
\bottomrule\noalign{}
\endlastfoot
H1001 & 2025-01-10 09:05 & 2025-01-11 11:20 & 2025-01-12 15:40 \\
\end{longtable}
}

The fragment makes duration calculations straightforward and avoids
repeated event-processing logic throughout the model.

Milestone tables are also known as accumulating snapshots in dimensional
modelling.

\subsection{Current or primary
version}\label{current-or-primary-version}

Many entities have multiple versions. The full version history may be
useful for audit or lineage, but the business often needs the current or
primary version.

In many systems, this logic is handled in the application tier and is
not obvious in the database. A fragment makes this explicit.

\textbf{\texttt{Sales.CustomerVersion}}

{\def\LTcaptype{none} % do not increment counter
\begin{longtable}[]{@{}lrl@{}}
\toprule\noalign{}
Customer ID & Version & Customer name \\
\midrule\noalign{}
\endhead
\bottomrule\noalign{}
\endlastfoot
C1001 & 1 & John Smith \\
C1001 & 2 & John A. Smith \\
\end{longtable}
}

A fragment can isolate the selection logic.

\textbf{\texttt{Sales.CurrentCustomerVersion}}

{\def\LTcaptype{none} % do not increment counter
\begin{longtable}[]{@{}lr@{}}
\toprule\noalign{}
Customer ID & Current version \\
\midrule\noalign{}
\endhead
\bottomrule\noalign{}
\endlastfoot
C1001 & 2 \\
\end{longtable}
}

Users can then retrieve the current state without reconstructing the
version-selection rule.

This is one case where fragment modelling can reduce immediate
expressiveness if naming is poor. Most users do not expect to retrieve a
row from one table and then consult another table to find whether it is
current. The data engineer should compensate by naming the main table
clearly, for example by using \texttt{version} in the table name or key
columns.

\subsection{Timeline and
end-of-period}\label{timeline-and-end-of-period}

When entities change over time, users often need consistent
point-in-time views. This can be error-prone for users when the
attributes are spread across multiple tables, because joins require
careful handling of time conditions. The data engineer can make these
results easily accessible through corresponding fragments.

The relevant fragments---timelines and end-of-period tables--- are
discussed in \hyperref[docs-creating-information-entity-tracking]{Entity
tracking}.

\subsection{Mapping}\label{mapping}

Business processes often do not record relationships at the grain
required for analysis.

Sometimes this granularity is impractical because processes do not
operate in the neat way expected by the business.

Sometimes the recording is imperfect, such as through nearest temporal
join. These relationships can be approximated using fuzzy logic.

These are often many-to-many relationships. Adding foreign key columns
to the core tables would disturb their grain.

Rather than including complex code in core tables, or disturbing their
grain, the mapping can be stored in separate fragments. These are
two-column tables with foreign keys pointing to each primary table. This
is the mapping table.

For example, the business may want to relate sales to production
batches, even though the source system does not record the link
directly.

\textbf{\texttt{Cake.SalesProductionMap}}

{\def\LTcaptype{none} % do not increment counter
\begin{longtable}[]{@{}ll@{}}
\toprule\noalign{}
Sales ID & Production ID \\
\midrule\noalign{}
\endhead
\bottomrule\noalign{}
\endlastfoot
S1001 & P5501 \\
S1002 & P5501 \\
S1003 & P5502 \\
\end{longtable}
}

Rather than disturbing the grain of either table, the relationship is
represented in its own fragment.

\subsection{Distribution weights}\label{distribution-weights}

Sometimes a value is recorded only in aggregate form. However, the
business wants to understand that value at the detail level and proposes
a rule for distributing it. This logic can be complex and should be
calculated in its own fragment. This is the distribution weights table.

Suppose a system records total effort for a work package.

\textbf{\texttt{Cake.WorkPackageEffort}}

{\def\LTcaptype{none} % do not increment counter
\begin{longtable}[]{@{}lr@{}}
\toprule\noalign{}
Work package ID & Total work hours \\
\midrule\noalign{}
\endhead
\bottomrule\noalign{}
\endlastfoot
WP1001 & 100 \\
WP1002 & 80 \\
\end{longtable}
}

The same work package contains multiple tasks, but the system does not
record actual effort per task.

\textbf{\texttt{Cake.Task}}

{\def\LTcaptype{none} % do not increment counter
\begin{longtable}[]{@{}lllll@{}}
\toprule\noalign{}
Task ID & Work package ID & Task type & Worker & Task start datetime \\
\midrule\noalign{}
\endhead
\bottomrule\noalign{}
\endlastfoot
T1001 & WP1001 & Design & Alex & 2025-01-10 09:00 \\
T1002 & WP1001 & Build & Priya & 2025-01-10 13:00 \\
T1003 & WP1001 & Test & Mei & 2025-01-11 10:00 \\
T2001 & WP1002 & Design & Alex & 2025-01-12 09:00 \\
T2002 & WP1002 & Build & Priya & 2025-01-12 14:00 \\
\end{longtable}
}

The business may define a rule for distributing the total effort across
tasks. The result should be stored as a fragment.

\textbf{\texttt{Cake.TaskWeight}}

{\def\LTcaptype{none} % do not increment counter
\begin{longtable}[]{@{}lr@{}}
\toprule\noalign{}
Task ID & Estimated work hours \\
\midrule\noalign{}
\endhead
\bottomrule\noalign{}
\endlastfoot
T1001 & 50 \\
T1002 & 30 \\
T1003 & 20 \\
T2001 & 35 \\
T2002 & 45 \\
\end{longtable}
}

The fragment does not claim that the source system directly recorded
these task-level hours. By placing the allocation in its own table, and
naming the column appropriately, the data engineer keeps the assumption
visible, testable, and replaceable.

In many cases, mapping tables and distribution weights may be combined
if their ambiguities stem from the same business phenomenon.

A common variation occurs when the detailed records contain repeated
item types, while the measure is recorded only at the header level.

Suppose a restaurant records the bill amount by table.

\textbf{\texttt{Restaurant.TableBill}}

{\def\LTcaptype{none} % do not increment counter
\begin{longtable}[]{@{}llr@{}}
\toprule\noalign{}
Table ID & Bill month & Bill amount \\
\midrule\noalign{}
\endhead
\bottomrule\noalign{}
\endlastfoot
T1001 & 2025-01 & 120.00 \\
T1002 & 2025-01 & 80.00 \\
\end{longtable}
}

The table also records drink orders.

\textbf{\texttt{Restaurant.DrinkOrder}}

{\def\LTcaptype{none} % do not increment counter
\begin{longtable}[]{@{}ll@{}}
\toprule\noalign{}
Table ID & Drink item \\
\midrule\noalign{}
\endhead
\bottomrule\noalign{}
\endlastfoot
T1001 & Beer \\
T1001 & Wine \\
T1001 & Wine \\
T1001 & Spirits \\
T1002 & Wine \\
T1002 & Wine \\
\end{longtable}
}

If the business assumes each drink order has equal value, the data
engineer can create a distribution-weight fragment.

\textbf{\texttt{Restaurant.DrinkOrderWeight}}

{\def\LTcaptype{none} % do not increment counter
\begin{longtable}[]{@{}llrr@{}}
\toprule\noalign{}
Table ID & Drink item & Item order count & Total order count \\
\midrule\noalign{}
\endhead
\bottomrule\noalign{}
\endlastfoot
T1001 & Beer & 1 & 4 \\
T1001 & Wine & 2 & 4 \\
T1001 & Spirits & 1 & 4 \\
T1002 & Wine & 2 & 2 \\
\end{longtable}
}

The final two columns can be used as a fractional weight. For example,
\texttt{T1001} has a bill amount of \texttt{120.00}, and wine accounts
for \texttt{2\ /\ 4} drink orders, so the estimated wine value is
\texttt{60.00}. This also simplifies whole-of-month measures at the item
level. A monthly estimate of wine sales can be calculated without
repeatedly scanning and counting the raw order rows.

\subsection{Hubs}\label{hubs}

A hub creates a consistent surrogate key for the same entity appearing
across multiple systems.

This is a concept from the Data Vault method of data warehousing, but it
is also useful as a fragment pattern. The hub centralises identity logic
that should not belong to any one subsystem.

For example, multiple systems may record staff by first name and last
name.

\textbf{\texttt{Payroll.Staff}}

{\def\LTcaptype{none} % do not increment counter
\begin{longtable}[]{@{}lll@{}}
\toprule\noalign{}
Payroll staff ID & First name & Last name \\
\midrule\noalign{}
\endhead
\bottomrule\noalign{}
\endlastfoot
P1001 & Jane & Smith \\
P1002 & John & Wong \\
P1003 & Alex & Spelic \\
\end{longtable}
}

\textbf{\texttt{Training.Staff}}

{\def\LTcaptype{none} % do not increment counter
\begin{longtable}[]{@{}lll@{}}
\toprule\noalign{}
Training staff ID & First name & Last name \\
\midrule\noalign{}
\endhead
\bottomrule\noalign{}
\endlastfoot
T884 & Jane & Smith \\
T921 & John & Wong \\
T312 & Mary & Zimm \\
\end{longtable}
}

A hub table assigns a consistent surrogate key to the shared identity.

\textbf{\texttt{Workforce.HubId}}

{\def\LTcaptype{none} % do not increment counter
\begin{longtable}[]{@{}rll@{}}
\toprule\noalign{}
Staff SK & First name & Last name \\
\midrule\noalign{}
\endhead
\bottomrule\noalign{}
\endlastfoot
1001 & Jane & Smith \\
1002 & John & Wong \\
1003 & Alex & Spelic \\
1004 & Mary & Zimm \\
\end{longtable}
}

This key can then be used for light anonymisation, relationship
modelling, or downstream tools such as Power BI.

The hub belongs to no single subsystem. It is a central fragment where
pairing and key assignment occur.

\section{The discipline of meaningful
fragments}\label{the-discipline-of-meaningful-fragments}

Fragment modelling is a discipline of decomposition.

The data engineer looks at a data product and asks:

\begin{itemize}
\tightlist
\item
  What are the entities?
\item
  What meanings are being created?
\item
  Which computations are complex enough to isolate?
\item
  Which relationships should not be forced into core tables?
\item
  Which pieces of information need to be reused?
\item
  Which dependencies should be delayed or reduced?
\end{itemize}

A meaningful fragment is justified when it creates business meaning,
isolates complexity, reduces dependency, supports reuse, or makes change
easier to manage.

In practice, anything can become a meaningful fragment. It does not need
to fall under a predefined category. If the information is
self-contained and its creation is justified by business meaning,
complex computation, or dependency management, it belongs in its own
place.

This is one of the hallmarks of mature data engineering: the ability to
see a data product not as a pile of tables, but as a set of minimal
informational components.

\section{Key ideas}\label{key-ideas-2}

\begin{quote}
\textbf{Key ideas.}

The second principle of data engineering---Instead of building giant
tables, create meaningful fragments.

A meaningful fragment is a narrow table with a specific informational
purpose.

Fragments isolate complex computation, reduce dependencies, support
reuse, and make change easier to manage.

Common fragment patterns include storytelling summaries, milestones,
current-version selectors, timelines, mappings, distribution weights,
and hubs.

Mature data engineering depends on seeing the data product as a set of
minimal informational components.
\end{quote}

\cleardoublepage

\thispagestyle{empty}
\phantomsection
\hypertarget{docs-presenting-insights}{}
\addcontentsline{toc}{part}{Presenting insights}
\vspace*{0.34\textheight}
\begin{center}
{\Huge\bfseries Presenting insights}
\end{center}
\cleardoublepage

\chapter{The craft of dimensional
modelling}\label{docs-presenting-insights-the-craft-of-dimensional-modelling}

\section{From information to insight}\label{from-information-to-insight}

Insight is information analysed in light of intent. The previous
section,
\hyperref[docs-creating-information-expressiveness-and-fragment-modelling]{Creating
information}, focused on building reusable blocks of information from
data.

This section, \textbf{Presenting insights}, asks the next question.

How should that information be assembled into a product that business
users can explore, understand, and trust?

Modern data engineers may present insights through a variety of ways.
Yet one method has remained central: \textbf{dimensional modelling}.

In the Microsoft ecosystem, Power BI semantic models make dimensional
modelling especially important because they provide the layer through
which many business users interact with data.

A dimensional model in Power BI is therefore best understood as a
product interface.

\section{The craft}\label{the-craft}

Building a data pipeline is largely an engineering exercise. It requires
correctness, reproducibility, performance, and maintainability.

Building a dimensional model also requires those things. But the work
becomes much easier when the upstream tables are ready for assembly:
meaningful fragments become facts; expressive reference tables become
dimensions; binary flags simplify measures. The technical work of
dimensional modelling is greatly simplified when it is built on a strong
foundation.

The central difficulty of dimensional modelling lies elsewhere.

When building a pipeline, the data engineer is mainly concerned with
informational value: how data can be transformed into faithful,
reusable, business-relevant information. When building a dimensional
model, the data engineer must also anticipate how users will think, ask,
filter, drill, and otherwise use the model.

The model must not only be correct. It must feel natural to use and
intuitively aligned to business intent.

This is why dimensional modelling is a craft.

In practice, this craft appears in modelling decisions, many of which
are not questions of correctness alone: which tables should be visible,
which columns should be hidden, which measures should be defined, which
relationships should be active, and how the model should guide the
user's attention.

The hard question is therefore not only: "How do I build the dimensional
model?"

It is also:

\begin{quote}
How should the dimensional model behave to anticipate user needs?
\end{quote}

Three guiding questions organise the chapters in this section:

\begin{itemize}
\tightlist
\item
  How do we know what the user will do with the model?
\item
  What does a good dimensional model look and feel like?
\item
  What modelling tools are available to anticipate a range of user
  behaviours?
\end{itemize}

\cleardoublepage

\chapter{Anticipating
questions}\label{docs-presenting-insights-anticipating-questions}

\section{Designing for questions not yet
asked}\label{designing-for-questions-not-yet-asked}

A good dimensional model does not merely answer the questions already
written in requirements. It preserves enough business structure to
answer all the questions users are likely to ask during the model's
lifetime.

This expectation can surprise data engineers who come from an
application-development background. In application development, the work
often begins by defining requirements with stakeholders and building
accordingly. Even when development is iterative, it is still usually
organised around delivering stated requirements.

That approach has its place, but analytics behaves differently.

In data work, answering one question often creates the next question.
Users discover what they want to ask by interacting with the data. They
may begin with a simple requirement, such as ``show quality-control
failures by month'', but once the data is visible, they will ask
follow-up questions: Which products failed? Which batches? Which
suppliers? Which customers were affected? Did shipping outcomes change
afterward?

A data project focused only on stated requirements is therefore likely
to frustrate everyone. If stakeholders have little interest, the stated
requirements may be shallow. If stakeholders are engaged, the
requirements will keep changing as the data reveals new possibilities.

The data engineer cannot predict every future question. But good design
can preserve the structure from which future questions can be answered.

This is the purpose of anticipating questions.

Power BI dimensional models are well suited to this task because facts
represent business processes, dimensions represent business information,
and relationships determine which questions can be answered through
filtering.

Instead of ticking off stated requirements, the data engineer should
identify the business processes, the information known to those
processes, and the information shared across processes. These become the
facts, dimensions, and relationships of the dimensional model.

To develop a model that anticipates questions and captures the breadth
of information, the data engineer can remember a simple three-step
process:

\begin{enumerate}
\def\labelenumi{\arabic{enumi}.}
\tightlist
\item
  All the facts---business processes
\item
  All the dimensions---business information
\item
  All the relationships---real-world relationships mirrored in the model
\end{enumerate}

Each step has an artefact to help the data engineer and business
stakeholders check the comprehensiveness of the model.

\section{All the facts}\label{all-the-facts}

The first step to building a dimensional model that can anticipate all
reasonable questions is to ensure it captures all the relevant business
processes. For ease of memory, this is the step to capture ``all the
facts.''

This step can be non-trivial for two reasons.

First, in a complex business, many teams work together to complete a
single goal. These teams often operate across different business
processes. The data engineer and project team are usually working with a
particular stakeholder group that focuses on specific processes rather
than the full set. For example, some may focus on manufacturing, some on
ordering, and others on quality control.

Second, different stakeholders often take different views of the same
process. Operational staff may focus on detailed steps, while others
take a broader, strategic view. Speaking with any one group always
results in a skewed perspective.

In other words, the data engineer faces two challenges---how much to
include, and how zoomed out the view should be.

Consider the following example of business processes:

\begin{enumerate}
\def\labelenumi{\arabic{enumi}.}
\tightlist
\item
  Research and development
\item
  Sourcing materials
\item
  Manufacture
\item
  Quality control
\item
  Orders
\item
  Shipping
\item
  Customer feedback
\end{enumerate}

For the question ``what are the different business processes to
include?'', the data engineer should be as comprehensive as possible
during horizon scanning. This means being persistent in asking ``where
does this come from?'' and ``where is this going to?'' until the
end-to-end is covered.

A useful rule of thumb is to trace the process beyond the immediate
team: where did the work come from, and where does it go next? For
end-to-end process discovery, it can help to ask where the work first
enters the organisation and where its effect eventually leaves.

Once all the processes are identified, it becomes easier for
stakeholders to collectively decide which ones to include. Usually,
there is a shared sense of what is core and what is peripheral. This is
a different question from ``what are your requirements?''

Requirements tend to narrow the perspective.

For example, a manager of quality control may focus narrowly on their
part of the process. If the project is driven by this
stakeholder---perhaps because quality control initiated the
engagement---the requirements may be framed as ``how much is spent on
quality control?'' or ``how many products are failing?'' Such a view is
unhelpfully narrow.

Instead, the data engineer can help by identifying the full business
process. The manager may then recognise that, even though reporting on
the cost of manufacture is not a ``must-have requirement'' for quality
control, it is still a core business process.

Research and development, by contrast, may be considered peripheral.

For the question ``how zoomed out should the view be?'', the answer is
to work at a level that makes sense for decision-making. It does not
help to model the lowest-level details such as sending emails or looking
up documents. On the other hand, it is also not helpful to see
manufacturing and quality control as one large step.

Zooming out often involves grouping sub-processes into a single process,
or denormalising detail rows into header-level structures that expose a
broader grain.

For example, the quality control process may involve multiple testing
steps for different criteria. Rather than creating one fact table for
quality samples and another for quality sample tests, it may be more
expressive to model the process as a single fact called
\texttt{\textquotesingle{}Quality\ control\textquotesingle{}}, with the
lower-level testing detail embedded where appropriate.

Power BI's flexibility with its DAX engine supports this approach. In
simple cases, a \texttt{distinctcount} can recover the embedded grain,
such as counting headers by their ID. In more complex cases, DAX can
still retrieve the embedded grain, as explained in
\hyperref[docs-presenting-insights-designing-measures]{Designing
measures}.

This question often challenges engineers trained in traditional
dimensional modelling who are learning Power BI for the first time. The
classical approach tends to split facts by grain rather than by business
process. This can obscure the business view.

The artefact for this step is the linear process diagram. A linear
process diagram is a business process diagram using only linear steps
with no branching, cycles, or decision trees. Readers find linear
diagrams much easier to understand than diagrams with arrows pointing in
all directions. The constraint of linearity forces the business analyst
to abstract the process to a level that is useful for themselves and
others. This is especially true for executive audiences.

It can also be helpful to categorise the business processes into
higher-level categories.

Continuing with the example above:

\begin{itemize}
\tightlist
\item
  Design---Research and development
\item
  Production---Sourcing, Manufacturing, Quality control
\item
  Sales---Orders, Shipping, Customer feedback
\end{itemize}

\pandocbounded{\includegraphics[keepaspectratio]{book/epub-assets/diagram-006.png}}

\emph{Figure 1. Linear process diagram. Business processes are arranged
in chronological order and grouped into Design, Production, and Sales.}

Making linear process diagrams does not mean that business analysts
cannot create more complex diagrams for other purposes. However, these
are not linear process diagrams for model planning.

The identified business processes form the facts of the dimensional
model. Hence this step can be remembered as ``all the facts.'' The
processes selected for the model define its scope. For example, if the
decision is to focus on manufacturing onwards, then questions about
research and development or sourcing will fall outside the model's
scope.

Identifying the business processes is the hardest part of anticipating
questions. It involves making business decisions and requires a high
degree of judgement and practical experience. Once the processes are
identified, the next step of identifying the business information
follows more mechanically.

\section{All the dimensions}\label{all-the-dimensions}

The second step, after identifying the business processes in scope, is
to identify all the business information known to those processes. This
ensures that all reasonable questions about the processes included in
the model can be answered. For ease of memory, this is the step to
capture ``all the dimensions.''

The key distinction is between information a process captures and
information a process knows.

This step is more straightforward than identifying the processes.
Business stakeholders usually know what information is captured. While
there may be debate about the relative importance of one attribute over
another, the answer is often that they are all important. If they were
not, they would not be captured in the first place.

The nuance for the data engineer is to distinguish between what is
captured by a business process and what is known to it. The general rule
is that downstream processes inherit information from upstream ones.

Consider the following example. Suppose each process captures the
following information:

\begin{itemize}
\tightlist
\item
  Manufacturing---Product type, Manufacture details (date, batch number)
\item
  Quality control---Testing details (manufacture sample, date, criteria,
  results)
\item
  Orders---Product type, Customer, Order details (date, order units,
  sales amount, order number), Order status
\item
  Shipping---Shipping order (order number, line items), Shipping
  logistics (shipping method, shipping company), Shipping status (date,
  status)
\item
  Customer feedback---Customer, Product type, Feedback details (method,
  star rating, comments)
\end{itemize}

In this example, the order process captures product type, customer and
order units.

The shipping process does not capture these directly but knows them
through the order number. This is inheritance. However, not all
downstream processes inherit everything.

For example, customer feedback may not record the order number and may
only capture the product type. In that case, it does not know the order
date, even though it is downstream.

Once the information is identified, it can be presented using a
cumulative information diagram. This is a table with business
information as rows and business processes as columns, both listed in
chronological order. The diagram is cumulative because information
typically accumulates downstream through inheritance.

Rows are business information. Columns are business processes.
\texttt{C} means the process captures the information directly.
\texttt{I} means the process inherits the information from an upstream
process.

It is also helpful to annotate the diagram with where the information is
physically stored---which database or table holds each attribute. This
helps assess what can be reasonably answered and informs project
planning.

The example below shows what information is known to each business
process.

The linear process diagram describes the real-world business processes.
The cumulative information diagram builds on the linear process diagram
by showing how information flows through those processes and where it is
stored in the database world.

{\def\LTcaptype{none} % do not increment counter
\begin{longtable}[]{@{}
  >{\raggedright\arraybackslash}p{(\linewidth - 10\tabcolsep) * \real{0.1667}}
  >{\raggedright\arraybackslash}p{(\linewidth - 10\tabcolsep) * \real{0.1667}}
  >{\raggedright\arraybackslash}p{(\linewidth - 10\tabcolsep) * \real{0.1667}}
  >{\raggedright\arraybackslash}p{(\linewidth - 10\tabcolsep) * \real{0.1667}}
  >{\raggedright\arraybackslash}p{(\linewidth - 10\tabcolsep) * \real{0.1667}}
  >{\raggedright\arraybackslash}p{(\linewidth - 10\tabcolsep) * \real{0.1667}}@{}}
\toprule\noalign{}
\begin{minipage}[b]{\linewidth}\raggedright
Business information
\end{minipage} & \begin{minipage}[b]{\linewidth}\raggedright
Manufacture
\end{minipage} & \begin{minipage}[b]{\linewidth}\raggedright
Quality control
\end{minipage} & \begin{minipage}[b]{\linewidth}\raggedright
Orders
\end{minipage} & \begin{minipage}[b]{\linewidth}\raggedright
Shipping
\end{minipage} & \begin{minipage}[b]{\linewidth}\raggedright
Customer feedback
\end{minipage} \\
\midrule\noalign{}
\endhead
\bottomrule\noalign{}
\endlastfoot
Product type & C & I & C & I & C \\
Manufacturing details & C & I & & & \\
Testing details & & C & & & \\
Order status & & & C & I & \\
Order details & & & C & I & \\
Customer & & & C & I & C \\
Shipping order & & & & C & \\
Shipping logistics & & & & C & \\
Shipping status & & & & C & \\
Feedback details & & & & & C \\
Database storage & XYZ SQL database & ABC diagnostics database & XYZ SAP
ERP & XYZ SAP ERP & Salesforce analytics \\
\end{longtable}
}

In other words, the cumulative information diagram maps real-world
processes to database storage.

Since business information becomes dimensions in a dimensional model,
this step can be remembered as ``all the dimensions.''

\section{All the relationships}\label{all-the-relationships}

The final step is to ensure that the business processes and their known
information are reflected in the dimensional model through filtering
relationships. If all relevant processes are captured, all relevant
information is identified, and valid relationships are correctly
implemented, then the model can anticipate all reasonable questions.

A relationship is not only a technical join. In a Power BI model, a
relationship is a path by which a user's question can reach a fact.

The rule of thumb is simple. If a business process knows a piece of
information, whether by direct capture or by inheritance, then the
corresponding dimension should filter that fact table.

Continuing with the earlier example, suppose
\texttt{\textquotesingle{}Order\ ID\textquotesingle{}},
\texttt{\textquotesingle{}Order\ calendar\textquotesingle{}}, and
\texttt{\textquotesingle{}Customer\textquotesingle{}} are dimensions,
and \texttt{\textquotesingle{}Orders\textquotesingle{}} is the fact
table. These dimensions should directly filter
\texttt{\textquotesingle{}Orders\textquotesingle{}}. The
\texttt{\textquotesingle{}Shipping\textquotesingle{}} fact table should
also be filtered by
\texttt{\textquotesingle{}Order\ ID\textquotesingle{}}, since shipping
needs to know the order number. In this case,
\texttt{\textquotesingle{}Order\ ID\textquotesingle{}} is a conformed ID
dimension for both processes. The
\texttt{\textquotesingle{}Order\ calendar\textquotesingle{}} and
\texttt{\textquotesingle{}Customer\textquotesingle{}} dimensions should
also filter \texttt{\textquotesingle{}Shipping\textquotesingle{}},
because this process inherits them through the order number.

New data engineers often forget about inheritance. This creates a
frustrating experience for users who expect to retrieve all shipments
for a customer but find the model does not support the question due to a
missing relationship.

There are exceptions to inheritance. Sometimes an upstream process has a
zero-to-many relationship with a downstream one. In these cases,
including the upstream information as a filter may be problematic.

For example, suppose the quality control process usually tests
manufactured batches, but sometimes tests other aspects not tied to a
specific batch, such as testing factory equipment itself. In this case,
\texttt{\textquotesingle{}Manufacture\ ID\textquotesingle{}} should
filter \texttt{\textquotesingle{}Quality\ control\textquotesingle{}} so
users can look up testing for a batch. However, it may be better not to
implement a filtering relationship between
\texttt{\textquotesingle{}Manufacture\ calendar\textquotesingle{}} and
\texttt{\textquotesingle{}Quality\ control\textquotesingle{}}. In Power
BI, if a user selects a manufacture date, this relationship will also
filter \texttt{\textquotesingle{}Quality\ control\textquotesingle{}} and
affect any measures, even when the test was unrelated to a batch.

It may be sufficient to use a role-playing
\texttt{\textquotesingle{}Reporting\ calendar\textquotesingle{}} that
filters \texttt{\textquotesingle{}Manufacture\textquotesingle{}} on
manufacture date and
\texttt{\textquotesingle{}Quality\ control\textquotesingle{}} on testing
date. This avoids misleading results. If the relationship is
implemented, the report design must be careful not to mislead users with
incorrect numbers caused by unintended filters.

Usually, upstream processes tend not to know about downstream processes.
This means that \texttt{\textquotesingle{}Customer\textquotesingle{}}
should not filter
\texttt{\textquotesingle{}Manufacture\textquotesingle{}} because this
information does not exist at the time of manufacture. However, the data
engineer has the benefit of hindsight, so some information can be
brought back to the earlier process. For example, while the
quality-control results of a manufactured batch are not known at the
time of manufacture, the result can be associated to the batch number
afterwards.

To express this at the batch-number grain, multiple quality control
criteria may need to be rolled back to the batch for an overall pass or
fail. A dimension for
\texttt{\textquotesingle{}Batch\ quality\ outcome\textquotesingle{}} can
be used to filter
\texttt{\textquotesingle{}Manufacture\textquotesingle{}}. In general,
storytelling dimensions summarise the journey of an entity through the
whole business process and can be associated to all steps of the process
in hindsight.

The artefact for this step is the chronological bus. A bus matrix is a
table that lists dimensions as rows and facts as columns, showing which
dimensions filter which facts. In Power BI, this should reflect the
actual filtering relationships, including their cardinality. A
chronological bus matrix is one where the facts and dimensions are
sorted in chronological order of their occurrence in the business.

The chronological bus is the main planning artefact of the dimensional
model. It shows, at a glance, which information can filter which
business processes.

The diagram for our example is below. The symbol
\texttt{1\ -\textgreater{}\ *} is used to indicate one-to-many, while
\texttt{*\ -\textgreater{}\ *}, which is used for product type and
customer feedback, indicates a many-to-many relationship. The latter
means one customer-feedback record may relate to several products at
once.

{\def\LTcaptype{none} % do not increment counter
\begin{longtable}[]{@{}
  >{\raggedright\arraybackslash}p{(\linewidth - 10\tabcolsep) * \real{0.1667}}
  >{\raggedright\arraybackslash}p{(\linewidth - 10\tabcolsep) * \real{0.1667}}
  >{\raggedright\arraybackslash}p{(\linewidth - 10\tabcolsep) * \real{0.1667}}
  >{\raggedright\arraybackslash}p{(\linewidth - 10\tabcolsep) * \real{0.1667}}
  >{\raggedright\arraybackslash}p{(\linewidth - 10\tabcolsep) * \real{0.1667}}
  >{\raggedright\arraybackslash}p{(\linewidth - 10\tabcolsep) * \real{0.1667}}@{}}
\toprule\noalign{}
\begin{minipage}[b]{\linewidth}\raggedright
Dimension
\end{minipage} & \begin{minipage}[b]{\linewidth}\raggedright
Manufacture
\end{minipage} & \begin{minipage}[b]{\linewidth}\raggedright
Quality control
\end{minipage} & \begin{minipage}[b]{\linewidth}\raggedright
Orders
\end{minipage} & \begin{minipage}[b]{\linewidth}\raggedright
Shipping
\end{minipage} & \begin{minipage}[b]{\linewidth}\raggedright
Customer feedback
\end{minipage} \\
\midrule\noalign{}
\endhead
\bottomrule\noalign{}
\endlastfoot
Reporting calendar & 1 -\textgreater{} * & 1 -\textgreater{} * & 1
-\textgreater{} * & 1 -\textgreater{} * & 1 -\textgreater{} * \\
Manufacture ID & 1 -\textgreater{} * & 1 -\textgreater{} * & 1
-\textgreater{} * & 1 -\textgreater{} * & \\
Manufacture calendar & 1 -\textgreater{} * & & 1 -\textgreater{} * & 1
-\textgreater{} * & \\
Batch number & 1 -\textgreater{} * & 1 -\textgreater{} * & 1
-\textgreater{} * & 1 -\textgreater{} * & \\
Batch quality outcome & 1 -\textgreater{} * & 1 -\textgreater{} * & 1
-\textgreater{} * & 1 -\textgreater{} * & \\
Product type & 1 -\textgreater{} * & 1 -\textgreater{} * & 1
-\textgreater{} * & 1 -\textgreater{} * & * -\textgreater{} * \\
Testing ID & & 1 -\textgreater{} * & & & \\
Testing date & & 1 -\textgreater{} * & & & \\
Quality criteria & & 1 -\textgreater{} * & & & \\
Testing result & & 1 -\textgreater{} * & & & \\
Order ID & & & 1 -\textgreater{} * & 1 -\textgreater{} * & \\
Order calendar & & & 1 -\textgreater{} * & 1 -\textgreater{} * & \\
Customer & & & 1 -\textgreater{} * & 1 -\textgreater{} * & 1
-\textgreater{} * \\
Order status & & & 1 -\textgreater{} * & 1 -\textgreater{} * & \\
Shipping calendar & & & & 1 -\textgreater{} * & \\
Shipping company & & & & 1 -\textgreater{} * & \\
Shipping status & & & & 1 -\textgreater{} * & \\
Arrival date & & & & 1 -\textgreater{} * & \\
Feedback calendar & & & & & 1 -\textgreater{} * \\
Feedback sentiment (star ratings) & & & & & 1 -\textgreater{} * \\
\end{longtable}
}

The chronological bus can be used to evaluate what questions are
answered by the model.

In the above example, because
\texttt{\textquotesingle{}Customer\textquotesingle{}} is a conformed
dimension for \texttt{\textquotesingle{}Orders\textquotesingle{}},
\texttt{\textquotesingle{}Shipping\textquotesingle{}}, and
\texttt{\textquotesingle{}Customer\ feedback\textquotesingle{}}, it is
possible to identify customers by the volume of orders, shipping delays,
and feedback results.

Because \texttt{\textquotesingle{}Product\ type\textquotesingle{}} is a
conformed dimension on all processes, it is the attribute that the model
can answer most questions for, such as comparing quality-control results
with customer feedback. Note, however, the many-to-many relationship
between \texttt{\textquotesingle{}Product\ type\textquotesingle{}} and
\texttt{\textquotesingle{}Customer\ feedback\textquotesingle{}}
indicates that any assignment of feedback to product is approximate and
may double-count.

Because \texttt{\textquotesingle{}Manufacture\ ID\textquotesingle{}} is
a conformed ID dimension from
\texttt{\textquotesingle{}Manufacture\textquotesingle{}} to
\texttt{\textquotesingle{}Shipping\textquotesingle{}}, the model can
support traceback of faulty products that have been shipped.

On the other hand, it is not possible to directly analyse the effect of
shipping on customer feedback because shipping information does not
filter the customer-feedback fact.

Such analysis may be possible indirectly through
\texttt{\textquotesingle{}Reporting\ calendar\textquotesingle{}},
\texttt{\textquotesingle{}Product\ type\textquotesingle{}} and
\texttt{\textquotesingle{}Customer\textquotesingle{}}, which together
may narrow the transactions down to a segment that correlates shipping
status and shipping company with feedback sentiment. If this can be done
systematically, it can be introduced as precomputed information in the
data pipeline.

Ensuring that the model's relationships match the real-world flow of
information means the model behaves in ways users intuitively expect.
Downstream processes tend to inherit from upstream ones. By default,
upstream processes do not know downstream outcomes. When downstream
outcomes are brought back upstream, they should be deliberately
modelled.

As a result, the chronological bus is typically dense in the upper
right. This structure makes gaps easy to spot. These gaps may be
accidental, reflect a limitation in the source system, or be a
deliberate design choice.

In the example above,
\texttt{\textquotesingle{}Manufacture\ calendar\textquotesingle{}} does
not filter
\texttt{\textquotesingle{}Quality\ control\textquotesingle{}}. This is a
design decision. It does not filter
\texttt{\textquotesingle{}Customer\ feedback\textquotesingle{}} either,
which is a limitation because the information is not recorded. This
example does not have any accidental gaps. If there were, they would
stand out.

The chronological bus only works if the data engineer has modelled
information as dimensions, processes as facts, and listed them in
chronological order.

The same information becomes much less useful when sorted
alphabetically:

{\def\LTcaptype{none} % do not increment counter
\begin{longtable}[]{@{}
  >{\raggedright\arraybackslash}p{(\linewidth - 10\tabcolsep) * \real{0.1667}}
  >{\raggedright\arraybackslash}p{(\linewidth - 10\tabcolsep) * \real{0.1667}}
  >{\raggedright\arraybackslash}p{(\linewidth - 10\tabcolsep) * \real{0.1667}}
  >{\raggedright\arraybackslash}p{(\linewidth - 10\tabcolsep) * \real{0.1667}}
  >{\raggedright\arraybackslash}p{(\linewidth - 10\tabcolsep) * \real{0.1667}}
  >{\raggedright\arraybackslash}p{(\linewidth - 10\tabcolsep) * \real{0.1667}}@{}}
\toprule\noalign{}
\begin{minipage}[b]{\linewidth}\raggedright
Dimension
\end{minipage} & \begin{minipage}[b]{\linewidth}\raggedright
Customer feedback
\end{minipage} & \begin{minipage}[b]{\linewidth}\raggedright
Manufacture
\end{minipage} & \begin{minipage}[b]{\linewidth}\raggedright
Orders
\end{minipage} & \begin{minipage}[b]{\linewidth}\raggedright
Quality control
\end{minipage} & \begin{minipage}[b]{\linewidth}\raggedright
Shipping
\end{minipage} \\
\midrule\noalign{}
\endhead
\bottomrule\noalign{}
\endlastfoot
Arrival date & & & & & 1 -\textgreater{} * \\
Batch number & & 1 -\textgreater{} * & 1 -\textgreater{} * & 1
-\textgreater{} * & 1 -\textgreater{} * \\
Batch quality outcome & & 1 -\textgreater{} * & 1 -\textgreater{} * & 1
-\textgreater{} * & 1 -\textgreater{} * \\
Customer & 1 -\textgreater{} * & & 1 -\textgreater{} * & & 1
-\textgreater{} * \\
Feedback calendar & 1 -\textgreater{} * & & & & \\
Feedback sentiment (star ratings) & 1 -\textgreater{} * & & & & \\
Manufacture calendar & & 1 -\textgreater{} * & 1 -\textgreater{} * & & 1
-\textgreater{} * \\
Manufacture ID & & 1 -\textgreater{} * & 1 -\textgreater{} * & 1
-\textgreater{} * & 1 -\textgreater{} * \\
Order calendar & & & 1 -\textgreater{} * & & \\
Order ID & & & 1 -\textgreater{} * & & 1 -\textgreater{} * \\
Order status & & & 1 -\textgreater{} * & & 1 -\textgreater{} * \\
Product type & * -\textgreater{} * & 1 -\textgreater{} * & 1
-\textgreater{} * & 1 -\textgreater{} * & 1 -\textgreater{} * \\
Quality criteria & & & & 1 -\textgreater{} * & \\
Reporting calendar & 1 -\textgreater{} * & 1 -\textgreater{} * & 1
-\textgreater{} * & 1 -\textgreater{} * & 1 -\textgreater{} * \\
Shipping calendar & & & & & 1 -\textgreater{} * \\
Shipping company & & & & & 1 -\textgreater{} * \\
Shipping status & & & & & 1 -\textgreater{} * \\
Testing date & & & & 1 -\textgreater{} * & \\
Testing ID & & & & 1 -\textgreater{} * & \\
Testing result & & & & 1 -\textgreater{} * & \\
\end{longtable}
}

In comparison, this matrix is difficult to read. This remains true even
though the dimension names have been deliberately prefixed to group
similar attributes together, and ``Manufacture'', ``Order'' and
``Shipping'' happen to sort roughly chronologically. In any other
instance, the bus matrix would be incomprehensible.

Consequently, the three-step process of linear process diagram,
cumulative information diagram, and chronological bus is not an
accident. It is designed to make it simple for the data engineer to keep
track of information, visualise what questions can be answered, and thus
accomplish the goal of anticipating questions.

The bus matrix is also a chance to check naming. Dimensions should be
nouns because they represent business attributes or information. Fact
names should usually be process-oriented because facts represent
business processes.

The relationship between business information and business processes is
mirrored in the relationship between dimensions and facts. Hence this
step can be remembered as ``all the relationships.''

\section{Conclusion}\label{conclusion}

The three-step approach guides the data engineer to work on the
information captured in business processes, rather than the stated
requirements of ``I would like to see X by Y''. The focus is on
preserving information, not on answering specific questions.

The three steps and their artefacts can be summarised in a table like
this:

{\def\LTcaptype{none} % do not increment counter
\begin{longtable}[]{@{}
  >{\raggedright\arraybackslash}p{(\linewidth - 4\tabcolsep) * \real{0.3333}}
  >{\raggedright\arraybackslash}p{(\linewidth - 4\tabcolsep) * \real{0.3333}}
  >{\raggedright\arraybackslash}p{(\linewidth - 4\tabcolsep) * \real{0.3333}}@{}}
\toprule\noalign{}
\begin{minipage}[b]{\linewidth}\raggedright
Step
\end{minipage} & \begin{minipage}[b]{\linewidth}\raggedright
Content
\end{minipage} & \begin{minipage}[b]{\linewidth}\raggedright
Artefact
\end{minipage} \\
\midrule\noalign{}
\endhead
\bottomrule\noalign{}
\endlastfoot
All the facts & Identify all the business processes. & Linear process
diagram. Shows all the major steps in a business, set at a level that
makes sense to decision-making. \\
All the dimensions & Identify all the information captured or inherited
by the processes. & Cumulative information diagram. Shows all the
information relevant to the business, how it flows, and where it is
physically stored. Used to inform what can be reasonably answered. \\
All the relationships & Relate information to its processes in the
dimensional model. & Chronological bus. Shows the dimensional model that
is implemented. Used to identify which possible questions can or cannot
be answered, particularly any gaps. \\
\end{longtable}
}

Following these three steps ensures the full range of information is
preserved, and questions relating to this information can be answered
because the relationships exist in the model.

This breadth of information requires a variety of modelling techniques.
The later chapters in this section explain how filtering, measures,
measure families, and security patterns make the preserved structure
usable.

In future, when a user asks a different question, the model can answer
it if the relevant information has already been articulated in the
dimensional model.

This is how the data engineer anticipates questions.

\section{Key ideas}\label{key-ideas-3}

\begin{quote}
\textbf{Key ideas.}

A strong dimensional model preserves enough business structure to answer
all reasonable future questions.

Requirements tend to narrow attention; analytical modelling must
preserve the structure from which follow-up questions can be answered.

``All the facts'' means identifying the business processes that should
become fact tables.

``All the dimensions'' means identifying the information captured or
inherited by those processes.

``All the relationships'' means implementing the filtering paths that
mirror what each process knows.

The three planning artefacts are the linear process diagram, the
cumulative information diagram, and the chronological bus.

The chronological bus shows which questions the model can and cannot
answer.
\end{quote}

\cleardoublepage

\chapter{A good dimensional
model}\label{docs-presenting-insights-a-good-dimensional-model}

\section{What makes a good dimensional
model?}\label{what-makes-a-good-dimensional-model}

A dimensional model is successful when users can ask real business
questions and obtain the expected answers without guessing how the model
works.

This chapter describes a good dimensional model through two questions:

\begin{itemize}
\tightlist
\item
  What should a good model achieve?
\item
  What are the signs of a good model?
\end{itemize}

\section{What should a good model
achieve?}\label{what-should-a-good-model-achieve}

A good dimensional model in Power BI should:

\begin{itemize}
\tightlist
\item
  Resonate with the business view
\item
  Be intuitive and unambiguous---it should "just work"
\item
  Anticipate questions
\item
  Support both summary and detail
\item
  Perform quickly
\end{itemize}

\subsection{Resonate with the view from business
intent}\label{resonate-with-the-view-from-business-intent}

A dimensional model is where expressiveness is tested by use.

In earlier chapters, expressiveness meant creating information that
corresponds to the business world in a recognisable way. In a
dimensional model, this becomes even more important. Users must be able
to recognise the business through the model.

In some cases, technically correct representations may need to give way
to what better reflects how the business sees the world. For example,
strict grain consistency can sometimes produce a proliferation of
objects that are technically correct but difficult to understand. A
simpler model may better match business understanding.

Resonance begins with naming. Models, tables, columns, and measures
should use language that is immediately meaningful to the business.

This may also require creating perspectives the business has not yet
articulated. Familiarity can cause business areas to miss insights in
their own processes. The data engineer, working across domains, is often
well placed to offer a fresh view.

Metadata also plays a key role. Descriptions should speak the business
language, state the intent clearly, and, in complex cases, explain the
logic in near-pseudocode using business terms.

\subsection{Be intuitive and
unambiguous}\label{be-intuitive-and-unambiguous}

As a self-service data model, a Power BI model should let users drag and
drop fields to get answers. For this purpose, the model should just
work.

If users must memorise complex filtering rules, understand the map of
relationships, or second-guess which column to use, the model is not
intuitive. An intuitive model behaves in a way that matches expectation
without requiring a heavy manual.

A related issue is ambiguity. An ambiguous model gives users multiple
apparent ways to get the same information, or presents the same
information in ways open to interpretation.

For example, ambiguity can occur when a column has been denormalised
into multiple fact tables, leaving the user unsure which version to use.
It can also occur when a column name such as
\texttt{{[}Calendar\ date{]}} is too generic, making it unclear which
business date the column refers to.

An unambiguous model gives the user one obvious way to answer a
question. The model should either answer the question accurately or make
it apparent that it cannot.

\subsection{Anticipate questions}\label{anticipate-questions}

A good model does not merely satisfy the initial requirements. It
anticipates the full range of reasonable business questions.

This requires the data engineer to think beyond the immediate request
and consider what the model may need to answer over time. There is no
escaping this reality. If a reasonable question is not answered now, a
business area will eventually return and ask for it later.

In practice, this means a business process should either be outside the
model, or, if included, represented in a form that supports its expected
usage. All information captured by that process should be accessible in
a way that allows users to ask sensible questions of it.

This is subject of the chapter
\hyperref[docs-presenting-insights-anticipating-questions]{anticipating
questions}.

\subsection{Support both summary and
detail}\label{support-both-summary-and-detail}

One of the most powerful features of Power BI is the ability for users
to move interactively through data by cross-filtering and drill-through.

Users often begin with a high-level view. If something catches their
attention, they drill into the details to see exactly what happened.
This is an important part of how users build trust in the model.

The opposite path is also common. Users who operate at the detailed
level often begin with transactions for troubleshooting, but eventually
want to zoom out in search of broader patterns. This movement from
detail to pattern is a sign of trust and maturity.

A data engineer cannot ignore the need to support both high-level and
detailed views. While individual reports may emphasise one or the other,
the underlying model often needs to support both.

\subsection{Perform quickly}\label{perform-quickly}

Most queries should return in under a second. Complex or infrequent
queries should complete in under two seconds. Longer runtimes are
acceptable only in rare cases.

Performance is not merely a user-experience consideration. It can also
signal deeper issues in the articulation of information.

A slow model often indicates that key information has not been
pre-computed through
\hyperref[docs-creating-information-meaningful-fragments]{meaningful
fragments}. The model is then forced to compensate with work at query
time. In this sense, a slow model may be a concern of business
expressiveness, not only technical performance.

\section{What are the signs of a good
model?}\label{what-are-the-signs-of-a-good-model}

The following signs do not guarantee that the above objectives are met,
and some of the signs may occasionally be absent for good reasons. But a
strong Power BI dimensional model will usually show many of them. Their
absence or frequent violation would hint at a model that is difficult to
use.

The signs can be grouped into four areas:

\begin{itemize}
\tightlist
\item
  Names and metadata
\item
  Dimensions
\item
  Facts and measures
\item
  Relationships
\end{itemize}

\subsection{Names and metadata}\label{names-and-metadata}

Names should be business-centric rather than system-centric. It is a
common mistake for data engineers to use acronyms or terms from the
source system rather than names that reflect the business content
itself. This should be avoided unless the model is intentionally about
the system.

Names should also avoid technical implementation terms. Words such as
\texttt{dim}, \texttt{fact}, or \texttt{bridge} should not appear in
visible names. These are implementation details, not business concepts.

Names should be explicit. For example, \texttt{{[}Sales\ date{]}} is
better than \texttt{{[}Date{]}}. It may even be useful to repeat the
table meaning in the column name, such as changing
\texttt{Sales{[}Date{]}} to \texttt{Sales{[}Sales\ date{]}}. This is
because the column name, and not the table name, is what appears in
report visuals.

Visible names should not repeat across the model. When a user searches
for a name in the field list, it should return one result. If a column
appears more than once, it presents an immediate ambiguity.

If the repeated columns are genuinely the same, but all are visible, the
model may need to be reworked to remove duplication. If the columns are
different but share the same name by accident, they should be renamed
rather than relying on table context to resolve the ambiguity.

Dimensions represent information or attributes, while facts represent
business processes. As such, dimensions should be nouns, while business
process facts should be verbal nouns or carry a sense of action.

For example, in a business system with three processes---manufacture,
order, and shipping---the facts may be called \texttt{Manufacture},
\texttt{Order}, and \texttt{Shipping}. The dimensions may be called
\texttt{Product}, \texttt{Region}, \texttt{Manufacture\ date},
\texttt{Order\ date}, and \texttt{Shipping\ date}. The facts name the
business processes; the dimensions name the attributes those processes
capture.

Naming facts after business processes is harder than it sounds. It
requires a clear focus on business intent and enough familiarity with
Power BI to design the model with simplicity.

A common case is the transaction and reference table pair, such as
\texttt{Sales} and \texttt{RefSales}. One approach is to use
\texttt{Sales} as the fact name and \texttt{Sales\ description} as the
dimension name.

Finally, a strong indicator of a good dimensional model is the presence
of rich, business-centric descriptions in the hover text of tables,
columns, and measures.

\subsection{Dimensions}\label{dimensions}

In Power BI, dimensions are the window through which users interact with
the model. A good dimensional model consciously treats dimensions as the
user\textquotesingle s point of access. They are the biggest influence
on the model\textquotesingle s look and feel.

On the surface, dimension table names should be nouns that convey the
type of information available to the user.

Dimensions should also have a reasonable number of columns, often around
five to ten. There can be more if some columns are simple derivations of
others, such as text versions of binary flags.

Too few columns in a dimension may mean that it carries too narrow a
business meaning. For example, a dimension called \texttt{Gender} with a
single column \texttt{{[}Gender{]}} is likely too thin. A large number
of dimensions with very few columns indicates a fragmented rather than
consolidated view of the business. It may be more useful to combine
\texttt{Gender} with related attributes into a broader
\texttt{Demographic} dimension.

Too few columns may also indicate missing helpful attributes. For
example, a \texttt{Gender} table should include
\texttt{{[}Gender\ display\ order{]}} so the values can be sorted
usefully.

On the other hand, too many columns in a dimension may suggest that the
grouping of attributes is too complicated and may overwhelm users. It
also reflects a view of the business that is insufficiently nuanced. In
that case, the dimension may need to be broken into several more
manageable dimensions.

A dimensional model often includes two special types of dimensions that
indicate thoughtful design.

The first is the ID dimension. This is a table of the primary keys of a
business process. It allows users to retrieve detailed transactions for
a business entity by looking up a familiar business key.

The second is the storytelling dimension. It sits at the opposite end of
the spectrum. As explained in the Storytelling chapter, a storytelling
dimension categorises business entities into a digestible number of
journeys. It is often the front door of the model: the first high-level
lens through which users enter the data.

The presence of both ID dimensions and storytelling dimensions suggests
that the data engineer has considered a broad range of user needs. It is
not a guarantee of quality, but it is a useful sign.

The full set of useful dimensions is covered in the next chapter on the
\hyperref[docs-presenting-insights-dimensional-modelling-for-ux]{components
of a dimensional model}.

\subsection{Facts and measures}\label{facts-and-measures}

Power BI is designed so that users primarily interact with dimensions
and measures rather than facts. This is explained in greater depth in
the chapter on
\hyperref[docs-presenting-insights-filtering-behaviour]{filtering
behaviour}.

For now, it is enough to note that a good Power BI dimensional model
usually reduces the prominence of fact tables. In the ideal case, fact
tables are hidden from the user.

A related sign is the absence, or at least limited use, of degenerate
dimensions (attributes stored directly in the fact table rather than in
a separate dimension) in fact tables. A proliferation of degenerate
dimensions suggests that the model is working against the natural
filtering behaviour of Power BI. It also indicates that business
information has been left as miscellaneous attributes of business
processes rather than articulated as standalone properties that deserve
dimension tables.

Measures become more prominent as fact tables become less visible. A
good dimensional model invests heavily in measures.

As a rough rule of thumb, each fact table representing a business
process should usually have at least ten measures. This reflects the
reality that there are often many metrics needed to understand a
business process. For example, a sales process may require measures for
sales volume, turnaround time, profit, cost, and related percentages.

If the model has implemented the concept of "good" and "bad", this often
leads to sub-measures for each category. Percentages and other
derivatives quickly add more.

If users can grab ready-made measures and arrive at useful answers with
minimal effort, this is a strong sign of a good model. It shows that the
data engineer has thought carefully about the business process and
placed important information at the user\textquotesingle s fingertips.

Conversely, a lack of ready-to-use measures means that the data engineer
has not sufficiently engaged with the business intent of the user.

When a model contains many measures, measure management becomes
important. A good pattern is to place all measures into a single table
at the top of the field list. A suitable name for this table is simply
\texttt{Measure}.

Power BI supports this behaviour. If all columns in a table are hidden
and the table contains at least one measure, Power BI places the table
at the top of the field list. This indicates a model that elevates
measures to first-class status.

Measures should be organised by display folders. By default, these
folders should be grouped by business process. For example, a model with
the processes \texttt{Manufacture}, \texttt{Order}, and
\texttt{Shipping} could have folders called \texttt{"Manufacture"},
\texttt{"Order"}, and \texttt{"Shipping"}.

The prominence of measures, their business-centricity, and their
organisation through business processes are strong contributors to the
look and feel of a good data model.

Measures are also easier to define when the pipeline has already created
meaningful fragments at the right grain. A good model should not rely on
complex measures to compensate for missing upstream structure.

\subsection{Relationships}\label{relationships}

Relationships in a Power BI model define user interactivity.

Following the theme of using dimensions as the main point of
interaction, dimensions should filter facts and not the other way
around.

A good sign is that, if a piece of information is known to a business
process, the dimension representing that information should filter the
corresponding fact.

For example, consider a model with two processes: \texttt{Manufacture}
and \texttt{Shipping}.

The \texttt{Product} dimension should filter both facts if both
processes know the product. \texttt{Manufacture} captures the product
directly. \texttt{Shipping} inherits the product because shipped items
trace back to what was manufactured.

By contrast, the \texttt{Shipping\ date} dimension should filter only
the \texttt{Shipping} fact. It should not filter \texttt{Manufacture},
because shipping occurs at the shipping grain. One manufactured batch
may be shipped across multiple dates, so there is no single shipping
date that naturally belongs to the manufacture process.

This does not mean that each business process is filtered only by its
own attributes. For example, all fact tables may be filtered by
\texttt{Manufacture\ date} if the shipment tracks back to a manufacture
batch.

The criterion is whether the dimension belongs naturally to the grain of
the fact being filtered.

{\def\LTcaptype{none} % do not increment counter
\begin{longtable}[]{@{}lll@{}}
\toprule\noalign{}
Dimension & Manufacture & Shipping \\
\midrule\noalign{}
\endhead
\bottomrule\noalign{}
\endlastfoot
Product & 1 → * & 1 → * \\
Manufacture date & 1 → * & 1 → * \\
Shipping date & & 1 → * \\
\end{longtable}
}

Relationships are explored in greater depth in the chapter on
\hyperref[docs-presenting-insights-anticipating-questions]{anticipating
questions}.

Since dimensions are the user\textquotesingle s access point to the
model, the column on the filtered side of the relationship, usually in
the fact table, should be hidden to avoid ambiguity. The user should use
the column on the filtering side, usually in the dimension.

\section{Conclusion}\label{conclusion-1}

This chapter sketched what a good dimensional model looks like.

The purpose is not to provide a definitive checklist that guarantees
quality. It is to attune the data engineer to the look and feel of a
good model in Power BI by taking the perspective of the user.

A list of signs is helpful, but the best test is for the data engineer
to use the model personally and ask a range of real business questions.
This is part of conducting visual checks.

This task requires deep familiarity with the business and strong
curiosity about the subject matter. It cannot be reduced to ticking off
formal requirements.

The common theme is business centricity and intuitive use. A good model
implements answers explicitly so that they appear naturally. It does not
require the user to interpret technical structure or apply guesswork.

The Zen of Python is a useful guide to designing a dimensional model.
Its emphasis on explicit over implicit, and on having one---and
preferably only one---obvious way of doing something, applies directly
to Power BI model design.

\begin{quote}
\textbf{Key ideas.}

A good dimensional model lets users ask real business questions without
guessing.

A dimensional model is judged by the experience it creates for users,
not only by whether has the technically correct information.

It should resonate with the business view, feel intuitive and
unambiguous, anticipate reasonable questions, support both summary and
detail, and perform quickly.

The signs of a good model are business-friendly names and metadata,
useful dimensions, ready-to-use measures, and relationships that support
natural filtering.

A strong model makes answers appear naturally because the data engineer
has already made the hard design decisions.
\end{quote}

\cleardoublepage

\chapter{Dimensional modelling for
UX}\label{docs-presenting-insights-dimensional-modelling-for-ux}

\section{Facts and dimensions as
interaction}\label{facts-and-dimensions-as-interaction}

Dimensional modelling is usually taught as a distinction between table
types---facts record activity; dimensions provide context.

That distinction is correct, but not enough to achieve the expectations
of a \hyperref[docs-presenting-insights-a-good-dimensional-model]{good
dimensional model}. The more important distinction is between what the
user can touch and what responds---dimensions are what the user touches;
facts are what respond.

A model is good when the user can answer questions without guessing.
This chapter explains facts and dimensions from that point of view. That
is, dimensional modelling for user experience (UX).

\section{Understanding facts and
dimensions}\label{understanding-facts-and-dimensions}

Classical dimensional modelling revolves around two core concepts: facts
and dimensions.

In a technical view:

\begin{itemize}
\tightlist
\item
  Facts are transactions, events, or activities that occur in the
  business.
\item
  Dimensions are reference tables that provide context for those
  transactions.
\end{itemize}

Fact tables are fast-moving tables. They grow quickly, have high
cardinality, and record frequent business activity. Dimension tables are
slower-moving tables. They have lower cardinality and provide meaningful
categories for filtering, grouping, and comparison.

From the business view:

\begin{itemize}
\tightlist
\item
  Facts are business processes. They convey action, such as manufacture,
  order, shipping, inspection, payment, or case handling.
\item
  Dimensions are business objects, attributes, or categories. They
  convey context, such as product, region, customer, officer, status, or
  date.
\end{itemize}

A common explanation says that dimensions filter facts. This is true,
but it hides the practical point: filtering is a form of control.

\subsection{An analogy of robots}\label{an-analogy-of-robots}

To make the distinction concrete, think of dimensions as levers and
facts as robots.

If filtering is a form of control, then:

\begin{enumerate}
\def\labelenumi{\arabic{enumi}.}
\tightlist
\item
  Levers control robots.
\item
  Levers may control other levers.
\item
  Robots may occasionally control other robots.
\item
  Robots should not control levers.
\end{enumerate}

This analogy helps reinforce a key design principle: interactivity
should be driven by dimensions.

Levers are what users control. Robots are what respond. In the same way,
dimensions are the user's interface to the model, while facts contain
the machinery that responds to that interaction.

This is also why the controls should not sit on the robots themselves.
They should sit on the controller. In practical terms, users should
interact with dimensions to drive filtering, grouping, searching,
comparison, and other model effects---not with columns embedded in fact
tables.

This is the main problem with degenerate dimensions. They may work
technically, but they place controls inside the fact table. They make
the robot carry its own levers.

For this chapter, we stick with the simpler view of dimensions as
tables.

But the wider principle is more general. A dimension is not merely a
lookup table. It is anything in the model the user interacts with to
produce an effect.

In more advanced scenarios, a dimension may function more like a button
than a filter. The user clicks or selects it, and the model responds.
The fact tables contain the ingredients that make the response possible.

In this view, the physical relationship pattern is secondary. A
dimension may have one relationship, many relationships, or no direct
relationship to a fact table. It may work alone, or as part of a set of
tables. What matters is that the user interacts with it to produce a
controlled effect. If it does that, it is functioning as a dimension.

This abstract view is useful when tackling demanding usage scenarios.

\section{A repertoire of facts and
dimensions}\label{a-repertoire-of-facts-and-dimensions}

The categories below are not a complete taxonomy of dimensional
modelling. They are a practical repertoire for handling usage scenarios.
They name the roles tables can play when a model is designed for
interaction.

There are 3 types of facts and 9 types of dimensions.

Facts:

\begin{enumerate}
\def\labelenumi{\arabic{enumi}.}
\tightlist
\item
  Measurable fact
\item
  End-of-period fact
\item
  Annotation fact
\end{enumerate}

Dimensions:

\begin{enumerate}
\def\labelenumi{\arabic{enumi}.}
\tightlist
\item
  Business dimension
\item
  Role-playing dimension
\item
  Histogram dimension
\item
  Combination dimension
\item
  Choices dimension
\item
  Sankey dimension
\item
  Storytelling dimension
\item
  Search dimension, a sub-type being the ID dimension
\item
  Degenerate dimension
\end{enumerate}

\section{Facts}\label{facts}

Facts are the responsive machinery of the model. They contain business
activity, business state, or display detail that becomes visible when
the user interacts with dimensions.

\subsection{Measurable fact}\label{measurable-fact}

A \textbf{measurable fact} records business activity that can be
counted, summed, averaged, or otherwise aggregated.

Examples include inspections, sales, orders, payments, helpdesk cases,
audit events, or shipments. These are the facts people usually mean when
they talk about fact tables.

A measurable fact is designed to work through measures. Users should
normally not consume the raw fact table directly. They should interact
with dimensions, and the fact should respond through measures such as
inspection count, failure rate, total value, or average processing time.

\textbf{Example structure of
\texttt{\textquotesingle{}Inspection\textquotesingle{}}}

{\def\LTcaptype{none} % do not increment counter
\begin{longtable}[]{@{}
  >{\raggedleft\arraybackslash}p{(\linewidth - 10\tabcolsep) * \real{0.1739}}
  >{\raggedright\arraybackslash}p{(\linewidth - 10\tabcolsep) * \real{0.1304}}
  >{\raggedleft\arraybackslash}p{(\linewidth - 10\tabcolsep) * \real{0.1739}}
  >{\raggedleft\arraybackslash}p{(\linewidth - 10\tabcolsep) * \real{0.1739}}
  >{\raggedleft\arraybackslash}p{(\linewidth - 10\tabcolsep) * \real{0.1739}}
  >{\raggedleft\arraybackslash}p{(\linewidth - 10\tabcolsep) * \real{0.1739}}@{}}
\toprule\noalign{}
\begin{minipage}[b]{\linewidth}\raggedleft
Inspection SK
\end{minipage} & \begin{minipage}[b]{\linewidth}\raggedright
Inspection date
\end{minipage} & \begin{minipage}[b]{\linewidth}\raggedleft
Location SK
\end{minipage} & \begin{minipage}[b]{\linewidth}\raggedleft
Officer SK
\end{minipage} & \begin{minipage}[b]{\linewidth}\raggedleft
Inspection status SK
\end{minipage} & \begin{minipage}[b]{\linewidth}\raggedleft
Inspection duration minutes
\end{minipage} \\
\midrule\noalign{}
\endhead
\bottomrule\noalign{}
\endlastfoot
1001 & 2024-01-15 & 12 & 31 & 1 & 45 \\
1002 & 2024-01-16 & 18 & 44 & 2 & 70 \\
\end{longtable}
}

This fact can support measures that use functions such as
\texttt{COUNTROWS}, \texttt{SUMX}, and \texttt{COUNTX}. The user
interacts with dimensions such as
\texttt{\textquotesingle{}Location\textquotesingle{}},
\texttt{\textquotesingle{}Officer\textquotesingle{}}, and
\texttt{\textquotesingle{}Inspection\ status\textquotesingle{}}. The
measures respond from
\texttt{\textquotesingle{}Inspection\textquotesingle{}}.

This example is at the grain of a single inspection, but a measurable
fact does not have to be at event grain. It may already be aggregated
before it reaches the model. What matters is not the physical grain of
the table, but the role it plays: it supplies measurable material for
measures to calculate against and for dimensions to control.

\textbf{Example structure of
\texttt{\textquotesingle{}Inspection\ daily\ aggregates\textquotesingle{}}}

{\def\LTcaptype{none} % do not increment counter
\begin{longtable}[]{@{}
  >{\raggedright\arraybackslash}p{(\linewidth - 10\tabcolsep) * \real{0.1304}}
  >{\raggedleft\arraybackslash}p{(\linewidth - 10\tabcolsep) * \real{0.1739}}
  >{\raggedleft\arraybackslash}p{(\linewidth - 10\tabcolsep) * \real{0.1739}}
  >{\raggedleft\arraybackslash}p{(\linewidth - 10\tabcolsep) * \real{0.1739}}
  >{\raggedleft\arraybackslash}p{(\linewidth - 10\tabcolsep) * \real{0.1739}}
  >{\raggedleft\arraybackslash}p{(\linewidth - 10\tabcolsep) * \real{0.1739}}@{}}
\toprule\noalign{}
\begin{minipage}[b]{\linewidth}\raggedright
Inspection date
\end{minipage} & \begin{minipage}[b]{\linewidth}\raggedleft
Location SK
\end{minipage} & \begin{minipage}[b]{\linewidth}\raggedleft
Inspection status SK
\end{minipage} & \begin{minipage}[b]{\linewidth}\raggedleft
Inspection count
\end{minipage} & \begin{minipage}[b]{\linewidth}\raggedleft
Total inspection duration minutes
\end{minipage} & \begin{minipage}[b]{\linewidth}\raggedleft
Failed inspection count
\end{minipage} \\
\midrule\noalign{}
\endhead
\bottomrule\noalign{}
\endlastfoot
2024-01-15 & 12 & 1 & 18 & 720 & 0 \\
2024-01-15 & 18 & 2 & 7 & 490 & 7 \\
2024-01-16 & 12 & 1 & 22 & 860 & 0 \\
\end{longtable}
}

This is also a measurable fact. It is not at the grain of a single
inspection. It has already been aggregated by inspection date, location,
and inspection status before it reaches the model.

The measures still respond from the fact. A measure may sum
\texttt{{[}Inspection\ count{]}}, divide
\texttt{{[}Failed\ inspection\ count{]}} by
\texttt{{[}Inspection\ count{]}}, or calculate average duration from
\texttt{{[}Total\ inspection\ duration\ minutes{]}} and
\texttt{{[}Inspection\ count{]}}.

The physical grain is different, but the modelling role is the same:
\texttt{\textquotesingle{}Inspection\ daily\ aggregates\textquotesingle{}}
supplies measurable material for measures to calculate against and for
dimensions to control.

Measurable facts are the default facts for
\hyperref[docs-creating-information-entity-processing]{entity
processing} scenarios.

\subsection{End-of-period fact}\label{end-of-period-fact}

An \textbf{end-of-period fact} records the state of an entity at a
regular reporting point.

This pattern applies when the business needs to ask what existed, or
what state something was in, at the end of a period. Examples include an
account at the end of each month, an employee at the end of each pay
cycle, or a supplier at the end of each reporting quarter.

End-of-period facts are usually produced from
\hyperref[docs-creating-information-entity-tracking]{entity tracking}
pipelines. They let the model answer point-in-time questions without
requiring users to solve validity-window logic inside the report.

\textbf{Example structure of
\texttt{\textquotesingle{}Account\ end\ of\ month\textquotesingle{}}}

{\def\LTcaptype{none} % do not increment counter
\begin{longtable}[]{@{}
  >{\raggedleft\arraybackslash}p{(\linewidth - 8\tabcolsep) * \real{0.2105}}
  >{\raggedright\arraybackslash}p{(\linewidth - 8\tabcolsep) * \real{0.1579}}
  >{\raggedleft\arraybackslash}p{(\linewidth - 8\tabcolsep) * \real{0.2105}}
  >{\raggedleft\arraybackslash}p{(\linewidth - 8\tabcolsep) * \real{0.2105}}
  >{\raggedleft\arraybackslash}p{(\linewidth - 8\tabcolsep) * \real{0.2105}}@{}}
\toprule\noalign{}
\begin{minipage}[b]{\linewidth}\raggedleft
Account SK
\end{minipage} & \begin{minipage}[b]{\linewidth}\raggedright
Period end date
\end{minipage} & \begin{minipage}[b]{\linewidth}\raggedleft
Account status SK
\end{minipage} & \begin{minipage}[b]{\linewidth}\raggedleft
Account level SK
\end{minipage} & \begin{minipage}[b]{\linewidth}\raggedleft
Closing balance
\end{minipage} \\
\midrule\noalign{}
\endhead
\bottomrule\noalign{}
\endlastfoot
501 & 2024-01-31 & 1 & 2 & 8700.00 \\
501 & 2024-02-29 & 1 & 3 & 13700.00 \\
\end{longtable}
}

This fact supports questions such as:

\begin{itemize}
\tightlist
\item
  How many accounts were active at the end of February?
\item
  How many accounts were gold-level accounts at the end of the quarter?
\item
  What was the total closing balance for the selected period?
\end{itemize}

The important point is that the fact records state at a reporting
boundary. It turns a changing entity into something the user can query
at a stable point in time.

\subsection{Annotation fact}\label{annotation-fact}

An \textbf{annotation fact} records display detail attached to another
fact, usually where the detail is too fine-grained or too
high-cardinality to belong in the main fact table.

Examples include inspection comments, certificate details, notes,
attachments, or free-text descriptions. These details are often not
designed for aggregation. They are designed for display after the user
has already arrived at a transaction.

\textbf{Example structure of
\texttt{\textquotesingle{}Inspection\ comments\textquotesingle{}}}

{\def\LTcaptype{none} % do not increment counter
\begin{longtable}[]{@{}rrll@{}}
\toprule\noalign{}
Inspection SK & Comment sequence & Comment datetime & Comment text \\
\midrule\noalign{}
\endhead
\bottomrule\noalign{}
\endlastfoot
1001 & 1 & 2024-01-15 10:32 & Missing supporting document. \\
1001 & 2 & 2024-01-15 10:47 & Officer requested follow-up evidence. \\
\end{longtable}
}

The reason to move these details into a separate fact is grain.

One inspection may have many comments. If the comments are placed
directly in \texttt{\textquotesingle{}Inspection\textquotesingle{}}, the
data engineer must either change the grain of
\texttt{\textquotesingle{}Inspection\textquotesingle{}} or concatenate
the comments into a single field. Both options are undesirable.

Changing the grain bloats
\texttt{\textquotesingle{}Inspection\textquotesingle{}} by repeating
inspection-level information once for every comment. This increases
model size and load cost, but it also creates a deeper pipeline problem:
the inspection and its comments may change on different rhythms.
Inspection status may be updated while comments remain unchanged, or
comments may be appended while the inspection itself is stable.
Combining them forces change detection and partitioned loading to treat
these different behaviours as one table.

Concatenation preserves the grain of
\texttt{\textquotesingle{}Inspection\textquotesingle{}}, but destroys
the structure of the comments. The comments become a block of text
rather than records with their own sequence, datetime, author, category,
and display behaviour. The user may be able to read the comments, but
the model can no longer sort them, filter them, display them
individually, or apply separate security and formatting logic to them.

An annotation fact avoids both problems. It lets
\texttt{\textquotesingle{}Inspection\textquotesingle{}} remain at the
inspection grain, while
\texttt{\textquotesingle{}Inspection\ comments\textquotesingle{}}
carries the finer-grain detail.

Annotation facts are one of the rare cases where fact-to-fact filtering
can be appropriate. Here, the \texttt{{[}Inspection\ SK{]}} column would
the related column to the same in
\texttt{\textquotesingle{}Inspection\textquotesingle{}} fact. The
typical usage scenario is one where the business user has arrived at a
set of transaction records through a dimension and would like to see
additional details about that transaction. This is also why they are
filtered through another fact, not a dimension.

Often, a measure is still needed to display annotation detail in a
controlled way. For example, \texttt{SELECTEDVALUE} can ensure that a
comment or detail field is only displayed when the relevant transaction
has been selected.

An annotation fact is therefore not a second business process. It is an
adjunct fact. It carries detail that describes another fact.

\section{Dimensions}\label{dimensions-1}

Dimensions are the controlled surfaces of interaction. They give users
handles for filtering, grouping, searching, comparing, and storytelling.

A dimension is not merely a lookup table. In an interactive model, it is
a control.

\subsection{Business dimension}\label{business-dimension}

A \textbf{business dimension} is an ordinary control surface that lets
users filter and group facts using familiar business categories.

Examples include statuses, locations, officers, inspection types,
customers, products, regions, suppliers, and programs. These dimensions
often come from source-system reference tables, but the data engineer
may supplement them with additional columns that make business meaning
clearer.

\textbf{Example structure of
\texttt{\textquotesingle{}Inspection\ status\textquotesingle{}}}

{\def\LTcaptype{none} % do not increment counter
\begin{longtable}[]{@{}rlll@{}}
\toprule\noalign{}
Inspection status SK & Inspection status & Is completed & Is failed \\
\midrule\noalign{}
\endhead
\bottomrule\noalign{}
\endlastfoot
1 & Passed & true & false \\
2 & Failed & true & true \\
3 & In progress & false & false \\
\end{longtable}
}

This dimension filters
\texttt{\textquotesingle{}Inspection\textquotesingle{}}. A user can
place \texttt{{[}Inspection\ status{]}} on a slicer, group inspection
counts by \texttt{{[}Is\ failed{]}}, or compare completed and incomplete
work.

The important point is not only that
\texttt{\textquotesingle{}Inspection\ status\textquotesingle{}} is a
low-cardinality lookup table. The important point is that it is a
natural thing for the user to touch.

\subsection{Role-playing dimension}\label{role-playing-dimension}

A \textbf{role-playing dimension} is a dimension that takes different
meanings depending on which fact column it filters.

The most common example is a calendar. The same
\texttt{\textquotesingle{}Reporting\ calendar\textquotesingle{}} table
may be used against multiple date columns. It may filter orders by order
date, shipments by shipment date, and completions by completion date.

\textbf{Example structure of
\texttt{\textquotesingle{}Reporting\ calendar\textquotesingle{}}}

{\def\LTcaptype{none} % do not increment counter
\begin{longtable}[]{@{}llll@{}}
\toprule\noalign{}
Reporting date & Financial year & Month name & Month end date \\
\midrule\noalign{}
\endhead
\bottomrule\noalign{}
\endlastfoot
2024-01-15 & 2023-24 & January & 2024-01-31 \\
2024-02-15 & 2023-24 & February & 2024-02-29 \\
\end{longtable}
}

The same
\texttt{\textquotesingle{}Reporting\ calendar\textquotesingle{}} table
may filter a fact through different columns, such as
\texttt{{[}Order\ date{]}}, \texttt{{[}Shipment\ date{]}}, and
\texttt{{[}Completion\ date{]}}.

A user selecting January may be asking for orders placed in January,
shipments sent in January, or cases completed in January. Role-playing
dimensions play an important role in comparing different business
processes using a common denominator, as explained in the chapter
\hyperref[docs-creating-information-conforming-systems]{Conforming
systems}.

The model should make that role clear. If a role-playing dimension is
exposed to users, the field names, measure names, or hover text should
explain which role is being used.

A dimension is not role-playing merely because it filters many facts.
The related columns must have different meanings.

For example, suppose \texttt{{[}Manufacture\ date{]}} appears on
\texttt{\textquotesingle{}Manufacture\textquotesingle{}},
\texttt{\textquotesingle{}Order\textquotesingle{}}, and
\texttt{\textquotesingle{}Shipping\textquotesingle{}}, and
\texttt{\textquotesingle{}Reporting\ calendar\textquotesingle{}} relates
to each fact through \texttt{{[}Manufacture\ date{]}}. Even though the
calendar filters three facts, it is not role-playing. It means the same
thing in each case: manufacture date. It is better understood as a
business dimension, and may be better named
\texttt{\textquotesingle{}Manufacture\ calendar\textquotesingle{}}.

\subsection{Histogram dimension}\label{histogram-dimension}

A \textbf{histogram dimension} turns numeric values into a filterable
and visualisable dimension.

This pattern is useful when the user needs to filter or chart records by
a numeric value, such as days to resolve, hours to process, number of
attempts, or value bands.

Measures do not work as ordinary slicer fields. If users need to filter
by a numeric result, that result often needs to be expressed as a
dimension.

\textbf{Example structure of
\texttt{\textquotesingle{}Days\ to\ resolve\textquotesingle{}}}

{\def\LTcaptype{none} % do not increment counter
\begin{longtable}[]{@{}
  >{\raggedleft\arraybackslash}p{(\linewidth - 6\tabcolsep) * \real{0.2667}}
  >{\raggedleft\arraybackslash}p{(\linewidth - 6\tabcolsep) * \real{0.2667}}
  >{\raggedleft\arraybackslash}p{(\linewidth - 6\tabcolsep) * \real{0.2667}}
  >{\raggedright\arraybackslash}p{(\linewidth - 6\tabcolsep) * \real{0.2000}}@{}}
\toprule\noalign{}
\begin{minipage}[b]{\linewidth}\raggedleft
Days to resolve SK
\end{minipage} & \begin{minipage}[b]{\linewidth}\raggedleft
Minimum days
\end{minipage} & \begin{minipage}[b]{\linewidth}\raggedleft
Maximum days
\end{minipage} & \begin{minipage}[b]{\linewidth}\raggedright
Days to resolve band
\end{minipage} \\
\midrule\noalign{}
\endhead
\bottomrule\noalign{}
\endlastfoot
1 & 0 & 1 & 0 to 1 days \\
2 & 2 & 5 & 2 to 5 days \\
3 & 6 & 10 & 6 to 10 days \\
4 & 11 & 9999 & More than 10 days \\
\end{longtable}
}

The fact table can carry \texttt{{[}Days\ to\ resolve\ SK{]}}, or it can
carry a numeric value that maps to
\texttt{\textquotesingle{}Days\ to\ resolve\textquotesingle{}}.

Care is needed around boundaries. The data engineer must decide whether
boundaries are inclusive or exclusive and ensure that maximum and
minimum values are handled consistently.

As well as allowing the user to filter the model on a numeric value, the
dimension is useful for creating histogram visuals---hence the name. The
visual counts entities against a dimension of days, hours, attempts,
amounts, or other numeric bands.

This is especially useful for time-to-process metrics. Instead of
showing only an average processing time, the model can show the
distribution: how many cases took 0 to 1 days, 2 to 5 days, 6 to 10
days, or more than 10 days.

\subsection{Combination dimension}\label{combination-dimension}

A \textbf{combination dimension} represents combinations of related
binary properties at the entity grain.

This pattern is useful when an entity can have several related
properties at once, and the business wants to filter by combinations of
those properties.

For example, a business entity may pass through multiple systems. The
detail may be recorded at a lower grain, with one row per system
touchpoint. But the user often wants to filter the business entity, not
the system touchpoints.

\textbf{Example structure of
\texttt{\textquotesingle{}System\ combination\textquotesingle{}}}

{\def\LTcaptype{none} % do not increment counter
\begin{longtable}[]{@{}
  >{\raggedleft\arraybackslash}p{(\linewidth - 8\tabcolsep) * \real{0.2500}}
  >{\raggedright\arraybackslash}p{(\linewidth - 8\tabcolsep) * \real{0.1875}}
  >{\raggedright\arraybackslash}p{(\linewidth - 8\tabcolsep) * \real{0.1875}}
  >{\raggedright\arraybackslash}p{(\linewidth - 8\tabcolsep) * \real{0.1875}}
  >{\raggedright\arraybackslash}p{(\linewidth - 8\tabcolsep) * \real{0.1875}}@{}}
\toprule\noalign{}
\begin{minipage}[b]{\linewidth}\raggedleft
System combination SK
\end{minipage} & \begin{minipage}[b]{\linewidth}\raggedright
Has System A
\end{minipage} & \begin{minipage}[b]{\linewidth}\raggedright
Has System B
\end{minipage} & \begin{minipage}[b]{\linewidth}\raggedright
Has System C
\end{minipage} & \begin{minipage}[b]{\linewidth}\raggedright
System combination
\end{minipage} \\
\midrule\noalign{}
\endhead
\bottomrule\noalign{}
\endlastfoot
1 & false & false & false & No system \\
2 & true & false & false & System A \\
3 & false & true & false & System B \\
4 & false & false & true & System C \\
... & ... & ... & ... & ... \\
8 & true & true & true & System A, System B, System C \\
\end{longtable}
}

The fact table carries \texttt{{[}System\ combination\ SK{]}}. The user
interacts with the flags in
\texttt{\textquotesingle{}System\ combination\textquotesingle{}}.

A combination dimension pushes lower-grain detail back to the entity
level, then exposes the result as a controlled surface of interaction.

A common case is the header-detail pattern. A consignment may be the
header entity, while the cargo line items are the detail rows. Some line
items may be fresh produce. Others may be inanimate goods. The user
often wants to filter consignments by the types of cargo they contain,
not analyse each line item separately.

\textbf{Example structure of
\texttt{\textquotesingle{}Cargo\ combination\textquotesingle{}}}

{\def\LTcaptype{none} % do not increment counter
\begin{longtable}[]{@{}
  >{\raggedleft\arraybackslash}p{(\linewidth - 8\tabcolsep) * \real{0.2500}}
  >{\raggedright\arraybackslash}p{(\linewidth - 8\tabcolsep) * \real{0.1875}}
  >{\raggedright\arraybackslash}p{(\linewidth - 8\tabcolsep) * \real{0.1875}}
  >{\raggedright\arraybackslash}p{(\linewidth - 8\tabcolsep) * \real{0.1875}}
  >{\raggedright\arraybackslash}p{(\linewidth - 8\tabcolsep) * \real{0.1875}}@{}}
\toprule\noalign{}
\begin{minipage}[b]{\linewidth}\raggedleft
Cargo combination SK
\end{minipage} & \begin{minipage}[b]{\linewidth}\raggedright
Has fresh produce
\end{minipage} & \begin{minipage}[b]{\linewidth}\raggedright
Has inanimate cargo
\end{minipage} & \begin{minipage}[b]{\linewidth}\raggedright
Cargo combination
\end{minipage} & \begin{minipage}[b]{\linewidth}\raggedright
Has mixed cargo
\end{minipage} \\
\midrule\noalign{}
\endhead
\bottomrule\noalign{}
\endlastfoot
1 & false & false & No cargo & false \\
2 & true & false & Fresh produce only & false \\
3 & false & true & Inanimate cargo only & false \\
4 & true & true & Fresh produce and inanimate cargo & true \\
\end{longtable}
}

The consignment fact carries \texttt{{[}Cargo\ combination\ SK{]}}. The
user interacts with
\texttt{\textquotesingle{}Cargo\ combination\textquotesingle{}} to find
consignments containing fresh produce, inanimate cargo, or both.

Combination dimensions are especially useful for \texttt{and} logic.
They allow the user to ask for entities that have this property and that
property while keeping the model at the correct grain.

\subsection{Choices dimension}\label{choices-dimension}

A \textbf{choices dimension} is a low-cardinality, multi-valued
dimension that lets users filter for entities that have any selected
attribute.

Choices dimensions are related to combination dimensions, but the
interaction is different. A combination dimension is useful when the
user cares about combinations of properties. A choices dimension is
useful when the user wants any matching choice.

The simplest example is a drink preference question where a person may
choose tea, coffee, both, or neither. The fact table stays at the grain
of the person. It does not expand to one row per drink.

\textbf{Example structure of
\texttt{\textquotesingle{}Table\ order\textquotesingle{}}}

{\def\LTcaptype{none} % do not increment counter
\begin{longtable}[]{@{}rlr@{}}
\toprule\noalign{}
Table order ID & Order date & Drink choice group SK \\
\midrule\noalign{}
\endhead
\bottomrule\noalign{}
\endlastfoot
1 & 2024-01-15 & 2 \\
2 & 2024-01-15 & 3 \\
3 & 2024-01-16 & 3 \\
4 & 2024-01-16 & 1 \\
5 & 2024-01-17 & 4 \\
6 & 2024-01-17 & 4 \\
\end{longtable}
}

\textbf{Example structure of
\texttt{\textquotesingle{}Drink\ choice\textquotesingle{}}}

{\def\LTcaptype{none} % do not increment counter
\begin{longtable}[]{@{}rl@{}}
\toprule\noalign{}
Drink choice group SK & Drink choice \\
\midrule\noalign{}
\endhead
\bottomrule\noalign{}
\endlastfoot
1 & No drink \\
2 & Tea \\
3 & Coffee \\
4 & Tea \\
4 & Coffee \\
\end{longtable}
}

The fact table carries \texttt{{[}Drink\ choice\ group\ SK{]}}. The
choices dimension contains one row per selected choice within each
group.

If the user selects \texttt{Tea}, the model returns table orders
\texttt{1}, \texttt{5}, and \texttt{6}. If the user selects
\texttt{Coffee}, the model returns table orders \texttt{2}, \texttt{3},
\texttt{5}, and \texttt{6}. Table orders \texttt{5} and \texttt{6}
appear in both results because their group contains both choices. The
grain of \texttt{\textquotesingle{}Table\ order\textquotesingle{}} has
not changed.

Without a choices dimension, there are two common alternatives.

The first is to store each choice as a separate Boolean column in the
fact table, such as \texttt{{[}Has\ tea{]}} and
\texttt{{[}Has\ coffee{]}}. But it is awkward when the user wants to
find tables with tea \emph{or} coffee---the filter will pick rows that
have tea \emph{and} coffee instead.

The second is to change the grain of the fact table so that each order
appears once per drink choice. This makes the choice easy to filter, but
\texttt{\textquotesingle{}Table\ order\textquotesingle{}} is no longer
at the grain of an order. It is now at the grain of an order-choice.
That may be correct for some models, but often not.

A choices dimension avoids both problems. The fact remains at the order
grain, while the user can still select individual choices using
\texttt{or} logic. Thus, a choices dimension is a set-membership bridge
expressed as a dimension. It preserves the entity grain while still
letting the user interact with individual choices.

A common business example is data quality. A transaction row may have
multiple issues: missing date, missing officer, invalid status, or
missing location. A data quality report needs to show records that have
any selected issue, while still preserving one row per transaction.

\textbf{Example structure of
\texttt{\textquotesingle{}Transaction\textquotesingle{}}}

{\def\LTcaptype{none} % do not increment counter
\begin{longtable}[]{@{}rlr@{}}
\toprule\noalign{}
Transaction SK & Transaction date & Issue group SK \\
\midrule\noalign{}
\endhead
\bottomrule\noalign{}
\endlastfoot
1001 & 2024-01-15 & 2 \\
1002 & 2024-01-16 & 3 \\
1003 & 2024-01-17 & 4 \\
1004 & 2024-01-18 & 1 \\
\end{longtable}
}

\textbf{Example structure of
\texttt{\textquotesingle{}Data\ quality\ issue\textquotesingle{}}}

{\def\LTcaptype{none} % do not increment counter
\begin{longtable}[]{@{}rl@{}}
\toprule\noalign{}
Issue group SK & Data quality issue \\
\midrule\noalign{}
\endhead
\bottomrule\noalign{}
\endlastfoot
1 & No issue \\
2 & Missing date \\
2 & Missing officer \\
3 & Missing location \\
3 & Invalid status \\
4 & Missing date \\
4 & Missing officer \\
4 & Missing location \\
\end{longtable}
}

The fact table carries \texttt{{[}Issue\ group\ SK{]}}. The choices
dimension contains one row per issue within each group.

This allows a user to select \texttt{Missing\ date} and see all
transactions whose issue group includes \texttt{Missing\ date}. It also
allows the user to select several issues without forcing
\texttt{\textquotesingle{}Transaction\textquotesingle{}} to expand to
one row per issue.

Physically, a choices dimension is created by taking the possible sets
of choices from a finite list and assigning each set a group key. The
relationship between the fact and the choices dimension is on that group
key.

Choice and combination dimensions are closely related. A combination
dimension can often be unpivoted to create a choices dimension.

For example,
\texttt{\textquotesingle{}System\ combination\textquotesingle{}} may
contain \texttt{{[}Has\ System\ A{]}}, \texttt{{[}Has\ System\ B{]}},
and \texttt{{[}Has\ System\ C{]}}. This can be unpivoted into a
\texttt{\textquotesingle{}System\ touch\ point\textquotesingle{}}
choices dimension with values such as \texttt{System\ A},
\texttt{System\ B}, and \texttt{System\ C}.

\textbf{Example structure of
\texttt{\textquotesingle{}System\ touch\ point\textquotesingle{}}}

{\def\LTcaptype{none} % do not increment counter
\begin{longtable}[]{@{}rl@{}}
\toprule\noalign{}
System combination SK & System touch point \\
\midrule\noalign{}
\endhead
\bottomrule\noalign{}
\endlastfoot
1 & No system \\
2 & System A \\
3 & System B \\
4 & System C \\
5 & System A \\
5 & System B \\
6 & System A \\
6 & System C \\
7 & System B \\
7 & System C \\
8 & System A \\
8 & System B \\
8 & System C \\
\end{longtable}
}

The decision depends on the question. If the business wants
combinations, use a combination dimension. If the business wants any
selected value, use a choices dimension. Since they are lightweight, it
is always possible to have both if beneficial.

\subsection{Sankey dimension}\label{sankey-dimension}

A \textbf{Sankey dimension} represents possible paths through a sequence
of checkpoints.

This pattern is useful when the business wants to visualise movement
through a process. The source fact should not be reshaped just to create
a flow visual. Instead, each possible path can be represented as a
dimension.

For example, an entity may move through follow-on systems: System A,
System B, and System C. Some entities begin in System A and continue
through all systems. Some begin in System B. Some skip System B and move
directly from System A to System C.

A \texttt{\textquotesingle{}System\ path\textquotesingle{}} dimension
can represent each possible path as a set of edges.

\textbf{Example structure of
\texttt{\textquotesingle{}System\ path\textquotesingle{}}}

{\def\LTcaptype{none} % do not increment counter
\begin{longtable}[]{@{}rll@{}}
\toprule\noalign{}
System combination SK & From system & To system \\
\midrule\noalign{}
\endhead
\bottomrule\noalign{}
\endlastfoot
1 & Start & No system \\
2 & Start & System A \\
3 & Start & System B \\
4 & Start & System C \\
5 & Start & System A \\
5 & System A & System B \\
6 & Start & System A \\
6 & System A & System C \\
7 & Start & System B \\
7 & System B & System C \\
8 & Start & System A \\
8 & System A & System B \\
8 & System B & System C \\
\end{longtable}
}

The fact table carries \texttt{{[}System\ combination\ SK{]}}. The
Sankey dimension expands that path into the rows required by the Sankey
visual.

This lets the user see process flow without changing the grain of the
fact. The fact remains at the entity grain. The dimension supplies the
visual structure.

Notice that \texttt{{[}System\ combination\ SK{]}} is the same key used
by \texttt{\textquotesingle{}System\ combination\textquotesingle{}},
\texttt{\textquotesingle{}System\ touch\ point\textquotesingle{}}, and
\texttt{\textquotesingle{}System\ path\textquotesingle{}}.

This is the same underlying structure expressed for different usage
scenarios.

\texttt{\textquotesingle{}System\ combination\textquotesingle{}}
presents the systems as a combination of binary choices. It answers:
which exact combination of systems was involved?

\texttt{\textquotesingle{}System\ touch\ point\textquotesingle{}}
presents the same structure as a choices dimension. It answers: did the
entity involve any selected system?

\texttt{\textquotesingle{}System\ path\textquotesingle{}} presents the
same structure as a Sankey dimension. It answers: how did the entity
move through the systems?

Computationally these are low costs and low complexity translations.

The key is to recognise the underlying axes of binary choice. Once those
axes are identified, the data engineer can express them in different
dimensional forms depending on the interaction the user needs. Thus, the
data engineer is not merely arranging tables, but recognising the latent
structure of interaction in the business problem.

\subsection{Storytelling dimension}\label{storytelling-dimension}

A \textbf{storytelling dimension} groups entities into meaningful
business journeys.

This pattern is used when the raw process is too detailed for the user's
first point of entry. Instead of beginning with many statuses, events,
checkpoints, or paths, the model gives users a small set of meaningful
journeys.

Where a Sankey dimension shows movement, a storytelling dimension
translates these movements into manageable categories that are
business-meaningful.

For example, the
\hyperref[docs-creating-information-storytelling]{Storytelling} chapter
describes a help desk escalation dimension that combines tier
combinations, highest escalation, escalation path, escalation summary,
and display order. The point is not merely to record the path a case
took, but to express what that path means to the business.

A storytelling dimension is created when the data engineer steps back
from raw data and asks how the business understands the movement of the
entity.

In a self-service model, storytelling dimensions often act as the first
portal into the data. They divide the population into meaningful strata
and help users begin analysis before drilling into detail.

\subsection{Search dimension}\label{search-dimension}

A \textbf{search dimension} is a high-cardinality dimension that gives
users a controlled way to find, retrieve, drill through, or cross-filter
records.

Most dimensions are useful because they support aggregation. Search
dimensions are useful because they provide a way into the model.

In traditional dimensional modelling, a transaction is usually modelled
as a fact table, not as a dimension. Transaction identifiers are often
kept inside fact tables as degenerate dimensions. In an interactive
Power BI model, this is not always enough. Sometimes the user needs to
search for a specific transaction, use it for drillthrough, or
cross-filter related facts across multiple business processes.

In that case, a transaction-level value can become a search dimension.

An important special case is the \textbf{ID dimension}.

An ID dimension is a minimal search dimension that contains only the key
fields needed for search, usually the surrogate key and one or more
business identifiers.

\textbf{Example structure of
\texttt{\textquotesingle{}Inspection\ ID\textquotesingle{}}}

{\def\LTcaptype{none} % do not increment counter
\begin{longtable}[]{@{}rl@{}}
\toprule\noalign{}
Inspection SK & Inspection number \\
\midrule\noalign{}
\endhead
\bottomrule\noalign{}
\endlastfoot
1001 & INS-2024-0001 \\
1002 & INS-2024-0002 \\
\end{longtable}
}

This dimension lets a user retrieve records using an identifier they
already know. In traditional dimensional modelling, this pattern may be
treated as an exception. In interactive reporting, it is often practical
and necessary because users frequently arrive with a known business key.
They want to paste or search an ID and see the relevant records.

Search dimensions can also support high-cardinality text values that
users treat as search handles.

\textbf{Example structure of
\texttt{\textquotesingle{}Package\ search\textquotesingle{}}}

{\def\LTcaptype{none} % do not increment counter
\begin{longtable}[]{@{}rll@{}}
\toprule\noalign{}
Package SK & Package reference & Package description \\
\midrule\noalign{}
\endhead
\bottomrule\noalign{}
\endlastfoot
7001 & PKG-2024-0001 & Refrigerated seafood consignment \\
7002 & PKG-2024-0002 & Machinery parts and spare components \\
\end{longtable}
}

The modelling decision depends on the user's intent.

If the user expects to search using the field, model it as a search
dimension. If the user only expects to display the field after arriving
at a record, model it as an annotation fact.

Search dimensions can significantly increase model size. They should be
used deliberately, and they are not ideal for aggregation.

Due to Power BI filtering behaviour, search dimensions often need to be
used with a unit-record display measure. This is explained in the
\hyperref[docs-presenting-insights-filtering-behaviour]{Filtering
behaviour} and
\hyperref[docs-presenting-insights-designing-measures]{Designing
measures} chapters.

\subsection{Degenerate dimension}\label{degenerate-dimension}

A \textbf{degenerate dimension} is a dimension-like field retained
inside a fact table rather than lifted into a separate dimension table.

This usually occurs for one of two reasons:

\begin{itemize}
\tightlist
\item
  The field does not have enough business weight to justify a standalone
  dimension.
\item
  The field has high cardinality and would increase model size if lifted
  into a dimension.
\end{itemize}

\textbf{Example of a degenerate dimension
\texttt{{[}Inspected\ goods{]}} inside
\texttt{\textquotesingle{}Inspection\textquotesingle{}}}

{\def\LTcaptype{none} % do not increment counter
\begin{longtable}[]{@{}
  >{\raggedleft\arraybackslash}p{(\linewidth - 8\tabcolsep) * \real{0.2353}}
  >{\raggedright\arraybackslash}p{(\linewidth - 8\tabcolsep) * \real{0.1765}}
  >{\raggedright\arraybackslash}p{(\linewidth - 8\tabcolsep) * \real{0.1765}}
  >{\raggedleft\arraybackslash}p{(\linewidth - 8\tabcolsep) * \real{0.2353}}
  >{\raggedright\arraybackslash}p{(\linewidth - 8\tabcolsep) * \real{0.1765}}@{}}
\toprule\noalign{}
\begin{minipage}[b]{\linewidth}\raggedleft
Inspection SK
\end{minipage} & \begin{minipage}[b]{\linewidth}\raggedright
Inspection number
\end{minipage} & \begin{minipage}[b]{\linewidth}\raggedright
Inspection date
\end{minipage} & \begin{minipage}[b]{\linewidth}\raggedleft
Inspection status SK
\end{minipage} & \begin{minipage}[b]{\linewidth}\raggedright
Inspected goods
\end{minipage} \\
\midrule\noalign{}
\endhead
\bottomrule\noalign{}
\endlastfoot
1001 & INS-2024-0001 & 2024-01-15 & 1 & Refrigerated seafood
consignment \\
1002 & INS-2024-0002 & 2024-01-16 & 2 & Machinery parts and spare
components \\
1003 & INS-2024-0003 & 2024-01-16 & 1 & Fresh mangoes and tropical
fruit \\
1004 & INS-2024-0004 & 2024-01-17 & 3 & Timber pallets and packaging
materials \\
1005 & INS-2024-0005 & 2024-01-18 & 2 & Laboratory testing equipment \\
\end{longtable}
}

Here, \texttt{{[}Inspected\ goods{]}} is the degenerate dimension, and
can be used for searching. This works technically, but it can weaken the
interaction logic of the model. It makes the fact table behave like a
control surface.

Degenerate dimensions should therefore be treated with caution. They are
sometimes necessary, but they are often a sign that a field has not been
properly expressed as a business dimension, search dimension, or
annotation fact.

For example, if \texttt{Inspected\ goods} is only needed for display,
the column can remain hidden in the model and be exposed through a
measure using \texttt{SELECTEDVALUE}. If it is needed for search or
cross-filtering, it may be better expressed as an
\texttt{\textquotesingle{}Inspected\ goods\textquotesingle{}} search
dimension related by \texttt{{[}Inspection\ SK{]}}.

The question is not whether the field can remain in the fact table. The
question is whether the user needs to touch it. The three problems of
degenerate dimensions are discussed in the next chapter.

If the user needs to touch it, it probably belongs on the controller,
not on the robot.

\section{Pushing detail back to the
entity}\label{pushing-detail-back-to-the-entity}

One common theme in this chapter is the movement from information at a
finer grain than the entity back to the entity level, and then into a
dimension.

Combination dimensions, choice dimensions, Sankey dimensions, and
storytelling dimensions all follow this pattern.

\pandocbounded{\includegraphics[keepaspectratio]{book/epub-assets/diagram-007.png}}

\emph{Figure 1. Lower-grain detail is pushed back to the entity grain,
then exposed as a dimension value that users can interact with safely.}

This movement is powerful because it keeps the focus on the entity of
interest, both technically and business-wise.

Technically, the entity of interest remains the common denominator for
joins. The user gets a meaningful control surface without forcing the
model into lower-grain chaos.

Business-wise, user interactions remain centred on the entity the
business actually cares about.

The purpose of this chapter is not to define an exhaustive list of all
possible artefacts in a dimensional model. Ultimately, each table is
either a dimension or a fact, and the sub-categorisation does not
significantly change this.

Rather, the purpose of listing these elements is to help data engineers
shift to a mindset of seeing dimensions and facts through usage
scenarios, and to recognise that they can play functionally different
roles depending on the interactivity requirements.

In the end, the final test is this: can the business user answer a
question without needing to guess?

That is the heart of dimensional modelling in an interactive model.

\begin{quote}
\textbf{Key ideas.}

In an interactive model, dimensions are what users touch and facts are
what respond.

Facts contain business activity, business state, or display detail.

Dimensions provide controlled surfaces for filtering, grouping,
searching, comparing, and storytelling.

A fact table should not become the user's control panel. The controls
belong on the controller, not on the robot.

Some dimensions exist to push lower-grain complexity back to the entity
grain and expose it as something users can interact with safely.

A good dimensional model lets users interact with business activity
reliably without guessing how the model works.
\end{quote}

\cleardoublepage

\chapter{Filtering
behaviour}\label{docs-presenting-insights-filtering-behaviour}

\section{Filtering as interaction
design}\label{filtering-as-interaction-design}

Filtering design is what makes a Power BI model feel intuitive---or
frustrating.

Many engineers stop too early, treating the job as complete once the
correct answer is reachable through some combination of filters.

That is not enough.

A self-service model is not a puzzle for the user to solve. The user
should not have to discover the one valid route through the model. They
should be able to approach the model from different
directions---starting with a date, a product, an ID, a measure, or a
business process---and still receive sensible results at each step.

The data engineer's task is therefore not only to make the correct
answer possible. It is to anticipate how users will try to reach that
answer. More generally, it is to anticipate what buttons the user will
click, and what effects the model should produce.

A good dimensional model is designed around the interaction between
\hyperref[docs-presenting-insights-dimensional-modelling-for-ux]{buttons
and effects}. Filtering is the mechanism through which this interaction
takes place.

This chapter explains how filtering works so the data engineer can
anticipate user interaction and design the model's response
deliberately.

The chapter proceeds in two parts.

First, it explains the main ways filtering occurs in Power BI. These are
the basic mechanisms the data engineer needs to understand: relationship
filtering, common-fact filtering, non-blank measure filtering,
visual-level filtering, and measure-defined filtering.

Second, it applies those mechanisms to tricky modelling scenarios. These
include displaying unit records, coordinating multiple business
processes, cascading filters, aggregating dimension values,
filter-by-having, and dynamic Type I behaviour.

The purpose is not to memorise tricks. The purpose is to understand how
Power BI responds when users interact with the model, so the data
engineer can design those responses deliberately.

\section{Ways of filtering}\label{ways-of-filtering}

There are five common ways filtering appears in Power BI:

\begin{itemize}
\tightlist
\item
  Relationship filtering
\item
  Common-fact filtering
\item
  Non-blank measure filtering
\item
  Visual-level filtering
\item
  Measure-defined filtering
\end{itemize}

\subsection{Relationship filtering}\label{relationship-filtering}

Relationship filtering is the standard case of filtering.

When table X filters table Y, a user who selects a value from X narrows
the rows in Y based on the relationship between the two tables. In a
dimensional model, the most important case is when dimensions filter
facts.

The user touches the dimension. The fact responds.

\textbf{Example structure of
\texttt{\textquotesingle{}Product\textquotesingle{}}}

{\def\LTcaptype{none} % do not increment counter
\begin{longtable}[]{@{}rl@{}}
\toprule\noalign{}
Product SK & Product name \\
\midrule\noalign{}
\endhead
\bottomrule\noalign{}
\endlastfoot
1 & Cake \\
2 & Bread \\
3 & Pastry \\
\end{longtable}
}

\textbf{Example structure of
\texttt{\textquotesingle{}Sale\textquotesingle{}}}

{\def\LTcaptype{none} % do not increment counter
\begin{longtable}[]{@{}lr@{}}
\toprule\noalign{}
Sale ID & Product SK \\
\midrule\noalign{}
\endhead
\bottomrule\noalign{}
\endlastfoot
S1001 & 1 \\
S1002 & 1 \\
S1003 & 2 \\
\end{longtable}
}

If the user selects \texttt{Cake} from
\texttt{\textquotesingle{}Product\textquotesingle{}}, the model filters
\texttt{\textquotesingle{}Sale\textquotesingle{}} to \texttt{S1001} and
\texttt{S1002}. If the user selects \texttt{Bread}, it filters to
\texttt{S1003}. If the user selects \texttt{Pastry}, there are no
matching sales rows.

New engineers often ask, ``Why do we need dimensions at all? Why not
leave all the information in the fact table? It is simpler.'' That is,
why not rely solely on degenerate dimensions?

There are three common problems.

First, non-existence becomes impossible to express. If product were only
stored as a degenerate column inside
\texttt{\textquotesingle{}Sale\textquotesingle{}}, then \texttt{Pastry}
would not exist anywhere in the model. A user would not be able to know
that no pastries were sold.

Second, the business meaning becomes buried. Attributes such as region,
product, demographic group, or status are no longer presented as
coherent business objects. They become miscellaneous columns inside a
fact table. A user would have to read every single row to reverse
engineer the entity and its grain.

Third, cross-filtering and drillthrough across multiple business
processes become harder. A conformed dimension such as
\texttt{\textquotesingle{}Calendar\textquotesingle{}},
\texttt{\textquotesingle{}Product\textquotesingle{}}, or
\texttt{\textquotesingle{}Region\textquotesingle{}} can filter several
facts. A degenerate column inside one fact cannot naturally serve as a
shared point of interaction.

\subsection{Common-fact filtering}\label{common-fact-filtering}

Filtering also appears indirectly through common filtered tables.

This behaviour is often missed by new engineers, but it is essential to
the intuitive experience of Power BI. When two dimensions both filter
the same fact, Power BI can use the fact to determine which combinations
of dimension values are valid.

Consider a simple sales model.

\textbf{Example structure of
\texttt{\textquotesingle{}Product\textquotesingle{}}}

{\def\LTcaptype{none} % do not increment counter
\begin{longtable}[]{@{}rl@{}}
\toprule\noalign{}
Product SK & Product name \\
\midrule\noalign{}
\endhead
\bottomrule\noalign{}
\endlastfoot
1 & Cake \\
2 & Bread \\
3 & Pastry \\
\end{longtable}
}

\textbf{Example structure of
\texttt{\textquotesingle{}Region\textquotesingle{}}}

{\def\LTcaptype{none} % do not increment counter
\begin{longtable}[]{@{}rl@{}}
\toprule\noalign{}
Region SK & Region name \\
\midrule\noalign{}
\endhead
\bottomrule\noalign{}
\endlastfoot
1 & North \\
2 & South \\
3 & West \\
\end{longtable}
}

\textbf{Example structure of
\texttt{\textquotesingle{}Sale\textquotesingle{}}}

{\def\LTcaptype{none} % do not increment counter
\begin{longtable}[]{@{}lrrr@{}}
\toprule\noalign{}
Sale ID & Product SK & Region SK & Sales amount \\
\midrule\noalign{}
\endhead
\bottomrule\noalign{}
\endlastfoot
S1001 & 1 & 1 & 100 \\
S1002 & 1 & 2 & 150 \\
S1003 & 2 & 1 & 200 \\
\end{longtable}
}

If a visual contains
\texttt{\textquotesingle{}Product\textquotesingle{}{[}Product\ name{]}}
and
\texttt{\textquotesingle{}Region\textquotesingle{}{[}Region\ name{]}},
Power BI returns only product-region combinations that exist in
\texttt{\textquotesingle{}Sale\textquotesingle{}}.

The fact table acts as the common filtered table that makes the
dimension combination valid.

The result would be:

{\def\LTcaptype{none} % do not increment counter
\begin{longtable}[]{@{}ll@{}}
\toprule\noalign{}
Product name & Region name \\
\midrule\noalign{}
\endhead
\bottomrule\noalign{}
\endlastfoot
Cake & North \\
Cake & South \\
Bread & North \\
\end{longtable}
}

\texttt{Pastry} does not appear because there is no matching row in
\texttt{\textquotesingle{}Sale\textquotesingle{}}. \texttt{West} does
not appear because there is no matching row in
\texttt{\textquotesingle{}Sale\textquotesingle{}}. The missing
combinations are not shown because they do not exist in the common fact.

This common-fact behaviour is part of what makes a dimensional model
feel natural. Users can place dimension fields together, and the model
silently removes combinations that do not exist.

If two dimensions do not share any filtered fact table, Power BI cannot
determine the valid combinations between them. In that case, the visual
may return the Power BI error message:
\texttt{Can\textquotesingle{}t\ determine\ relationships\ between\ the\ fields.}

If multiple fact tables are filtered by both dimensions, Power BI shows
combinations where any fact table has a matching row. This behaviour is
powerful, but it also means the data engineer must think carefully about
which facts share which dimensions.

This ``upward'' filtering from facts to dimensions enables the intuitive
experience of a dimensional model. It is why dimensions alone can serve
as the interface, even when facts are hidden. Mastering this behaviour
is central to designing good Power BI models.

\subsection{Non-blank measure
filtering}\label{non-blank-measure-filtering}

The previous method relies on implicit filtering through relationships
and common facts. In simple models, this is often enough. In complex
models with several fact tables, it may not be.

Measures can also control what appears in a visual.

Suppose \texttt{\textquotesingle{}Product\textquotesingle{}} filters
\texttt{\textquotesingle{}Sale\textquotesingle{}}, and
\texttt{{[}Total\ sales\ amount{]}} sums the sales amount in
\texttt{\textquotesingle{}Sale\textquotesingle{}}.

\textbf{Visual before adding \texttt{{[}Total\ sales\ amount{]}}}

{\def\LTcaptype{none} % do not increment counter
\begin{longtable}[]{@{}l@{}}
\toprule\noalign{}
Product name \\
\midrule\noalign{}
\endhead
\bottomrule\noalign{}
\endlastfoot
Cake \\
Bread \\
Pastry \\
\end{longtable}
}

\textbf{Visual after adding \texttt{{[}Total\ sales\ amount{]}}}

{\def\LTcaptype{none} % do not increment counter
\begin{longtable}[]{@{}lr@{}}
\toprule\noalign{}
Product name & Total sales amount \\
\midrule\noalign{}
\endhead
\bottomrule\noalign{}
\endlastfoot
Cake & 250 \\
Bread & 200 \\
\end{longtable}
}

\texttt{Pastry} disappears because \texttt{{[}Total\ sales\ amount{]}}
is blank for that product in the current filter context.

This is non-blank measure filtering. When a measure is added to a table
or matrix visual, Power BI tends to retain rows where the measure
returns a non-blank value.

If several measures are added, the visual may retain rows where any
measure returns a value.

This mirrors fact table filtering, but gives the data engineer more
control. Instead of relying only on row existence in a fact table, the
data engineer can define the measure logic that determines which values
should remain visible.

\subsection{Visual-level filtering}\label{visual-level-filtering}

Sometimes a measure cannot be placed directly in the visual, or doing so
would interfere with the visual's layout.

For example, a slicer cannot simply display a helper measure as one of
its fields. A table visual may also need to be filtered by a measure
without showing that measure to the user.

In these cases, the report developer can apply a visual-level filter
using the measure value.

For example, \texttt{{[}Has\ sales{]}} may return \texttt{1} when there
is at least one sale in the current context. A slicer on
\texttt{\textquotesingle{}Product\textquotesingle{}{[}Product\ name{]}}
can then use \texttt{{[}Has\ sales{]}} as a visual-level filter so that
only products with sales appear.

This method is useful, but should be used with care. The evaluation cost
can be significant when the visual has a large number of possible values
and the measure must be calculated for each one.

\subsection{Measure-defined filtering}\label{measure-defined-filtering}

The previous methods control which dimension values appear in a visual.
Sometimes the data engineer needs to apply filters directly inside the
measure calculation itself.

This is measure-defined filtering.

Power BI supports this through DAX expressions such as
\texttt{userelationship} and \texttt{crossfilter}. These expressions
allow the data engineer to activate a relationship, modify a
relationship direction, apply a virtual relationship, or otherwise
change the filter context for the duration of the calculation.

Consider a sales model where
\texttt{\textquotesingle{}Sale\textquotesingle{}} has both
\texttt{{[}Order\ date{]}} and \texttt{{[}Delivery\ date{]}}.

\textbf{Example structure of
\texttt{\textquotesingle{}Calendar\textquotesingle{}}}

{\def\LTcaptype{none} % do not increment counter
\begin{longtable}[]{@{}ll@{}}
\toprule\noalign{}
Date & Month name \\
\midrule\noalign{}
\endhead
\bottomrule\noalign{}
\endlastfoot
2024-01-15 & January \\
2024-01-20 & January \\
2024-02-05 & February \\
\end{longtable}
}

\textbf{Example structure of
\texttt{\textquotesingle{}Sale\textquotesingle{}}}

{\def\LTcaptype{none} % do not increment counter
\begin{longtable}[]{@{}lrllr@{}}
\toprule\noalign{}
Sale ID & Product SK & Order date & Delivery date & Sales amount \\
\midrule\noalign{}
\endhead
\bottomrule\noalign{}
\endlastfoot
S1001 & 1 & 2024-01-15 & 2024-01-20 & 100 \\
S1002 & 1 & 2024-01-16 & 2024-02-05 & 150 \\
S1003 & 2 & 2024-02-01 & 2024-02-06 & 200 \\
\end{longtable}
}

Suppose \texttt{\textquotesingle{}Calendar\textquotesingle{}{[}Date{]}}
has an active relationship to
\texttt{\textquotesingle{}Sale\textquotesingle{}{[}Order\ date{]}} and
an inactive relationship to
\texttt{\textquotesingle{}Sale\textquotesingle{}{[}Delivery\ date{]}}.

A normal measure such as \texttt{{[}Total\ sales\ amount{]}} responds to
the active relationship. If the user selects January, the measure
returns sales ordered in January.

\begin{verbatim}
Total sales amount =
sum ( 'Sale'[Sales amount] )
\end{verbatim}

But the user may also need sales delivered in January. The data engineer
can define a second measure that activates the delivery-date
relationship only for that calculation.

\begin{verbatim}
Total delivered sales amount =
calculate (
    [Total sales amount],
    userelationship ( 'Calendar'[Date], 'Sale'[Delivery date] )
)
\end{verbatim}

With the example above, selecting January gives
\texttt{{[}Total\ sales\ amount{]}\ =\ 250}, because \texttt{S1001} and
\texttt{S1002} were ordered in January. The delivered-sales measure
gives \texttt{100}, because only \texttt{S1001} was delivered in
January.

The model relationship does not permanently change. The filtering path
changes only while \texttt{{[}Total\ delivered\ sales\ amount{]}} is
being evaluated.

Measure-defined filtering is powerful because it allows precise
behaviour without permanently changing the model. It is also dangerous
because the filtering rule becomes hidden inside the measure. Used well,
it solves specific interaction problems. Used carelessly, it creates a
model where results are technically correct but difficult to explain.

\section{Filtering scenarios}\label{filtering-scenarios}

The five filtering mechanisms above can be combined to support a wide
range of interactivity goals while preserving the standard dimensional
model pattern: dimensions filter facts using single-direction
relationships.

The scenarios below are not exhaustive. Their purpose is to show how
filtering design supports user interaction in real models.

\subsection{Displaying unit records for a single business
process}\label{displaying-unit-records-for-a-single-business-process}

A common requirement is to show one row per business transaction.

The obvious approach is to expose the fact table. For example, if the
transaction is a sale, the user might expect to drag fields directly
from \texttt{\textquotesingle{}Sale\textquotesingle{}}.

This is usually the wrong interface.

A better approach is to hide
\texttt{\textquotesingle{}Sale\textquotesingle{}} and let dimensions and
measures reconstruct the transaction view.

The point of hiding the fact table is not tidiness. It is to protect the
user experience. If users interact directly with the fact table, they
can produce technically valid but strange results. If they interact
through dimensions and measures, the model is much more likely to behave
sensibly no matter which direction they approach it from.

\subsubsection{The model setup}\label{the-model-setup}

Suppose there is a business process represented by the
\texttt{\textquotesingle{}Sale\textquotesingle{}} fact table.

The setup is as follows:

\begin{itemize}
\tightlist
\item
  A fact table \texttt{\textquotesingle{}Sale\textquotesingle{}}
\item
  An ID dimension \texttt{\textquotesingle{}Sales\ ID\textquotesingle{}}
  with \texttt{{[}Sales\ order\ number{]}}
\item
  A calendar dimension
  \texttt{\textquotesingle{}Sales\ calendar\textquotesingle{}}
\item
  A business dimension
  \texttt{\textquotesingle{}Sales\ product\textquotesingle{}}
\item
  A measure \texttt{{[}Total\ sales\ amount{]}} that sums the sales
  amount in \texttt{\textquotesingle{}Sale\textquotesingle{}}
\end{itemize}

\textbf{Example structure of
\texttt{\textquotesingle{}Sale\textquotesingle{}}}

{\def\LTcaptype{none} % do not increment counter
\begin{longtable}[]{@{}llrr@{}}
\toprule\noalign{}
Sales order number & Sales date & Product SK & Sales amount \\
\midrule\noalign{}
\endhead
\bottomrule\noalign{}
\endlastfoot
SO1001 & 2024-01-15 & 1 & 100 \\
SO1002 & 2024-01-16 & 2 & 150 \\
\end{longtable}
}

\textbf{Example structure of
\texttt{\textquotesingle{}Sales\ ID\textquotesingle{}}}

{\def\LTcaptype{none} % do not increment counter
\begin{longtable}[]{@{}l@{}}
\toprule\noalign{}
Sales order number \\
\midrule\noalign{}
\endhead
\bottomrule\noalign{}
\endlastfoot
SO1001 \\
SO1002 \\
\end{longtable}
}

\textbf{Example structure of
\texttt{\textquotesingle{}Sales\ calendar\textquotesingle{}}}

{\def\LTcaptype{none} % do not increment counter
\begin{longtable}[]{@{}ll@{}}
\toprule\noalign{}
Sales date & Month name \\
\midrule\noalign{}
\endhead
\bottomrule\noalign{}
\endlastfoot
2024-01-15 & January \\
2024-01-16 & January \\
\end{longtable}
}

\textbf{Example structure of
\texttt{\textquotesingle{}Sales\ product\textquotesingle{}}}

{\def\LTcaptype{none} % do not increment counter
\begin{longtable}[]{@{}rl@{}}
\toprule\noalign{}
Product SK & Product name \\
\midrule\noalign{}
\endhead
\bottomrule\noalign{}
\endlastfoot
1 & Cake \\
2 & Bread \\
\end{longtable}
}

\subsubsection{Reconstructing the transaction
row}\label{reconstructing-the-transaction-row}

The user should reach the sale through:

\begin{itemize}
\tightlist
\item
  \texttt{\textquotesingle{}Sales\ ID\textquotesingle{}{[}Sales\ order\ number{]}}
\item
  \texttt{\textquotesingle{}Sales\ calendar\textquotesingle{}{[}Sales\ date{]}}
\item
  \texttt{\textquotesingle{}Sales\ product\textquotesingle{}{[}Product\ name{]}}
\item
  \texttt{{[}Total\ sales\ amount{]}}
\end{itemize}

Together, these fields give the user the full sale record without
exposing the fact table itself.

This works because Power BI returns only combinations of dimension
values that correspond to existing rows in
\texttt{\textquotesingle{}Sale\textquotesingle{}} (common-fact
filtering). Since
\texttt{\textquotesingle{}Sales\ ID\textquotesingle{}{[}Sales\ order\ number{]}}
uniquely identifies the sales transaction, adding
\texttt{\textquotesingle{}Sales\ ID\textquotesingle{}{[}Sales\ order\ number{]}},
with either
\texttt{\textquotesingle{}Sales\ calendar\textquotesingle{}{[}Sales\ date{]}}
or
\texttt{\textquotesingle{}Sales\ product\textquotesingle{}{[}Product\ name{]}}
downfilters \texttt{\textquotesingle{}Sales\textquotesingle{}} to the
transactions that exist.

\subsubsection{Why use a measure for the
amount?}\label{why-use-a-measure-for-the-amount}

\texttt{\textquotesingle{}Sale\textquotesingle{}} has an amount recorded
in \texttt{{[}Sales\ amount{]}}. A simplistic approach is to expose
\texttt{\textquotesingle{}Sale\textquotesingle{}{[}Sales\ amount{]}} to
the user as though it were an ordinary field. Instead, a better approach
is to expose \texttt{{[}Total\ sales\ amount{]}}.

Because
\texttt{\textquotesingle{}Sales\ ID\textquotesingle{}{[}Sales\ order\ number{]}}
uniquely identifies one row, \texttt{{[}Total\ sales\ amount{]}} returns
the sum of that one row in the current context---that is, the same
numeric value as
\texttt{\textquotesingle{}Sale\textquotesingle{}{[}Sales\ amount{]}} at
the transaction level (non-blank measure filtering).

At the sales order level, the measure returns the value of one
transaction. At the product level, the same measure aggregates all
matching transactions. The user does not need to choose between a
transaction amount and an aggregate amount. The measure behaves
correctly in both situations. For example:

{\def\LTcaptype{none} % do not increment counter
\begin{longtable}[]{@{}lr@{}}
\toprule\noalign{}
Current context & Result \\
\midrule\noalign{}
\endhead
\bottomrule\noalign{}
\endlastfoot
\texttt{SO1001} & \texttt{{[}Total\ sales\ amount{]}\ =\ 100} \\
\texttt{January} & \texttt{{[}Total\ sales\ amount{]}\ =\ 250} \\
\end{longtable}
}

This sounds complex in the backend, but it is seamless for the user. In
the ideal setup, \texttt{\textquotesingle{}Sale\textquotesingle{}} and
\texttt{\textquotesingle{}Sale\textquotesingle{}{[}Sales\ amount{]}} are
hidden. The user sees \texttt{{[}Total\ sales\ amount{]}}, not both the
measure and the raw fact column.

This is far superior to exposing both. At the transaction level, both
may show the same number, which encourages users to treat them as
interchangeable. But in charts, aggregations, and time intelligence,
only the measure behaves correctly. That inconsistency is confusing.
Something appears to work in one visual, then fails in another.

With this setup, \texttt{{[}Total\ sales\ amount{]}} behaves
consistently across situations. The user gets sensible results without
needing to remember special cases.

\subsubsection{Why this remains
intuitive}\label{why-this-remains-intuitive}

The benefit is that the model behaves sensibly from multiple directions.
The user can start with the sales order number, the date, the product,
or the measure. In each case, Power BI's filtering behaviour narrows the
model toward valid sales records.

First, suppose the user selects
\texttt{\textquotesingle{}Sales\ ID\textquotesingle{}{[}Sales\ order\ number{]}},
then
\texttt{\textquotesingle{}Sales\ calendar\textquotesingle{}{[}Sales\ date{]}}.

Power BI immediately gives the valid combination of sales order numbers
and sales dates because the dimension values are filtered by their
common fact. This returns the desired transaction record because the
business key selects the correct grain.

Second, suppose the user selects
\texttt{\textquotesingle{}Sales\ ID\textquotesingle{}{[}Sales\ order\ number{]}},
then \texttt{{[}Total\ sales\ amount{]}}.

This also gives the transaction record. This time, it is not only
because of the common fact. It is also because
\texttt{{[}Total\ sales\ amount{]}} returns a non-blank value for the
sales order number. If there are sales order numbers without a sales
amount, such as unconfirmed sales, those order numbers will not appear.
For most use cases, that is the desired outcome.

Third, suppose the user selects \texttt{{[}Total\ sales\ amount{]}},
then
\texttt{\textquotesingle{}Sales\ product\textquotesingle{}{[}Product\ name{]}}.

The model first returns the total sales amount for the whole model as a
single number, then breaks it down by product. When the user adds
\texttt{\textquotesingle{}Sales\ ID\textquotesingle{}{[}Sales\ order\ number{]}},
the table expands to unit records.

The experience is correct at every step.

This is the power of a properly designed dimensional model focused on
dimensions and measures. Regardless of the sequence of field selection,
the model produces sensible business results.

This seamless experience would not occur if the user selected the
degenerate
\texttt{\textquotesingle{}Sale\textquotesingle{}{[}Sales\ amount{]}}
instead of \texttt{{[}Total\ sales\ amount{]}}. The user would first see
a distinct list of sales amount values, then a confusing combination of
sales amount and product, and only after adding
\texttt{\textquotesingle{}Sales\ ID\textquotesingle{}{[}Sales\ order\ number{]}}
would the table make sense.

The experience would be jarring rather than seamless.

Even this simple use case requires careful modelling work. The model
must cater for a large number of possible user interaction sequences.
These possibilities are handled elegantly when the data engineer focuses
on dimensions and measures, and hides facts.

In this example,
\texttt{\textquotesingle{}Sales\ ID\textquotesingle{}{[}Sales\ order\ number{]}}
is a single-column business key. If the business key has multiple
columns, all columns need to be used.

\subsection{Displaying unit records for two business
processes}\label{displaying-unit-records-for-two-business-processes}

The single-process case is difficult enough. The two-process case
introduces a new problem: different business processes may share some
controls but not others.

Continuing the previous example, suppose sales can be refunded.

The model contains:

\begin{itemize}
\tightlist
\item
  A fact table \texttt{\textquotesingle{}Sale\textquotesingle{}}
\item
  A fact table \texttt{\textquotesingle{}Refund\textquotesingle{}}
\item
  A shared ID dimension
  \texttt{\textquotesingle{}Sales\ ID\textquotesingle{}}
\item
  A shared product dimension
  \texttt{\textquotesingle{}Sales\ product\textquotesingle{}}
\item
  A sales calendar dimension
  \texttt{\textquotesingle{}Sales\ calendar\textquotesingle{}}
\item
  A refund calendar dimension
  \texttt{\textquotesingle{}Refund\ calendar\textquotesingle{}}
\item
  A measure \texttt{{[}Total\ sales\ amount{]}}
\item
  A measure \texttt{{[}Total\ refund\ amount{]}}
\end{itemize}

\textbf{Example structure of
\texttt{\textquotesingle{}Sale\textquotesingle{}}}

{\def\LTcaptype{none} % do not increment counter
\begin{longtable}[]{@{}llrr@{}}
\toprule\noalign{}
Sales order number & Sales date & Product SK & Sales amount \\
\midrule\noalign{}
\endhead
\bottomrule\noalign{}
\endlastfoot
SO1001 & 2024-01-15 & 1 & 100 \\
SO1002 & 2024-01-16 & 2 & 150 \\
SO1003 & 2024-01-17 & 1 & 200 \\
\end{longtable}
}

\textbf{Example structure of
\texttt{\textquotesingle{}Refund\textquotesingle{}}}

{\def\LTcaptype{none} % do not increment counter
\begin{longtable}[]{@{}llr@{}}
\toprule\noalign{}
Sales order number & Refund date & Refund amount \\
\midrule\noalign{}
\endhead
\bottomrule\noalign{}
\endlastfoot
SO1001 & 2024-01-20 & 50 \\
SO1001 & 2024-01-20 & 50 \\
SO1003 & 2024-01-22 & 75 \\
\end{longtable}
}

The dimensions relate to the two facts differently.

Sales and refunds share a sales order and product. They do not share the
same process date. Sales occur on sales dates; refunds occur on refund
dates. Sales are usually one row per order; refunds may have multiple
rows per order.

{\def\LTcaptype{none} % do not increment counter
\begin{longtable}[]{@{}llll@{}}
\toprule\noalign{}
Dimension & \texttt{\textquotesingle{}Sale\textquotesingle{}} &
\texttt{\textquotesingle{}Refund\textquotesingle{}} & Meaning \\
\midrule\noalign{}
\endhead
\bottomrule\noalign{}
\endlastfoot
\texttt{\textquotesingle{}Sales\ ID\textquotesingle{}} & 1 → * & 1 → * &
Shared sales order identifier \\
\texttt{\textquotesingle{}Sales\ product\textquotesingle{}} & 1 → * & 1
→ * & Product sold and later refunded \\
\texttt{\textquotesingle{}Sales\ calendar\textquotesingle{}} & 1 → * & &
Date of sale \\
\texttt{\textquotesingle{}Refund\ calendar\textquotesingle{}} & & 1 → *
& Date of refund \\
\end{longtable}
}

The data engineer therefore has to anticipate which process the user is
trying to display. Showing sales, showing refunds, and showing sales
with refunds are three different interaction problems.

This is the source of the complexity.

Some controls are shared.
\texttt{\textquotesingle{}Sales\ ID\textquotesingle{}} and
\texttt{\textquotesingle{}Sales\ product\textquotesingle{}} can filter
both sales and refunds.

Some controls are process-specific.
\texttt{\textquotesingle{}Sales\ calendar\textquotesingle{}} filters
sales. \texttt{\textquotesingle{}Refund\ calendar\textquotesingle{}}
filters refunds.

The model therefore behaves differently depending on what the user is
trying to display:

\begin{itemize}
\tightlist
\item
  Sales transactions
\item
  Refund transactions
\item
  Sales and refund information together
\end{itemize}

\subsection{Displaying sales}\label{displaying-sales}

When displaying sales records, everything works as if there were only
the \texttt{\textquotesingle{}Sale\textquotesingle{}} fact table.

This is because every refund comes from a sale.
\texttt{\textquotesingle{}Sale\textquotesingle{}} is therefore a
superset of \texttt{\textquotesingle{}Refund\textquotesingle{}}. Adding
\texttt{\textquotesingle{}Refund\textquotesingle{}} to the model does
not interfere with the display of sales transactions.

The superset relationship matters. If there were sales order numbers in
\texttt{\textquotesingle{}Refund\textquotesingle{}} that did not exist
in \texttt{\textquotesingle{}Sale\textquotesingle{}}, then selecting
\texttt{\textquotesingle{}Sales\ calendar\textquotesingle{}{[}Sales\ date{]}}
and
\texttt{\textquotesingle{}Sales\ ID\textquotesingle{}{[}Sales\ order\ number{]}}
would not correctly identify the list of sales transactions.

Single-direction filters also matter. If the report canvas has a filter
on \texttt{\textquotesingle{}Refund\ calendar\textquotesingle{}}, that
filter does not interfere with the presentation of sales transactions
because \texttt{\textquotesingle{}Refund\ calendar\textquotesingle{}}
filters \texttt{\textquotesingle{}Refund\textquotesingle{}}, not
\texttt{\textquotesingle{}Sale\textquotesingle{}}.

This would not be true if the model used bidirectional filters
carelessly. The refund date could unexpectedly affect sales visuals.

\subsection{Displaying refunds}\label{displaying-refunds}

When displaying refund records, the user is interested in
\texttt{{[}Total\ refund\ amount{]}}, but they may also need the sales
context, such as
\texttt{\textquotesingle{}Sales\ ID\textquotesingle{}{[}Sales\ order\ number{]}}
and
\texttt{\textquotesingle{}Sales\ product\textquotesingle{}{[}Product\ name{]}}.

While the outcome---displaying refund---is the same, the user may come
from two angles: starting with refund information, or starting from
sales information. Both need to work. In a properly designed dimensional
model, that is the case.

If the user selects \texttt{{[}Total\ refund\ amount{]}} first, the
model returns the total refund amount for the whole dataset (non-blank
measure filtering). Adding sales attributes breaks that number down.
Adding the relevant ID dimension returns unit-level records.

If the user selects sales attributes first, Power BI initially displays
the list of sales orders and products, regardless of whether they were
refunded. When the user adds either
\texttt{\textquotesingle{}Refund\ calendar\textquotesingle{}{[}Refund\ date{]}}
or \texttt{{[}Total\ refund\ amount{]}}, the table filters down to
refund transactions. The former because of common-fact filtering, the
latter because of non-blank measure filtering. And both happens because
those fields are linked to
\texttt{\textquotesingle{}Refund\textquotesingle{}}.

Whether the user starts with the sales attributes or the refund-specific
fields, the dimensional model gives a sensible path.

However, the situation is more complicated than it first appears.

Suppose a sales order of \texttt{\$100} (\texttt{SO1001} in the example)
is refunded in two partial refunds of \texttt{\$50} on the same day.
Because there is no primary key for refunds, the refund for that sales
order may appear as a single transaction with
\texttt{{[}Total\ refund\ amount{]}\ =\ 100}, rather than two rows of
\texttt{\$50}.

For many scenarios, showing the total refund amount, regardless of the
partial breakdowns, may be sufficient. In some cases, it is not.

When the user must see each refund action as a separate transaction, the
data engineer needs a key that distinguishes each refund transaction.

One option is to create \texttt{{[}Refund\ SK{]}}, a unique integer for
each refund, and expose it as
\texttt{\textquotesingle{}Refund\ ID\textquotesingle{}{[}Refund\ SK{]}}.
This is usually not preferable:

\begin{itemize}
\tightlist
\item
  \texttt{{[}Refund\ SK{]}} is high cardinality, increasing model size
  through the \texttt{\textquotesingle{}Refund\ ID\textquotesingle{}}
  dimension and relationship columns.
\item
  \texttt{{[}Refund\ SK{]}} has no natural business meaning and creates
  an additional learning burden for users.
\end{itemize}

A better approach is often to create \texttt{{[}Refund\ number{]}}, a
sequence number for each refund action within a sales order. The
composite key becomes \texttt{{[}Sales\ order\ number{]}} and
\texttt{{[}Refund\ number{]}}.

\texttt{{[}Refund\ number{]}} can be exposed through
\texttt{\textquotesingle{}Refund\ ID\textquotesingle{}{[}Refund\ number{]}}.
This is a partial ID dimension because it supplies one part of the full
key needed to identify a refund.

This dimension is low cardinality, has minimal model-size impact, and is
natural for users.

From the user's perspective:

\begin{itemize}
\tightlist
\item
  \texttt{{[}Total\ refund\ amount{]}} gives the refund value for the
  transaction.
\item
  Adding
  \texttt{\textquotesingle{}Refund\ ID\textquotesingle{}{[}Refund\ number{]}}
  breaks the sale into separate refund actions.
\end{itemize}

This is natural at every step.

The use of \texttt{{[}Total\ refund\ amount{]}}, rather than the
degenerate
\texttt{\textquotesingle{}Refund\textquotesingle{}{[}Refund\ amount{]}},
is essential. Suppose the user selects the raw value
\texttt{\textquotesingle{}Refund\textquotesingle{}{[}Refund\ amount{]}}
instead. In the absence of a refund primary key, Power BI may display
one row of \texttt{\$50} for the sales order rather than the correct
total of \texttt{\$100}.

This is another example of why degenerate numeric columns are
problematic and why fact tables are better hidden.

\subsection{Displaying sales and refunds
together}\label{displaying-sales-and-refunds-together}

The user may want to see sales and refund business processes
simultaneously. This is not one problem. There are three common
intentions:

\begin{itemize}
\tightlist
\item
  Show the sales and refund processes side by side.
\item
  Show sales transactions with supplementary refund detail.
\item
  Show refund transactions with the original sales amount.
\end{itemize}

Each intention has a different grain and therefore a different filtering
problem.

\subsubsection{Side-by-side visuals}\label{side-by-side-visuals}

Displaying sales and refunds side by side on the same report page is the
simplest case.

Suppose there is a table visual for sales with:

\begin{itemize}
\tightlist
\item
  \texttt{\textquotesingle{}Sales\ ID\textquotesingle{}{[}Sales\ order\ number{]}}
\item
  \texttt{\textquotesingle{}Sales\ product\textquotesingle{}{[}Product\ name{]}}
\item
  \texttt{\textquotesingle{}Sales\ calendar\textquotesingle{}{[}Sales\ date{]}}
\item
  \texttt{{[}Total\ sales\ amount{]}}
\end{itemize}

And a table visual for refunds with:

\begin{itemize}
\tightlist
\item
  \texttt{\textquotesingle{}Sales\ ID\textquotesingle{}{[}Sales\ order\ number{]}}
\item
  \texttt{\textquotesingle{}Sales\ product\textquotesingle{}{[}Product\ name{]}}
\item
  \texttt{\textquotesingle{}Refund\ calendar\textquotesingle{}{[}Refund\ date{]}}
\item
  \texttt{{[}Total\ refund\ amount{]}}
\end{itemize}

These two tables work well because each visual keeps its own
business-process grain.

\textbf{Shared dimensions behave well.}

If there is a slicer on
\texttt{\textquotesingle{}Sales\ product\textquotesingle{}{[}Product\ name{]}},
both visuals filter down to the selected product.

If there is a slicer on
\texttt{\textquotesingle{}Sales\ ID\textquotesingle{}{[}Sales\ order\ number{]}},
the user can look up a specific sales order and retrieve details of both
the sale and any refund.

A user can also click on a row in the sales table and cross-filter to
any refund. This works because
\texttt{\textquotesingle{}Sales\ ID\textquotesingle{}} is a conformed
dimension for both \texttt{\textquotesingle{}Sale\textquotesingle{}} and
\texttt{\textquotesingle{}Refund\textquotesingle{}}.

For the business user, this is extremely convenient. It is one reason
the ID dimension is one of the most powerful dimensions and an indicator
of a good model.

The same experience cannot be achieved if
\texttt{{[}Sales\ order\ number{]}} is used only as a degenerate
dimension inside \texttt{\textquotesingle{}Sale\textquotesingle{}}. If
degenerate dimensions are used, the filtering experience becomes
frustrating because the shared point of interaction is missing.

\textbf{Process-specific dimensions are harder.}

The user experience is seamless when operating on conformed dimensions
such as \texttt{\textquotesingle{}Sales\ ID\textquotesingle{}} and
\texttt{\textquotesingle{}Sales\ product\textquotesingle{}}.

It becomes more complicated with the process-specific calendar
dimensions.

Users expect to slice by time. But neither calendar works perfectly in
every context:

\begin{itemize}
\tightlist
\item
  A slicer on
  \texttt{\textquotesingle{}Refund\ calendar\textquotesingle{}} has no
  effect on the sales table because it does not filter
  \texttt{\textquotesingle{}Sale\textquotesingle{}}.
\item
  A slicer on
  \texttt{\textquotesingle{}Sales\ calendar\textquotesingle{}} correctly
  filters the sales table, but in the refund table, it returns refunds
  for products sold on that date---not refunds that occurred on that
  date.
\end{itemize}

In simple examples, users may interpret this correctly. In real-world
scenarios with multiple dates, this nuance is often difficult to
understand. In the worst case, it silently misleads users.

One solution is to create a role-playing
\texttt{\textquotesingle{}Reporting\ calendar\textquotesingle{}} that
links \texttt{{[}Reporting\ date{]}} to both
\texttt{\textquotesingle{}Sale\textquotesingle{}{[}Sales\ date{]}} and
\texttt{\textquotesingle{}Refund\textquotesingle{}{[}Refund\ date{]}}.

When used as a slicer, both tables return rows for the selected date.
This works because the selection filters the sales and refund facts
directly, and the visuals then filter out combinations without matching
fact rows.

The chain of filtering behind the scenes creates a seamless experience
for the user.

\subsubsection{Sales with refund
details}\label{sales-with-refund-details}

If the user wants a single table of sales transactions with refund
amount, the visual should remain at the sales grain.

Use the sales table columns and add
\texttt{{[}Total\ refund\ amount{]}}.

The table visual then has two measures:

\begin{itemize}
\tightlist
\item
  \texttt{{[}Total\ sales\ amount{]}}
\item
  \texttt{{[}Total\ refund\ amount{]}}
\end{itemize}

The result is:

\begin{itemize}
\tightlist
\item
  All sales transactions display with
  \texttt{{[}Total\ sales\ amount{]}}.
\item
  Rows with refunds show a value in
  \texttt{{[}Total\ refund\ amount{]}}.
\item
  Rows without refunds show a blank refund amount.
\end{itemize}

This works because Power BI displays rows where any measure is
non-blank.

In addition, if a refund exists without a sales amount, the row stands
out with a blank \texttt{{[}Total\ sales\ amount{]}} and a non-blank
\texttt{{[}Total\ refund\ amount{]}}.

However, the user cannot simply add
\texttt{\textquotesingle{}Refund\ calendar\textquotesingle{}{[}Refund\ date{]}}
to this sales-grain visual.

The technical reason is that
\texttt{\textquotesingle{}Refund\ calendar\textquotesingle{}} relates
only to \texttt{\textquotesingle{}Refund\textquotesingle{}}, not
\texttt{\textquotesingle{}Sale\textquotesingle{}}. The business reason
is that one sale can have multiple refunds on different days. It is
therefore not logically correct to add refund date to a visual whose
intent is one row per sale.

If displaying refund dates is necessary, the data engineer can create a
measure such as \texttt{{[}Sales\ refund\ dates{]}} that returns the
unique refund date if there is one, or concatenates distinct dates if
there are several.

\subsubsection{Refunds with sales
details}\label{refunds-with-sales-details}

The reverse scenario does not mirror cleanly.

If the user creates a refund table and adds
\texttt{{[}Total\ sales\ amount{]}}, Power BI may return incorrect
results by cross-joining sales and refunds. This happens because
\texttt{\textquotesingle{}Refund\ calendar\textquotesingle{}{[}Refund\ date{]}}
relates only to \texttt{\textquotesingle{}Refund\textquotesingle{}} and
does not filter \texttt{\textquotesingle{}Sale\textquotesingle{}}. There
is no valid relationship path for \texttt{{[}Total\ sales\ amount{]}} to
work naturally.

Users may expect to drag \texttt{{[}Total\ sales\ amount{]}} into the
refund table and see the original sales amount. Instead, they may get a
cross-join of rows.

This is one of the situations where the default dimensional model is not
intuitive out of the box and requires additional tuning.

There are several solutions while staying within dimensional modelling:

\begin{itemize}
\tightlist
\item
  Define \texttt{{[}Total\ sales\ amount{]}} to behave differently
  depending on context using DAX functions such as \texttt{treatas}.
\item
  Denormalise the sales amount into
  \texttt{\textquotesingle{}Refund\textquotesingle{}} as a value such as
  \texttt{{[}Sales\ amount\ before\ refund{]}}, then expose a specific
  measure for that use case.
\end{itemize}

The first option gives the most seamless experience for the user, but
adds complexity and performance cost to a core business measure. It is
often not advisable to complicate a central measure for a narrow
scenario.

The second option adds a more narrowly defined measure, such as
\texttt{{[}Sales\ amount\ before\ refund{]}}. This is more explicit,
easier to explain, and usually keeps DAX complexity lower overall.

The measure is unlikely to be useful everywhere, but it is useful in the
refund context.

\subsection{Displaying unit records for three or more business
processes}\label{displaying-unit-records-for-three-or-more-business-processes}

The techniques above rely on conformed dimensions and the existence of
fact rows to filter. The ID dimension plays a special role.

This does not scale neatly.

With three or more facts, it becomes difficult to know which
combinations of dimension values still have at least one row in any
fact. Managing multiple ID dimensions is also complicated because some
are conformed across some facts but not others. Many processes also lack
a ready business key.

A practical solution is to control display with measures.

The data engineer can create one display measure per business process:

\begin{itemize}
\tightlist
\item
  \texttt{{[}Display\ sales\ transaction{]}}
\item
  \texttt{{[}Display\ refund\ transaction{]}}
\item
  \texttt{{[}Display\ ...\ transaction{]}}
\end{itemize}

Each measure returns \texttt{1} when there is at least one row in its
fact under the current filter context.

These measures should be placed in a display folder named
\texttt{Display\ unit\ records}.

A refinement is to return a value only when the relevant ID is in scope.
For example, \texttt{{[}Display\ sales\ transaction{]}} may return a
value only when
\texttt{\textquotesingle{}Sales\ ID\textquotesingle{}{[}Sales\ order\ number{]}}
is being used. This can be done with functions such as
\texttt{ISINSCOPE}.

Adding one of these measures to a table visual filters dimension values
to those with a matching fact row. This is non-blank measure filtering.

The user can hide the helper column by renaming it and shrinking its
width to zero.

An alternative is to use the measure as a visual-level filter. This is
convenient but may have a noticeable performance penalty on large facts.

This has a learning curve, but it can be helpful. A model with more than
four or five business processes is already likely to be complicated. For
a user dealing with this level of complexity, the additional complexity
of display measures is minimal. In fact, an explicit menu of display
measures can clarify which business process the user is trying to
display.

\subsection{Cascading filters}\label{cascading-filters}

Cascading filters refer to the behaviour where selecting a value in one
filter narrows the choices in another so that only valid options remain.

For example, if a report page has a slicer for
\texttt{\textquotesingle{}Sales\ product\textquotesingle{}{[}Product\ name{]}}
and another for
\texttt{\textquotesingle{}Sales\ calendar\textquotesingle{}{[}Sales\ date{]}},
then choosing a product should limit the dates to dates where that
product was sold. Choosing a date should limit products to products sold
on that date.

This feels obvious to users, but it is not always automatic.

Relationship filtering requires a chain of relationships from one table
to the other. In this case, there is no direct relationship path from
\texttt{\textquotesingle{}Sales\ product\textquotesingle{}} to
\texttt{\textquotesingle{}Sales\ calendar\textquotesingle{}}. Both
filter \texttt{\textquotesingle{}Sale\textquotesingle{}}, but neither
directly filters the other.

Common-fact filtering and non-blank measure filtering usually apply
inside a single visual where dimension values and measures are evaluated
together. They do not automatically make two slicers filter each other.

The solution is to use visual-level filtering. For example, a measure
such as \texttt{{[}Has\ sales{]}} can return \texttt{1} if there is at
least one row in \texttt{\textquotesingle{}Sale\textquotesingle{}} under
the current filter context. The slicer can then use
\texttt{{[}Has\ sales{]}} as a visual-level filter.

A simple measure might be conceptually:

\begin{verbatim}
Has sales =
if (
    countrows ( 'Sale' ) > 0,
    1
)
\end{verbatim}

The slicer for
\texttt{\textquotesingle{}Sales\ product\textquotesingle{}{[}Product\ name{]}}
can then be filtered to values where \texttt{{[}Has\ sales{]}} is not
blank. The same can be done for
\texttt{\textquotesingle{}Sales\ calendar\textquotesingle{}{[}Sales\ date{]}}.

The solution relies on a common fact table and a measure that passes
information back to the dimension against the default filter direction.

This is a case where the dimension is on the receiving end of the
button-and-effect relationship. The data engineer can support this
without using bidirectional relationships across the whole model.

\subsection{Aggregating dimension
values}\label{aggregating-dimension-values}

Cascading filters aim to filter one dimension by another. Sometimes the
goal is to aggregate dimension values by another dimension.

Following the sales and refunds example, suppose the user wants to
display a concatenated list of distinct products refunded on each refund
date.

If \texttt{{[}Product\ name{]}} were already included in
\texttt{\textquotesingle{}Refund\textquotesingle{}} as a raw value, the
measure could take the distinct values in
\texttt{\textquotesingle{}Refund\textquotesingle{}{[}Product\ name{]}}
and use \texttt{CONCATENATEX}.

However, this approach is often not available or desirable.

For example:

\begin{itemize}
\tightlist
\item
  \texttt{\textquotesingle{}Sales\ product\textquotesingle{}{[}Product\ name{]}}
  may be linked to \texttt{\textquotesingle{}Refund\textquotesingle{}}
  through \texttt{{[}Product\ SK{]}}.
\item
  The user may want to aggregate another column from
  \texttt{\textquotesingle{}Sales\ product\textquotesingle{}}, such as
  \texttt{{[}Product\ type{]}}.
\item
  It is not practical to include all such columns in the fact table.
\end{itemize}

An inexperienced data engineer might denormalise the required column
into the fact table, or even precompute string concatenation in the data
layer so the column can be used like a degenerate dimension.

Sometimes this is pragmatic. In general, it is not desirable. It is
heavy work to modify a large fact table for every dimension column that
users may want to aggregate. Pre-aggregation is also limited because it
cannot respond flexibly to user filter context.

The better solution is to notice that the information is already in the
model. It only needs to be retrieved under the right filter context.

Measure-defined filtering provides the answer.

The DAX pattern is to concatenate values from
\texttt{\textquotesingle{}Sales\ product\textquotesingle{}{[}Product\ name{]}}
while using \texttt{calculate} and \texttt{crossfilter} to temporarily
allow \texttt{\textquotesingle{}Refund\ calendar\textquotesingle{}} to
filter \texttt{\textquotesingle{}Sales\ product\textquotesingle{}}.

Conceptually:

\begin{verbatim}
Refunded products =
calculate (
    concatenatex (
        values ( 'Sales product'[Product name] ),
        'Sales product'[Product name],
        ", "
    ),
    crossfilter ( 'Sales product'[Product SK], 'Refund'[Product SK], both )
)
\end{verbatim}

The exact DAX depends on the model, but the principle is the same:
temporarily change the filter path for this calculation only.

This avoids compromising the whole model with bidirectional
relationships or denormalised columns. The model remains clean, and the
special behaviour is contained inside the measure that needs it.

\subsection{Filter by having}\label{filter-by-having}

Sometimes the user does not want to filter a dimension by the value it
currently has. They want to filter by whether the entity has ever had an
attribute.

This is filter by having.

Filter by having means: filter the entity by whether it has, or once
had, a matching attribute, rather than filtering only the rows where the
attribute currently appears.

Consider an organisation's HR system. The model contains:

\begin{itemize}
\tightlist
\item
  A dimension \texttt{\textquotesingle{}Employee\textquotesingle{}} that
  stores details such as name, email, date of birth, and role.
\item
  \texttt{\textquotesingle{}Employee\textquotesingle{}} is Type II, so
  it stores historical changes to employee attributes.
\item
  The key columns are \texttt{{[}Employee\ ID{]}},
  \texttt{{[}Start\ datetime{]}}, and \texttt{{[}End\ datetime{]}}.
\item
  \texttt{{[}Employee\ SK{]}} is used as the surrogate key for Power BI
  relationships.
\item
  An end-of-month fact table
  \texttt{\textquotesingle{}Employee\ end\ of\ month\textquotesingle{}}
  stores the employee's attributes at the end of each month.
\end{itemize}

\textbf{Example structure of
\texttt{\textquotesingle{}Employee\textquotesingle{}}}

{\def\LTcaptype{none} % do not increment counter
\begin{longtable}[]{@{}
  >{\raggedleft\arraybackslash}p{(\linewidth - 10\tabcolsep) * \real{0.2105}}
  >{\raggedright\arraybackslash}p{(\linewidth - 10\tabcolsep) * \real{0.1579}}
  >{\raggedright\arraybackslash}p{(\linewidth - 10\tabcolsep) * \real{0.1579}}
  >{\raggedright\arraybackslash}p{(\linewidth - 10\tabcolsep) * \real{0.1579}}
  >{\raggedright\arraybackslash}p{(\linewidth - 10\tabcolsep) * \real{0.1579}}
  >{\raggedright\arraybackslash}p{(\linewidth - 10\tabcolsep) * \real{0.1579}}@{}}
\toprule\noalign{}
\begin{minipage}[b]{\linewidth}\raggedleft
Employee SK
\end{minipage} & \begin{minipage}[b]{\linewidth}\raggedright
Employee ID
\end{minipage} & \begin{minipage}[b]{\linewidth}\raggedright
Start date
\end{minipage} & \begin{minipage}[b]{\linewidth}\raggedright
End date
\end{minipage} & \begin{minipage}[b]{\linewidth}\raggedright
Employee name
\end{minipage} & \begin{minipage}[b]{\linewidth}\raggedright
Role
\end{minipage} \\
\midrule\noalign{}
\endhead
\bottomrule\noalign{}
\endlastfoot
1 & E1001 & 2023-01-01 & 2023-06-30 & Alice Chen & Data analyst \\
2 & E1001 & 2023-07-01 & 9999-12-31 & Alice Chen & Senior analyst \\
3 & E1002 & 2023-01-01 & 9999-12-31 & Ben Smith & Finance officer \\
\end{longtable}
}

\textbf{Example structure of
\texttt{\textquotesingle{}Employee\ end\ of\ month\textquotesingle{}}}

{\def\LTcaptype{none} % do not increment counter
\begin{longtable}[]{@{}rlr@{}}
\toprule\noalign{}
Employee SK & Period end date & FTE \\
\midrule\noalign{}
\endhead
\bottomrule\noalign{}
\endlastfoot
1 & 2023-06-30 & 1.0 \\
2 & 2023-07-31 & 1.0 \\
2 & 2023-08-31 & 1.0 \\
3 & 2023-08-31 & 1.0 \\
\end{longtable}
}

\texttt{\textquotesingle{}Employee\textquotesingle{}} is a dimension
filtering the fact
\texttt{\textquotesingle{}Employee\ end\ of\ month\textquotesingle{}}.

{\def\LTcaptype{none} % do not increment counter
\begin{longtable}[]{@{}
  >{\raggedright\arraybackslash}p{(\linewidth - 4\tabcolsep) * \real{0.3333}}
  >{\raggedright\arraybackslash}p{(\linewidth - 4\tabcolsep) * \real{0.3333}}
  >{\raggedright\arraybackslash}p{(\linewidth - 4\tabcolsep) * \real{0.3333}}@{}}
\toprule\noalign{}
\begin{minipage}[b]{\linewidth}\raggedright
Dimension
\end{minipage} & \begin{minipage}[b]{\linewidth}\raggedright
\texttt{\textquotesingle{}Employee\ end\ of\ month\textquotesingle{}}
\end{minipage} & \begin{minipage}[b]{\linewidth}\raggedright
Meaning
\end{minipage} \\
\midrule\noalign{}
\endhead
\bottomrule\noalign{}
\endlastfoot
\texttt{\textquotesingle{}Employee\textquotesingle{}} & 1 → * via
\texttt{{[}Employee\ SK{]}} & Point-in-time employee attributes \\
\end{longtable}
}

In this model, filtering
\texttt{\textquotesingle{}Employee\textquotesingle{}{[}Role{]}} filters
the fact through \texttt{{[}Employee\ SK{]}}. This is correct for
point-in-time reporting. It returns the employee-month rows where the
selected Type II employee row is active.

That is not always the desired search behaviour.

The user wants to search an employee's history by personal attributes
such as name or role.

Since \texttt{\textquotesingle{}Employee\textquotesingle{}} is Type II,
a direct search returns only the end-of-period states that match the
selected point-in-time attribute values. If the user filters
\texttt{\textquotesingle{}Employee\textquotesingle{}{[}Role{]}} to
\texttt{Data\ analyst}, the model returns only rows where the employee
was a data analyst. In the example above, it returns Alice's June row,
but not Alice's later history as a senior analyst.

That is correct Type II behaviour, but it is not the desired search
behaviour.

The desired behaviour is different. The user wants to find employees who
have ever had \texttt{Data\ analyst} as a role, then return their full
history.

New data engineers often try to solve this by restructuring the fact
table. This leads to complex code and confusing artefacts.

A simpler solution is to note that \texttt{{[}Employee\ ID{]}} is the
identifier of the "true" employee across time, and add a search
dimension keyed on this column.

Create \texttt{\textquotesingle{}Employee\ search\textquotesingle{}} as
a duplicate \texttt{\textquotesingle{}Employee\textquotesingle{}}. It
relates to \texttt{\textquotesingle{}Employee\textquotesingle{}} on
\texttt{{[}Employee\ ID{]}} in a snowflake configuration. The
relationship is many-to-many and single direction.

\textbf{Example structure of
\texttt{\textquotesingle{}Employee\ search\textquotesingle{}}}

{\def\LTcaptype{none} % do not increment counter
\begin{longtable}[]{@{}lll@{}}
\toprule\noalign{}
Employee ID & Employee name & Role \\
\midrule\noalign{}
\endhead
\bottomrule\noalign{}
\endlastfoot
E1001 & Alice Chen & Data analyst \\
E1001 & Alice Chen & Senior analyst \\
E1002 & Ben Smith & Finance officer \\
\end{longtable}
}

The revised model looks like this.

{\def\LTcaptype{none} % do not increment counter
\begin{longtable}[]{@{}
  >{\raggedright\arraybackslash}p{(\linewidth - 6\tabcolsep) * \real{0.2500}}
  >{\raggedright\arraybackslash}p{(\linewidth - 6\tabcolsep) * \real{0.2500}}
  >{\raggedright\arraybackslash}p{(\linewidth - 6\tabcolsep) * \real{0.2500}}
  >{\raggedright\arraybackslash}p{(\linewidth - 6\tabcolsep) * \real{0.2500}}@{}}
\toprule\noalign{}
\begin{minipage}[b]{\linewidth}\raggedright
Dimension
\end{minipage} & \begin{minipage}[b]{\linewidth}\raggedright
\texttt{\textquotesingle{}Employee\textquotesingle{}}
\end{minipage} & \begin{minipage}[b]{\linewidth}\raggedright
\texttt{\textquotesingle{}Employee\ end\ of\ month\textquotesingle{}}
\end{minipage} & \begin{minipage}[b]{\linewidth}\raggedright
Meaning
\end{minipage} \\
\midrule\noalign{}
\endhead
\bottomrule\noalign{}
\endlastfoot
\texttt{\textquotesingle{}Employee\ search\textquotesingle{}} & * → *
via \texttt{{[}Employee\ ID{]}} & & Search by any historical employee
attribute \\
\texttt{\textquotesingle{}Employee\textquotesingle{}} & & 1 → * via
\texttt{{[}Employee\ SK{]}} & Point-in-time employee attributes \\
\end{longtable}
}

Using \texttt{\textquotesingle{}Employee\ search\textquotesingle{}}, the
user can look up any historical attribute of the employee, retrieve
\texttt{{[}Employee\ ID{]}}, and return all rows in
\texttt{\textquotesingle{}Employee\textquotesingle{}} that match that
ID.

When the user selects \texttt{Data\ analyst} in
\texttt{\textquotesingle{}Employee\ search\textquotesingle{}{[}Role{]}},
the model first finds \texttt{{[}Employee\ ID{]}\ =\ E1001}. It then
filters \texttt{\textquotesingle{}Employee\textquotesingle{}} to all
rows for \texttt{E1001}, not only the row where the role was
\texttt{Data\ analyst}.

The user searched by a historical attribute, but the model returned the
full entity history.

Under this view,
\texttt{\textquotesingle{}Employee\ search\textquotesingle{}} and
\texttt{\textquotesingle{}Employee\textquotesingle{}} work together as
one logical dimension that acts as a search interface for the model.

This pattern generalises. The data engineer can add
\texttt{\textquotesingle{}Compliance\ search\textquotesingle{}},
\texttt{\textquotesingle{}Qualification\ search\textquotesingle{}}, or
other derived attribute search dimensions as additional branches of the
snowflake.

From the user's perspective, these are search tables that allow
searching employees by their history.

Unlike tampering with fact tables, this solution is non-destructive. The
fact table has not changed. The ordinary point-in-time path remains
intact. The search dimension adds a new interface path for a different
user intention.

\subsection{Dynamic Type I}\label{dynamic-type-i}

A common workforce reporting request sounds like this: \emph{"When I
choose August, show me the last 12 months using the organisation
structure as at August."}

This sounds simple, but it is not the default behaviour of a Type II
dimension.

Type II dimensions show an entity at its point-in-time attribute. Type I
dimensions show the latest attribute value. Neither meets the user
intention. What the user would like is Dynamic Type I---latest attribute
value as at a user-selected time.

Consider employees and organisation units. The fact table
\texttt{\textquotesingle{}Employee\ end\ of\ month\textquotesingle{}}
returns employee metrics at the end of each month. The dimension
\texttt{\textquotesingle{}Organisation\ unit\textquotesingle{}} is Type
II because organisation structures change over time.

\textbf{Example structure of
\texttt{\textquotesingle{}Organisation\ unit\textquotesingle{}}}

{\def\LTcaptype{none} % do not increment counter
\begin{longtable}[]{@{}
  >{\raggedleft\arraybackslash}p{(\linewidth - 10\tabcolsep) * \real{0.2105}}
  >{\raggedright\arraybackslash}p{(\linewidth - 10\tabcolsep) * \real{0.1579}}
  >{\raggedright\arraybackslash}p{(\linewidth - 10\tabcolsep) * \real{0.1579}}
  >{\raggedright\arraybackslash}p{(\linewidth - 10\tabcolsep) * \real{0.1579}}
  >{\raggedright\arraybackslash}p{(\linewidth - 10\tabcolsep) * \real{0.1579}}
  >{\raggedright\arraybackslash}p{(\linewidth - 10\tabcolsep) * \real{0.1579}}@{}}
\toprule\noalign{}
\begin{minipage}[b]{\linewidth}\raggedleft
Team unit SK
\end{minipage} & \begin{minipage}[b]{\linewidth}\raggedright
Team unit ID
\end{minipage} & \begin{minipage}[b]{\linewidth}\raggedright
Start date
\end{minipage} & \begin{minipage}[b]{\linewidth}\raggedright
End date
\end{minipage} & \begin{minipage}[b]{\linewidth}\raggedright
Team name
\end{minipage} & \begin{minipage}[b]{\linewidth}\raggedright
Group name
\end{minipage} \\
\midrule\noalign{}
\endhead
\bottomrule\noalign{}
\endlastfoot
1 & T1001 & 2023-01-01 & 2023-06-30 & Data engineering & Group A \\
2 & T1001 & 2023-07-01 & 9999-12-31 & Data engineering & Group B \\
3 & T1002 & 2023-01-01 & 9999-12-31 & Finance operations & Group A \\
\end{longtable}
}

\textbf{Example structure of
\texttt{\textquotesingle{}Employee\ end\ of\ month\textquotesingle{}}}

{\def\LTcaptype{none} % do not increment counter
\begin{longtable}[]{@{}rllr@{}}
\toprule\noalign{}
Team unit SK & Team unit ID & Period end date & Commencements \\
\midrule\noalign{}
\endhead
\bottomrule\noalign{}
\endlastfoot
1 & T1001 & 2023-06-30 & 2 \\
2 & T1001 & 2023-07-31 & 1 \\
2 & T1001 & 2023-08-31 & 1 \\
3 & T1002 & 2023-08-31 & 1 \\
\end{longtable}
}

In the ordinary Type II model,
\texttt{\textquotesingle{}Employee\ end\ of\ month\textquotesingle{}}
relates to
\texttt{\textquotesingle{}Organisation\ unit\textquotesingle{}} through
\texttt{{[}Team\ unit\ SK{]}}.

{\def\LTcaptype{none} % do not increment counter
\begin{longtable}[]{@{}
  >{\raggedright\arraybackslash}p{(\linewidth - 4\tabcolsep) * \real{0.3333}}
  >{\raggedright\arraybackslash}p{(\linewidth - 4\tabcolsep) * \real{0.3333}}
  >{\raggedright\arraybackslash}p{(\linewidth - 4\tabcolsep) * \real{0.3333}}@{}}
\toprule\noalign{}
\begin{minipage}[b]{\linewidth}\raggedright
Dimension
\end{minipage} & \begin{minipage}[b]{\linewidth}\raggedright
\texttt{\textquotesingle{}Employee\ end\ of\ month\textquotesingle{}}
\end{minipage} & \begin{minipage}[b]{\linewidth}\raggedright
Meaning
\end{minipage} \\
\midrule\noalign{}
\endhead
\bottomrule\noalign{}
\endlastfoot
\texttt{\textquotesingle{}Organisation\ unit\textquotesingle{}} & 1 → *
via \texttt{{[}Team\ unit\ SK{]}} & Point-in-time organisation
hierarchy \\
\end{longtable}
}

In the example, Team \texttt{T1001} belonged to \texttt{Group\ A} in
June, then moved to \texttt{Group\ B} in July.

If the user selects August and asks for commencements in the last twelve
months, they may expect all activity for Team \texttt{T1001} to appear
under \texttt{Group\ B}, because that is the organisation structure as
at August. The correct answer is a total of 4 commencements.

The ordinary Type II relationship does not do that. It reports each fact
row using the organisation structure that was valid at the time of the
row. June remains under \texttt{Group\ A}; July and August appear under
\texttt{Group\ B}. This returns 2 commencements, the incorrect result.

Thus, Dynamic Type I changes the lens. It says: use the selected
reporting date to decide which organisation-unit row is current, then
view the relevant historical facts through that selected-date structure.

\subsubsection{Solving Dynamic Type}\label{solving-dynamic-type}

The problem is notoriously complex to solve by amending the fact table.
It requires associating the organisation unit's full history with each
fact row. That approach is complex, fragile, and expensive.

The key is to recognise that the model already has the history in
\texttt{\textquotesingle{}Organisation\ unit\textquotesingle{}}. The
information is already in the model, and what\textquotesingle s missing
is to surface it. The better solution is to surface it at query time.

There are two steps.

First, broaden the relationship for evaluation.

The problem with the existing relationship is that it is too
"narrow"---locking the dimension into a smaller set of data than
intuitive to business sense. The first step is to "unlock" this
relationship between the organisation unit to the activity rows that are
relevant to it.

This means relating each fact row not only to a single point-in-time row
in \texttt{\textquotesingle{}Organisation\ unit\textquotesingle{}}, but
to all rows for that \texttt{{[}Team\ unit\ ID{]}}. This requires a
many-to-many, single-direction relationship on
\texttt{{[}Team\ unit\ ID{]}}. Each fact row now has the full history of
its team available during calculation.

The relationship becomes:

{\def\LTcaptype{none} % do not increment counter
\begin{longtable}[]{@{}
  >{\raggedright\arraybackslash}p{(\linewidth - 4\tabcolsep) * \real{0.3333}}
  >{\raggedright\arraybackslash}p{(\linewidth - 4\tabcolsep) * \real{0.3333}}
  >{\raggedright\arraybackslash}p{(\linewidth - 4\tabcolsep) * \real{0.3333}}@{}}
\toprule\noalign{}
\begin{minipage}[b]{\linewidth}\raggedright
Dimension
\end{minipage} & \begin{minipage}[b]{\linewidth}\raggedright
\texttt{\textquotesingle{}Employee\ end\ of\ month\textquotesingle{}}
\end{minipage} & \begin{minipage}[b]{\linewidth}\raggedright
Meaning
\end{minipage} \\
\midrule\noalign{}
\endhead
\bottomrule\noalign{}
\endlastfoot
\texttt{\textquotesingle{}Organisation\ unit\textquotesingle{}} & * → *
via \texttt{{[}Team\ unit\ ID{]}} & Relates all rows based on the
persistent team unit identity \\
\end{longtable}
}

Second, select the valid row at query time.

The consequence of the relationship is that, though the data is
available in the filter context, too many rows are selected for display.
We now need to narrow this again.

In the measure, we take the latest date in the current context and
filter \texttt{\textquotesingle{}Organisation\ unit\textquotesingle{}}
to the row whose validity covers that date:

\begin{verbatim}
[Start date] <= selected date < [End date]
\end{verbatim}

For a given \texttt{{[}Team\ unit\ ID{]}}, this should return one row
because the primary key is \texttt{{[}Team\ unit\ ID{]}} and
\texttt{{[}Start\ date{]}}.

Under this filter, the table behaves like a one-to-many join for that
evaluation. The measure returns values only for the valid
organisation-unit row. Other rows go blank and drop out through
non-blank measure filtering.

For August 2023, \texttt{T1001} resolves to the row where
\texttt{{[}Group\ name{]}\ =\ Group\ B}. The measure can then report
prior-period activity using the August 2023 organisation view.

This creates a problem. The original
\texttt{\textquotesingle{}Organisation\ unit\textquotesingle{}} no
longer behaves as before. This could be problematic. If the original
behaviour needs to be preserved, there are two broad implementation
options.

\subsubsection{Option A---Keep one dimension and switch in
measures}\label{option-akeep-one-dimension-and-switch-in-measures}

The first option is to do the switching in the measure.

In this scenario, we keep the normal one-to-many, single-direction
relationship on \texttt{{[}Team\ unit\ SK{]}}, and then add a
many-to-many, single-direction relationship on
\texttt{{[}Team\ unit\ ID{]}}. It remains dormant by default.

In measures that need dynamic Type I, deactivate the
\texttt{{[}Team\ unit\ SK{]}} path and allow the
\texttt{{[}Team\ unit\ ID{]}} path to apply for the calculation. Then
filter to the valid row as above.

In this implementation,
\texttt{\textquotesingle{}Organisation\ unit\textquotesingle{}} works as
Type II by default, and the user chooses the dynamic Type I perspective
through measures.

\subsubsection{Option B---Add a new
dimension}\label{option-badd-a-new-dimension}

The second option is to switch through dimensions. This is done by
duplicating
\texttt{\textquotesingle{}Organisation\ unit\textquotesingle{}} as
\texttt{\textquotesingle{}End\ of\ period\ organisation\ unit\textquotesingle{}},
and adding \texttt{{[}Team\ unit\ ID{]}} to
\texttt{\textquotesingle{}Employee\ end\ of\ month\textquotesingle{}}.

The fact now carries:

\begin{itemize}
\tightlist
\item
  \texttt{{[}Team\ unit\ SK{]}} for the normal relationship to
  \texttt{\textquotesingle{}Organisation\ unit\textquotesingle{}}
\item
  \texttt{{[}Team\ unit\ ID{]}} for the many-to-many relationship to
  \texttt{\textquotesingle{}End\ of\ period\ organisation\ unit\textquotesingle{}}
\end{itemize}

The revised evaluation model is therefore:

{\def\LTcaptype{none} % do not increment counter
\begin{longtable}[]{@{}
  >{\raggedright\arraybackslash}p{(\linewidth - 4\tabcolsep) * \real{0.3333}}
  >{\raggedright\arraybackslash}p{(\linewidth - 4\tabcolsep) * \real{0.3333}}
  >{\raggedright\arraybackslash}p{(\linewidth - 4\tabcolsep) * \real{0.3333}}@{}}
\toprule\noalign{}
\begin{minipage}[b]{\linewidth}\raggedright
Dimension
\end{minipage} & \begin{minipage}[b]{\linewidth}\raggedright
\texttt{\textquotesingle{}Employee\ end\ of\ month\textquotesingle{}}
\end{minipage} & \begin{minipage}[b]{\linewidth}\raggedright
Meaning
\end{minipage} \\
\midrule\noalign{}
\endhead
\bottomrule\noalign{}
\endlastfoot
\texttt{\textquotesingle{}Organisation\ unit\textquotesingle{}} & 1 → *
via \texttt{{[}Team\ unit\ SK{]}} & Ordinary Type II path \\
\texttt{\textquotesingle{}End\ of\ period\ organisation\ unit\textquotesingle{}}
& * → * via \texttt{{[}Team\ unit\ ID{]}} & Dynamic Type I evaluation
path \\
\end{longtable}
}

Measures that need dynamic Type I behaviour can be wrapped.

For example, rename \texttt{{[}Employee\ commencements{]}} to
\texttt{{[}\_Employee\ commencements{]}}. Then define
\texttt{{[}Employee\ commencements{]}} so that it evaluates
\texttt{{[}\_Employee\ commencements{]}} inside \texttt{calculate},
filtering
\texttt{\textquotesingle{}End\ of\ period\ organisation\ unit\textquotesingle{}}
to the valid row as at the latest date in context.

In this option, the user chooses perspective by choosing the dimension:

\begin{itemize}
\tightlist
\item
  \texttt{\textquotesingle{}Organisation\ unit\textquotesingle{}} gives
  Type II.
\item
  \texttt{\textquotesingle{}End\ of\ period\ organisation\ unit\textquotesingle{}}
  gives dynamic Type I.
\end{itemize}

Under both options, the original dimension keeps its normal behaviour.
Existing measures continue to work. Both options are non-destructive.

The distinction is whether the user chooses Type II or dynamic Type I
through different dimensions, or through different measures.

The important point is that Dynamic Type I does not rewrite history in
the fact table. It changes the reporting lens at query time.

\section{Buttons and effects}\label{buttons-and-effects}

The take-away from above is not exhaustive list of filtering techniques.
That would be a book in itself.

Their purpose is to illustrate a philosophy of seeing Power BI as an
interface that users touch to produce effects.

In this philosophy, every table tends to have one of two roles:

\begin{itemize}
\tightlist
\item
  Tables that define the interface as buttons to click, usually
  dimensions
\item
  Tables that define the content that drives the effect, usually facts
\end{itemize}

A good model keeps these roles separate.

Every table should have one clear purpose. Tables should not mix
interface and content unless there is a deliberate reason. This is why
an indicator of quality is that fact tables are hidden.

The first scenarios showed how dimensions present information from facts
through the difficult case of displaying transaction records. They
illustrate that fact tables, which are the natural focus for new data
engineers, are not the artefacts to expose to users.

Trying to expose fact tables and degenerate dimensions causes issues.
Dimensions should act as the entry point to the data.

The same philosophy explains why degenerate dimensions and bidirectional
relationships often frustrate users. They blur the line between tables
that act as buttons and tables that create effects.

The examples of cascading filters and aggregating dimension values show
that the data engineer can handle cases where a dimension is on the
receiving end of a button-and-effect relationship without compromising
the model with degenerate dimensions or broad bidirectional
relationships.

The final examples generalise the idea of a dimension to a set of
tables. A dimension is not simply a lookup table. If the user sees it as
intuitive and gets the desired outcome unambiguously and quickly, that
is what matters.

In this view, a dimension's chief purpose is to be the user's interface
to the model.

This is the philosophy of separation of concerns from software design.
It maintains a clear distinction between the interface layer and the
content layer.

Power BI supports this through the five filtering mechanisms described
at the start of the chapter. This perspective pushes the data engineer
to see a data model as software and to design a clean interface that is
business-centric and intuitive for the user.

\section{Key ideas}\label{key-ideas-4}

\begin{quote}
\textbf{Key ideas.}

Filtering design is what makes a Power BI model feel intuitive---or
frustrating.

A self-service model is not a puzzle for the user to solve. The correct
answer should not merely be reachable; it should be reachable through
likely paths of interaction.

Filtering is the mechanism that turns model structure into user
experience.

Dimensions usually provide the buttons: the things users select, slice,
search, group, and compare.

Facts and measures provide the effects: the activity, state, and values
that respond.

Power BI filtering commonly occurs through relationships, common facts,
non-blank measures, visual-level filters, and measure-defined filtering.

Degenerate dimensions and broad bidirectional relationships often blur
the distinction between buttons and effects.

Good filtering design anticipates how users will interact with the model
and makes sensible behaviour occur naturally.
\end{quote}

\cleardoublepage

\chapter{Designing
measures}\label{docs-presenting-insights-designing-measures}

\section{Measures as compression and
re-expression}\label{measures-as-compression-and-re-expression}

A fact table may contain thousands, millions, or billions of rows. A
measure compresses those rows into a business expression, then unpacks
that expression again in user context.

A good measure therefore does not merely calculate correctly. It states
clearly what business meaning is being compressed, and it behaves
predictably as that meaning is re-expressed through the user's selected
dimensions.

Power BI's implicit calculations can undermine this design. They
encourage users to treat raw columns as ready-made measures, even when
those columns do not express a clear business answer. For this reason,
\textbf{Discourage implicit measures} should usually be turned on.

This chapter covers three areas:

\begin{itemize}
\tightlist
\item
  signs of good measures;
\item
  interface roles of measures;
\item
  technical patterns for building measures.
\end{itemize}

The next chapter,
\hyperref[docs-presenting-insights-measure-of-measures]{Measure of
measures}, introduces a pattern for managing structured families of
measures.

\section{Signs of good measures}\label{signs-of-good-measures}

There are two signs of good measures:

\begin{itemize}
\tightlist
\item
  business centricity;
\item
  technical simplicity.
\end{itemize}

\subsection{Business centricity}\label{business-centricity}

A measure is not merely a calculation. It is a business expression made
available to the user for unpacking.

Every aspect of a measure should reflect business meaning. This includes
its definition, name, description, display folder, and placement
alongside related measures.

\subsubsection{Alignment to business
reality}\label{alignment-to-business-reality}

A measure should align with real-world events.

A common mistake is to count a database key and assume that the result
is a meaningful business metric.

For example, a data engineer might define
\texttt{{[}Inspection\ count{]}} as a distinct count of
\texttt{{[}Inspection\ ID{]}} in the
\texttt{\textquotesingle{}Inspection\textquotesingle{}} fact table. That
may be fine if \texttt{{[}Inspection\ ID{]}} corresponds to a real
inspection. But in many systems, \texttt{{[}Inspection\ ID{]}} is only a
system-generated record key. It may exist for retrieval, workflow, or
database convenience.

If the system later changes so that one real inspection is split into
ten records, the measure changes even though inspection effort has not
changed. Conversely, if ten records are merged into one, the measure
falls even though no real-world activity has disappeared.

That is a measure misaligned with reality.

A better measure quantifies something with real-world business meaning.
Depending on the scenario, the data engineer may measure:

\begin{itemize}
\tightlist
\item
  inspection hours;
\item
  entities inspected.
\end{itemize}

The right choice depends on the business question.

Even \texttt{inspection} itself can be unreliable as the basis for a
measure. If new inspection types are introduced that are much lighter or
heavier than the old ones, then \texttt{{[}Number\ of\ inspections{]}}
may drift in meaning over time.

By contrast, \texttt{{[}Total\ inspection\ hours{]}} remains close to
concrete reality. An hour is still an hour, even when new inspections
get added.

This does not mean \texttt{{[}Number\ of\ inspections{]}} is always
wrong. The point is that the data engineer should understand what the
measure is really compressing, and where possible, stay closer to
concrete reality than abstraction.

\subsubsection{Business specificity}\label{business-specificity}

Measures should be specific to the business question they answer.

A common misstep is to create a single generic measure and expect users
to combine it with filters to answer every question.

For example, the model may contain a generic measure:

\begin{verbatim}
Employee count =
distinctcount ( 'Employee end of month'[Employee ID] )
\end{verbatim}

The data engineer may then expect users to filter manually for active
employees, separated employees, employees on leave, current employees,
and other cases.

This turns the model into a puzzle. Users must remember which filters to
apply, which dates matter, and which status values are valid for each
business question.

A better model presents specific measures that answer common business
questions directly.

For example:

\begin{itemize}
\tightlist
\item
  \texttt{{[}Active\ employees\ end\ of\ period{]}}---employees with an
  active employment status at the end of the selected period.
\item
  \texttt{{[}Separated\ employees\ end\ of\ period{]}}---employees with
  a separated employment status at the end of the selected period.
\item
  \texttt{{[}Current\ employees{]}}---the latest value of
  \texttt{{[}Active\ employees\ end\ of\ period{]}}.
\item
  \texttt{{[}Current\ employees\ on\ long\ service\ leave{]}}---current
  employees who are on long service leave.
\end{itemize}

These measures are not merely variants of
\texttt{{[}Employee\ count{]}}. They are specific business expressions
that make common questions directly available.

A technical measure such as \texttt{{[}\_Employee\ count{]}} and
\texttt{{[}\_Employee\ count\ end\ of\ period{]}} may still be useful
internally. If it is not useful to users directly, it should be hidden
to support business-facing measures without cluttering the interface.

\subsubsection{Names}\label{names}

Measure names should be explicit, grammatical, and able to stand alone.

Measures can be dragged anywhere on the canvas. Their names should carry
enough context to remain meaningful outside their original folder or
visual.

Prefer:

\begin{verbatim}
[Active employees end of period]
\end{verbatim}

over:

\begin{verbatim}
[Count]
\end{verbatim}

Avoid non-standard abbreviations. Saving a few characters does not help
the user or the next engineer.

Names should also support search. Users often find measures by typing
keywords in the field list. If the business calls something
\texttt{commencements}, the measure name may include
\texttt{commencements}. Anticipating the user search should be part of
the measure design.

\subsubsection{Descriptions}\label{descriptions}

Descriptions should explain business meaning.

The description should not simply restate the DAX. Users are not asking
for implementation trivia, nor should they be expected to accept complex
logic without explanation. They want to know what the measure means and
whether they can trust it.

For simple measures, a description may be short:

\begin{verbatim}
Counts employees with an active employment status at the end of the selected reporting period.
\end{verbatim}

For more complex measures, near-pseudocode in business language may be
appropriate:

\begin{verbatim}
Counts employees who were active at the end of the selected period and whose leave status indicates long service leave. Uses the latest reporting period in the current user context.
\end{verbatim}

The description is part of the interface. It is often the first
explanation the user sees, and an opportunity to build trust.

\subsubsection{Display folders}\label{display-folders}

Display folders help measures appear in business-facing groups.

Measures can be grouped by business process, such as:

\begin{itemize}
\tightlist
\item
  \texttt{Workforce}
\item
  \texttt{Inspections}
\item
  \texttt{Sales}
\item
  \texttt{Refunds}
\end{itemize}

Technical or report-specific measures can be placed in folders such as:

\begin{itemize}
\tightlist
\item
  \texttt{Dashboard}
\item
  \texttt{Display\ unit\ records}
\end{itemize}

\subsubsection{Proximity to related
measures}\label{proximity-to-related-measures}

The meaning of a measure is influenced by nearby measures.

Related measures should be named so that they appear near each other.

For example:

\begin{itemize}
\tightlist
\item
  \texttt{{[}Active\ employees\ end\ of\ period{]}}
\item
  \texttt{{[}Separated\ employees\ end\ of\ period{]}}
\item
  \texttt{{[}Current\ employees{]}}
\item
  \texttt{{[}Current\ employees\ on\ long\ service\ leave{]}}
\end{itemize}

The presence of \texttt{{[}Separated\ employees\ end\ of\ period{]}}
reinforces the meaning of
\texttt{{[}Active\ employees\ end\ of\ period{]}}. The presence of
\texttt{{[}Current\ employees{]}} clarifies that
\texttt{{[}Active\ employees\ end\ of\ period{]}} is a point-in-time
measure.

These compare-and-contrast relationships help users form a clearer
mental model of the data.

There are two practical forms of proximity:

\begin{itemize}
\tightlist
\item
  alphabetical placement within a display folder;
\item
  keyword-based lookup in the field list.
\end{itemize}

A data engineer can take advantage of the proximity effect during
measure naming and placement.

\subsection{Technical simplicity}\label{technical-simplicity}

Measures should be technically simple.

This does not mean every measure must be trivial. It means complex
business logic should usually be pushed upstream into the data pipeline.

A measure should aggregate, filter, compare, or display information that
is already well prepared.

There are two important foundations for technical simplicity:

\begin{itemize}
\tightlist
\item
  precomputing complex information;
\item
  preparing binary flags.
\end{itemize}

\subsubsection{Precomputing complex
information}\label{precomputing-complex-information}

Complex business logic should usually be computed in the data layer.

This frees measures from carrying too much responsibility. A measure
should not have to reconstruct a complicated business state every time a
user clicks a slicer.

For example, suppose \texttt{{[}Is\ on\ long\ service\ leave{]}}
requires analysing leave transactions, effective dates, overlapping
periods, and employee status. That logic should usually be calculated in
the pipeline and stored in a fragment such as \texttt{HR.EmployeeLeave}
or a related reference table.

Then the measure can remain simple:

\begin{verbatim}
Current employees on long service leave =
calculate (
    [Current employees],
    keepfilters ( 'Employee leave'[Is on long service leave] )
)
\end{verbatim}

The complexity has been moved to the data layer, where it can be tested
and reused.

Storytelling dimensions are an important example of this principle. When
complex detail has been reduced into business-facing categories, paths,
summaries, or special cases, measures can refer to those prepared
structures instead of repeating the interpretation in DAX.

Precomputation can be overdone. If the data engineer precomputes too
much, the model may lose interactivity. A pre-aggregated string may no
longer respond correctly to filters. An average stored at the wrong
grain may produce wrong results when aggregated again.

The rule is not ``precompute everything.'' The rule is:

\begin{quote}
Precompute business interpretation when doing so preserves interactivity
and simplifies the model.
\end{quote}

The data engineer should understand the filtering behaviour of Power BI
well enough to know which logic belongs in the pipeline and which logic
belongs in a measure.

\subsubsection{Preparing binary flags}\label{preparing-binary-flags}

Binary flags are true-or-false columns that express business logic
clearly.

Their most common place is in reference tables and dimensions. For
example, an \texttt{\textquotesingle{}Employee\ leave\textquotesingle{}}
dimension may contain \texttt{{[}Is\ on\ long\ service\ leave{]}} to
indicate whether the employee is on long service leave at the end of the
reporting period.

Binary flags simplify measures.

Instead of writing a measure that filters text values directly:

\begin{verbatim}
Current employees on long service leave =
calculate (
    [Current employees],
    keepfilters ( 'Employee leave'[Leave type] in { "Long service leave", "Extended long service leave" } )
)
\end{verbatim}

The measure can use the prepared binary flag:

\begin{verbatim}
Current employees on long service leave =
calculate (
    [Current employees],
    keepfilters ( 'Employee leave'[Is on long service leave] )
)
\end{verbatim}

The reverse can be expressed with \texttt{not}:

\begin{verbatim}
Current employees not on long service leave =
calculate (
    [Current employees],
    keepfilters ( not 'Employee leave'[Is on long service leave] )
)
\end{verbatim}

Binary flags have two advantages.

First, they are syntactically checked. If a developer mistypes
\texttt{Long\ service\ leave} as text, the formula may still be valid
but return the wrong result. If the developer mistypes the column name
\texttt{{[}Is\ on\ long\ service\ leave{]}}, the formula fails visibly.

Second, binary flags make the business logic readable. The measure says
what it means.

Binary flags also participate cleanly in \texttt{and} and \texttt{or}
logic. For example:

\begin{verbatim}
Current employees on leave or acting =
calculate (
    [Current employees],
    keepfilters (
        'Employee status'[Is on leave]
            || 'Employee status'[Is acting]
    )
)
\end{verbatim}

Using binary flags makes measures faster to write, easier to maintain,
and easier to understand. The business logic has been named and surfaced
in the model.

\section{Interface roles of measures}\label{interface-roles-of-measures}

Measures do not all play the same role in the user experience. Some
measures answer business questions directly. Some display contextual
information. Some control whether rows appear. Some support the
presentation of a dashboard.

These are interface roles. There are four common roles:

\begin{itemize}
\tightlist
\item
  aggregating measures;
\item
  dimensional measures;
\item
  filtering measures;
\item
  dashboard measures.
\end{itemize}

\subsection{Aggregating measures}\label{aggregating-measures}

Aggregating measures are the measures most people think of first.

Examples include:

\begin{itemize}
\tightlist
\item
  \texttt{{[}Total\ sales\ amount{]}}
\item
  \texttt{{[}Total\ inspection\ hours{]}}
\item
  \texttt{{[}Average\ inspection\ duration{]}}
\item
  \texttt{{[}Active\ employees\ end\ of\ period{]}}
\item
  \texttt{{[}Median\ time\ to\ resolve{]}}
\end{itemize}

For measurable facts, aggregating measures may use functions such as
\texttt{sum}, \texttt{countrows}, \texttt{distinctcount},
\texttt{averagex}, or \texttt{medianx}.

For end-of-period facts, they often need to identify the relevant
reporting period before aggregating. For example,
\texttt{{[}Current\ employees{]}} may evaluate
\texttt{{[}Active\ employees\ end\ of\ period{]}} at the latest period
in the current filter context.

Aggregating measures should usually be business-facing. They are the
main way facts become visible to users.

\subsection{Dimensional measures}\label{dimensional-measures}

A dimensional measure turns dimension values into a measure result.

This sounds unusual because dimensions are normally used for grouping or
filtering. But sometimes a user wants a dimension value displayed as a
measure.

For example, in
\hyperref[docs-presenting-insights-filtering-behaviour]{Filtering
behaviour}, the user may want to show all products refunded on a
particular day. The product names come from
\texttt{\textquotesingle{}Sales\ product\textquotesingle{}}, but the
result is displayed as a measure.

A simplified example is:

\begin{verbatim}
Refunded products =
calculate (
    concatenatex (
        values ( 'Sales product'[Product name] ),
        'Sales product'[Product name],
        ", "
    ),
    crossfilter ( 'Sales product'[Product SK], 'Refund'[Product SK], both )
)
\end{verbatim}

This is not an aggregating measure in the normal sense. It is not
summing or counting fact rows. It is retrieving dimension values under a
particular filter context and turning them into display text.

Dimensional measures can be confused with dimension columns.

A dimension column such as
\texttt{\textquotesingle{}Sales\ product\textquotesingle{}{[}Product\ name{]}}
can be used in slicers, rows, columns, and cross-filtering.

A dimensional measure such as \texttt{{[}Refunded\ products{]}} cannot
do the same thing. It is a displayed result. It does not provide the
same interactive control.

Its chief purpose is display.

\subsection{Filtering measures}\label{filtering-measures}

Filtering measures exist to control what appears.

As explained in
\hyperref[docs-presenting-insights-filtering-behaviour]{Filtering
behaviour}, Power BI often retains visual rows where at least one
measure returns a non-blank value. A measure can therefore be used as a
filtering device.

Filtering measures typically return \texttt{1} or blank.

Examples include:

\begin{itemize}
\tightlist
\item
  \texttt{{[}Has\ sales{]}}
\item
  \texttt{{[}Has\ refund{]}}
\item
  \texttt{{[}Display\ sales\ transaction{]}}
\item
  \texttt{{[}Display\ refund\ transaction{]}}
\end{itemize}

A simple filtering measure might be:

\begin{verbatim}
Has sales =
if (
    countrows ( 'Sale' ) > 0,
    1
)
\end{verbatim}

This can be added to a visual or used as a visual-level filter to show
only products, dates, or IDs that have sales under the current filter
context.

Display unit-record measures are a special case. They are used to
control whether transaction records should appear. They may include
functions such as \texttt{isinscope} to ensure that rows appear only
when the relevant ID dimension is present.

Filtering measures are technical, but they still need business meaning.
Their descriptions should explain how they are intended to be used.

\subsection{Dashboard measures}\label{dashboard-measures}

Dashboard measures support report presentation.

They may provide:

\begin{itemize}
\tightlist
\item
  dynamic titles;
\item
  selected-value labels;
\item
  data currency messages;
\item
  conditional formatting values;
\item
  colour rules;
\item
  icons;
\item
  arrows;
\item
  warning messages.
\end{itemize}

For example:

\begin{verbatim}
Selected product label =
if (
    hasonevalue ( 'Sales product'[Product name] ),
    selectedvalue ( 'Sales product'[Product name] ),
    "Multiple products"
)
\end{verbatim}

Dashboard measures often inspect user context. They commonly use
functions such as \texttt{selectedvalue}, \texttt{hasonevalue},
\texttt{isfiltered}, and \texttt{isinscope}.

\section{Technical patterns for
measures}\label{technical-patterns-for-measures}

Most measures fall into a small number of technical patterns.

The first three fall under the standard case:

\begin{itemize}
\tightlist
\item
  base measures;
\item
  derived measures;
\item
  context-aware measures.
\end{itemize}

These patterns cover most ordinary scenarios.

Some measures become difficult because the model contains a more
demanding semantic problem. The final part of this section discusses
three advanced scenarios:

\begin{itemize}
\tightlist
\item
  polymorphic keys;
\item
  embedded grain;
\item
  unsupported relationships.
\end{itemize}

These are not measure types. They are situations where the data engineer
must be especially clear about grain, meaning, and filter context.

The techniques in this section assume that the model has been prepared
properly. They work best when the data engineer has preserved business
keys, built useful dimensions, hidden facts, prepared binary flags, and
created meaningful fragments in the pipeline.

If those foundations are missing, the same DAX patterns may become
difficult or impossible to apply. A measure cannot recover a grain if
the key for that grain has been discarded. It cannot express a business
category if the category has never been modelled. It cannot follow a
relationship path if the relevant identifier has not been preserved.

\subsection{Base measures}\label{base-measures}

Base measures interact directly with fact tables.

They are the first layer of measure design. They should usually be
simple, explicit, and reusable.

For example:

\begin{verbatim}
Total sales amount =
sum ( 'Sale'[Sales amount] )
\end{verbatim}

\begin{verbatim}
Sales transaction count =
countrows ( 'Sale' )
\end{verbatim}

\begin{verbatim}
Total inspection hours =
sum ( 'Inspection'[Inspection duration hours] )
\end{verbatim}

For end-of-period facts, base measures often include logic to select a
reporting period.

For example, a current-state measure may select the latest period in
context:

\begin{verbatim}
Employees end of period =
var latest_period =
    max ( 'Reporting calendar'[Period end date] )
return
    calculate (
        distinctcount ( 'Employee end of month'[Employee ID] ),
        keepfilters ( 'Reporting calendar'[Period end date] = latest_period )
    )
\end{verbatim}

The exact formula depends on the calendar design, but the pattern is
common: identify the relevant period, then evaluate the base measure
there.

For annotation facts, base measures may use \texttt{concatenatex} to
display detail.

For example:

\begin{verbatim}
Inspection comments =
concatenatex (
    'Inspection comments',
    'Inspection comments'[Comment text],
    " , ",
    'Inspection comments'[Comment sequence number]
)
\end{verbatim}

Base measures should carry as little business complexity as possible.
They should aggregate prepared facts and use prepared dimensions. For
example, failed inspections should be:

\begin{verbatim}
Failed inspection count =
calculate (
    [Inspection count],
    keepfilters ( 'Inspection outcome'[Is failed inspection] )
)
\end{verbatim}

\subsection{Derived measures}\label{derived-measures}

Derived measures build on base measures.

They include ratios, rates, comparisons, rolling periods, and
time-intelligence measures.

For example:

\begin{verbatim}
Inspection failure rate =
divide ( [Failed inspection count], [Inspection count] )
\end{verbatim}

Derived measures should usually be built from other measures rather than
repeating base logic.

Population comparison measures are another important derived pattern.
These measures escape part of the filter context to compare the current
value with a larger population.

For example:

\begin{verbatim}
National inspection count =
calculate (
    [Inspection count],
    removefilters ( 'Region' )
)
\end{verbatim}

The meaning should be explicit. The user should understand which filters
are removed and which remain.

\subsection{Context-aware measures}\label{context-aware-measures}

Context-aware measures inspect the user's current selection and respond
accordingly.

Dashboard measures are often context-aware, but aggregating and
filtering measures can be context-aware too.

Common functions include:

\begin{itemize}
\tightlist
\item
  \texttt{selectedvalue}
\item
  \texttt{hasonevalue}
\item
  \texttt{isfiltered}
\item
  \texttt{isinscope}
\item
  \texttt{values}
\item
  \texttt{isblank}
\end{itemize}

A report title may use:

\begin{verbatim}
Sales title =
"Sales for "
    & if (
        hasonevalue ( 'Sales product'[Product name] ),
        selectedvalue ( 'Sales product'[Product name] ),
        "all products"
    )
\end{verbatim}

A display measure may require an ID dimension before returning a value:

\begin{verbatim}
Display sales transaction =
if (
    isinscope ( 'Sales ID'[Sales order number] )
        && countrows ( 'Sale' ) > 0,
    1
)
\end{verbatim}

Context-aware measures are useful because Power BI users interact with
the model in many sequences. A measure can inspect the current context
and choose an appropriate response.

But context-aware logic should remain explainable. If a measure behaves
differently in different visuals, its description should make that
behaviour clear.

\subsection{Advanced scenario: polymorphic
keys}\label{advanced-scenario-polymorphic-keys}

Polymorphic keys occur when the meaning of one key column depends on
another column.

For example, a business process may inspect travellers and bags. Both
are recorded in the same
\texttt{\textquotesingle{}Inspection\textquotesingle{}} fact table. A
column such as \texttt{{[}Inspection\ item\ type{]}} has values such as
\texttt{Traveller} and \texttt{Bag}. Another column,
\texttt{{[}Inspected\ item\ SK{]}}, stores the inspected item key.

The meaning of \texttt{{[}Inspected\ item\ SK{]}} depends on
\texttt{{[}Inspection\ item\ type{]}}.

\textbf{Example structure of
\texttt{\textquotesingle{}Inspection\textquotesingle{}}}

{\def\LTcaptype{none} % do not increment counter
\begin{longtable}[]{@{}rlr@{}}
\toprule\noalign{}
Inspection SK & Inspection item type & Inspected item SK \\
\midrule\noalign{}
\endhead
\bottomrule\noalign{}
\endlastfoot
1 & Traveller & 1001 \\
2 & Traveller & 1002 \\
3 & Bag & 1001 \\
4 & Bag & 1003 \\
\end{longtable}
}

A simple distinct count of \texttt{{[}Inspected\ item\ SK{]}} gives the
wrong answer. It treats traveller \texttt{1001} and bag \texttt{1001} as
the same item.

The correct logic must count distinct items within each item type, then
add the results.

In a well-designed model, the determining column should usually be
expressed in a dimension such as
\texttt{\textquotesingle{}Inspection\ type\textquotesingle{}}. The
measure can then group by the determining column.

Conceptually:

\begin{verbatim}
Inspected item count =
var inspection_per_type =
    summarize (
        'Inspection type', -- group by dimension
        'Inspection type'[Inspection item type], -- on the polymorphic resolving column
        "Inspection count",
            distinctcountnoblank ( 'Inspection'[Inspected item SK] ) -- aggregating fact per segment
    )
return
    sumx ( inspection_per_type, [Inspection count] ) -- adds up the result from each segment
\end{verbatim}

If additional filters are needed, \texttt{summarize} can be wrapped
inside \texttt{calculatetable} with \texttt{keepfilters}.

New data engineers often solve this by concatenating item type and item
key into a single artificial column. That can work, but it changes the
model to solve a calculation problem. The approach above keeps the model
stable and handles the polymorphism inside the measure.

The idea is similar to SQL: group by the resolving column, aggregate
each segment, then aggregate the result.

\subsection{Advanced scenario: embedded
grain}\label{advanced-scenario-embedded-grain}

Embedded grain occurs when a table is physically stored at one grain,
but contains enough repeated keys or attributes to recover another
grain.

Power BI is more flexible than traditional dimensional modelling. A
single fact table can sometimes contain detail rows with header-level
values repeated on each row. This can be practical, but it creates a
measure problem.

For example, suppose
\texttt{\textquotesingle{}Inspection\textquotesingle{}} records
inspections at the traveller level, but also contains bag results.

\textbf{Example structure of
\texttt{\textquotesingle{}Inspection\textquotesingle{}}}

{\def\LTcaptype{none} % do not increment counter
\begin{longtable}[]{@{}rrrr@{}}
\toprule\noalign{}
Inspection SK & Traveller SK & Bag SK & Inspection duration minutes \\
\midrule\noalign{}
\endhead
\bottomrule\noalign{}
\endlastfoot
1 & 1001 & 501 & 45 \\
1 & 1001 & 502 & 45 \\
3 & 1002 & 503 & 30 \\
\end{longtable}
}

The inspection duration for \texttt{SK\ =\ 1} is repeated on bag rows.
If the data engineer averages or sums directly over
\texttt{\textquotesingle{}Inspection\textquotesingle{}}, the row may be
aggregated twice.

The solution is to reconstruct the intended grain before aggregating.

For example, to calculate median inspection duration per traveller:

\begin{verbatim}
Median time to inspect traveller =
var inspection_time_per_sk =
    summarize (
        'Inspection',
        'Inspection'[Inspection SK],
        'Inspection'[Inspection duration minutes]
    )
return
    medianx (
        inspection_time_per_sk,
        [Inspection duration minutes]
    )
\end{verbatim}

In SQL terms, this is similar to using \texttt{select\ distinct} on the
columns that define the intended grain, then aggregating over that
result.

ID dimensions can help because they preserve the keys needed to
reconstruct the grain. This is another reason ID dimensions are powerful
in Power BI models.

The broader lesson is that measures must know the grain of the thing
being measured. If the grain is embedded inside a denormalised fact, the
measure may need to recover it explicitly.

\subsection{Advanced scenario: unsupported
relationships}\label{advanced-scenario-unsupported-relationships}

Unsupported relationships occur when the standard single-direction
dimension-to-fact relationship structure does not naturally support the
required calculation.

The answer is not always to change the model. Sometimes the data
engineer can use DAX to create the required filter path for one measure.

The most useful tools are:

\begin{itemize}
\tightlist
\item
  \texttt{crossfilter}
\item
  \texttt{treatas}
\item
  \texttt{userelationship}
\end{itemize}

\texttt{crossfilter} temporarily changes relationship direction for a
calculation. As shown in
\hyperref[docs-presenting-insights-filtering-behaviour]{Filtering
behaviour}, it can be used to retrieve dimension values through a
relationship path that normally flows the wrong way.

\texttt{treatas} applies values from one table as though they filtered
another table. It is powerful when the information is already in the
model but no physical relationship exists.

For example, a refund table may need to retrieve the original sale
amount through sales order numbers. If the default relationship path
does not support the calculation, \texttt{treatas} may be used to apply
the refund's sales order numbers to
\texttt{\textquotesingle{}Sale\textquotesingle{}} for that measure.

\texttt{userelationship} activates an inactive relationship for the
duration of a calculation.

For example, suppose
\texttt{\textquotesingle{}Calendar\textquotesingle{}{[}Date{]}} has an
active relationship to
\texttt{\textquotesingle{}Sale\textquotesingle{}{[}Order\ date{]}} and
an inactive relationship to
\texttt{\textquotesingle{}Sale\textquotesingle{}{[}Delivery\ date{]}}. A
delivered-sales measure can activate the delivery-date relationship:

\begin{verbatim}
Total delivered sales amount =
calculate (
    [Total sales amount],
    userelationship ( 'Calendar'[Date], 'Sale'[Delivery date] )
)
\end{verbatim}

These functions are powerful because they allow specific measures to
behave differently without changing the whole model.

They should be used deliberately. If many measures require complex
relationship overrides, the model may need redesign. But when the
exception is narrow and well understood, measure-defined filtering is
often cleaner than adding confusing bidirectional relationships or
distorting facts.

\section{Key ideas}\label{key-ideas-5}

\begin{quote}
\textbf{Key ideas.}

Measures compress fact-table content into business meaning, then unpack
that meaning in user context.

Good measures should be business-centric and technically simple.

Business-facing measures should answer common questions directly rather
than forcing users to memorise filter combinations.

Complex business logic should usually be prepared in the data pipeline,
reference data, binary flags, or storytelling dimensions.

Measures can play different interface roles: aggregating, dimensional,
filtering, or dashboard-oriented.

Business keys help with difficult measure challenges by preserving
recoverable meaning across time, grain, and modelling complexity.

If measure logic becomes difficult to explain, the data engineer should
consider whether the model or pipeline should be improved.
\end{quote}

\cleardoublepage

\chapter{Measure of
measures}\label{docs-presenting-insights-measure-of-measures}

\section{Measures as a modelled
structure}\label{measures-as-a-modelled-structure}

The previous chapter described measures as the face of facts. A measure
compresses fact-table content into business meaning, then unpacks that
meaning in user context.

This chapter extends that idea.

When the measures themselves have dimensional structure, that structure
should be modelled.

The main idea is that a family of measures may have business dimensions
of its own. Measures may vary by process stage, metric type, or
reporting perspective. If these dimensions exist only inside measure
names, the structure is hidden from the model.

A measure of measures makes this structure explicit.

This should be distinguished from other Power BI approaches that look
similar on the surface: create a disconnected table of measure names,
let the user choose a measure, and use a \texttt{switch} measure to
return the selected result. That can be useful, but it is not the main
idea here.

Instead, \textbf{measure of measures} refers to the modelling pattern of
making the underlying structure of measures explicit for user
interaction. The \texttt{switch} measure is only one component of that
pattern.

\section{The problem}\label{the-problem}

Consider a company that manages three operational stages:

\begin{itemize}
\tightlist
\item
  Manufacture
\item
  Orders
\item
  Shipping
\end{itemize}

Each stage has three operational metrics:

\begin{itemize}
\tightlist
\item
  Process volume
\item
  Median time to process
\item
  Volume processed on schedule
\end{itemize}

This creates nine base measures.

These measures have a clear two-axis structure:

\begin{itemize}
\tightlist
\item
  Process stage
\item
  Metric
\end{itemize}

{\def\LTcaptype{none} % do not increment counter
\begin{longtable}[]{@{}
  >{\raggedright\arraybackslash}p{(\linewidth - 4\tabcolsep) * \real{0.3333}}
  >{\raggedright\arraybackslash}p{(\linewidth - 4\tabcolsep) * \real{0.3333}}
  >{\raggedright\arraybackslash}p{(\linewidth - 4\tabcolsep) * \real{0.3333}}@{}}
\toprule\noalign{}
\begin{minipage}[b]{\linewidth}\raggedright
Process stage
\end{minipage} & \begin{minipage}[b]{\linewidth}\raggedright
Metric
\end{minipage} & \begin{minipage}[b]{\linewidth}\raggedright
Base measure
\end{minipage} \\
\midrule\noalign{}
\endhead
\bottomrule\noalign{}
\endlastfoot
Manufacture & Process volume &
\texttt{{[}Manufacture\ process\ volume{]}} \\
Manufacture & Median time to process &
\texttt{{[}Manufacture\ median\ time\ to\ process\ (days){]}} \\
Manufacture & Volume processed on schedule &
\texttt{{[}Manufacture\ volume\ processed\ on\ schedule{]}} \\
Orders & Process volume & \texttt{{[}Orders\ process\ volume{]}} \\
Orders & Median time to process &
\texttt{{[}Orders\ median\ time\ to\ process\ (days){]}} \\
Orders & Volume processed on schedule &
\texttt{{[}Orders\ volume\ processed\ on\ schedule{]}} \\
Shipping & Process volume & \texttt{{[}Shipping\ process\ volume{]}} \\
Shipping & Median time to process &
\texttt{{[}Shipping\ median\ time\ to\ process\ (days){]}} \\
Shipping & Volume processed on schedule &
\texttt{{[}Shipping\ volume\ processed\ on\ schedule{]}} \\
\end{longtable}
}

That structure matters to business intent. A manager may want to compare
process volume across stages. Another may want to see whether time to
process is deteriorating across years. Another may want to view
operational performance by stage and metric in a single matrix.

If the structure exists only in measure names, the user cannot interact
with this structure directly. The report builder has to manually place
many separate measures into visuals and repeat derived logic.

Moreover, the complexity quickly explodes. For example, if two
additional concepts need to be measured:

\begin{itemize}
\tightlist
\item
  Percentage processed on schedule
\item
  National comparison
\end{itemize}

Calculating percentage processed on schedule for each process stage
leads to 3 additional measures and 12 total. If national versions are
created for all 12, the total becomes 24.

\section{The measure of measures
pattern}\label{the-measure-of-measures-pattern}

The measure of measures pattern introduces four elements:

\begin{enumerate}
\def\labelenumi{\arabic{enumi}.}
\tightlist
\item
  A measure table
\item
  A switch measure
\item
  Derived measures
\item
  A formatting calculation group
\end{enumerate}

\subsection{Step 1---Create a measure
table}\label{step-1create-a-measure-table}

The measure table is a small table that lists the measures and annotates
them with business attributes.

In this example, create a table called
\texttt{\textquotesingle{}Operational\ metric\textquotesingle{}}.

\textbf{Example structure of
\texttt{\textquotesingle{}Operational\ metric\textquotesingle{}}}

{\def\LTcaptype{none} % do not increment counter
\begin{longtable}[]{@{}
  >{\raggedleft\arraybackslash}p{(\linewidth - 10\tabcolsep) * \real{0.1905}}
  >{\raggedright\arraybackslash}p{(\linewidth - 10\tabcolsep) * \real{0.1429}}
  >{\raggedright\arraybackslash}p{(\linewidth - 10\tabcolsep) * \real{0.1429}}
  >{\raggedleft\arraybackslash}p{(\linewidth - 10\tabcolsep) * \real{0.1905}}
  >{\raggedright\arraybackslash}p{(\linewidth - 10\tabcolsep) * \real{0.1429}}
  >{\raggedleft\arraybackslash}p{(\linewidth - 10\tabcolsep) * \real{0.1905}}@{}}
\toprule\noalign{}
\begin{minipage}[b]{\linewidth}\raggedleft
Operational metric SK
\end{minipage} & \begin{minipage}[b]{\linewidth}\raggedright
Measure name
\end{minipage} & \begin{minipage}[b]{\linewidth}\raggedright
Process stage
\end{minipage} & \begin{minipage}[b]{\linewidth}\raggedleft
Process stage display order
\end{minipage} & \begin{minipage}[b]{\linewidth}\raggedright
Metric
\end{minipage} & \begin{minipage}[b]{\linewidth}\raggedleft
Metric display order
\end{minipage} \\
\midrule\noalign{}
\endhead
\bottomrule\noalign{}
\endlastfoot
1 & Manufacture process volume & Manufacture & 1 & Process volume & 1 \\
2 & Manufacture median time to process (days) & Manufacture & 1 & Median
time to process & 2 \\
3 & Manufacture volume processed on schedule & Manufacture & 1 & Volume
processed on schedule & 3 \\
4 & Orders process volume & Orders & 2 & Process volume & 1 \\
5 & Orders median time to process (days) & Orders & 2 & Median time to
process & 2 \\
6 & Orders volume processed on schedule & Orders & 2 & Volume processed
on schedule & 3 \\
7 & Shipping process volume & Shipping & 3 & Process volume & 1 \\
8 & Shipping median time to process (days) & Shipping & 3 & Median time
to process & 2 \\
9 & Shipping volume processed on schedule & Shipping & 3 & Volume
processed on schedule & 3 \\
\end{longtable}
}

This table is not a fact table. It is also not a normal business
dimension over transactional data. It is a dimension over the measures
themselves.

With the table, the model now knows that
\texttt{{[}Manufacture\ process\ volume{]}},
\texttt{{[}Orders\ process\ volume{]}}, and
\texttt{{[}Shipping\ process\ volume{]}} are the same kind of metric
applied to different process stages.

The display order columns are important. Alphabetical order is rarely
the same as business order. For example, \texttt{{[}Process\ stage{]}}
is in order of manufacture, orders, and shipping.

\subsection{Step 2---Create a switch
measure}\label{step-2create-a-switch-measure}

The switch measure returns the correct base measure depending on the
selected row in
\texttt{\textquotesingle{}Operational\ metric\textquotesingle{}}.

A simplified version of the measure \texttt{{[}Operational\ metric{]}}
looks like this:

\begin{verbatim}
Operational metric =
var selected_metric =
    selectedvalue ( 'Operational metric'[Measure name] )
return
    switch (
        selected_metric,
        "Manufacture process volume", [Manufacture process volume],
        "Manufacture median time to process (days)", [Manufacture median time to process (days)],
        "Manufacture volume processed on schedule", [Manufacture volume processed on schedule],
        "Orders process volume", [Orders process volume],
        "Orders median time to process (days)", [Orders median time to process (days)],
        "Orders volume processed on schedule", [Orders volume processed on schedule],
        "Shipping process volume", [Shipping process volume],
        "Shipping median time to process (days)", [Shipping median time to process (days)],
        "Shipping volume processed on schedule", [Shipping volume processed on schedule]
    )
\end{verbatim}

The measure table provides the selection. The switch measure provides
the value.

The ordinary underlying measures can be called base measures. They
calculate directly from facts or from ordinary dimensional context. The
switch measure is the first measure that operates over the structure of
those base measures.

\subsection{Step 3---Define derived
measures}\label{step-3define-derived-measures}

Once the switch measure exists, the data engineer can define derived
measures over the structure of the measure table.

Suppose the business wants percentage processed on schedule. Without a
measure of measures, the data engineer may create one percentage measure
per process stage:

\begin{itemize}
\tightlist
\item
  \texttt{{[}Manufacture\ percentage\ processed\ on\ schedule{]}}
\item
  \texttt{{[}Orders\ percentage\ processed\ on\ schedule{]}}
\item
  \texttt{{[}Shipping\ percentage\ processed\ on\ schedule{]}}
\end{itemize}

With the measure table, the data engineer can define a more general
derived measure.

For example, the model may contain these derived measures:

\begin{verbatim}
Process volume =
calculate (
    [Operational metric],
    keepfilters ( 'Operational metric'[Metric] = "Process volume" )
)
\end{verbatim}

\begin{verbatim}
Volume processed on schedule =
calculate (
    [Operational metric],
    keepfilters ( 'Operational metric'[Metric] = "Volume processed on schedule" )
)
\end{verbatim}

The percentage measure can then be written once:

\begin{verbatim}
Percentage processed on schedule =
divide (
    [Volume processed on schedule],
    [Process volume]
)
\end{verbatim}

A national comparison can also be written once:

\begin{verbatim}
National operational metric =
calculate (
    [Operational metric],
    removefilters ( 'Region' )
)
\end{verbatim}

The same measure works across process stages because
\texttt{{[}Process\ stage{]}} remains in the filter context.

For example, when the process stage is filtered to \texttt{Manufacture},
\texttt{{[}Process\ volume{]}} resolves to
\texttt{{[}Manufacture\ process\ volume{]}}, and
\texttt{{[}Volume\ processed\ on\ schedule{]}} resolves to
\texttt{{[}Manufacture\ volume\ processed\ on\ schedule{]}}.

When the process stage is filtered to \texttt{Orders}, the same measure
resolves to the order measures.

\subsection{Step 4---Handle formatting}\label{step-4handle-formatting}

A switch measure does not automatically preserve the original format of
each underlying measure.

In practice, this can be handled with a calculation group or another
dynamic format-string pattern.

A simplified format string expression may look like this:

\begin{verbatim}
switch (
    selectedvalue ( 'Operational metric'[Measure name] ),
    "Manufacture process volume", "#,0",
    "Orders process volume", "#,0",
    "Shipping process volume", "#,0",
    "Manufacture median time to process (days)", "#,0.0",
    "Orders median time to process (days)", "#,0.0",
    "Shipping median time to process (days)", "#,0.0",
    "Manufacture volume processed on schedule", "#,0",
    "Orders volume processed on schedule", "#,0",
    "Shipping volume processed on schedule", "#,0",
    selectedmeasureformatstring ()
)
\end{verbatim}

The exact implementation depends on the model and the calculation group
design. The important point is that dynamic measure selection also
requires deliberate formatting design.

Without this step, the measure of measures may return correct numbers
but display them poorly.

\begin{quote}
\textbf{Note.}

In practice, long \texttt{switch} measures and matching format-string
expressions should usually be generated by code when this pattern is
used regularly.
\end{quote}

\section{Benefits}\label{benefits}

The measure of measures pattern has several benefits.

\subsection{Simplified DAX management}\label{simplified-dax-management}

Instead of repeating similar logic across many measures, the data
engineer can define reusable derived measures.

This reduces duplication and lowers maintenance cost. In some cases, it
is the only viable solution without creating a tangle of DAX definitions
or nightmare reports.

\subsection{Explicit business
structure}\label{explicit-business-structure}

The measure table expresses the structure of the metrics directly.

Process stage, metric, display order, and other annotations become part
of the model rather than being hidden inside measure names.

For example, a matrix can place \texttt{{[}Metric{]}} on rows,
\texttt{{[}Process\ stage{]}} on columns, and
\texttt{{[}Operational\ metric{]}} as the value.

{\def\LTcaptype{none} % do not increment counter
\begin{longtable}[]{@{}lrrr@{}}
\toprule\noalign{}
Metric & Manufacture & Orders & Shipping \\
\midrule\noalign{}
\endhead
\bottomrule\noalign{}
\endlastfoot
Process volume & 12,400 & 11,850 & 11,200 \\
Median time to process & 2.1 & 1.4 & 3.8 \\
Volume processed on schedule & 11,780 & 10,960 & 9,520 \\
\end{longtable}
}

The numbers are returned by different underlying measures, but the row
and column structure comes from
\texttt{\textquotesingle{}Operational\ metric\textquotesingle{}}.

This visual is difficult to build with separate measures.

\subsection{More powerful reporting}\label{more-powerful-reporting}

The third benefit is powerful reporting. A visual can place
\texttt{{[}Process\ stage{]}} and \texttt{{[}Metric{]}} on rows,
\texttt{{[}Reporting\ year{]}} on columns, and
\texttt{{[}Operational\ metric{]}} as the value.

{\def\LTcaptype{none} % do not increment counter
\begin{longtable}[]{@{}llrrr@{}}
\toprule\noalign{}
Process stage & Metric & 2022 & 2023 & 2024 \\
\midrule\noalign{}
\endhead
\bottomrule\noalign{}
\endlastfoot
Manufacture & Process volume & 10,800 & 11,600 & 12,400 \\
Manufacture & Median time to process & 2.4 & 2.2 & 2.1 \\
Manufacture & Volume processed on schedule & 9,900 & 10,850 & 11,780 \\
Orders & Process volume & 10,200 & 11,100 & 11,850 \\
Orders & Median time to process & 1.7 & 1.5 & 1.4 \\
Orders & Volume processed on schedule & 9,300 & 10,100 & 10,960 \\
Shipping & Process volume & 9,800 & 10,600 & 11,200 \\
Shipping & Median time to process & 3.2 & 3.5 & 3.8 \\
Shipping & Volume processed on schedule & 8,950 & 9,350 & 9,520 \\
\end{longtable}
}

This gives the business a compact view of many measures across time. For
example, shipping volume is increasing, but median time to process is
also deteriorating. Conditional formatting can then highlight
deterioration or improvement.

There are many ways to use the structure expressed by
\texttt{\textquotesingle{}Operational\ metric\textquotesingle{}}. The
\texttt{{[}Process\ stage{]}} column can be used for small multiples in
line charts, creating a series of visuals in a controlled and automated
manner. It can also serve as a legend to compare metrics directly, such
as identifying lags or discrepancies in \texttt{{[}Process\ volume{]}}
between manufacture, orders, and shipping.

If each process is managed by a different team,
\texttt{{[}Process\ stage{]}} can be used as a report page filter to
tailor the view to each manager's concern.

Used appropriately, the measure of measures enables powerful reports
that would otherwise be impractical or impossible.

\subsection{Avoiding awkward fact-table
unions}\label{avoiding-awkward-fact-table-unions}

Inexperienced data engineers sometimes try to solve the same problem by
forcing fact tables into an awkward union so that different process
stages can appear as rows in a single table.

That may be appropriate in some models, but not always. If the processes
are genuinely different and already have good measures, it may be better
to structure the measures rather than reshape the facts.

The measure of measures allows users to operate on the measures
directly. It is a form of
\hyperref[docs-creating-information-conforming-systems]{horizontal
integration} when vertical integration of the underlying facts is not
appropriate.

\section{Dangers}\label{dangers}

The measure of measures pattern is powerful, but it should be used
judiciously.

\subsection{Performance}\label{performance}

Large \texttt{switch} measures can perform poorly, especially when the
selected field is not the same column used by the switch.

For example, performance may be worse when the visual filters on
\texttt{{[}Metric{]}} or \texttt{{[}Process\ stage{]}}, but the switch
branches on \texttt{{[}Measure\ name{]}}.

With many branches, performance should be tested in practice.

\subsection{Dynamism for its own sake}\label{dynamism-for-its-own-sake}

Dynamism can be overused.

In theory, any list of measures can be placed into a switch. That does
not mean it should be done.

The criterion is business expressiveness. The pattern is appropriate
when the measures form a coherent structure: process stages, business
lines, quality criteria, product families, or similar.

\subsection{Hiding important information behind
clicks}\label{hiding-important-information-behind-clicks}

A dynamic report is not automatically a good report.

If users must keep clicking before they can see what matters, the report
may be poorly designed. A report with ten selectable metrics is often
less intuitive than a report that shows the important comparisons
directly.

The measure of measures should be used to create explicitly expressive
visuals, not to hide information behind unnecessary interaction.

\subsection{Ambiguous aggregation across
measures}\label{ambiguous-aggregation-across-measures}

The measure table often has a composite key. In the example above,
\texttt{{[}Process\ stage{]}} and \texttt{{[}Metric{]}} together
identify a specific measure.

If the user selects \texttt{{[}Metric{]}\ =\ Process\ volume} without
selecting \texttt{{[}Process\ stage{]}}, the switch measure returns
blank. This is often correct. The default behaviour prevents summing
over items that should not be summed.

If aggregation across measures is required, it should be explicitly
defined and tailored to the metric.

\section{Relationship to switch
measures}\label{relationship-to-switch-measures}

Many discussions of \texttt{switch} measures in Power BI focus on
technical convenience: allowing users to choose a measure from a slicer,
reducing visual clutter, or reusing a visual for different metrics.

Those uses can be valid, but they are not the main point here.

A measure of measures is not a slicer trick. It is a modelling technique
for expressing business intent. The measure table allows the data
engineer to:

\begin{itemize}
\tightlist
\item
  make business metrics explicit and navigable;
\item
  build structured visuals by metric and process;
\item
  define derived measures without duplication.
\end{itemize}

Without anchor in business meaning, dynamic measure selection can easily
become frustrating. The model may become clever, but not clearer.

\section{Key ideas}\label{key-ideas-6}

\begin{quote}
\textbf{Key ideas.}

Sometimes the measures themselves have dimensional structure, and that
structure should be modelled.

Measures can themselves be modelled as a business object.

A measure of measures is useful when measures have a coherent structure,
such as process stage by metric.

The important idea is not merely that a \texttt{switch} measure can
change what is displayed. The important idea is that the measure family
may have dimensional structure worth modelling.

The pattern uses a measure table, a switch measure, derived measures,
and usually a formatting calculation group.

The measure table makes the structure of the measures explicit rather
than hiding it in measure names.

Derived measures can often be written once because the process or metric
structure is carried by the measure table.

The pattern should be used for business expressiveness, not dynamism for
its own sake.

Aggregating across measures must be defined deliberately. Blank results
are often a useful safeguard against accidental nonsense.
\end{quote}

\cleardoublepage

\chapter{Interactive row level
security}\label{docs-presenting-insights-interactive-row-level-security}

\section{Security without losing
context}\label{security-without-losing-context}

Row level security (RLS) in Power BI means exposing to users only the
rows they are allowed to see. This is implemented through a DAX filter
attached to the user's security role.

The important subtlety is that an RLS filter is not visible to DAX as an
ordinary, inspectable filter. Functions such as \texttt{isfiltered()} or
\texttt{isinscope()} cannot detect that RLS is active. A formula simply
evaluates over the rows available to the user after security has been
applied.

This makes restriction straightforward but comparison difficult. If a
regional user is restricted to regional rows, then ordinary measures
automatically evaluate over that regional population. A national
comparison cannot be recovered by asking DAX to notice and remove the
RLS filter. The wider population must be deliberately modelled.

A common solution is to summarise the wider population into aggregate
tables. This can reduce disclosure risk, and it may be the right
approach for highly sensitive data. However, summary tables do not
preserve unit-record interactivity. They also introduce their own
modelling complexity, because every supported comparison must be
anticipated in the summary design.

This chapter focuses on approaches that preserve unit-record
interactivity while still allowing comparison with a wider population.

This leads to two distinct tasks in implementing RLS:

\begin{itemize}
\tightlist
\item
  limiting what users can see;
\item
  preserving population context.
\end{itemize}

\section{Limiting what users can see}\label{limiting-what-users-can-see}

RLS first needs an access-control pattern. There are two common
approaches:

\begin{itemize}
\tightlist
\item
  static roles;
\item
  dynamic roles.
\end{itemize}

This choice determines how the permitted population is assigned to each
user.

\subsection{Static roles}\label{static-roles}

Static roles use one RLS role per access category.

For example, the model may have separate roles for \texttt{North},
\texttt{South}, \texttt{West}, and \texttt{Unlimited}.

The \texttt{North} role may apply this filter to
\texttt{\textquotesingle{}Region\textquotesingle{}}:

\begin{verbatim}
'Region'[Region name] = "North"
\end{verbatim}

The \texttt{South} role may apply:

\begin{verbatim}
'Region'[Region name] = "South"
\end{verbatim}

The \texttt{West} role may apply:

\begin{verbatim}
'Region'[Region name] = "West"
\end{verbatim}

The \texttt{Unlimited} role has no filter. Users assigned to this role
can see all rows allowed by the model.

This is simple to understand and easy to test. Each role has its own DAX
filter, and users are assigned to the role that matches their access.

The weakness is maintenance. Static roles work when the access
categories are few and stable. They become awkward when users have many
different access combinations, when access changes often, or when
permissions need to follow organisational hierarchy.

A clear naming convention is important. \texttt{Limited} and
\texttt{Unlimited} are often clearer than \texttt{Restricted} and
\texttt{Unrestricted}, because \texttt{restricted} can mean either
access to restricted data or restriction from data.

\subsection{Dynamic roles}\label{dynamic-roles}

Dynamic roles use a security mapping table and the user's identity,
usually through \texttt{userprincipalname()}.

Instead of creating one role per access category, the model has a role
that looks up which rows the current user is allowed to see. The mapping
table may connect users to regions, teams, programs, customers, or other
secured entities.

For example, a simple user-region mapping table may look like this.

\textbf{Example structure of
\texttt{\textquotesingle{}User\ access\textquotesingle{}}}

{\def\LTcaptype{none} % do not increment counter
\begin{longtable}[]{@{}ll@{}}
\toprule\noalign{}
User principal name & Region \\
\midrule\noalign{}
\endhead
\bottomrule\noalign{}
\endlastfoot
alice@example.com & North \\
ben@example.com & South \\
clara@example.com & West \\
dana@example.com & North \\
dana@example.com & South \\
dana@example.com & West \\
\end{longtable}
}

The RLS rule can then filter the model by the regions assigned to the
current user.

A DAX expression may look like this:

\begin{verbatim}
'User access'[User principal name] = userprincipalname ()
\end{verbatim}

In practice, the filter often applies through relationships from the
access table to a secured dimension or directly to the fact table.

Dynamic roles are more flexible and auditable. They allow access rules
to be maintained in data rather than hard-coded into many RLS roles.
They are especially useful when permissions are numerous, changing,
hierarchical, or managed outside Power BI.

The trade-off is that the access table becomes part of the security
infrastructure. It must be governed, tested, and documented. If the
mapping logic is complex, the table should include metadata explaining
why each assignment exists.

From a performance perspective, the filter should be applied as close as
possible to the sensitive data. If the sensitive data resides in fact
tables, the filter should ideally reach those fact tables directly. If
the sensitive data resides in dimensions, the filter may apply through a
snowflake relationship.

Multiple copies of the access table may be needed if filtering applies
across multiple dimensions. In such cases, care must be taken to avoid
conflicting filters that accidentally combine as an unintended
\texttt{and} condition.

Whether access is static or dynamic, the result is the same: each user
sees only their permitted rows. The harder problem is what happens when
the user needs to compare those rows with a wider population.

\section{Preserving population
context}\label{preserving-population-context}

RLS removes rows from the user's visible population. If the user needs
comparison against a wider population, that wider population must remain
available through another modelling path.

For example, a regional team may need to compare its sales performance
against national sales. The comparison may need to remain filterable by
product type, sales period, customer segment, or other dimensions.

A measure such as this cannot recover rows removed by RLS:

\begin{verbatim}
National sales volume =
calculate (
    [Sales volume],
    removefilters ( 'Sales team' )
)
\end{verbatim}

\texttt{removefilters()} can remove ordinary filter context. It cannot
bypass security. If RLS has already restricted the user to regional
rows, then the national rows are not available to the measure.

There are two approaches to preserving population context while
retaining unit-record interactivity:

\begin{itemize}
\tightlist
\item
  anonymous facts;
\item
  pseudonymous dimensions.
\end{itemize}

\subsection{Anonymous facts}\label{anonymous-facts}

The anonymous approach duplicates the fact table and removes sensitive
columns and sensitive relationship paths.

For example, suppose the model contains:

\begin{itemize}
\tightlist
\item
  a fact table \texttt{\textquotesingle{}Sales\textquotesingle{}};
\item
  a sensitive dimension
  \texttt{\textquotesingle{}Sales\ team\textquotesingle{}};
\item
  a measure \texttt{{[}Sales\ volume{]}}.
\end{itemize}

The restricted user should see only their permitted rows in
\texttt{\textquotesingle{}Sales\textquotesingle{}}, but should also
compare against the wider population.

The data engineer can create a second fact table:

\begin{itemize}
\tightlist
\item
  \texttt{\textquotesingle{}Sales\ anonymous\textquotesingle{}}
\end{itemize}

This table contains the population sales rows, but removes columns and
relationships that identify the sensitive team or user. The relationship
to \texttt{\textquotesingle{}Sales\ team\textquotesingle{}} is removed.
Any sensitive columns are removed. The table is therefore not filtered
by the same RLS path.

A population measure can then be defined over the anonymous fact:

\begin{verbatim}
Population sales volume =
sum ( 'Sales anonymous'[Sales volume] )
\end{verbatim}

The user-facing comparison measure might be:

\begin{verbatim}
Regional sales proportion =
divide (
    [Sales volume],
    [Population sales volume]
)
\end{verbatim}

The ordinary \texttt{{[}Sales\ volume{]}} measure reflects the user's
restricted rows. \texttt{{[}Population\ sales\ volume{]}} reflects the
wider anonymous population. Both can still respond to safe shared
dimensions such as calendar, product, or region grouping where
appropriate.

For example, if both \texttt{\textquotesingle{}Sales\textquotesingle{}}
and \texttt{\textquotesingle{}Sales\ anonymous\textquotesingle{}} relate
to \texttt{\textquotesingle{}Calendar\textquotesingle{}} and
\texttt{\textquotesingle{}Product\textquotesingle{}}, the user can still
compare their regional sales with population sales by month and product
type.

The anonymous approach preserves interactivity by keeping the population
fact at the grain needed for analysis, while removing the sensitive
columns and relationship paths that would identify restricted entities.

When there are many pairs of user and population measures, maintaining
the measures can become tedious. The measure of measures pattern can
help by creating one switch measure for restricted measures and another
for population measures. Derived comparison measures can then be defined
over the two switch measures rather than repeated for every metric.

\subsection{Pseudonymous dimensions}\label{pseudonymous-dimensions}

The pseudonymous approach assumes sensitive information is contained in
dimensions.

For example, the sensitive information may be in
\texttt{\textquotesingle{}Sales\ team\textquotesingle{}}, which contains
team names, staff names, locations, or other identifying attributes.

Instead of duplicating the fact table, the data engineer duplicates the
dimension rows through a union. One set of rows contains the real
values. The other set contains masked values.

\textbf{Example structure of
\texttt{\textquotesingle{}Sales\ team\textquotesingle{}}}

{\def\LTcaptype{none} % do not increment counter
\begin{longtable}[]{@{}rllll@{}}
\toprule\noalign{}
Sales team SK & Sales team ID & Team name & Region & Is masked \\
\midrule\noalign{}
\endhead
\bottomrule\noalign{}
\endlastfoot
1 & T001 & North Metro Team & North & false \\
2 & T002 & South Coast Team & South & false \\
101 & T001 & (masked) & North & true \\
102 & T002 & (masked) & South & true \\
\end{longtable}
}

The access table determines which doubled dimension rows each user
principal may see. Since each \texttt{{[}Sales\ team\ ID{]}} has both a
masked and unmasked row, the access table grants access to the
appropriate \texttt{{[}Sales\ team\ SK{]}} for each user.

For example, one user may see the unmasked row for \texttt{T001} but the
masked row for \texttt{T002}. Another user may have the reverse access.
A national manager may see unmasked rows for all sales teams. The access
table may look like this:

\textbf{Example structure of
\texttt{\textquotesingle{}User\ sales\ team\ access\textquotesingle{}}}

{\def\LTcaptype{none} % do not increment counter
\begin{longtable}[]{@{}lr@{}}
\toprule\noalign{}
User principal name & Sales team SK \\
\midrule\noalign{}
\endhead
\bottomrule\noalign{}
\endlastfoot
alice@example.com & 1 \\
alice@example.com & 102 \\
ben@example.com & 101 \\
ben@example.com & 2 \\
clara@example.com & 1 \\
clara@example.com & 2 \\
\end{longtable}
}

This user access table filters the
\texttt{\textquotesingle{}Sales\ team\textquotesingle{}} dimension on
\texttt{{[}Sales\ team\ SK{]}}. The RLS rule is applied to
\texttt{\textquotesingle{}User\ sales\ team\ access\textquotesingle{}}
using \texttt{{[}User\ principal\ name{]}}. The filtered access table
then filters \texttt{\textquotesingle{}Sales\ team\textquotesingle{}}
through \texttt{{[}Sales\ team\ SK{]}}, leaving only the masked or
unmasked rows the user is allowed to see.

The relationship from the fact table to
\texttt{\textquotesingle{}Sales\ team\textquotesingle{}}, based on
\texttt{{[}Sales\ team\ ID{]}}, becomes many-to-many. Since
\texttt{{[}Sales\ team\ SK{]}} has multiple values in the access table,
it is also many-to-many.

The relationships can be summarised as a bus matrix.

{\def\LTcaptype{none} % do not increment counter
\begin{longtable}[]{@{}lll@{}}
\toprule\noalign{}
Dimension / access table &
\texttt{\textquotesingle{}Sales\ team\textquotesingle{}} &
\texttt{\textquotesingle{}Sales\textquotesingle{}} \\
\midrule\noalign{}
\endhead
\bottomrule\noalign{}
\endlastfoot
\texttt{\textquotesingle{}User\ sales\ team\ access\textquotesingle{}} &
* → * via \texttt{{[}Sales\ team\ SK{]}} & \\
\texttt{\textquotesingle{}Sales\ team\textquotesingle{}} & & * → * via
\texttt{{[}Sales\ team\ ID{]}} \\
\end{longtable}
}

Following the previous example, these would be the view for each user.

\textbf{Alice's view}

{\def\LTcaptype{none} % do not increment counter
\begin{longtable}[]{@{}llr@{}}
\toprule\noalign{}
Team name & Region & Goals reached \\
\midrule\noalign{}
\endhead
\bottomrule\noalign{}
\endlastfoot
North Metro Team & North & 84 \\
(masked) & South & 57 \\
\end{longtable}
}

\textbf{Ben's view}

{\def\LTcaptype{none} % do not increment counter
\begin{longtable}[]{@{}llr@{}}
\toprule\noalign{}
Team name & Region & Goals reached \\
\midrule\noalign{}
\endhead
\bottomrule\noalign{}
\endlastfoot
(masked) & North & 84 \\
South Coast Team & South & 57 \\
\end{longtable}
}

\textbf{Clara's view}

{\def\LTcaptype{none} % do not increment counter
\begin{longtable}[]{@{}llr@{}}
\toprule\noalign{}
Team name & Region & Goals reached \\
\midrule\noalign{}
\endhead
\bottomrule\noalign{}
\endlastfoot
North Metro Team & North & 84 \\
South Coast Team & South & 57 \\
\end{longtable}
}

In this case, each user can only see the full details of the teams they
are granted access to. The fact table is unaffected because RLS is
selecting visible dimension rows rather than removing fact rows
directly. This means population measures can still use the same fact
table through \texttt{removefilters()} where the security design permits
it.

Care must be taken when masking. In some cases, the primary key itself
is sensitive. For example, if the business key is a social security
number or another identifying value, it should not be exposed as the
relationship key. A surrogate key or non-sensitive mapping key is
required.

The pseudonymous approach is lighter than duplicating facts, but it is
also riskier. Because the same underlying facts remain available, an
informed user may infer identities from distinctive combinations of
values.

\section{Choosing an approach}\label{choosing-an-approach}

The anonymous and pseudonymous approaches solve the same problem in
different ways.

The anonymous approach duplicates the population facts and removes
sensitive columns and relationship paths. It is safer and easier to
reason about, but it increases model size and requires parallel
population measures.

The pseudonymous approach doubles dimension rows and masks sensitive
values. It keeps the model smaller and preserves the same fact table.
However, because the same underlying fact rows remain available, it
exposes more detail than the anonymous approach.

It is most useful when the identity of the row is not itself the main
concern, but particular attributes need to be hidden. For example, a
user may be allowed to see all traveller records and analyse their
activity, but should not see sensitive values such as passport numbers,
phone numbers, or personal identifiers.

Both approaches can preserve unit-level interactivity. This is
intentional. The goal is to keep dimensional filtering, drillthrough,
and detailed comparison available while limiting exposure of sensitive
attributes.

This also means both approaches carry inference risk. A user may deduce
identities from distinctive combinations of values, even when direct
identifiers are removed or masked.

If that risk is unacceptable, the required design is no longer merely
anonymous or pseudonymous RLS. It becomes a summary-table problem, with
specific applications of mitigation strategies. This is outside the
scope of this chapter.

The choice between anonymous population facts and pseudonymous
dimensions is not purely technical. It reflects the organisation's
appetite for risk, governance requirements, security obligations, and
expected user behaviour.

\section{Key ideas}\label{key-ideas-7}

\begin{quote}
\textbf{Key ideas.}

RLS should limit access without destroying analytical context.

DAX cannot detect RLS as an ordinary filter. Measures evaluate over the
rows available after security has been applied.

\texttt{removefilters()} cannot recover rows removed directly by RLS.
Population context must be preserved through modelling design.

RLS has two distinct tasks: limiting what users can see, and preserving
population context.

Anonymous population facts preserve wider population context by
duplicating facts and removing sensitive columns and relationship paths.

Pseudonymous dimensions preserve population context by masking sensitive
dimension values while keeping the model interactive.

Both anonymous and pseudonymous approaches preserve unit-record
interactivity, but both carry inference risk.
\end{quote}

\cleardoublepage

\thispagestyle{empty}
\phantomsection
\hypertarget{docs-quality-reliability}{}
\addcontentsline{toc}{part}{Quality \& reliability}
\vspace*{0.34\textheight}
\begin{center}
{\Huge\bfseries Quality \& reliability}
\end{center}
\cleardoublepage

\chapter{The foundations of
trust}\label{docs-quality-reliability-the-foundations-of-trust}

\section{Quality and reliability}\label{quality-and-reliability}

Data is a fragment of reality captured by process.

Because of this, data is an
\hyperref[docs-foundations-what-is-data]{imperfect projection} of the
business world into the data world. Data quality issues arise from the
gap between the business world and this projection.

For example, key information required for decisions may not be captured.
It may be captured in free-text format. It may be captured at the wrong
grain.

When business users encounter these gaps, they experience a mismatch
between the data product and the business world they already know. In
many cases, this means they cannot trust the product for
decision-making.

No amount of blaming the source system changes this reality. From the
business's point of view, the product simply cannot be used.

In this perspective, one of the data engineer's jobs is to close the gap
left by the projection, to the satisfaction of the decision-maker. This
is how the data engineer adds to data \textbf{quality}.

Beyond the imperfections of the projection itself, the mechanical
processes of a data pipeline may introduce further errors.

An example is schema change. Users may experience this when entire
tables become empty or when inconsistencies arise between business
processes.

On some occasions, this can be disastrous. For example, if the cost
forecast has been updated to the latest month but the actual expenditure
has not, the final reconciliation may contain an unacceptable
discrepancy. However, there may be no obvious signs of such an error.

The issue is not with the source data, but with mechanical
\textbf{reliability}: the pipeline has produced a data product whose
parts are no longer aligned and has nothing to signal this divergence.

Reliability is the discipline of making the data product stable against
predictable forms of breakage.

\section{Trust requires both}\label{trust-requires-both}

Quality keeps users from guessing what the data means. Reliability keeps
users from guessing whether the output worked.

When both are available, users can use the product without
second-guessing. Thus, quality and reliability are the foundations of
trust.

{\def\LTcaptype{none} % do not increment counter
\begin{longtable}[]{@{}
  >{\raggedright\arraybackslash}p{(\linewidth - 4\tabcolsep) * \real{0.3333}}
  >{\raggedright\arraybackslash}p{(\linewidth - 4\tabcolsep) * \real{0.3333}}
  >{\raggedright\arraybackslash}p{(\linewidth - 4\tabcolsep) * \real{0.3333}}@{}}
\toprule\noalign{}
\begin{minipage}[b]{\linewidth}\raggedright
Foundation
\end{minipage} & \begin{minipage}[b]{\linewidth}\raggedright
Question
\end{minipage} & \begin{minipage}[b]{\linewidth}\raggedright
User failure mode
\end{minipage} \\
\midrule\noalign{}
\endhead
\bottomrule\noalign{}
\endlastfoot
Quality & Does this data product represent the business reality well
enough for the decision? & ``This does not match the world I know.'' \\
Reliability & Does the pipeline consistently produce the data product
without hidden breakage? & ``This number might be broken.'' \\
\end{longtable}
}

\section{The third principle: anticipate
errors}\label{the-third-principle-anticipate-errors}

This section introduces the third principle of data engineering:
anticipate errors.

This principle follows from the nature of quality and reliability. Data
is an imperfect projection of complex reality, and pipelines are
mechanical systems that operate in a changing world.

Because of these changes, a trustworthy data product cannot depend on
everything continuing as expected. The data engineer does not stop with
what works now, but anticipate errors that may occur.

\section{Key ideas}\label{key-ideas-8}

\begin{quote}
\textbf{Key ideas.}

Trust is produced when a data product is fit for decision.

Quality means the data product represents business reality well enough
for its intended use.

Reliability means the pipeline consistently produces the data product
without hidden failure.

Quality keeps users from guessing what the data means.

Reliability keeps users from guessing whether the data worked.

The third principle of data engineering is to anticipate errors.
\end{quote}

\cleardoublepage

\chapter{Quality
metadata}\label{docs-quality-reliability-quality-metadata}

\section{When metadata is missing}\label{when-metadata-is-missing}

Metadata is the place where meaning is either carried by the data
product or abandoned to the user.

Data quality issues reflect the gap between the business world and the
data world. Metadata helps bridge that gap by adding the interpretative
context that connects the user to the business reality behind the data.

A source column called \texttt{{[}Valid{]}} may look harmless until
someone has to use it.

Valid for what? Valid now, or valid at creation? Valid according to the
source system, the business rule, or the reporting model? Does
\texttt{false} mean expired, rejected, superseded, or never approved?

Renaming \texttt{{[}Valid{]}} to \texttt{{[}Is\ current\ name{]}} is not
cosmetic. It changes the user's ability to understand the data.

This is why metadata is not an add-on task. It is part of modelling. No
transformation can help a user who cannot tell what the result means.

In complex cases, writing metadata first is a way to lay out the logic
in plain language before implementation.

There are three basic kinds of metadata discussed in this chapter:

{\def\LTcaptype{none} % do not increment counter
\begin{longtable}[]{@{}
  >{\raggedright\arraybackslash}p{(\linewidth - 4\tabcolsep) * \real{0.3333}}
  >{\raggedright\arraybackslash}p{(\linewidth - 4\tabcolsep) * \real{0.3333}}
  >{\raggedright\arraybackslash}p{(\linewidth - 4\tabcolsep) * \real{0.3333}}@{}}
\toprule\noalign{}
\begin{minipage}[b]{\linewidth}\raggedright
Metadata
\end{minipage} & \begin{minipage}[b]{\linewidth}\raggedright
Question it answers
\end{minipage} & \begin{minipage}[b]{\linewidth}\raggedright
Danger when compromised
\end{minipage} \\
\midrule\noalign{}
\endhead
\bottomrule\noalign{}
\endlastfoot
Names & What is this thing? & Unrecognisable \\
Business descriptions & What does this mean, and when is it unsafe to
use? & Incomprehensible \\
Database keys & How does this relate to reality and to other artefacts?
& Lost \\
\end{longtable}
}

Without good names, artefacts become unrecognisable. Without
descriptions, artefacts become incomprehensible. Without keys, they
become lost: disconnected from reality and from the other artefacts that
give them meaning.

\section{Names answer: what is this
thing?}\label{names-answer-what-is-this-thing}

Names make artefacts recognisable.

A table, column, measure, calculation group, or display folder may be
technically visible but still unrecognisable in business terms. If a
user cannot recognise what business thing an artefact represents, the
artefact cannot be used.

This is why names should be business-centric and written in plain
language. They should not use idiosyncratic acronyms or abbreviations
for the developer's convenience. There is no reason to use
\texttt{{[}Employee\ cnt{]}} when the meaning is
\texttt{{[}Employee\ count{]}}. If the infrastructure supports spaces in
column names, they should be used.

A recognisable name can also contain useful signals. \texttt{Is} and
\texttt{Has} can indicate binary columns, such as
\texttt{{[}Is\ non-compliant{]}} or \texttt{{[}Has\ active\ case{]}}.
Columns ending with \texttt{date} and \texttt{datetime} can signal
precision. Terms such as \texttt{current}, \texttt{historical},
\texttt{earliest}, \texttt{latest}, \texttt{start}, and \texttt{end}
should be used consistently.

These signals are helpful only when they clarify business meaning. They
should not advertise technical structure at the expense of user
understanding.

This is why a Power BI dimensional model should not prefix tables with
\texttt{dim} or \texttt{fact}. Such prefixes describe the modelling
technique, not the user-facing meaning.

Recognisability also depends on consistency. If the same business
concept appears under different names across the warehouse, the user
cannot tell whether the concepts are genuinely different or merely named
differently.

The data engineer should therefore pay particular attention to conformed
concepts. Before naming something new, check whether the same concept
already appears elsewhere under a different name.

For example, when adding \texttt{Cake.RefQualityControl}, it is
important to check whether \texttt{RefTesting}, \texttt{RefAudit}, or
\texttt{RefInspection} are already being used to mean the same thing. If
so, it may be better to reuse an existing name.

Names should also be explicit and unambiguous because ambiguity creates
false recognition. The user thinks they know what the artefact means,
but may be wrong.

An explicit column name means using
\texttt{{[}Inspection\ first\ of\ month{]}} rather than
\texttt{{[}First\ of\ month{]}}.

An unambiguous name means ensuring the word has only one meaning for the
user. For example, does \texttt{{[}Last\ inspection\ date{]}} mean the
previous inspection date or the final inspection date? Does
\texttt{restricted} mean less access, or access to restricted data?

Names should also be able to stand alone. Data tables and columns can
often be taken out of context. For key concepts, it may be worthwhile
for names to carry their own context.

For example, \texttt{{[}Is\ non-compliant{]}} may be better as
\texttt{{[}Is\ financially\ non-compliant{]}}. This may involve
repeating some of the table meaning in the column name.

When in doubt, the data engineer should err on the side of explicitness,
because it is easier to shorten a name in a Power BI visual than for a
user to creatively extend it.

Another compromise to recognisability is filler language.

Words such as \texttt{type}, \texttt{group}, \texttt{summary}, and
\texttt{details} are often used by developers without adding meaning.
They can function like ``um'' and ``ah'' in a sentence: they fill space
but obscure meaning.

If a table is called \texttt{summary}, it should usually be renamed to
reflect what is being summarised. \texttt{Details} should be avoided
unless it is truly a finer-grain expansion of a header row.

These principles apply just as much to code. Temporary tables and CTEs
should be named with business meaning, not cryptic abbreviations like
\texttt{tmp}. In complex code, column renaming should happen early, on
first contact, and persist throughout transformation.

\section{Descriptions answer: what does this
mean?}\label{descriptions-answer-what-does-this-mean}

Descriptions make artefacts comprehensible.

A name may tell the user what an artefact is, but the name cannot carry
the whole interpretation. The user may still need to know what business
reality is represented, how the value was derived, what assumptions were
made, when the artefact is unsafe, and how it differs from related
concepts.

This is why business descriptions are part of the meaning rather than an
add-on. They enhance data quality by aligning artefacts to business
reality through interpretative context.

Too often, descriptions merely restate the artefact name. For example,
\texttt{Cake.Sales} may be described as ``A list of cakes sold'', or
\texttt{{[}Is\ non-compliant{]}} as ``True if found non-compliant.''
These descriptions do nothing more than take up storage space. At worst,
they harm trust by conveying shallowness.

Like names, descriptions should be business-centric and written in plain
language. But unlike names, they should be substantial.

A good description should explain:

\begin{itemize}
\tightlist
\item
  what business reality is represented;
\item
  how the value was derived;
\item
  what assumptions were made;
\item
  when the artefact is unsafe or limited;
\item
  how it relates to similar concepts.
\end{itemize}

A useful description carries the interpretation that would otherwise
have to be supplied by tribal knowledge or guesswork.

A column such as \texttt{{[}Is\ non-compliant{]}} should describe what
non-compliance means in business terms, not merely say ``if the result
code is \texttt{F}.''

Data engineers and stakeholders often slip into using system names as
shorthand for business concepts. Descriptions should refer to business
entities and business processes, not merely the systems that implement
them.

There will be cases where system names or technical implementation
details are necessary for transparency. These should be included as
supplementary information, not as standalone descriptions.

For example, the description for \texttt{{[}Is\ non-compliant{]}} may
be:

\begin{quote}
True if the transaction failed to meet the minimum regulatory
requirements for importation, based on the department's published
compliance criteria. This includes both automatic and manual
assessments. Transactions flagged by automated rules are included
regardless of whether they were later manually overridden. This column
does not include cases where compliance was not assessed. The result
code is derived from the \texttt{ResultCode} field in the source system,
where \texttt{F} indicates failure.
\end{quote}

The same principle applies to transformation logic. The transformation
should be stated in business terms. Technical statements can be included
as elaboration, but not as a substitute for meaning.

For example:

\begin{quote}
This column is derived from the inspection result and officer notes. It
is true if the inspection outcome was \texttt{Spoiled\ cake} or
\texttt{Not\ tasty}, and the inspection was conducted by a certified
officer. It excludes inspections conducted during training. The logic is
implemented in the \texttt{InspectionOutcomeFlag} measure using a
combination of result codes and officer certification status.
\end{quote}

Limitations should also be framed in terms of business interpretation.

``The column can have a null value'' is not a useful limitation by
itself.

A better statement is:

\begin{quote}
The inspector may not record a result if there are no issues. This means
it is not possible to distinguish between a successful inspection and an
inspection not performed.
\end{quote}

Clear statements of limitation are essential for transparency and
reliable application of business insight. They tell the user not only
what the artefact means, but when it becomes dangerous to rely on it.

This is where descriptions protect users from misuse. They do not merely
explain what an artefact means in normal cases; they mark the boundary
beyond which that meaning breaks down.

\section{Keys answer: how does this
relate?}\label{keys-answer-how-does-this-relate}

Database keys include primary keys, foreign keys, and unique keys. They
are also metadata.

Keys stop records from becoming lost.

As explained in
\hyperref[docs-creating-information-mapping-the-data-world]{Mapping the
data world}, the primary key is what links a database record to its
real-world business entity. For the user, they are the anchor between a
table and the entity it claims to represent.

Moreover, a business domain is rarely three simple tables. It is usually
a swamp of entities, events, statuses, and inferred relationships.
Without keys, the user may see the records but have no reliable way to
navigate them.

In this sense, keys are a map for the warehouse. They show which row
represents which business entity, which records belong together, and how
one fragment of reality relates to another.

Keys also recognise that relationships are part of meaning. The
relationship between one table and another tells the user something
about both tables that neither table can say in isolation.

For example, a foreign key between \texttt{Cake.Sales} and
\texttt{Cake.RefProfitability} immediately signals that some sales may
be profitable and some may not.

Each organisation implements keys differently. Some store them only in
conceptual diagrams or third-party tools. The disadvantage is that they
cannot be easily read or queried.

The best approach is to house them as standalone tables in the
warehouse. This makes them visible and open to querying like every other
table in the warehouse.

Ultimately, keys are not just technical constraints. They are statements
of identity and relationship.

This makes them a non-negotiable part of metadata.

\section{Metadata should be data}\label{metadata-should-be-data}

Data without metadata is data without context.

The primary purpose of metadata is not governance, tagging,
discoverability, or data exchange. These are all important, but they are
secondary effects.

Metadata should be understood in the context of the data engineer's
task: shaping data in light of business intent. From this perspective,
the primary purpose of metadata is to add interpretative context that
aligns the data world to its underlying business reality.

Without this context, the data engineer is abandoning the user to guess.

Finally, whether metadata describes names, descriptions, or keys, it
should be stored and made available as data.

Instead of recording metadata only in diagrams, specialist tools, or
database constructs, it should also be available as tables. Treating
metadata as data enhances the ability to distribute it, automate it, and
surface it to appropriate tools through APIs.

\section{Key ideas}\label{key-ideas-9}

\begin{quote}
\textbf{Key ideas.}

Without metadata, everyone is guessing.

Metadata gives data its interpretative context.

Without good names, artefacts become unrecognisable.

Without descriptions, artefacts become incomprehensible.

Without keys, artefacts become lost.

Metadata should be stored and treated as data.
\end{quote}

\cleardoublepage

\chapter{Three approaches to data
quality}\label{docs-quality-reliability-three-approaches-to-data-quality}

\section{Quality means fitness for
intent}\label{quality-means-fitness-for-intent}

Most people see data quality issues as failures of data capture to
reflect reality.

That view is useful, but incomplete.

In practice, errors in the database are rarely neutral. They reflect the
incentives and pressures of the process that produced the data. In
general, errors tend to be biased toward the advantage of the person
providing the information, or toward the path of least resistance in the
business process.

Consider personal income.

Both government social services and the tax office may be concerned with
underestimation of income, but the direction and consequence of error
differ. For tax, a taxpayer is unlikely to complain if they accidentally
report less income, because they may receive a higher tax return or pay
less tax. For social services, the same person is more likely to notice
an underestimate, because they may receive less support than expected.

A bank credit rater faces a different concern again. It may care more
about overestimates of income, because overestimated income may lead to
excessive lending risk.

The same data element---personal income---therefore has different
quality implications depending on the business intent. Social services
may actively manage underestimates to meet legislative requirements.
Credit raters may actively manage overestimates to reduce risk. The tax
office may design its controls around yet another risk profile.

All three may prefer perfect accuracy. But none takes the same approach
when under constraint.

\begin{quote}
There is no accuracy without margin of error, no margin of error without
a risk threshold, and no risk threshold without an articulation of
business intent.
\end{quote}

Consequently, data quality is not merely data versus reality. It is data
versus business intent.

A dogmatic insistence on reflecting the real world perfectly is costly
and often impractical. It may also conflict with privacy, timeliness,
proportionality, or operational burden. The data engineer's job is not
to pursue perfect data in the abstract, but to help the business decide
what quality means for the decision at hand.

With that in mind, common data quality issues arise because digital
systems and business processes were not built to honour the full intent
of the business. This may be due to poor design, such as missing data
type constraints. It may be due to bias in the business process, such as
different incentives for reporting income. It may also be due to a
mismatch of purpose, such as operational forms captured for one
regulatory purpose but later analysed for another.

The data engineer therefore needs a practical set of interventions. This
chapter introduces thre broad approaches, with further methods for
different scenarios.

{\def\LTcaptype{none} % do not increment counter
\begin{longtable}[]{@{}ll@{}}
\toprule\noalign{}
Situation & Approach \\
\midrule\noalign{}
\endhead
\bottomrule\noalign{}
\endlastfoot
Business judgement is needed & Human curation \\
Business intent can be formalised & Precise rules \\
Business intent can only be approximated & Fuzzy logic \\
\end{longtable}
}

Each approach contains recurring methods.

\textbf{Human curation}

{\def\LTcaptype{none} % do not increment counter
\begin{longtable}[]{@{}
  >{\raggedright\arraybackslash}p{(\linewidth - 2\tabcolsep) * \real{0.5000}}
  >{\raggedright\arraybackslash}p{(\linewidth - 2\tabcolsep) * \real{0.5000}}@{}}
\toprule\noalign{}
\begin{minipage}[b]{\linewidth}\raggedright
Method
\end{minipage} & \begin{minipage}[b]{\linewidth}\raggedright
Use when...
\end{minipage} \\
\midrule\noalign{}
\endhead
\bottomrule\noalign{}
\endlastfoot
Data annotation & Records need to be classified, mapped, enriched, or
corrected by a business expert. \\
Applying assumptions & A practical assumption can handle most cases, but
violations need expert review. \\
Data quality reports & Issues are frequent, numerous, or easier to
remediate in bulk. \\
\end{longtable}
}

\textbf{Precise rules}

{\def\LTcaptype{none} % do not increment counter
\begin{longtable}[]{@{}
  >{\raggedright\arraybackslash}p{(\linewidth - 2\tabcolsep) * \real{0.5000}}
  >{\raggedright\arraybackslash}p{(\linewidth - 2\tabcolsep) * \real{0.5000}}@{}}
\toprule\noalign{}
\begin{minipage}[b]{\linewidth}\raggedright
Method
\end{minipage} & \begin{minipage}[b]{\linewidth}\raggedright
Use when...
\end{minipage} \\
\midrule\noalign{}
\endhead
\bottomrule\noalign{}
\endlastfoot
Defining analytical concepts & Operational details need to be turned
into usable business concepts, such as good/bad, milestones, or
conformed categories. \\
Defining primary keys & Rows need a clear identity that connects them to
real-world entities. \\
Defining the primary record & Multiple records appear to refer to one
underlying entity. \\
Defining relationships & Analytically important relationships were not
recorded by the source system. \\
\end{longtable}
}

\textbf{Fuzzy logic}

Fuzzy logic is handled through iterative approximation.

{\def\LTcaptype{none} % do not increment counter
\begin{longtable}[]{@{}
  >{\raggedright\arraybackslash}p{(\linewidth - 2\tabcolsep) * \real{0.5000}}
  >{\raggedright\arraybackslash}p{(\linewidth - 2\tabcolsep) * \real{0.5000}}@{}}
\toprule\noalign{}
\begin{minipage}[b]{\linewidth}\raggedright
Step
\end{minipage} & \begin{minipage}[b]{\linewidth}\raggedright
Question
\end{minipage} \\
\midrule\noalign{}
\endhead
\bottomrule\noalign{}
\endlastfoot
Loose--tight iteration & Can the pattern be tuned to reduce both false
matches and false rejects? \\
Random validation & Does the pattern hold up when checked against
unbiased samples? \\
Monitoring for drift & Does the pattern continue to behave over time? \\
\end{longtable}
}

One theme runs through all three approaches: assumptions must be
monitored. If a data engineer closes a quality gap by applying
judgement, rules, or approximation, the data product should also monitor
whether that intervention continues to behave as intended.

This is covered in more detail in
\hyperref[docs-quality-reliability-tests-and-assumptions]{Tests and
assumptions}.

\section{Human curation: when judgement is
needed}\label{human-curation-when-judgement-is-needed}

If data quality issues arise from a gap between recorded data and
business intent, one simple remedy is to allow business experts to
intervene directly.

Human curation closes the quality gap by allowing business judgement to
enter the data product. Its main monitoring need is to detect new
records that require human attention.

It is appropriate when business judgement is needed and the data
engineer should not pretend the answer can be fully inferred by the
system.

This section covers three scenarios where human curation is appropriate:

\begin{enumerate}
\def\labelenumi{\arabic{enumi}.}
\tightlist
\item
  Data annotation
\item
  Applying assumptions
\item
  Data quality report
\end{enumerate}

\subsection{Data annotation}\label{data-annotation}

Data annotation is useful when records need to be classified, mapped,
enriched, or corrected by a business expert.

For example, a digital system may record store locations as addresses.
These addresses may be sufficient for operational use, but not enough
for regional reporting. A business expert may later annotate each store
with its sales region.

This annotation occurs in the data pipeline, not the source system, and
is commonly known as reference data management.

\textbf{Example structure of
\texttt{\textquotesingle{}Store\textquotesingle{}}}

{\def\LTcaptype{none} % do not increment counter
\begin{longtable}[]{@{}lll@{}}
\toprule\noalign{}
Store ID & Store address & Sales region (annotated) \\
\midrule\noalign{}
\endhead
\bottomrule\noalign{}
\endlastfoot
S001 & 12 Smith St, Northvale & North \\
S002 & 88 Market Rd, Southport & South \\
S003 & 45 Central Ave, Westhaven & West \\
S004 & 9 Harbour Rd, Eastbank & New store not annotated \\
\end{longtable}
}

The value \texttt{New\ store\ not\ annotated} flags the need for human
attention.

Another example is harmonising employee identities across systems. In
large organisations, the same employee may appear under different
accounts for different business processes. A business expert who
understands the enterprise view may annotate these accounts to map them
to a canonical representation. This is known as master data management.

Human annotation extends beyond reference and master data. It applies
wherever a human can interpret records outside the operational context.

There are two legitimate views of data: one fast-paced and operational,
suited to digital systems; the other reflective and analytical, suited
to business interpretation. Both may require human input. It is not
reasonable to expect operational systems to carry the full load of
analytical meaning.

Traditionally, data warehousing has treated human curation as an
exception. But if human editing can be implemented transparently, with
control and auditability, it should be treated as a standard business
tool.

Whether the curation concerns reference data, master data, or other
forms of annotation, the data engineer must monitor for incoming records
that require curation. Examples include uncategorised store locations,
unmapped employee accounts, or new values that do not yet belong to a
known category.

\subsection{Applying assumptions}\label{applying-assumptions}

Applying assumptions is useful when a practical assumption can handle
most cases, but violations need expert review.

Consider the recording of dates. If users manually enter dates, they may
accidentally type \texttt{2300} instead of \texttt{2030}. If the digital
system lacks validation, such errors can end up in the database. Even a
single mistake can distort a line chart or skew time-based analysis.

One way to address this is to \emph{assume} that future dates must be
within 50 years. Dates outside this range are treated as invalid and
converted to blank. Offending records should also be flagged so that a
business expert can review and correct the source data.

Before applying the assumption, the data may look like this.

\textbf{Example structure before flagging}

{\def\LTcaptype{none} % do not increment counter
\begin{longtable}[]{@{}rl@{}}
\toprule\noalign{}
Record ID & Event date \\
\midrule\noalign{}
\endhead
\bottomrule\noalign{}
\endlastfoot
1 & 2030-06-15 \\
2 & 2300-06-15 \\
3 & 2025-04-02 \\
\end{longtable}
}

After applying the assumption, the invalid date is blanked out and the
row is flagged.

\textbf{Example structure after flagging}

{\def\LTcaptype{none} % do not increment counter
\begin{longtable}[]{@{}rll@{}}
\toprule\noalign{}
Record ID & Event date & Is flagged for invalid date \\
\midrule\noalign{}
\endhead
\bottomrule\noalign{}
\endlastfoot
1 & 2030-06-15 & false \\
2 & & true \\
3 & 2025-04-02 & false \\
\end{longtable}
}

This preserves analytical usefulness while making the intervention
visible. The invalid date no longer distorts time-based analysis, but
the issue remains available for remediation.

This approach works best when violations are rare and the business
expert can correct the issue before the next batch load. It is suitable
for irregular data quality issues that do not require bulk remediation.

If data quality issues are frequent or systemic, a data quality report
may be more appropriate.

\subsection{Data quality report}\label{data-quality-report}

Data quality reports are suitable when issues are frequent, numerous, or
easier to treat in bulk on a periodic basis.

An effective implementation is to create a
\hyperref[docs-presenting-insights-dimensional-modelling-for-ux]{combination
or choice dimension}. A combination dimension may include one column per
issue and describe each transaction. For example, a dimension called
\texttt{\textquotesingle{}Sales\ data\ quality\textquotesingle{}} might
include \texttt{{[}Is\ missing\ sales\ amount{]}},
\texttt{{[}Is\ invalid\ sales\ date{]}},
\texttt{{[}Is\ unknown\ customer{]}}, and so on.

\textbf{Example structure of
\texttt{\textquotesingle{}Sales\ data\ quality\textquotesingle{}}}

{\def\LTcaptype{none} % do not increment counter
\begin{longtable}[]{@{}
  >{\raggedleft\arraybackslash}p{(\linewidth - 8\tabcolsep) * \real{0.2500}}
  >{\raggedright\arraybackslash}p{(\linewidth - 8\tabcolsep) * \real{0.1875}}
  >{\raggedright\arraybackslash}p{(\linewidth - 8\tabcolsep) * \real{0.1875}}
  >{\raggedright\arraybackslash}p{(\linewidth - 8\tabcolsep) * \real{0.1875}}
  >{\raggedright\arraybackslash}p{(\linewidth - 8\tabcolsep) * \real{0.1875}}@{}}
\toprule\noalign{}
\begin{minipage}[b]{\linewidth}\raggedleft
Sales data quality SK
\end{minipage} & \begin{minipage}[b]{\linewidth}\raggedright
Is missing sales amount
\end{minipage} & \begin{minipage}[b]{\linewidth}\raggedright
Is invalid sales date
\end{minipage} & \begin{minipage}[b]{\linewidth}\raggedright
Is unknown customer
\end{minipage} & \begin{minipage}[b]{\linewidth}\raggedright
Data quality category
\end{minipage} \\
\midrule\noalign{}
\endhead
\bottomrule\noalign{}
\endlastfoot
1 & false & false & false & No known issue \\
2 & true & false & false & Missing amount \\
3 & false & true & false & Invalid date \\
4 & false & false & true & Unknown customer \\
... & ... & ... & ... & ... \\
8 & true & true & true & Multiple issues \\
\end{longtable}
}

The fact table can then refer to the data quality dimension.

\textbf{Example structure of
\texttt{\textquotesingle{}Sales\textquotesingle{}}}

{\def\LTcaptype{none} % do not increment counter
\begin{longtable}[]{@{}
  >{\raggedright\arraybackslash}p{(\linewidth - 8\tabcolsep) * \real{0.1765}}
  >{\raggedright\arraybackslash}p{(\linewidth - 8\tabcolsep) * \real{0.1765}}
  >{\raggedright\arraybackslash}p{(\linewidth - 8\tabcolsep) * \real{0.1765}}
  >{\raggedleft\arraybackslash}p{(\linewidth - 8\tabcolsep) * \real{0.2353}}
  >{\raggedleft\arraybackslash}p{(\linewidth - 8\tabcolsep) * \real{0.2353}}@{}}
\toprule\noalign{}
\begin{minipage}[b]{\linewidth}\raggedright
Sales ID
\end{minipage} & \begin{minipage}[b]{\linewidth}\raggedright
Sales date
\end{minipage} & \begin{minipage}[b]{\linewidth}\raggedright
Customer ID
\end{minipage} & \begin{minipage}[b]{\linewidth}\raggedleft
Sales amount
\end{minipage} & \begin{minipage}[b]{\linewidth}\raggedleft
Sales data quality SK
\end{minipage} \\
\midrule\noalign{}
\endhead
\bottomrule\noalign{}
\endlastfoot
S1001 & 2025-04-01 & C001 & 120.00 & 1 \\
S1002 & 2300-04-01 & C002 & 95.00 & 3 \\
S1003 & 2025-04-03 & C999 & 80.00 & 4 \\
S1004 & 2025-04-04 & C004 & & 2 \\
\end{longtable}
}

Coupled with the ID dimension, this makes it easy to build a Power BI
report that surfaces transactions requiring remediation by a business
expert.

A data quality dimension also supports systematic analysis by enabling
statistics on the number and type of issues.

Reports are only effective if they embed in the user's workflow. A
decontextualised report becomes forgotten and unused. The same applies
to data quality reports. The business should have a workflow trigger
that prompts engagement with the report.

\section{Precise rules: when intent can be
formalised}\label{precise-rules-when-intent-can-be-formalised}

Translating business knowledge into precise rules allows it to be
automated and reduces reliance on human intervention.

One way to look at these rules is to see them as ``fixing issues'' in
the data. A more useful way is to see them as rules that bridge recorded
data to business intent.

Precise rules close data quality gaps by formalising business intent
into repeatable logic. Their main monitoring need is to detect rule
violations, edge cases, and changes in the business process that make
the rule unsafe.

Precise rules are appropriate when the business can say what should
happen clearly enough for the data engineer to implement it.

This section covers four recurring types of precise rule:

\begin{enumerate}
\def\labelenumi{\arabic{enumi}.}
\tightlist
\item
  Defining analytical concepts
\item
  Defining primary keys
\item
  Defining the primary record
\item
  Defining relationships
\end{enumerate}

\subsection{Defining analytical
concepts}\label{defining-analytical-concepts}

Defining analytical concepts closes the quality gap between operational
detail and business interpretation.

Since business insight is information analysed in light of business
intent, a direct way of improving data quality is to define the
analytical concept that the business needs.

Defining analytical concepts takes
\hyperref[docs-foundations-data-and-organisations]{leadership and
negotiation skill}. The data engineer is often well placed to broker
between stakeholder groups by experimenting with the data and visually
communicating possible outcomes.

\subsubsection{Good and bad entities}\label{good-and-bad-entities-1}

Good and bad entity definitions close the quality gap between detailed
operational outcomes and the business's need for a usable judgement.

A business process may collect detailed information such as audit
results, inspection results, incidents, sentiment, or sales volume.
However, these details may not yet be formalised into a systematic
judgement.

Sometimes numeric values such as star ratings need to be translated into
good and bad using thresholds. More commonly, a business process records
information at a lower grain that needs to be rolled up to the entity of
interest.

For example, a manufacturing batch may have multiple quality-control
criteria. The business may decide that failing one critical criterion,
or multiple non-critical criteria, defines failure at the batch level.

\textbf{Example structure of
\texttt{\textquotesingle{}Batch\ quality\ criterion\textquotesingle{}}}

{\def\LTcaptype{none} % do not increment counter
\begin{longtable}[]{@{}
  >{\raggedright\arraybackslash}p{(\linewidth - 6\tabcolsep) * \real{0.2500}}
  >{\raggedright\arraybackslash}p{(\linewidth - 6\tabcolsep) * \real{0.2500}}
  >{\raggedright\arraybackslash}p{(\linewidth - 6\tabcolsep) * \real{0.2500}}
  >{\raggedright\arraybackslash}p{(\linewidth - 6\tabcolsep) * \real{0.2500}}@{}}
\toprule\noalign{}
\begin{minipage}[b]{\linewidth}\raggedright
Batch ID
\end{minipage} & \begin{minipage}[b]{\linewidth}\raggedright
Quality criterion
\end{minipage} & \begin{minipage}[b]{\linewidth}\raggedright
Quality result
\end{minipage} & \begin{minipage}[b]{\linewidth}\raggedright
Is critical criterion
\end{minipage} \\
\midrule\noalign{}
\endhead
\bottomrule\noalign{}
\endlastfoot
B1001 & Temperature control & Pass & true \\
B1001 & Packaging integrity & Pass & false \\
B1001 & Labelling accuracy & Fail & false \\
B1002 & Temperature control & Fail & true \\
B1002 & Packaging integrity & Pass & false \\
\end{longtable}
}

The analytical concept can then be expressed at the batch level.

\textbf{Example structure of
\texttt{\textquotesingle{}Batch\ quality\ outcome\textquotesingle{}}}

{\def\LTcaptype{none} % do not increment counter
\begin{longtable}[]{@{}
  >{\raggedright\arraybackslash}p{(\linewidth - 6\tabcolsep) * \real{0.2143}}
  >{\raggedleft\arraybackslash}p{(\linewidth - 6\tabcolsep) * \real{0.2857}}
  >{\raggedleft\arraybackslash}p{(\linewidth - 6\tabcolsep) * \real{0.2857}}
  >{\raggedright\arraybackslash}p{(\linewidth - 6\tabcolsep) * \real{0.2143}}@{}}
\toprule\noalign{}
\begin{minipage}[b]{\linewidth}\raggedright
Batch ID
\end{minipage} & \begin{minipage}[b]{\linewidth}\raggedleft
Critical criteria failed
\end{minipage} & \begin{minipage}[b]{\linewidth}\raggedleft
Non-critical criteria failed
\end{minipage} & \begin{minipage}[b]{\linewidth}\raggedright
Batch quality outcome
\end{minipage} \\
\midrule\noalign{}
\endhead
\bottomrule\noalign{}
\endlastfoot
B1001 & 0 & 1 & Pass \\
B1002 & 1 & 0 & Fail \\
\end{longtable}
}

This gives the user an immediate view of the process at the level of
business judgement.

It is not reasonable to expect operational systems to always define
these concepts. Their primary job is to execute workflows while
preserving records. The data engineer plays a role by adding the
analytical lens.

This implementation is studied in greater depth in
\hyperref[docs-creating-information-entity-processing]{Entity
processing}.

\subsubsection{Milestones}\label{milestones}

Milestones close the quality gap between messy workflow events and
meaningful process control.

A business process may have many detailed steps. Some may loop back.
Some may start and stop multiple times. Some may occur repeatedly for
the same business entity. This mass of events can bury insight in
operational detail.

A useful analytical view is to define a small number of major
checkpoints that can be used to measure performance. These checkpoints
are milestones.

The difficulty is choosing the right level of abstraction. A milestone
can be too narrow, in which case it merely reproduces operational noise.
It can also be too broad, in which case it merges several different
responsibilities into one measure.

A practical rule is that each milestone should have one main control
point or responsible owner. If a milestone mixes responsibilities, then
a delay becomes difficult to interpret. The business can see that
something is late, but not who is able to act on it.

For example, an HR recruitment process may contain dozens of operational
steps: position approved, advertisement drafted, advertisement cleared,
advertisement published, applications received, applications closed,
eligibility checked, longlist completed, shortlist completed, interviews
scheduled, interviews completed, referee checks completed, delegate
approval received, offer issued, offer accepted, and commencement
recorded.

In practice, such an operational process may contain more than forty
recorded steps. That level of detail is necessary for workflow
execution, but too noisy for business decision-making.

The solution is to define key milestones.

\textbf{Example structure of
\texttt{\textquotesingle{}Recruitment\ event\textquotesingle{}}}

{\def\LTcaptype{none} % do not increment counter
\begin{longtable}[]{@{}lll@{}}
\toprule\noalign{}
Recruitment ID & Event datetime & Event type \\
\midrule\noalign{}
\endhead
\bottomrule\noalign{}
\endlastfoot
R1001 & 2025-04-01 09:00 & Advertisement drafted \\
R1001 & 2025-04-03 11:30 & Advertisement cleared \\
R1001 & 2025-04-04 09:00 & Advertisement published \\
R1001 & 2025-04-18 17:00 & Applications closed \\
R1001 & 2025-04-21 10:00 & Eligibility checked \\
R1001 & 2025-04-23 16:00 & Shortlist completed \\
R1001 & 2025-04-29 09:00 & Interviews scheduled \\
R1001 & 2025-05-02 16:30 & Interviews completed \\
R1001 & 2025-05-07 10:00 & Delegate approved \\
R1001 & 2025-05-09 14:00 & Offer accepted \\
R1002 & 2025-04-03 09:00 & Advertisement drafted \\
R1002 & 2025-04-05 10:00 & Advertisement returned for changes \\
R1002 & 2025-04-08 13:00 & Advertisement cleared \\
R1002 & 2025-04-09 09:00 & Advertisement published \\
R1002 & 2025-04-23 17:00 & Applications closed \\
R1002 & 2025-04-29 15:00 & Eligibility checked \\
R1002 & 2025-05-05 15:00 & Shortlist completed \\
R1002 & 2025-05-14 16:00 & Interviews completed \\
\end{longtable}
}

The data engineer can turn the event history into a milestone fragment.

\textbf{Example structure of
\texttt{\textquotesingle{}Recruitment\ milestone\textquotesingle{}}}

{\def\LTcaptype{none} % do not increment counter
\begin{longtable}[]{@{}
  >{\raggedright\arraybackslash}p{(\linewidth - 10\tabcolsep) * \real{0.1500}}
  >{\raggedright\arraybackslash}p{(\linewidth - 10\tabcolsep) * \real{0.1500}}
  >{\raggedright\arraybackslash}p{(\linewidth - 10\tabcolsep) * \real{0.1500}}
  >{\raggedright\arraybackslash}p{(\linewidth - 10\tabcolsep) * \real{0.1500}}
  >{\raggedleft\arraybackslash}p{(\linewidth - 10\tabcolsep) * \real{0.2000}}
  >{\raggedleft\arraybackslash}p{(\linewidth - 10\tabcolsep) * \real{0.2000}}@{}}
\toprule\noalign{}
\begin{minipage}[b]{\linewidth}\raggedright
Recruitment ID
\end{minipage} & \begin{minipage}[b]{\linewidth}\raggedright
Advertisement published datetime
\end{minipage} & \begin{minipage}[b]{\linewidth}\raggedright
Interview completed datetime
\end{minipage} & \begin{minipage}[b]{\linewidth}\raggedright
Offer accepted datetime
\end{minipage} & \begin{minipage}[b]{\linewidth}\raggedleft
Days from advertisement to interview
\end{minipage} & \begin{minipage}[b]{\linewidth}\raggedleft
Days from interview to offer
\end{minipage} \\
\midrule\noalign{}
\endhead
\bottomrule\noalign{}
\endlastfoot
R1001 & 2025-04-04 09:00 & 2025-05-02 16:30 & 2025-05-09 14:00 & 28.3 &
6.9 \\
R1002 & 2025-04-09 09:00 & 2025-05-14 16:00 & & 35.3 & \\
\end{longtable}
}

The milestone fragment does not replace the event history. It gives the
business a stable analytical view of the process.

In this example, the three milestones correspond to different control
points. Advertisement publication depends on the area preparing and
clearing the advertisement. Interview completion depends on the
assessment process. Offer acceptance depends on post-interview approval
and candidate response. If these were collapsed into one broad measure
such as \texttt{{[}Days\ to\ recruit{]}}, the business would know the
process was slow but not where the delay occurred.

To deal with loops and repeated events, the data engineer can define the
earliest or latest relevant occurrence. For example,
\texttt{{[}Advertisement\ published\ datetime{]}} may use the first
publication event, while \texttt{{[}Offer\ accepted\ datetime{]}} may
use the final accepted offer event.

This pattern is studied in
\hyperref[docs-creating-information-meaningful-fragments]{Meaningful
fragments}.

\subsubsection{Conformed dimensions}\label{conformed-dimensions}

Conformed dimensions close the quality gap between local system
categories and enterprise intent.

A large organisation may have many processes, and many of them share
similar concepts under different names or codes. A conformed dimension
reconciles these local concepts into a shared view.

For example, one system may store country using two-letter country
codes, while another system may store country using its own operational
codes.

\textbf{Example structure of \texttt{Sales.Customer}}

{\def\LTcaptype{none} % do not increment counter
\begin{longtable}[]{@{}ll@{}}
\toprule\noalign{}
Customer ID & Customer country code \\
\midrule\noalign{}
\endhead
\bottomrule\noalign{}
\endlastfoot
C001 & AU \\
C002 & NZ \\
C003 & US \\
\end{longtable}
}

\textbf{Example structure of \texttt{Shipping.Delivery}}

{\def\LTcaptype{none} % do not increment counter
\begin{longtable}[]{@{}ll@{}}
\toprule\noalign{}
Delivery ID & Delivery country code \\
\midrule\noalign{}
\endhead
\bottomrule\noalign{}
\endlastfoot
D001 & AUS \\
D002 & NZL \\
D003 & USA \\
\end{longtable}
}

Both systems refer to the same business concept: country. However, the
codes are not the same. If the data engineer leaves them as separate
local values, enterprise reporting becomes harder than it needs to be.

The data engineer can create a standard country reference table.

\textbf{Example structure of
\texttt{\textquotesingle{}Country\textquotesingle{}}}

{\def\LTcaptype{none} % do not increment counter
\begin{longtable}[]{@{}rlll@{}}
\toprule\noalign{}
Country SK & Country name & ISO alpha-2 code & ISO alpha-3 code \\
\midrule\noalign{}
\endhead
\bottomrule\noalign{}
\endlastfoot
1 & Australia & AU & AUS \\
2 & New Zealand & NZ & NZL \\
3 & United States & US & USA \\
\end{longtable}
}

The local system codes can then be mapped to the conformed dimension.

\textbf{Example structure of
\texttt{\textquotesingle{}Country\ code\ map\textquotesingle{}}}

{\def\LTcaptype{none} % do not increment counter
\begin{longtable}[]{@{}llr@{}}
\toprule\noalign{}
Source system & Source country code & Country SK \\
\midrule\noalign{}
\endhead
\bottomrule\noalign{}
\endlastfoot
Sales & AU & 1 \\
Sales & NZ & 2 \\
Sales & US & 3 \\
Shipping & AUS & 1 \\
Shipping & NZL & 2 \\
Shipping & USA & 3 \\
\end{longtable}
}

This allows sales and shipping data to be analysed through the same
country dimension, even though the source systems use different local
codes.

When done appropriately, a conformed view can empower decision-makers at
the most senior levels of the organisation.

Conformed dimensions are studied in
\hyperref[docs-creating-information-reference-data]{Reference data}.

\subsection{Defining primary keys}\label{defining-primary-keys}

Defining primary keys closes the quality gap between database rows and
real-world entities by making their link explicit.

Primary keys serve as the
\hyperref[docs-creating-information-mapping-the-data-world]{link between
data records and their counterparts in the real world}. Unfortunately,
some business processes do not rigorously define primary keys.

This can lead to slippery definitions in the database. In other cases,
the primary key is defined at the application layer but remains
invisible at the data layer, leaving the business user without a clear
way to interpret the data.

The data engineer should establish the primary key where it is not
meaningfully defined.

The primary key articulates the data engineer's view of how the data row
should be interpreted as a business entity.

Most primary keys can be traced by examining how a business process
creates, retrieves, and updates a data record. Three common patterns are
sequence numbers, version numbers, and temporality.

\subsubsection{Sequence numbers}\label{sequence-numbers}

Sequence numbers close the quality gap where the source system stores a
list, but the data product needs stable business identity for each row.

Sequence numbers are useful when a header has multiple detail rows and
the detail table lacks a meaningful primary key.

For example, a sales order may have multiple items. The order is stored
in \texttt{Sales}, with \texttt{{[}Order\ number{]}} as the primary key.
The items are stored in \texttt{SalesItems}, which may be a
miscellaneous list without a meaningful key.

If each product can appear only once, then \texttt{{[}Order\ number{]}}
and \texttt{{[}Product\ ID{]}} may suffice as a primary key. But if a
product can appear more than once, this fails.

\textbf{Example structure of \texttt{SalesItems} before sequence
numbering}

{\def\LTcaptype{none} % do not increment counter
\begin{longtable}[]{@{}llrr@{}}
\toprule\noalign{}
Order number & Product ID & Sales item ID & Quantity \\
\midrule\noalign{}
\endhead
\bottomrule\noalign{}
\endlastfoot
O1001 & P100 & 501 & 1 \\
O1001 & P100 & 502 & 2 \\
O1001 & P200 & 503 & 1 \\
\end{longtable}
}

\texttt{{[}Order\ number{]}}, \texttt{{[}Product\ ID{]}} cannot identify
a row because \texttt{P100} appears twice for the same order.

A simple solution is to treat the list of items as a sequence and create
an artificial column \texttt{{[}Sales\ item\ sequence\ number{]}},
forming the key \texttt{{[}Order\ number{]}},
\texttt{{[}Sales\ item\ sequence\ number{]}}.

\textbf{Example structure of \texttt{SalesItems} after sequence
numbering}

{\def\LTcaptype{none} % do not increment counter
\begin{longtable}[]{@{}
  >{\raggedright\arraybackslash}p{(\linewidth - 8\tabcolsep) * \real{0.1667}}
  >{\raggedleft\arraybackslash}p{(\linewidth - 8\tabcolsep) * \real{0.2222}}
  >{\raggedright\arraybackslash}p{(\linewidth - 8\tabcolsep) * \real{0.1667}}
  >{\raggedleft\arraybackslash}p{(\linewidth - 8\tabcolsep) * \real{0.2222}}
  >{\raggedleft\arraybackslash}p{(\linewidth - 8\tabcolsep) * \real{0.2222}}@{}}
\toprule\noalign{}
\begin{minipage}[b]{\linewidth}\raggedright
Order number
\end{minipage} & \begin{minipage}[b]{\linewidth}\raggedleft
Sales item sequence number
\end{minipage} & \begin{minipage}[b]{\linewidth}\raggedright
Product ID
\end{minipage} & \begin{minipage}[b]{\linewidth}\raggedleft
Sales item ID
\end{minipage} & \begin{minipage}[b]{\linewidth}\raggedleft
Quantity
\end{minipage} \\
\midrule\noalign{}
\endhead
\bottomrule\noalign{}
\endlastfoot
O1001 & 1 & P100 & 501 & 1 \\
O1001 & 2 & P100 & 502 & 2 \\
O1001 & 3 & P200 & 503 & 1 \\
\end{longtable}
}

Sometimes a system will implement a meaningless key such as
\texttt{{[}Sales\ item\ ID{]}}, a simple integer used for UI or database
constraints. Even in these cases, the data engineer may still define a
sequence number and use \texttt{{[}Order\ number{]}},
\texttt{{[}Sales\ item\ sequence\ number{]}} as the business-facing
primary key. \texttt{{[}Sales\ item\ ID{]}} can be retained for joins.

In general, sequence numbers are effective wherever there is a
miscellaneous list of line items within a header. Care must be taken to
ensure the sequence is deterministic by breaking ties---using a
surrogate key like \texttt{{[}Sales\ item\ ID{]}} as a sort order is a
reliable approach.

\subsubsection{Version numbers}\label{version-numbers}

Version numbers close the quality gap where the source system records
change, but does not make the continuity of the entity clear.

Version numbers are useful for
\hyperref[docs-creating-information-entity-processing]{entities whose
changes should be preserved as new versions}.

For example, a customer order may be submitted, revised, and resubmitted
before it is fulfilled. The business may still regard this as the same
order, but each revision changes the content that was true at a
particular point in time.

Digital systems are not always consistent in how they manage versions. A
common but problematic approach is to allow the entity ID to change each
time the record is revised and record the relationship using
\texttt{{[}Previous\ order\ ID{]}} or
\texttt{{[}Superseded\ order\ ID{]}}. This may work for rendering a web
UI but can cause confusion in analysis.

Where versioning is lost or muddled, the data engineer can restore
clarity using a consistent pattern such as \texttt{{[}Entity\ ID{]}},
\texttt{{[}Version\ number{]}}.

The order can be represented with a stable \texttt{{[}Order\ ID{]}} and
a changing \texttt{{[}Order\ version\ number{]}}.

\textbf{Example structure of \texttt{Sales.Order}}

{\def\LTcaptype{none} % do not increment counter
\begin{longtable}[]{@{}llrl@{}}
\toprule\noalign{}
Order ID & Customer ID & Current order version number & Order status \\
\midrule\noalign{}
\endhead
\bottomrule\noalign{}
\endlastfoot
O1001 & C042 & 3 & Submitted \\
O1002 & C087 & 1 & Fulfilled \\
\end{longtable}
}

\textbf{Example structure of \texttt{Sales.OrderVersion}}

{\def\LTcaptype{none} % do not increment counter
\begin{longtable}[]{@{}
  >{\raggedright\arraybackslash}p{(\linewidth - 10\tabcolsep) * \real{0.1429}}
  >{\raggedleft\arraybackslash}p{(\linewidth - 10\tabcolsep) * \real{0.1905}}
  >{\raggedright\arraybackslash}p{(\linewidth - 10\tabcolsep) * \real{0.1429}}
  >{\raggedright\arraybackslash}p{(\linewidth - 10\tabcolsep) * \real{0.1429}}
  >{\raggedleft\arraybackslash}p{(\linewidth - 10\tabcolsep) * \real{0.1905}}
  >{\raggedleft\arraybackslash}p{(\linewidth - 10\tabcolsep) * \real{0.1905}}@{}}
\toprule\noalign{}
\begin{minipage}[b]{\linewidth}\raggedright
Order ID
\end{minipage} & \begin{minipage}[b]{\linewidth}\raggedleft
Order version number
\end{minipage} & \begin{minipage}[b]{\linewidth}\raggedright
Version datetime
\end{minipage} & \begin{minipage}[b]{\linewidth}\raggedright
Product ID
\end{minipage} & \begin{minipage}[b]{\linewidth}\raggedleft
Quantity
\end{minipage} & \begin{minipage}[b]{\linewidth}\raggedleft
Order value
\end{minipage} \\
\midrule\noalign{}
\endhead
\bottomrule\noalign{}
\endlastfoot
O1001 & 1 & 2025-04-01 09:15 & P100 & 10 & 250.00 \\
O1001 & 2 & 2025-04-01 11:40 & P100 & 12 & 300.00 \\
O1001 & 3 & 2025-04-02 08:30 & P200 & 12 & 420.00 \\
O1002 & 1 & 2025-04-03 14:10 & P300 & 5 & 175.00 \\
\end{longtable}
}

The primary key for \texttt{Sales.OrderVersion} is
\texttt{{[}Order\ ID{]}}, \texttt{{[}Order\ version\ number{]}}.

\subsubsection{Temporality}\label{temporality}

Temporality closes the quality gap where the business needs historical
interpretation but the source system only presents the current state.

Temporality refers to tracking changes over time. In data warehousing,
this is commonly handled through Type II tracking.

Some business processes are designed to handle only the current event,
without tracking history. But when the business is interested in change
over time, the
\hyperref[docs-creating-information-entity-tracking]{entity is mutable}.

Mutable entities should usually include a time component in their
business key. The data engineer may define the business key as
\texttt{{[}Employee\ ID{]}}, \texttt{{[}Start\ datetime{]}}, with
\texttt{{[}End\ datetime{]}} marking the end of the validity period.

In implementation, it is necessary to add a surrogate key as well.

\textbf{Example structure of \texttt{Employee.TeamAssignment}}

{\def\LTcaptype{none} % do not increment counter
\begin{longtable}[]{@{}
  >{\raggedleft\arraybackslash}p{(\linewidth - 8\tabcolsep) * \real{0.2500}}
  >{\raggedright\arraybackslash}p{(\linewidth - 8\tabcolsep) * \real{0.1875}}
  >{\raggedright\arraybackslash}p{(\linewidth - 8\tabcolsep) * \real{0.1875}}
  >{\raggedright\arraybackslash}p{(\linewidth - 8\tabcolsep) * \real{0.1875}}
  >{\raggedright\arraybackslash}p{(\linewidth - 8\tabcolsep) * \real{0.1875}}@{}}
\toprule\noalign{}
\begin{minipage}[b]{\linewidth}\raggedleft
Team assignment SK
\end{minipage} & \begin{minipage}[b]{\linewidth}\raggedright
Employee ID
\end{minipage} & \begin{minipage}[b]{\linewidth}\raggedright
Start datetime
\end{minipage} & \begin{minipage}[b]{\linewidth}\raggedright
End datetime
\end{minipage} & \begin{minipage}[b]{\linewidth}\raggedright
Team name
\end{minipage} \\
\midrule\noalign{}
\endhead
\bottomrule\noalign{}
\endlastfoot
1 & E1001 & 2024-01-01 00:00 & 2025-03-15 09:30 & Finance \\
2 & E1001 & 2025-03-15 09:30 & 9999-12-31 00:00 & Analytics \\
3 & E1002 & 2024-07-01 00:00 & 9999-12-31 00:00 & Operations \\
\end{longtable}
}

The business key is \texttt{{[}Employee\ ID{]}},
\texttt{{[}Start\ datetime{]}}. \texttt{{[}End\ datetime{]}} marks the
end of the validity period. \texttt{{[}Team\ assignment\ SK{]}} provides
a simple surrogate key for joins, relationships, and downstream
modelling.

Recovering temporality depends on how history is stored. Sometimes it is
available in audit tables. Sometimes it must be reconstructed with help
from business experts.

\subsection{Defining the primary
record}\label{defining-the-primary-record}

Defining the primary record closes the quality gap between multiple
records and one underlying entity.

This occurs when the database has primary keys, but different primary
keys relate to the same underlying entity. It is common in business
processes that gather decentralised observations.

Consider the example of whale sightings. A database may collect
observations from citizen scientists, stored in
\texttt{Whale.Observation}, with \texttt{{[}User\ ID{]}},
\texttt{{[}Observation\ number{]}} identifying each user's submissions.
A surrogate key such as \texttt{{[}Observation\ ID{]}} may also exist
for system retrieval.

Multiple citizens may report sightings of the same whale. As a result,
multiple \texttt{{[}Observation\ ID{]}} values relate to the same
real-world entity.

\textbf{Example structure of \texttt{Whale.Observation}}

{\def\LTcaptype{none} % do not increment counter
\begin{longtable}[]{@{}
  >{\raggedleft\arraybackslash}p{(\linewidth - 10\tabcolsep) * \real{0.2000}}
  >{\raggedright\arraybackslash}p{(\linewidth - 10\tabcolsep) * \real{0.1500}}
  >{\raggedleft\arraybackslash}p{(\linewidth - 10\tabcolsep) * \real{0.2000}}
  >{\raggedright\arraybackslash}p{(\linewidth - 10\tabcolsep) * \real{0.1500}}
  >{\raggedright\arraybackslash}p{(\linewidth - 10\tabcolsep) * \real{0.1500}}
  >{\raggedright\arraybackslash}p{(\linewidth - 10\tabcolsep) * \real{0.1500}}@{}}
\toprule\noalign{}
\begin{minipage}[b]{\linewidth}\raggedleft
Observation ID
\end{minipage} & \begin{minipage}[b]{\linewidth}\raggedright
User ID
\end{minipage} & \begin{minipage}[b]{\linewidth}\raggedleft
Observation number
\end{minipage} & \begin{minipage}[b]{\linewidth}\raggedright
Species
\end{minipage} & \begin{minipage}[b]{\linewidth}\raggedright
GPS area
\end{minipage} & \begin{minipage}[b]{\linewidth}\raggedright
Observation date
\end{minipage} \\
\midrule\noalign{}
\endhead
\bottomrule\noalign{}
\endlastfoot
101 & U01 & 1 & Blue whale & A17 & 2025-05-01 \\
102 & U02 & 1 & Blue whale & A17 & 2025-05-01 \\
103 & U03 & 1 & Humpback whale & B04 & 2025-05-01 \\
\end{longtable}
}

Suppose the business has a rule such as:

\begin{quote}
A whale species at one GPS proximity should only appear once within a
day.
\end{quote}

This rule can be used to group observations and identify a
representative record.

The data engineer can express this by identifying the primary
observation. For example, the pipeline may select the first observation
of the day at a location and store it as
\texttt{{[}Primary\ observation\ ID{]}} in a table called
\texttt{Whale.PrimaryObservation}.

\textbf{Example structure of \texttt{Whale.PrimaryObservation}}

{\def\LTcaptype{none} % do not increment counter
\begin{longtable}[]{@{}rr@{}}
\toprule\noalign{}
Observation ID & Primary observation ID \\
\midrule\noalign{}
\endhead
\bottomrule\noalign{}
\endlastfoot
101 & 101 \\
102 & 101 \\
103 & 103 \\
\end{longtable}
}

\texttt{{[}Observation\ ID{]}} is the key for the original record.
\texttt{{[}Primary\ observation\ ID{]}} identifies the chosen
representative.

Subsequent transformation can focus on the
\texttt{{[}Primary\ observation\ ID{]}} grain. The same key can also be
used in Power BI to return the true count of whale sightings.

Care must be taken to resolve race conditions. If two users submit
observations at the same time, only one record should be selected as
primary. This can be resolved deterministically using a surrogate key
such as \texttt{{[}Observation\ ID{]}}.

\subsection{Defining relationships}\label{defining-relationships}

Defining relationships closes the quality gap where analytically
important relationships were not recorded by the source system.

Digital systems are usually good at recording relationships needed by
operational workflows. They are less reliable at recording relationships
that are useful for analysis.

This is especially common when the system is designed for operational
efficiency rather than analytical understanding.

The data engineer can sometimes close this quality gap by defining
relationships that the source system did not record.

Two useful techniques are:

\begin{itemize}
\tightlist
\item
  nearest temporal joins;
\item
  mapping tables.
\end{itemize}

\subsubsection{Nearest temporal join}\label{nearest-temporal-join}

Nearest temporal joins close the quality gap between related events
whose relationship was not explicitly recorded.

Sometimes two sets of events are related, but the relationship is not
recorded as a database key.

If business knowledge suggests that one set of events is expected to
precede another, the relationship can sometimes be recovered by
identifying the nearest preceding event. This technique is known as a
nearest temporal join.

Consider a building fire-safety process.

Building owners submit fire-safety certificates periodically, or when
the building's fire-safety arrangements change. Fire-safety officers
conduct inspections separately. Some inspections may be scheduled.
Others may be triggered by complaints, incidents, risk reviews, or
follow-up action.

The two processes are related, but they are not one-to-one. One
certificate may be followed by several inspections. Several certificates
may be submitted before the next inspection. Some inspections may have
no recent certificate at all.

The inspection system records the building and inspection date, but it
does not record which certificate was in force at the time of
inspection. This makes the relationship analytically important but
structurally missing.

The fire-safety certificates are stored in
\texttt{FireSafety.Certificate}.

\textbf{Example structure of \texttt{FireSafety.Certificate}}

{\def\LTcaptype{none} % do not increment counter
\begin{longtable}[]{@{}
  >{\raggedleft\arraybackslash}p{(\linewidth - 8\tabcolsep) * \real{0.2500}}
  >{\raggedright\arraybackslash}p{(\linewidth - 8\tabcolsep) * \real{0.1875}}
  >{\raggedright\arraybackslash}p{(\linewidth - 8\tabcolsep) * \real{0.1875}}
  >{\raggedright\arraybackslash}p{(\linewidth - 8\tabcolsep) * \real{0.1875}}
  >{\raggedright\arraybackslash}p{(\linewidth - 8\tabcolsep) * \real{0.1875}}@{}}
\toprule\noalign{}
\begin{minipage}[b]{\linewidth}\raggedleft
Fire safety certificate SK
\end{minipage} & \begin{minipage}[b]{\linewidth}\raggedright
Building ID
\end{minipage} & \begin{minipage}[b]{\linewidth}\raggedright
Certificate submitted datetime
\end{minipage} & \begin{minipage}[b]{\linewidth}\raggedright
Certificate type
\end{minipage} & \begin{minipage}[b]{\linewidth}\raggedright
Declared risk level
\end{minipage} \\
\midrule\noalign{}
\endhead
\bottomrule\noalign{}
\endlastfoot
1 & B001 & 2025-04-01 09:00 & Annual certificate & Low \\
2 & B001 & 2025-05-01 09:00 & Revised certificate & Medium \\
3 & B002 & 2025-04-10 14:30 & High-risk occupancy & High \\
4 & B002 & 2025-04-18 16:00 & Remediation certificate & Medium \\
\end{longtable}
}

The inspections are stored separately in \texttt{FireSafety.Inspection}.

\textbf{Example structure of \texttt{FireSafety.Inspection}}

{\def\LTcaptype{none} % do not increment counter
\begin{longtable}[]{@{}
  >{\raggedleft\arraybackslash}p{(\linewidth - 8\tabcolsep) * \real{0.2500}}
  >{\raggedright\arraybackslash}p{(\linewidth - 8\tabcolsep) * \real{0.1875}}
  >{\raggedright\arraybackslash}p{(\linewidth - 8\tabcolsep) * \real{0.1875}}
  >{\raggedright\arraybackslash}p{(\linewidth - 8\tabcolsep) * \real{0.1875}}
  >{\raggedright\arraybackslash}p{(\linewidth - 8\tabcolsep) * \real{0.1875}}@{}}
\toprule\noalign{}
\begin{minipage}[b]{\linewidth}\raggedleft
Inspection SK
\end{minipage} & \begin{minipage}[b]{\linewidth}\raggedright
Building ID
\end{minipage} & \begin{minipage}[b]{\linewidth}\raggedright
Inspection datetime
\end{minipage} & \begin{minipage}[b]{\linewidth}\raggedright
Inspection trigger
\end{minipage} & \begin{minipage}[b]{\linewidth}\raggedright
Inspection outcome
\end{minipage} \\
\midrule\noalign{}
\endhead
\bottomrule\noalign{}
\endlastfoot
101 & B001 & 2025-04-18 10:00 & Scheduled & Pass \\
102 & B001 & 2025-05-20 11:30 & Complaint & Fail \\
103 & B002 & 2025-04-15 09:00 & Risk review & Fail \\
104 & B002 & 2025-04-30 13:15 & Follow-up & Pass \\
\end{longtable}
}

Although both tables contain \texttt{{[}Building\ ID{]}}, there is no
direct relationship between a specific inspection and a specific
fire-safety certificate. The source system does not say which
certificate the inspection should be read against.

To support analysis, the data engineer can create a nearest temporal
join.

For example, the rule might be:

\begin{quote}
An inspection relates to the most recent fire-safety certificate
submitted for the same building before the inspection.
\end{quote}

The result can be stored in \texttt{FireSafety.InspectionCertificate}.

\textbf{Example structure of \texttt{FireSafety.InspectionCertificate}}

{\def\LTcaptype{none} % do not increment counter
\begin{longtable}[]{@{}
  >{\raggedleft\arraybackslash}p{(\linewidth - 4\tabcolsep) * \real{0.3333}}
  >{\raggedleft\arraybackslash}p{(\linewidth - 4\tabcolsep) * \real{0.3333}}
  >{\raggedleft\arraybackslash}p{(\linewidth - 4\tabcolsep) * \real{0.3333}}@{}}
\toprule\noalign{}
\begin{minipage}[b]{\linewidth}\raggedleft
Inspection SK
\end{minipage} & \begin{minipage}[b]{\linewidth}\raggedleft
Fire safety certificate SK
\end{minipage} & \begin{minipage}[b]{\linewidth}\raggedleft
Days from certificate to inspection
\end{minipage} \\
\midrule\noalign{}
\endhead
\bottomrule\noalign{}
\endlastfoot
101 & 1 & 17 \\
102 & 2 & 19 \\
103 & 3 & 5 \\
104 & 4 & 12 \\
\end{longtable}
}

This table re-establishes an analytical relationship between the
inspection and the fire-safety certificate that preceded it.

The relationship is not ``true'' in the same way as a source-system
foreign key. It is an inferred relationship created from business logic.
For that reason, it should be documented and tested.

The data engineer should monitor cases where the assumption may be
unsafe, such as:

\begin{itemize}
\tightlist
\item
  inspections with no preceding fire-safety certificate;
\item
  inspections where the most recent certificate is too old;
\item
  buildings with multiple certificates close together;
\item
  certificate types that should not be associated with inspections;
\item
  changes in the inspection process that alter the expected timing.
\end{itemize}

For example, the business may decide that an inspection should only be
linked to a fire-safety certificate if it occurred within 60 days. If
the most recent certificate is older than that, the inspection should be
flagged for review rather than automatically linked.

Nearest temporal joins are especially valuable at the macro level.

At unit level, the problem is often that two records should be related
but no foreign key was recorded. At macro level, the problem is broader:
two business processes are related, but they operate in bulk, at
different grains, and without unit-level traceability.

In the fire-safety example, certificate submission and inspection are
not the same business process. They are separate processes that become
analytically meaningful only when the data engineer defines a
relationship between them.

This is the discipline of the pattern: use time to restore a
relationship only where business knowledge says the relationship is
meaningful, then monitor the assumptions that make it safe.

\subsubsection{Mapping tables}\label{mapping-tables}

Mapping tables close the quality gap between entities that must be
related analytically but are not related operationally.

Mapping tables are introduced when the original system does not record
relationships needed for analysis.

A common scenario is a header table with two separate detail tables at
different grains. Both detail tables relate to the header, but not to
each other.

The data engineer may introduce a mapping table between them based on a
business rule.

Consider a club restaurant with data captured in three tables:

\begin{itemize}
\tightlist
\item
  \texttt{Club.TableSitting}---records of customers seated at the
  restaurant, one per sitting, keyed by
  \texttt{{[}Table\ sitting\ ID{]}}
\item
  \texttt{Club.TableCustomer}---club members identified by
  \texttt{{[}Member\ ID{]}}, linked to \texttt{{[}Table\ sitting\ ID{]}}
\item
  \texttt{Club.TableFoodOrder}---menu items ordered per
  \texttt{{[}Table\ sitting\ ID{]}}, with food items identified by
  \texttt{{[}Food\ item\ ID{]}}
\end{itemize}

\textbf{Example structure of \texttt{Club.TableCustomer}}

{\def\LTcaptype{none} % do not increment counter
\begin{longtable}[]{@{}ll@{}}
\toprule\noalign{}
Table sitting ID & Member ID \\
\midrule\noalign{}
\endhead
\bottomrule\noalign{}
\endlastfoot
T001 & M001 \\
T001 & M002 \\
\end{longtable}
}

\textbf{Example structure of \texttt{Club.TableFoodOrder}}

{\def\LTcaptype{none} % do not increment counter
\begin{longtable}[]{@{}lllr@{}}
\toprule\noalign{}
Table sitting ID & Food order ID & Food item ID & Food item cost \\
\midrule\noalign{}
\endhead
\bottomrule\noalign{}
\endlastfoot
T001 & FO001 & F100 & 12.00 \\
T001 & FO002 & F200 & 18.00 \\
\end{longtable}
}

The system does not record which member ordered which item. For
analysis, the business may wish to associate members with food items.
However, this relationship is not captured in the source system.

To support analysis, the data engineer may introduce a mapping table
called \texttt{Club.CustomerFoodOrderMap}.

The logic may be based on a business rule such as:

\begin{quote}
Every member at the table shares the cost of all items ordered.
\end{quote}

\textbf{Example structure of \texttt{Club.CustomerFoodOrderMap}}

{\def\LTcaptype{none} % do not increment counter
\begin{longtable}[]{@{}lllr@{}}
\toprule\noalign{}
Table sitting ID & Member ID & Food order ID & Allocation weight \\
\midrule\noalign{}
\endhead
\bottomrule\noalign{}
\endlastfoot
T001 & M001 & FO001 & 0.5 \\
T001 & M002 & FO001 & 0.5 \\
T001 & M001 & FO002 & 0.5 \\
T001 & M002 & FO002 & 0.5 \\
\end{longtable}
}

Because the relationship is inferred rather than recorded, it should be
tested. Techniques for validating mapping logic are described in
\hyperref[docs-quality-reliability-tests-and-assumptions]{Tests and
assumptions}.

Variants of this are studied in
\hyperref[docs-creating-information-meaningful-fragments]{Meaningful
fragments}.

\section{Fuzzy logic: when intent can only be
approximated}\label{fuzzy-logic-when-intent-can-only-be-approximated}

There are cases where precise rules do not apply.

Fuzzy logic closes data quality gaps by approximating business intent
where exact rules are not available. Its main monitoring need is to
detect drift in match rates, validation results, or alignment with a
known reference.

A common example is natural language processing. When working with free
text, it is often impossible to define exact rules that meet business
intent.

Fuzzy logic is appropriate when the business intent is real but cannot
be captured precisely. In many scenarios, something is better than
nothing, and a perfect match with reality does not justify the cost.

This aligns with the view that data quality is not data versus reality,
but data versus business intent.

Examples of fuzzy logic include:

\begin{itemize}
\tightlist
\item
  optical character recognition to extract information from images;
\item
  large language models to infer sentiment from user text;
\item
  entity resolution techniques to identify the underlying entity behind
  multiple records;
\item
  pattern extraction from free text, such as phone numbers, invoice
  numbers, or addresses.
\end{itemize}

Fuzzy logic often pushes into the domain of data science, where
statistical and mathematical disciplines bring rigour to its
application. However, there are cases where a data engineer can apply a
simple and practical approach.

The following is a three-step process, a kind of ``poor man's data
science'':

\begin{enumerate}
\def\labelenumi{\arabic{enumi}.}
\tightlist
\item
  Find the pattern through loose--tight iteration.
\item
  Check the pattern with random validation.
\item
  Monitor the pattern by watching for drift.
\end{enumerate}

This approach can be highly effective in common scenarios that require
extracting information from free text. A data engineer should be
comfortable considering this approach and ready to pivot to a more
sophisticated data science method when necessary.

\subsection{Step 1---Loose--tight
iteration}\label{step-1loosetight-iteration}

Loose--tight iteration closes the quality gap where a pattern is real
enough to use, but too messy to define perfectly at the start.

Every pattern has matches and rejects. The matches are records that meet
the criteria. The rejects are records that do not.

The data engineer tunes the pattern by adjusting the criteria:

\begin{itemize}
\tightlist
\item
  tightening to reduce over-matching;
\item
  loosening to recover missed records.
\end{itemize}

The goal is not to maximise matches. The goal is to reduce mistakes on
both sides.

Consider the example of extracting phone numbers from a free-text field.
Phone numbers may appear in formats such as \texttt{(02)\ 1234\ 5678},
\texttt{0412345678}, \texttt{12345678}, or \texttt{02-1234-5678}.

The challenge is that the field may also contain other numbers, such as
invoice numbers, dates, quantities, account numbers, or address numbers.

\textbf{Example structure of
\texttt{\textquotesingle{}Customer\ note\textquotesingle{}}}

{\def\LTcaptype{none} % do not increment counter
\begin{longtable}[]{@{}rl@{}}
\toprule\noalign{}
Customer note ID & Customer note \\
\midrule\noalign{}
\endhead
\bottomrule\noalign{}
\endlastfoot
1 & Please call me on (02) 1234 5678 after 3pm. \\
2 & Mobile is 0412345678. Invoice 889912 is unrelated. \\
3 & Address is 12 Hill Street. No phone provided. \\
4 & Try 02-9876-5432 or office 98765432. \\
5 & Order 12345678 was delayed. \\
6 & Phone number changed to 0412 345 678 on 2025-04-01. \\
\end{longtable}
}

A rough first pattern might look for any eight-digit number.

This will produce matches, but it will also over-match.

\textbf{Example match set after rough pattern}

{\def\LTcaptype{none} % do not increment counter
\begin{longtable}[]{@{}rlll@{}}
\toprule\noalign{}
Customer note ID & Extracted value & Correct extraction? & Comment \\
\midrule\noalign{}
\endhead
\bottomrule\noalign{}
\endlastfoot
2 & 0412345678 & true & Mobile number \\
4 & 98765432 & true & Local phone number \\
5 & 12345678 & false & Order number, not phone number \\
\end{longtable}
}

The pattern found real phone numbers, but also picked up an order
number. This is a false match. The pattern is too loose.

The data engineer may tighten the pattern by looking for phone-like
context, such as nearby words like \texttt{phone}, \texttt{mobile},
\texttt{call}, or \texttt{office}, or by excluding known phrases such as
\texttt{order} and \texttt{invoice}.

After tightening, the match set may improve, but some genuine phone
numbers may now be missed.

\textbf{Example reject set after tightened pattern}

{\def\LTcaptype{none} % do not increment counter
\begin{longtable}[]{@{}
  >{\raggedleft\arraybackslash}p{(\linewidth - 6\tabcolsep) * \real{0.3077}}
  >{\raggedright\arraybackslash}p{(\linewidth - 6\tabcolsep) * \real{0.2308}}
  >{\raggedright\arraybackslash}p{(\linewidth - 6\tabcolsep) * \real{0.2308}}
  >{\raggedright\arraybackslash}p{(\linewidth - 6\tabcolsep) * \real{0.2308}}@{}}
\toprule\noalign{}
\begin{minipage}[b]{\linewidth}\raggedleft
Customer note ID
\end{minipage} & \begin{minipage}[b]{\linewidth}\raggedright
Customer note
\end{minipage} & \begin{minipage}[b]{\linewidth}\raggedright
Should have matched?
\end{minipage} & \begin{minipage}[b]{\linewidth}\raggedright
Comment
\end{minipage} \\
\midrule\noalign{}
\endhead
\bottomrule\noalign{}
\endlastfoot
1 & Please call me on (02) 1234 5678 after 3pm. & true & Landline with
area code \\
3 & Address is 12 Hill Street. No phone provided. & false & No phone
number \\
6 & Phone number changed to 0412 345 678 on 2025-04-01. & true & Mobile
with spaces \\
\end{longtable}
}

The rejected records reveal false rejects. The pattern is now too tight,
or it does not recognise enough valid phone-number formats.

The data engineer then loosens the pattern to recognise area codes,
mobile numbers with spaces, and hyphenated numbers while still avoiding
obvious order numbers, invoice numbers, dates, and addresses.

A later iteration may produce a more useful extraction.

\textbf{Example extraction after loose--tight iteration}

{\def\LTcaptype{none} % do not increment counter
\begin{longtable}[]{@{}rll@{}}
\toprule\noalign{}
Customer note ID & Extracted phone number & Extraction status \\
\midrule\noalign{}
\endhead
\bottomrule\noalign{}
\endlastfoot
1 & 0212345678 & Matched \\
2 & 0412345678 & Matched \\
3 & & No phone detected \\
4 & 0298765432; 98765432 & Matched \\
5 & & Rejected as likely order number \\
6 & 0412345678 & Matched \\
\end{longtable}
}

The important point is not the exact pattern. The important point is the
discipline of iteration.

Finding the best pattern is an iterative process:

\begin{enumerate}
\def\labelenumi{\arabic{enumi}.}
\tightlist
\item
  Define a rough pattern.
\item
  Inspect the match set for false matches.
\item
  Inspect the reject set for false rejects.
\item
  Adjust and repeat.
\end{enumerate}

Each adjustment changes the match set and reject set. It is usually
better to tighten or loosen one thing at a time, then inspect the
result.

The key idea is that focusing only on matches creates a blind spot. The
rejected records are equally significant. A useful pattern is found by
examining both incorrect matches and incorrect rejects.

\subsection{Step 2---Random validation}\label{step-2random-validation}

Random validation closes the quality gap where approximate logic looks
plausible, but has not been checked against independent judgement.

Once the pattern has been tuned, it should be validated by a business
expert.

This is done by randomly sampling records from both the match set and
the reject set.

The sample must be unbiased. It should not focus only on records where
the engineer has low confidence, nor only on records that seem easy to
validate. The point is to test the pattern across its full range of
behaviour.

For example, the data engineer may extract phone numbers from 50,000
customer notes. The pattern may return 18,000 notes with phone numbers
and 32,000 notes without phone numbers. The validation sample should
include records from both groups.

\textbf{Example validation sample}

{\def\LTcaptype{none} % do not increment counter
\begin{longtable}[]{@{}
  >{\raggedleft\arraybackslash}p{(\linewidth - 10\tabcolsep) * \real{0.2105}}
  >{\raggedright\arraybackslash}p{(\linewidth - 10\tabcolsep) * \real{0.1579}}
  >{\raggedright\arraybackslash}p{(\linewidth - 10\tabcolsep) * \real{0.1579}}
  >{\raggedright\arraybackslash}p{(\linewidth - 10\tabcolsep) * \real{0.1579}}
  >{\raggedright\arraybackslash}p{(\linewidth - 10\tabcolsep) * \real{0.1579}}
  >{\raggedright\arraybackslash}p{(\linewidth - 10\tabcolsep) * \real{0.1579}}@{}}
\toprule\noalign{}
\begin{minipage}[b]{\linewidth}\raggedleft
Sample ID
\end{minipage} & \begin{minipage}[b]{\linewidth}\raggedright
Source set
\end{minipage} & \begin{minipage}[b]{\linewidth}\raggedright
Customer note
\end{minipage} & \begin{minipage}[b]{\linewidth}\raggedright
Pattern result
\end{minipage} & \begin{minipage}[b]{\linewidth}\raggedright
Expert judgement
\end{minipage} & \begin{minipage}[b]{\linewidth}\raggedright
Validation outcome
\end{minipage} \\
\midrule\noalign{}
\endhead
\bottomrule\noalign{}
\endlastfoot
1 & Match set & Mobile is 0412345678. & Phone found & Phone present &
Correct match \\
2 & Match set & Order 12345678. Please follow up. & Phone found & No
phone present & False match \\
3 & Reject set & Call through reception on 98765432. & No phone found &
Phone present & False reject \\
4 & Reject set & Address is 12 Hill Street. & No phone found & No phone
present & Correct reject \\
\end{longtable}
}

The validation result can then be summarised.

\textbf{Example validation summary}

{\def\LTcaptype{none} % do not increment counter
\begin{longtable}[]{@{}lr@{}}
\toprule\noalign{}
Validation outcome & Count \\
\midrule\noalign{}
\endhead
\bottomrule\noalign{}
\endlastfoot
Correct match & 86 \\
False match & 6 \\
Correct reject & 92 \\
False reject & 16 \\
\end{longtable}
}

This tells the data engineer how the approximation behaves. It may be
good enough for some purposes and not good enough for others.

If the business only needs a rough indication of whether contact details
are likely to be present, this may be acceptable. If the extracted phone
number will be used for direct customer contact, the tolerance for false
matches and false rejects will be much lower.

This is why fuzzy logic still depends on business intent. The question
is not whether the pattern is perfect. The question is whether the
pattern is good enough for the decision or action it supports.

This validation process should be repeated periodically. Even if the
pattern was correct at the time of implementation, changes in business
processes or user behaviour may cause it to degrade.

Regular validation helps ensure the pattern continues to serve its
intended purpose.

\subsection{Step 3---Monitoring for
drift}\label{step-3monitoring-for-drift}

Monitoring for drift closes the quality gap where an approximation was
once good enough, but may stop behaving as intended over time.

In practice, random validation cannot be performed continuously. Drift
monitoring provides a lightweight alternative.

The idea is to identify a statistic that should remain relatively stable
over time. If the statistic changes significantly, the pattern may no
longer be behaving as expected.

For example, suppose a free-text field is used to extract phone numbers.
Not all extracted numbers will match the customer database, but a
certain percentage---say 85\%---typically do. This percentage reflects
the stability of the pattern.

The data engineer can monitor this statistic over time.

\textbf{Example drift monitoring}

{\def\LTcaptype{none} % do not increment counter
\begin{longtable}[]{@{}
  >{\raggedright\arraybackslash}p{(\linewidth - 6\tabcolsep) * \real{0.2000}}
  >{\raggedleft\arraybackslash}p{(\linewidth - 6\tabcolsep) * \real{0.2667}}
  >{\raggedleft\arraybackslash}p{(\linewidth - 6\tabcolsep) * \real{0.2667}}
  >{\raggedleft\arraybackslash}p{(\linewidth - 6\tabcolsep) * \real{0.2667}}@{}}
\toprule\noalign{}
\begin{minipage}[b]{\linewidth}\raggedright
Month
\end{minipage} & \begin{minipage}[b]{\linewidth}\raggedleft
Notes processed
\end{minipage} & \begin{minipage}[b]{\linewidth}\raggedleft
Notes with extracted phone number
\end{minipage} & \begin{minipage}[b]{\linewidth}\raggedleft
Extracted numbers matching customer records
\end{minipage} \\
\midrule\noalign{}
\endhead
\bottomrule\noalign{}
\endlastfoot
2025-01 & 12,400 & 4,820 & 86\% \\
2025-02 & 13,100 & 5,040 & 85\% \\
2025-03 & 12,850 & 4,910 & 84\% \\
2025-04 & 13,300 & 5,110 & 85\% \\
2025-05 & 13,050 & 7,900 & 61\% \\
\end{longtable}
}

A sudden drop from the usual range may indicate that the pattern is no
longer behaving as expected.

Users may have started entering phone numbers in a new format. The
source system may have changed its note template. Another number, such
as a case number or ticket number, may have started appearing near
phone-like words. The customer database may also have changed, reducing
the apparent match rate.

This is a form of assumption monitoring. The assumption is that the
extracted values will continue to resemble the known population. If the
assumption fails, the pattern may need to be revisited.

Drift monitoring is especially useful when the business has a relatively
static reference point, such as a customer table, a list of known codes,
or a set of standard formats. However, any stable statistic can be used.

Useful drift statistics may include:

\begin{itemize}
\tightlist
\item
  percentage of records matched;
\item
  percentage of extracted values that match a known reference table;
\item
  number of distinct extracted values;
\item
  proportion of records with multiple extracted values;
\item
  proportion of records requiring manual review;
\item
  proportion of null or rejected records.
\end{itemize}

\section{Key ideas}\label{key-ideas-10}

\begin{quote}
\textbf{Key ideas.}

Data quality is fitness for business intent, not perfect correspondence
with reality.

Different intents create different quality requirements.

Human curation is appropriate when business judgement is needed.

Precise rules are appropriate when business intent can be formalised.

Fuzzy logic is appropriate when business intent can only be
approximated.

All three approaches depend on monitoring the assumptions that make the
intervention safe.
\end{quote}

\cleardoublepage

\chapter{Tests and
assumptions}\label{docs-quality-reliability-tests-and-assumptions}

\section{Anticipating errors}\label{anticipating-errors}

A data product is not trustworthy merely because it is correct today. It
is trustworthy when users are not caught unawares by eventual failures.

This is the third principle of data engineering:

\begin{quote}
instead of stopping with what works now, anticipate errors that may
occur.
\end{quote}

The principle follows from two facts.

First, the world changes. What is true today may not be true tomorrow.
New data may arrive outside previously conceived parameters. Business
processes may evolve. Source systems may change. Another developer may
alter existing code.

Second, engineers make mistakes. Complex transformations create many
opportunities for small errors to become consequential.

The danger is not only that errors occur---errors are inevitable.

The deeper danger is that errors occur silently.

A silent error allows a data product to remain polished and apparently
authoritative while no longer being safe to use.

Therefore:

\begin{quote}
Data engineering is not only building the transformation that works. It
is building the means by which failure becomes visible.
\end{quote}

This chapter covers two such mechanisms:

{\def\LTcaptype{none} % do not increment counter
\begin{longtable}[]{@{}
  >{\raggedright\arraybackslash}p{(\linewidth - 4\tabcolsep) * \real{0.3333}}
  >{\raggedright\arraybackslash}p{(\linewidth - 4\tabcolsep) * \real{0.3333}}
  >{\raggedright\arraybackslash}p{(\linewidth - 4\tabcolsep) * \real{0.3333}}@{}}
\toprule\noalign{}
\begin{minipage}[b]{\linewidth}\raggedright
Mechanism
\end{minipage} & \begin{minipage}[b]{\linewidth}\raggedright
Question
\end{minipage} & \begin{minipage}[b]{\linewidth}\raggedright
Output
\end{minipage} \\
\midrule\noalign{}
\endhead
\bottomrule\noalign{}
\endlastfoot
Tests & Did two independent calculations produce the same result? & Pass
or fail \\
Monitored assumptions & Has something appeared that requires human
attention? & A list of violating records \\
\end{longtable}
}

A test performs the same calculation in two different ways and checks
whether the results agree.

A monitored assumption checks whether a condition presumed to hold has
been violated, and returns the records requiring attention.

Tests protect the correctness of implemented logic. Monitored
assumptions protect the safety of the conditions around that logic.

Both should run regularly, such as once per pipeline batch. Their
purpose is not to eliminate failure entirely. Their purpose is to
prevent avoidable failure from remaining hidden.

\section{Tests: checking the same result two
ways}\label{tests-checking-the-same-result-two-ways}

A data test is strongest when expected and actual results are calculated
independently.

The premise of testing is this: when the same problem is solved in two
different ways, and both methods arrive at the same result, confidence
is justified.

In data engineering, a test has two parts:

{\def\LTcaptype{none} % do not increment counter
\begin{longtable}[]{@{}
  >{\raggedright\arraybackslash}p{(\linewidth - 2\tabcolsep) * \real{0.5000}}
  >{\raggedright\arraybackslash}p{(\linewidth - 2\tabcolsep) * \real{0.5000}}@{}}
\toprule\noalign{}
\begin{minipage}[b]{\linewidth}\raggedright
Test part
\end{minipage} & \begin{minipage}[b]{\linewidth}\raggedright
Meaning
\end{minipage} \\
\midrule\noalign{}
\endhead
\bottomrule\noalign{}
\endlastfoot
Expected result & A calculation of what the result should be, preferably
from a different source or logic path. \\
Actual result & The result produced by the pipeline, model, or measure
being tested. \\
\end{longtable}
}

The test passes if expected and actual results match.

In a data pipeline, the expected part of a test may be calculated from
raw or partly transformed data. The actual part uses the transformed
output.

In a Power BI semantic model, the expected part of a test may be
calculated from the underlying data source, such as SQL or a data lake.
The actual part may be calculated using DAX. This approach is effective
for validating complex DAX measures, especially measures that rely on
embedded grain or complex DAX.

Power BI tests are easy to create. The data engineer can drag and drop
to build a table visual with dimensions and measures. The Performance
Analyzer in Power BI Desktop reveals the DAX query behind the visual.
This can serve as the actual portion of the test. The engineer then
compares the DAX evaluation with a calculation from the underlying data
source.

Tests lose effectiveness when the expected query is a copy of the
pipeline code that produced the actual result. This entanglement reduces
independence and narrows the surface of possible errors caught by the
test.

A thoughtful test maximises the difference in logic between expected and
actual results.

The way to write a good test is to ask:

\begin{quote}
Where is the most complex or fragile part of the transformation, and how
can this result be calculated in a different way?
\end{quote}

The following examples are not exhaustive. They illustrate the mindset
of checking weak points by calculating the same result through a
different path.

\subsection{Row counts}\label{row-counts}

Row count tests check whether records have appeared, disappeared, or
duplicated unexpectedly.

The simplest test is to compare row counts before and after
transformation. The expected query counts rows from the raw data with
appropriate filters. The actual query counts rows from the transformed
output.

\textbf{Example row count test}

{\def\LTcaptype{none} % do not increment counter
\begin{longtable}[]{@{}lrr@{}}
\toprule\noalign{}
Test & Expected result & Actual result \\
\midrule\noalign{}
\endhead
\bottomrule\noalign{}
\endlastfoot
Count sales records for April 2025 & 10,240 & 10,240 \\
\end{longtable}
}

Even small variations can make row count tests more effective.

For example, instead of only counting all rows, the data engineer may
group by a reference value.

\textbf{Example grouped row count test}

{\def\LTcaptype{none} % do not increment counter
\begin{longtable}[]{@{}lrr@{}}
\toprule\noalign{}
Sales region & Expected sales rows & Actual sales rows \\
\midrule\noalign{}
\endhead
\bottomrule\noalign{}
\endlastfoot
North & 2,450 & 2,450 \\
South & 3,100 & 3,100 \\
West & 1,920 & 1,918 \\
New store not annotated & 12 & 14 \\
\end{longtable}
}

The total row count may still match because one group is undercounted
while another is overcounted. Grouped counts reveal the problem, helping
catch update errors, join errors, and mapping issues.

The expected and actual calculations should use different logical paths.
For example, the expected query may group raw records using source
fields, while the actual query groups transformed records using curated
dimensions.

Row count tests are especially important for incremental loads.

If the extract relies on an architectural column such as
\texttt{{[}Row\ update\ datetime{]}}, the test should use a different
datetime column, such as \texttt{{[}Staff\ update\ datetime{]}},
\texttt{{[}Submission\ datetime{]}}, or another business-centric column.
This separation avoids self-confirmation.

Selecting a sample of recent records is also useful. A recent-record
sample can check creation, update, and deletion behaviour across many
columns, not only row counts.

\subsection{Checksums on key results}\label{checksums-on-key-results}

Checksum tests check whether an important summary result still balances.

They are useful when a transformation rolls detailed records up into a
business concept.

Consider the example of manufacturing batch quality.

A business process may record individual quality-control criteria. The
data engineer then rolls these up to define the overall quality outcome
for each batch.

\textbf{Example structure of
\texttt{\textquotesingle{}Batch\ quality\ criterion\textquotesingle{}}}

{\def\LTcaptype{none} % do not increment counter
\begin{longtable}[]{@{}
  >{\raggedright\arraybackslash}p{(\linewidth - 6\tabcolsep) * \real{0.2500}}
  >{\raggedright\arraybackslash}p{(\linewidth - 6\tabcolsep) * \real{0.2500}}
  >{\raggedright\arraybackslash}p{(\linewidth - 6\tabcolsep) * \real{0.2500}}
  >{\raggedright\arraybackslash}p{(\linewidth - 6\tabcolsep) * \real{0.2500}}@{}}
\toprule\noalign{}
\begin{minipage}[b]{\linewidth}\raggedright
Batch ID
\end{minipage} & \begin{minipage}[b]{\linewidth}\raggedright
Quality criterion
\end{minipage} & \begin{minipage}[b]{\linewidth}\raggedright
Quality result
\end{minipage} & \begin{minipage}[b]{\linewidth}\raggedright
Is critical criterion
\end{minipage} \\
\midrule\noalign{}
\endhead
\bottomrule\noalign{}
\endlastfoot
B1001 & Temperature control & Pass & true \\
B1001 & Packaging integrity & Pass & false \\
B1001 & Labelling accuracy & Fail & false \\
B1002 & Temperature control & Fail & true \\
B1002 & Packaging integrity & Pass & false \\
\end{longtable}
}

The transformed output may look like this.

\textbf{Example structure of
\texttt{\textquotesingle{}Batch\ quality\ outcome\textquotesingle{}}}

{\def\LTcaptype{none} % do not increment counter
\begin{longtable}[]{@{}
  >{\raggedright\arraybackslash}p{(\linewidth - 6\tabcolsep) * \real{0.2143}}
  >{\raggedleft\arraybackslash}p{(\linewidth - 6\tabcolsep) * \real{0.2857}}
  >{\raggedleft\arraybackslash}p{(\linewidth - 6\tabcolsep) * \real{0.2857}}
  >{\raggedright\arraybackslash}p{(\linewidth - 6\tabcolsep) * \real{0.2143}}@{}}
\toprule\noalign{}
\begin{minipage}[b]{\linewidth}\raggedright
Batch ID
\end{minipage} & \begin{minipage}[b]{\linewidth}\raggedleft
Critical criteria failed
\end{minipage} & \begin{minipage}[b]{\linewidth}\raggedleft
Non-critical criteria failed
\end{minipage} & \begin{minipage}[b]{\linewidth}\raggedright
Batch quality outcome
\end{minipage} \\
\midrule\noalign{}
\endhead
\bottomrule\noalign{}
\endlastfoot
B1001 & 0 & 1 & Pass \\
B1002 & 1 & 0 & Fail \\
\end{longtable}
}

A checksum can test the core logic without copying the transformation.

For example, every batch should fall into one and only one of two
groups: pass or fail.

\textbf{Example checksum test}

{\def\LTcaptype{none} % do not increment counter
\begin{longtable}[]{@{}lrr@{}}
\toprule\noalign{}
Test & Expected result & Actual result \\
\midrule\noalign{}
\endhead
\bottomrule\noalign{}
\endlastfoot
Count distinct batches from criterion table & 2 & 2 \\
Count batches marked pass or fail in outcome table & 2 & 2 \\
Count batches with any critical failure from criterion table & 1 & 1 \\
Count batches marked fail in outcome table & 1 & 1 \\
\end{longtable}
}

This test is not a copy of the transformation. It checks whether the
transformed result still balances against the source logic.

A more comprehensive version can group by values such as
\texttt{{[}Product\ type{]}}, \texttt{{[}Manufacture\ site{]}}, or
\texttt{{[}Inspection\ month{]}}. This increases the chance of catching
errors that affect only a segment of the data.

\subsection{Bypassing mapping tables}\label{bypassing-mapping-tables}

Mapping tables are often introduced where the source system did not
record an analytical relationship.

A good test for a mapping table is to compare results with and without
the mapping table.

Consider the restaurant example from the previous chapter. The source
system records table sittings, table customers, and food orders. It does
not record which customer ordered which food item.

The data engineer introduces \texttt{Club.CustomerFoodOrderMap} to
allocate food orders across members seated at the table.

\textbf{Example structure of \texttt{Club.TableCustomer}}

{\def\LTcaptype{none} % do not increment counter
\begin{longtable}[]{@{}ll@{}}
\toprule\noalign{}
Table sitting ID & Member ID \\
\midrule\noalign{}
\endhead
\bottomrule\noalign{}
\endlastfoot
T001 & M001 \\
T001 & M002 \\
\end{longtable}
}

\textbf{Example structure of \texttt{Club.TableFoodOrder}}

{\def\LTcaptype{none} % do not increment counter
\begin{longtable}[]{@{}lllr@{}}
\toprule\noalign{}
Table sitting ID & Food order ID & Food item ID & Food item cost \\
\midrule\noalign{}
\endhead
\bottomrule\noalign{}
\endlastfoot
T001 & FO001 & F100 & 12.00 \\
T001 & FO002 & F200 & 18.00 \\
\end{longtable}
}

\textbf{Example structure of \texttt{Club.CustomerFoodOrderMap}}

{\def\LTcaptype{none} % do not increment counter
\begin{longtable}[]{@{}lllr@{}}
\toprule\noalign{}
Table sitting ID & Member ID & Food order ID & Allocation weight \\
\midrule\noalign{}
\endhead
\bottomrule\noalign{}
\endlastfoot
T001 & M001 & FO001 & 0.5 \\
T001 & M002 & FO001 & 0.5 \\
T001 & M001 & FO002 & 0.5 \\
T001 & M002 & FO002 & 0.5 \\
\end{longtable}
}

Regardless of how complex the mapping logic becomes, some totals should
remain the same. The total cost of food ordered at the table should not
change after allocation.

\textbf{Example mapping checksum}

{\def\LTcaptype{none} % do not increment counter
\begin{longtable}[]{@{}
  >{\raggedright\arraybackslash}p{(\linewidth - 4\tabcolsep) * \real{0.2727}}
  >{\raggedleft\arraybackslash}p{(\linewidth - 4\tabcolsep) * \real{0.3636}}
  >{\raggedleft\arraybackslash}p{(\linewidth - 4\tabcolsep) * \real{0.3636}}@{}}
\toprule\noalign{}
\begin{minipage}[b]{\linewidth}\raggedright
Table sitting ID
\end{minipage} & \begin{minipage}[b]{\linewidth}\raggedleft
Expected food cost from \texttt{Club.TableFoodOrder}
\end{minipage} & \begin{minipage}[b]{\linewidth}\raggedleft
Actual allocated food cost from \texttt{Club.CustomerFoodOrderMap}
\end{minipage} \\
\midrule\noalign{}
\endhead
\bottomrule\noalign{}
\endlastfoot
T001 & 30.00 & 30.00 \\
\end{longtable}
}

This is a powerful way to validate mapping logic. The test does not need
to repeat every detail of the mapping. It checks whether the mapping
preserves an invariant.

More specific tests can then check the business rule itself.

For example, if the business rule assumes food cost is evenly
distributed among seated members, the test can check that allocation
weights sum to \texttt{1} per food order.

\textbf{Example allocation-weight test}

{\def\LTcaptype{none} % do not increment counter
\begin{longtable}[]{@{}
  >{\raggedright\arraybackslash}p{(\linewidth - 4\tabcolsep) * \real{0.2727}}
  >{\raggedleft\arraybackslash}p{(\linewidth - 4\tabcolsep) * \real{0.3636}}
  >{\raggedleft\arraybackslash}p{(\linewidth - 4\tabcolsep) * \real{0.3636}}@{}}
\toprule\noalign{}
\begin{minipage}[b]{\linewidth}\raggedright
Food order ID
\end{minipage} & \begin{minipage}[b]{\linewidth}\raggedleft
Expected allocation weight total
\end{minipage} & \begin{minipage}[b]{\linewidth}\raggedleft
Actual allocation weight total
\end{minipage} \\
\midrule\noalign{}
\endhead
\bottomrule\noalign{}
\endlastfoot
FO001 & 1.0 & 1.0 \\
FO002 & 1.0 & 1.0 \\
\end{longtable}
}

Grouping by reference columns such as \texttt{{[}Food\ item\ ID{]}},
\texttt{{[}Food\ item\ type{]}}, or \texttt{{[}Table\ sitting\ ID{]}}
can make the test more sensitive.

\subsection{Checking boundary cases}\label{checking-boundary-cases}

Boundary cases are common sources of error.

A typical example is a rolling-window measure. Any implementation of a
rolling window should be tested at boundary points.

Suppose there is a Power BI measure \texttt{{[}Total\ sales\ volume{]}},
and a derived rolling 12-month version
\texttt{{[}Total\ sales\ volume\ past\ 12\ months{]}}.

The rolling measure is easy to get subtly wrong. It may include the
wrong start date, exclude the current date, mishandle leap years, or
behave incorrectly at the beginning of the available history.

One useful test is to compare the rolling result against known calendar
boundaries.

For example, if \texttt{{[}Total\ sales\ volume\ past\ 12\ months{]}} is
evaluated at 31 December 2025, it should match total sales volume for
the 2025 calendar year.

\textbf{Example boundary test}

{\def\LTcaptype{none} % do not increment counter
\begin{longtable}[]{@{}
  >{\raggedright\arraybackslash}p{(\linewidth - 4\tabcolsep) * \real{0.3333}}
  >{\raggedright\arraybackslash}p{(\linewidth - 4\tabcolsep) * \real{0.3333}}
  >{\raggedright\arraybackslash}p{(\linewidth - 4\tabcolsep) * \real{0.3333}}@{}}
\toprule\noalign{}
\begin{minipage}[b]{\linewidth}\raggedright
Evaluation date
\end{minipage} & \begin{minipage}[b]{\linewidth}\raggedright
Expected result
\end{minipage} & \begin{minipage}[b]{\linewidth}\raggedright
Actual rolling 12-month result
\end{minipage} \\
\midrule\noalign{}
\endhead
\bottomrule\noalign{}
\endlastfoot
2025-12-31 & Total sales from 2025-01-01 to 2025-12-31 & Total sales
volume past 12 months at 2025-12-31 \\
2026-06-30 & Total sales from 2025-07-01 to 2026-06-30 & Total sales
volume past 12 months at 2026-06-30 \\
\end{longtable}
}

This checks that the measure behaves properly at calendar-year and
financial-year boundaries.

However, this test would still pass if
\texttt{{[}Total\ sales\ volume\ past\ 12\ months{]}} were mistakenly
implemented as a copy of \texttt{{[}Total\ sales\ volume{]}} in a model
that is already filtered to one year.

To guard against this, the data engineer should also test turning
points, such as the first and last day of a month.

\textbf{Example turning-point test}

{\def\LTcaptype{none} % do not increment counter
\begin{longtable}[]{@{}ll@{}}
\toprule\noalign{}
Evaluation date & Boundary being tested \\
\midrule\noalign{}
\endhead
\bottomrule\noalign{}
\endlastfoot
2025-02-28 & End of February \\
2025-03-01 & Start of March \\
2025-06-30 & End of financial year \\
2025-07-01 & Start of financial year \\
\end{longtable}
}

A rolling-window measure that passes at calendar-year boundaries,
financial-year boundaries, and interim month turning points provides
much stronger assurance than one tested only at convenient dates.

\subsection{Using subject matter
knowledge}\label{using-subject-matter-knowledge}

Subject matter knowledge can introduce independence into a test.

Normally, the data engineer should avoid building unverified assumptions
into the pipeline. This prevents blind spots and reduces the risk of
silent errors.

However, tests are one place where business knowledge can be safely used
to strengthen robustness. If the business knowledge is wrong, the test
will reveal that. If the business knowledge was once right but stops
being right, the test will reveal a change in the process.

For example, a business expert may know that every table sitting must
have at least one customer. The data engineer should not blindly assume
this in the pipeline, because load errors may still occur. But the
knowledge can be used in a test.

\textbf{Example subject-matter test}

{\def\LTcaptype{none} % do not increment counter
\begin{longtable}[]{@{}
  >{\raggedright\arraybackslash}p{(\linewidth - 4\tabcolsep) * \real{0.3333}}
  >{\raggedright\arraybackslash}p{(\linewidth - 4\tabcolsep) * \real{0.3333}}
  >{\raggedright\arraybackslash}p{(\linewidth - 4\tabcolsep) * \real{0.3333}}@{}}
\toprule\noalign{}
\begin{minipage}[b]{\linewidth}\raggedright
Business knowledge
\end{minipage} & \begin{minipage}[b]{\linewidth}\raggedright
Expected calculation
\end{minipage} & \begin{minipage}[b]{\linewidth}\raggedright
Actual calculation
\end{minipage} \\
\midrule\noalign{}
\endhead
\bottomrule\noalign{}
\endlastfoot
Every table sitting has at least one customer & Count table sittings
using an inner join to \texttt{Club.TableCustomer} & Count all rows in
\texttt{Club.TableSitting} \\
\end{longtable}
}

If the results differ, either the data has an error, the pipeline has an
error, or the business knowledge was incomplete.

Business knowledge can also validate roll-up logic.

Consider a help desk system where cases escalate from Tier 1 to Tier 4.
Escalations are recorded in \texttt{Helpdesk.Escalation}, keyed by
\texttt{{[}Case\ ID{]}}, with \texttt{{[}Tier{]}} indicating escalation
level. The table \texttt{Helpdesk.CaseEscalation} rolls this information
up to case grain. A reference table \texttt{Helpdesk.RefCaseEscalation}
includes \texttt{{[}Highest\ escalation{]}} to summarise each case.

Tests can use business knowledge to validate the roll-up.

{\def\LTcaptype{none} % do not increment counter
\begin{longtable}[]{@{}
  >{\raggedright\arraybackslash}p{(\linewidth - 4\tabcolsep) * \real{0.3333}}
  >{\raggedright\arraybackslash}p{(\linewidth - 4\tabcolsep) * \real{0.3333}}
  >{\raggedright\arraybackslash}p{(\linewidth - 4\tabcolsep) * \real{0.3333}}@{}}
\toprule\noalign{}
\begin{minipage}[b]{\linewidth}\raggedright
Business knowledge
\end{minipage} & \begin{minipage}[b]{\linewidth}\raggedright
Expected
\end{minipage} & \begin{minipage}[b]{\linewidth}\raggedright
Actual
\end{minipage} \\
\midrule\noalign{}
\endhead
\bottomrule\noalign{}
\endlastfoot
Every case starts at Tier 1 & Count distinct \texttt{{[}Case\ ID{]}} in
\texttt{Helpdesk.Escalation} where \texttt{{[}Tier{]}\ =\ "Tier\ 1"} &
Count rows in \texttt{Helpdesk.CaseEscalation} \\
Tier 4 is the maximum escalation & Count distinct
\texttt{{[}Case\ ID{]}} in \texttt{Helpdesk.Escalation} where
\texttt{{[}Tier{]}\ =\ "Tier\ 4"} & Count rows in
\texttt{Helpdesk.CaseEscalation} where
\texttt{{[}Highest\ escalation{]}\ =\ "Tier\ 4"} \\
Cases do not skip the existence of earlier tiers & Count Tier 4 cases
with no Tier 1 record & Should be zero \\
\end{longtable}
}

In general, existence and uniqueness conditions known to business
experts can be used to bypass or vary the original pipeline logic.

This is one of the strongest ways to write tests because it introduces a
source of independence outside the code itself.

\section{Monitored assumptions: surfacing records that require
attention}\label{monitored-assumptions-surfacing-records-that-require-attention}

Data engineers must make assumptions.

However, assumptions may fail over time. Sometimes the failure is
unanticipated. At other times, the failure is anticipated but still
needs to be detected.

A monitored assumption is a query that returns records if and only if
human attention is required.

If no rows are returned, the assumption remains valid. If rows are
returned, each row is a violation that needs attention.

The key to monitoring assumptions is recognising that an assumption
exists.

The following are common cases.

{\def\LTcaptype{none} % do not increment counter
\begin{longtable}[]{@{}
  >{\raggedright\arraybackslash}p{(\linewidth - 2\tabcolsep) * \real{0.5000}}
  >{\raggedright\arraybackslash}p{(\linewidth - 2\tabcolsep) * \real{0.5000}}@{}}
\toprule\noalign{}
\begin{minipage}[b]{\linewidth}\raggedright
Assumption
\end{minipage} & \begin{minipage}[b]{\linewidth}\raggedright
Returns rows when...
\end{minipage} \\
\midrule\noalign{}
\endhead
\bottomrule\noalign{}
\endlastfoot
Source data is up to date & The latest expected batch has not
arrived. \\
Reference data is complete & New values need mapping. \\
Known values remain stable & A new hard-coded or unhandled value
appears. \\
Data quality rules still hold & Invalid dates, duplicate keys, or
out-of-range values appear. \\
Fuzzy logic has not drifted & Match rates or validation results move
outside tolerance. \\
\end{longtable}
}

\subsection{Is source data up to date?}\label{is-source-data-up-to-date}

Source data cannot always be assumed to arrive on time.

Any number of failures may delay the arrival of new data: an upstream
system may be unavailable, an extract may fail, a file may not arrive,
or a source team may change a schedule without warning.

A monitored assumption can check whether the latest expected data has
arrived.

\textbf{Example source freshness assumption}

{\def\LTcaptype{none} % do not increment counter
\begin{longtable}[]{@{}
  >{\raggedright\arraybackslash}p{(\linewidth - 6\tabcolsep) * \real{0.2500}}
  >{\raggedright\arraybackslash}p{(\linewidth - 6\tabcolsep) * \real{0.2500}}
  >{\raggedright\arraybackslash}p{(\linewidth - 6\tabcolsep) * \real{0.2500}}
  >{\raggedright\arraybackslash}p{(\linewidth - 6\tabcolsep) * \real{0.2500}}@{}}
\toprule\noalign{}
\begin{minipage}[b]{\linewidth}\raggedright
Source table
\end{minipage} & \begin{minipage}[b]{\linewidth}\raggedright
Expected latest date
\end{minipage} & \begin{minipage}[b]{\linewidth}\raggedright
Actual latest date
\end{minipage} & \begin{minipage}[b]{\linewidth}\raggedright
Requires attention
\end{minipage} \\
\midrule\noalign{}
\endhead
\bottomrule\noalign{}
\endlastfoot
\texttt{Sales.Order} & 2025-05-01 & 2025-05-01 & false \\
\texttt{Sales.Payment} & 2025-05-01 & 2025-04-28 & true \\
\end{longtable}
}

If the query returns \texttt{Sales.Payment}, the issue needs attention.
The data product may be incomplete even though the pipeline has
technically run.

This assumption is especially important when multiple source systems are
combined. A report may look normal while one process has updated and
another has not.

\subsection{Is reference data
complete?}\label{is-reference-data-complete}

Reference data is vital to a quality warehouse.

It enables different systems to map local values to conformed values.
But reference mappings require maintenance as new source values arrive.

For example, one system may start sending a new country code that is not
yet mapped.

\textbf{Example structure of
\texttt{\textquotesingle{}Country\ code\ map\textquotesingle{}}}

{\def\LTcaptype{none} % do not increment counter
\begin{longtable}[]{@{}llr@{}}
\toprule\noalign{}
Source system & Source country code & Country SK \\
\midrule\noalign{}
\endhead
\bottomrule\noalign{}
\endlastfoot
Sales & AU & 1 \\
Sales & NZ & 2 \\
Sales & US & 3 \\
Shipping & AUS & 1 \\
Shipping & NZL & 2 \\
Shipping & USA & 3 \\
\end{longtable}
}

A monitored assumption can check whether any source country codes have
arrived without a mapping.

\textbf{Example assumption output}

{\def\LTcaptype{none} % do not increment counter
\begin{longtable}[]{@{}llr@{}}
\toprule\noalign{}
Source system & Source country code & Records affected \\
\midrule\noalign{}
\endhead
\bottomrule\noalign{}
\endlastfoot
Shipping & SGP & 14 \\
\end{longtable}
}

The query should return only the values that require human attention. In
this case, someone needs to decide whether \texttt{SGP} maps to an
existing country record or whether the country reference data need to be
extended.

\subsection{Are there unanticipated
values?}\label{are-there-unanticipated-values}

Hard-coded values are sometimes unavoidable.

This is acceptable when a source column contains a small number of
stable values. For example, a status column may have contained
\texttt{Open}, \texttt{Closed}, and \texttt{Cancelled} for many years.

However, if a new value appears, the transformation logic may no longer
be safe.

\textbf{Example status assumption}

{\def\LTcaptype{none} % do not increment counter
\begin{longtable}[]{@{}lr@{}}
\toprule\noalign{}
Source status & Records affected \\
\midrule\noalign{}
\endhead
\bottomrule\noalign{}
\endlastfoot
Pending review & 37 \\
\end{longtable}
}

This assumption does not necessarily mean the pipeline is broken. It
means the business process has produced a value not anticipated by the
transformation. The data engineer or business owner must decide how it
should be handled.

This is especially important when a transformation uses \texttt{case}
logic, hard-coded mappings, or manually curated categories.

\subsection{Are there data quality
issues?}\label{are-there-data-quality-issues}

Source data may contain issues from unvalidated collection.

These can include duplicate keys, dates outside expected ranges,
negative values where only positive values make sense, or records that
violate business rules.

In some cases, the pipeline may treat the issue to preserve analytical
usefulness. For example, a future date outside the allowed range may be
converted to blank and flagged.

\textbf{Example assumption output for invalid dates}

{\def\LTcaptype{none} % do not increment counter
\begin{longtable}[]{@{}rll@{}}
\toprule\noalign{}
Record ID & Source date & Issue \\
\midrule\noalign{}
\endhead
\bottomrule\noalign{}
\endlastfoot
2 & 2300-06-15 & Date exceeds allowed future range \\
\end{longtable}
}

Similarly, duplicate records may be removed from the main analytical
table and written to a rejected-records table.

\textbf{Example assumption output for duplicate submissions}

{\def\LTcaptype{none} % do not increment counter
\begin{longtable}[]{@{}llll@{}}
\toprule\noalign{}
Submission ID & Customer ID & Submission date & Rejection reason \\
\midrule\noalign{}
\endhead
\bottomrule\noalign{}
\endlastfoot
S1003 & C001 & 2025-04-01 & Duplicate customer submission for date \\
\end{longtable}
}

These are monitored assumptions because the business needs to know that
the issue occurred. The pipeline may be able to continue, but the
underlying source issue still requires attention.

\subsection{Is there data drift?}\label{is-there-data-drift}

Statistical and fuzzy logic often rely on assumptions about the
distribution or pattern of input data.

These assumptions can change over time. This is commonly known as data
drift.

Data engineers may encounter this when applying fuzzy patterns to
extract information from free-text fields. A pattern may work well when
users write notes in one style, then fail when a source form changes or
users begin entering data differently.

For example, suppose a pattern extracts phone numbers from customer
notes. The data engineer may monitor the proportion of extracted numbers
that match known customer phone records.

\textbf{Example drift monitoring}

{\def\LTcaptype{none} % do not increment counter
\begin{longtable}[]{@{}
  >{\raggedright\arraybackslash}p{(\linewidth - 6\tabcolsep) * \real{0.2000}}
  >{\raggedleft\arraybackslash}p{(\linewidth - 6\tabcolsep) * \real{0.2667}}
  >{\raggedleft\arraybackslash}p{(\linewidth - 6\tabcolsep) * \real{0.2667}}
  >{\raggedleft\arraybackslash}p{(\linewidth - 6\tabcolsep) * \real{0.2667}}@{}}
\toprule\noalign{}
\begin{minipage}[b]{\linewidth}\raggedright
Month
\end{minipage} & \begin{minipage}[b]{\linewidth}\raggedleft
Notes processed
\end{minipage} & \begin{minipage}[b]{\linewidth}\raggedleft
Notes with extracted phone number
\end{minipage} & \begin{minipage}[b]{\linewidth}\raggedleft
Extracted numbers matching customer records
\end{minipage} \\
\midrule\noalign{}
\endhead
\bottomrule\noalign{}
\endlastfoot
2025-01 & 12,400 & 4,820 & 86\% \\
2025-02 & 13,100 & 5,040 & 85\% \\
2025-03 & 12,850 & 4,910 & 84\% \\
2025-04 & 13,300 & 5,110 & 85\% \\
2025-05 & 13,050 & 7,900 & 61\% \\
\end{longtable}
}

The May result may indicate drift. Users may have started entering phone
numbers in a new format. Another number may have started appearing near
phone-like words. The customer reference table may also have changed.

A monitored assumption should return rows, periods, or segments where
the statistic moves outside tolerance.

The important point is that fuzzy logic should not be treated as
set-and-forget. Its assumptions must be monitored.

\section{Tests and assumptions
together}\label{tests-and-assumptions-together}

Tests and monitored assumptions are closely related, but they protect
different things.

Tests detect whether implemented logic has failed. Monitored assumptions
detect whether the world has changed in a way that makes the logic
unsafe.

A condition can often be expressed either as a test or as a monitored
assumption by rephrasing.

For example, row counts before and after transformation should usually
be expressed as tests. A discrepancy may indicate a serious error in the
pipeline.

New country codes requiring annotation are better expressed as monitored
assumptions. Their appearance does not necessarily mean the pipeline has
failed. It means the world has produced something that requires human
attention.

A useful rule of thumb is:

{\def\LTcaptype{none} % do not increment counter
\begin{longtable}[]{@{}
  >{\raggedright\arraybackslash}p{(\linewidth - 4\tabcolsep) * \real{0.3333}}
  >{\raggedright\arraybackslash}p{(\linewidth - 4\tabcolsep) * \real{0.3333}}
  >{\raggedright\arraybackslash}p{(\linewidth - 4\tabcolsep) * \real{0.3333}}@{}}
\toprule\noalign{}
\begin{minipage}[b]{\linewidth}\raggedright
Condition
\end{minipage} & \begin{minipage}[b]{\linewidth}\raggedright
Prefer
\end{minipage} & \begin{minipage}[b]{\linewidth}\raggedright
Reason
\end{minipage} \\
\midrule\noalign{}
\endhead
\bottomrule\noalign{}
\endlastfoot
Critical correctness condition & Test & Failure suggests the product may
be wrong. \\
Slow-moving business condition & Monitored assumption & Failure suggests
human attention is needed. \\
New values or mappings & Monitored assumption & The product needs
curation, not necessarily shutdown. \\
Core reconciliation or row preservation & Test & Failure suggests the
transformation may have broken. \\
Statistical or fuzzy pattern stability & Monitored assumption & Failure
suggests logic may need review. \\
\end{longtable}
}

\section{Conclusion}\label{conclusion-2}

Tests and assumptions accelerate delivery rather than slow it down. They
reduce the cost of change. They allow data engineers to modify complex
transformations with confidence. They make it safer to refactor code,
adjust business logic, optimise performance, and release improvements
without relying on hope.

If something is easy to test, it should be tested because the cost is
low. If something is hard to test, it definitely should be tested
because the logic is complex. Either way, it should be tested.

Thoughtful tests are best, but even simple tests are valuable. Many
errors in data products are not profound misunderstandings. They are
careless mistakes made during rapid development. Simple tests and
assumptions can catch these before they reach the user.

New engineers naturally spend more time choosing patterns and building
their implementation. With experience, design becomes more mechanical
and rapid. Mature engineers dealing with greater complexity spend
proportionately more time anticipating failure: defining tests,
monitoring assumptions, and making sure that the data product will
reveal when it is no longer safe.

\section{Key ideas}\label{key-ideas-11}

\begin{quote}
\textbf{Key ideas.}

Data engineering is not only building the transformation that works. It
is building the means by which failure becomes visible.

The third principle of data engineering is to anticipate errors.

Tests check the same result in two independent ways.

Tests are strongest when expected and actual calculations are not
entangled.

Monitored assumptions return records that require human attention.

Tests detect failures in implemented logic.

Monitored assumptions detect changes that make the logic unsafe.
\end{quote}

\cleardoublepage

\chapter{Fault
tolerance}\label{docs-quality-reliability-fault-tolerance}

\section{Containing failure}\label{containing-failure}

Good engineering contains failure before it spreads.

A brittle pipeline treats every error as catastrophic. A fault-tolerant
pipeline distinguishes between errors that should stop everything,
errors that should stop one table, and errors that should be isolated to
a small set of records.

The goal is not to ignore errors. The goal is to keep the rest of the
data product usable while surfacing failure clearly.

On the other hand, surfacing failure is not the same as producing
endless alerts.

If the same warning appears every day and no one knows what to do with
it, the warning becomes background noise rather than a warning.

This is the cry-wolf problem of monitoring. Repeated, unactionable
alerts train people to ignore the pipeline.

For this reason, fault tolerance requires judgement. Some errors should
abort a load. Some occasional errors should be isolated into reject
tables while the rest of the pipeline proceeds. Some recurring errors
should be given targeted handling so that the pipeline explicitly
represents the known issue rather than repeatedly rediscovering it.

This is fault tolerance.

Fault tolerance is therefore an extension of the third principle of data
engineering: anticipate errors.

This chapter introduces three common fault patterns:

{\def\LTcaptype{none} % do not increment counter
\begin{longtable}[]{@{}
  >{\raggedright\arraybackslash}p{(\linewidth - 4\tabcolsep) * \real{0.3333}}
  >{\raggedright\arraybackslash}p{(\linewidth - 4\tabcolsep) * \real{0.3333}}
  >{\raggedright\arraybackslash}p{(\linewidth - 4\tabcolsep) * \real{0.3333}}@{}}
\toprule\noalign{}
\begin{minipage}[b]{\linewidth}\raggedright
Fault pattern
\end{minipage} & \begin{minipage}[b]{\linewidth}\raggedright
Core question
\end{minipage} & \begin{minipage}[b]{\linewidth}\raggedright
Example failure
\end{minipage} \\
\midrule\noalign{}
\endhead
\bottomrule\noalign{}
\endlastfoot
Uniqueness & Does one real-world entity correspond to the right number
of records? & Duplicate submissions, duplicated incremental loads, or
two people sharing one identifier. \\
Existence & Are required records and values present where they should
be? & A detail row without a header, a missing sales amount, or a
deleted source record not removed downstream. \\
Stability & Are small real-world changes producing proportionate
data-world changes? & One source update timestamp causes every
historical row to reload. \\
\end{longtable}
}

These patterns illustrate the mindset of building a fault-tolerant
pipeline.

\section{Uniqueness}\label{uniqueness}

Uniqueness faults occur when the relationship between real-world
entities and database records breaks down.

If the basic task of data engineering is to represent real-world objects
and processes in the data world, then uniqueness is one of the most
fundamental expectations.

In the simplest case, one real-world entity should correspond to one
database record.

Uniqueness is violated when one real-world entity appears as multiple
database records. For example, a staff member may accidentally enter the
same transaction twice.

The reverse can also occur. Two distinct real-world entities may be
forced to occupy the same database record. For example, if staff are
identified only by first and last name, two staff members with the same
name may collapse into one identifier value.

Uniqueness violations can arise from business process failures, such as
duplicate entry. They can also arise from mechanical failures, such as
incremental load logic incorrectly loading the same record twice.

The data engineer must express uniqueness expectations in the warehouse.
They are physical expressions of business intent.

If the constraint is violated, the pipeline has several possible
responses:

{\def\LTcaptype{none} % do not increment counter
\begin{longtable}[]{@{}
  >{\raggedright\arraybackslash}p{(\linewidth - 2\tabcolsep) * \real{0.5000}}
  >{\raggedright\arraybackslash}p{(\linewidth - 2\tabcolsep) * \real{0.5000}}@{}}
\toprule\noalign{}
\begin{minipage}[b]{\linewidth}\raggedright
Response
\end{minipage} & \begin{minipage}[b]{\linewidth}\raggedright
Use when...
\end{minipage} \\
\midrule\noalign{}
\endhead
\bottomrule\noalign{}
\endlastfoot
Abort the load & The violation suggests the table is unsafe to
publish. \\
Reject the violating records & The invalid records can be isolated while
the rest of the table remains usable. \\
Deduplicate automatically & The business rule for choosing the accepted
record is clear and safe. \\
\end{longtable}
}

For example, a source table may contain duplicate submissions.

\textbf{Example structure before deduplication}

{\def\LTcaptype{none} % do not increment counter
\begin{longtable}[]{@{}lllr@{}}
\toprule\noalign{}
Submission ID & Customer ID & Submission date & Submission amount \\
\midrule\noalign{}
\endhead
\bottomrule\noalign{}
\endlastfoot
S1001 & C001 & 2025-04-01 & 120.00 \\
S1002 & C002 & 2025-04-01 & 95.00 \\
S1003 & C001 & 2025-04-01 & 120.00 \\
S1004 & C003 & 2025-04-02 & 80.00 \\
\end{longtable}
}

If the business expects one submission per customer per date, the
duplicate can be separated from the main analytical output.

\textbf{Example structure of accepted records}

{\def\LTcaptype{none} % do not increment counter
\begin{longtable}[]{@{}lllr@{}}
\toprule\noalign{}
Submission ID & Customer ID & Submission date & Submission amount \\
\midrule\noalign{}
\endhead
\bottomrule\noalign{}
\endlastfoot
S1001 & C001 & 2025-04-01 & 120.00 \\
S1002 & C002 & 2025-04-01 & 95.00 \\
S1004 & C003 & 2025-04-02 & 80.00 \\
\end{longtable}
}

\textbf{Example structure of rejected records}

{\def\LTcaptype{none} % do not increment counter
\begin{longtable}[]{@{}
  >{\raggedright\arraybackslash}p{(\linewidth - 8\tabcolsep) * \real{0.1875}}
  >{\raggedright\arraybackslash}p{(\linewidth - 8\tabcolsep) * \real{0.1875}}
  >{\raggedright\arraybackslash}p{(\linewidth - 8\tabcolsep) * \real{0.1875}}
  >{\raggedleft\arraybackslash}p{(\linewidth - 8\tabcolsep) * \real{0.2500}}
  >{\raggedright\arraybackslash}p{(\linewidth - 8\tabcolsep) * \real{0.1875}}@{}}
\toprule\noalign{}
\begin{minipage}[b]{\linewidth}\raggedright
Submission ID
\end{minipage} & \begin{minipage}[b]{\linewidth}\raggedright
Customer ID
\end{minipage} & \begin{minipage}[b]{\linewidth}\raggedright
Submission date
\end{minipage} & \begin{minipage}[b]{\linewidth}\raggedleft
Submission amount
\end{minipage} & \begin{minipage}[b]{\linewidth}\raggedright
Rejection reason
\end{minipage} \\
\midrule\noalign{}
\endhead
\bottomrule\noalign{}
\endlastfoot
S1003 & C001 & 2025-04-01 & 120.00 & Duplicate customer submission for
date \\
\end{longtable}
}

The accepted records remain available for users. The duplicate record is
not silently lost but surfaced for attention. This approach localises
the fault.

If uniqueness violations are rare, a reject-record pattern may be
sufficient. If violations are frequent, the table needs targeted
handling. Repeated alerts can create a cry-wolf effect, dulling
attention and undermining the purpose of monitoring.

In such cases, the data engineer should implement a specific treatment
for the table. This may include deterministic deduplication, a
remediation workflow, or a data quality report.

\section{Existence}\label{existence}

Existence faults occur when required records or values are missing, or
when records remain in the data world after their real-world counterpart
has disappeared.

Existence is violated when a real-world entity exists but does not
appear in the database. For example, a completed sale may be missing its
sales amount.

Existence is also violated when a database record exists but the
real-world entity no longer does. For example, a record may fail to be
deleted during an incremental load.

Some existence expectations are simple. A completed sales record must
have a sales amount. This can be expressed through a not-null
constraint.

\textbf{Example structure of
\texttt{\textquotesingle{}Sales\textquotesingle{}} with missing value}

{\def\LTcaptype{none} % do not increment counter
\begin{longtable}[]{@{}lllr@{}}
\toprule\noalign{}
Sales ID & Sales date & Customer ID & Sales amount \\
\midrule\noalign{}
\endhead
\bottomrule\noalign{}
\endlastfoot
S1001 & 2025-04-01 & C001 & 120.00 \\
S1002 & 2025-04-02 & C002 & \\
S1003 & 2025-04-03 & C003 & 95.00 \\
\end{longtable}
}

If \texttt{{[}Sales\ amount{]}} is required for completed sales, then
the missing value should be treated as a fault. Depending on the use
case, the row may be rejected, flagged, or excluded from measures that
require a sales amount.

More complex existence faults occur between related tables.

For example, a sales item should not exist without a corresponding sales
header.

\textbf{Example structure of \texttt{Sales.Sales}}

{\def\LTcaptype{none} % do not increment counter
\begin{longtable}[]{@{}lll@{}}
\toprule\noalign{}
Sales ID & Sales date & Customer ID \\
\midrule\noalign{}
\endhead
\bottomrule\noalign{}
\endlastfoot
S1001 & 2025-04-01 & C001 \\
S1002 & 2025-04-02 & C002 \\
\end{longtable}
}

\textbf{Example structure of \texttt{Sales.SalesItems}}

{\def\LTcaptype{none} % do not increment counter
\begin{longtable}[]{@{}lllr@{}}
\toprule\noalign{}
Sales item ID & Sales ID & Product ID & Quantity \\
\midrule\noalign{}
\endhead
\bottomrule\noalign{}
\endlastfoot
SI001 & S1001 & P100 & 1 \\
SI002 & S1001 & P200 & 2 \\
SI003 & S999 & P300 & 1 \\
\end{longtable}
}

In this example, \texttt{SI003} refers to \texttt{S999}, but the sales
header does not exist.

This may occur because the header failed to load, the item arrived
earlier than the header, or the source system produced an invalid
relationship.

The data engineer can detect this by left-joining the detail table to
the header table, then checking whether the header key is missing.

Conceptually, the check looks like this.

\textbf{Example structure after checking for header existence}

{\def\LTcaptype{none} % do not increment counter
\begin{longtable}[]{@{}lllrl@{}}
\toprule\noalign{}
Sales item ID & Sales ID & Product ID & Quantity & Header sales ID \\
\midrule\noalign{}
\endhead
\bottomrule\noalign{}
\endlastfoot
SI001 & S1001 & P100 & 1 & S1001 \\
SI002 & S1001 & P200 & 2 & S1001 \\
SI003 & S999 & P300 & 1 & \\
\end{longtable}
}

The existence expectation can then be expressed as a not-null
requirement on \texttt{{[}Header\ sales\ ID{]}}. If
\texttt{{[}Header\ sales\ ID{]}} is blank, the detail row has no
corresponding header.

When the not-null check fails, the pipeline can send the violating rows
into a reject table and allow the rest of the detail table to proceed.

\textbf{Example structure of accepted sales items}

{\def\LTcaptype{none} % do not increment counter
\begin{longtable}[]{@{}lllr@{}}
\toprule\noalign{}
Sales item ID & Sales ID & Product ID & Quantity \\
\midrule\noalign{}
\endhead
\bottomrule\noalign{}
\endlastfoot
SI001 & S1001 & P100 & 1 \\
SI002 & S1001 & P200 & 2 \\
\end{longtable}
}

\textbf{Example structure of rejected sales items}

{\def\LTcaptype{none} % do not increment counter
\begin{longtable}[]{@{}lllrl@{}}
\toprule\noalign{}
Sales item ID & Sales ID & Product ID & Quantity & Rejection reason \\
\midrule\noalign{}
\endhead
\bottomrule\noalign{}
\endlastfoot
SI003 & S999 & P300 & 1 & Missing sales header \\
\end{longtable}
}

The rest of the table can continue to load, while the missing
relationship is surfaced.

This is the same pattern as uniqueness handling: detect the violation,
isolate the unsafe rows, preserve them for attention, and allow the safe
records to continue.

Existence faults are especially important in batch and incremental
pipelines. Tables may be extracted at different times. A header may
arrive in one batch and detail records in another. In traditional
warehousing this is often discussed as late-arriving data.

The data engineer should not assume that every related record will
always arrive in the expected order.

For existence, the best mindset is to design the pipeline as if it were
streaming.

Even in a batch pipeline, this is a useful design mindset:

\begin{itemize}
\tightlist
\item
  tables may load continuously, or at least several times a day;
\item
  some tables may fail to load from time to time;
\item
  related records may arrive out of order;
\item
  records may catch up in a later batch.
\end{itemize}

In this mindset, temporary existence violations are expected. The
pipeline should handle them safely.

However, safe handling must not become silent handling.

If a missing header catches up in the next batch, the issue may resolve.
If it does not catch up, tests or monitored assumptions should surface
the prolonged discrepancy.

\section{Stability}\label{stability}

Stability faults occur when small real-world changes produce
disproportionate changes in the data world.

A stable pipeline changes in proportion to the world it represents.

If ten sales changed in the source system, the pipeline should not
update ten million downstream records. If one reference value changed,
half the warehouse should not change as a side effect.

Stability means that ordinary changes should produce ordinary effects.

This is important because the real world rarely rewrites all historical
records at once. But the data world can suffer from massive accidental
changes that do not correspond to real-world change.

For example, a source system update may refresh
\texttt{{[}Row\ update\ datetime{]}} for every historical row. If the
pipeline relies on \texttt{{[}Row\ update\ datetime{]}} for incremental
extraction, it may attempt to reload the entire history.

Another example is a non-deterministic ranking. If a row number is
assigned using an incomplete sort order, the selected ``first'' row may
change between runs even when the underlying data has not changed.

A third example is a reporting column such as
\texttt{{[}Today’s\ date{]}}. If this is stored on every row, every row
changes every day even though the underlying business entity has not
changed.

The data engineer should therefore design pipelines so that changes are
commensurate with real-world change.

One way to do this is to avoid unstable elements:

{\def\LTcaptype{none} % do not increment counter
\begin{longtable}[]{@{}
  >{\raggedright\arraybackslash}p{(\linewidth - 2\tabcolsep) * \real{0.3429}}
  >{\raggedright\arraybackslash}p{(\linewidth - 2\tabcolsep) * \real{0.6571}}@{}}
\toprule\noalign{}
\begin{minipage}[b]{\linewidth}\raggedright
Unstable element
\end{minipage} & \begin{minipage}[b]{\linewidth}\raggedright
Why it is risky
\end{minipage} \\
\midrule\noalign{}
\endhead
\bottomrule\noalign{}
\endlastfoot
\texttt{{[}Today’s\ date{]}} stored on every row & Causes every row to
change every day. \\
Non-deterministic ranking & Causes rows to change between runs without
real-world change. \\
Incomplete tie-breaking logic & Makes selected records unstable. \\
Single-row control tables & Causes downstream volatility when
append-only history would be safer. \\
Wide miscellaneous tables & Amplifies small changes because unrelated
attributes share one row. \\
\end{longtable}
}

Another way is to use stability thresholds.

For example, the pipeline may abort or quarantine a table load if more
than 5\% of records are about to be deleted.

\textbf{Example stability threshold}

{\def\LTcaptype{none} % do not increment counter
\begin{longtable}[]{@{}
  >{\raggedright\arraybackslash}p{(\linewidth - 8\tabcolsep) * \real{0.1667}}
  >{\raggedleft\arraybackslash}p{(\linewidth - 8\tabcolsep) * \real{0.2222}}
  >{\raggedleft\arraybackslash}p{(\linewidth - 8\tabcolsep) * \real{0.2222}}
  >{\raggedleft\arraybackslash}p{(\linewidth - 8\tabcolsep) * \real{0.2222}}
  >{\raggedright\arraybackslash}p{(\linewidth - 8\tabcolsep) * \real{0.1667}}@{}}
\toprule\noalign{}
\begin{minipage}[b]{\linewidth}\raggedright
Table
\end{minipage} & \begin{minipage}[b]{\linewidth}\raggedleft
Existing rows
\end{minipage} & \begin{minipage}[b]{\linewidth}\raggedleft
Rows proposed for deletion
\end{minipage} & \begin{minipage}[b]{\linewidth}\raggedleft
Deletion percentage
\end{minipage} & \begin{minipage}[b]{\linewidth}\raggedright
Action
\end{minipage} \\
\midrule\noalign{}
\endhead
\bottomrule\noalign{}
\endlastfoot
\texttt{Sales.Sales} & 1,000,000 & 2,500 & 0.25\% & Proceed \\
\texttt{Sales.Customer} & 250,000 & 180,000 & 72.00\% & Abort and
alert \\
\end{longtable}
}

The threshold does not imply that the change is wrong. It says the
proposed change is large enough to require attention before it is
allowed to spread.

One important habit is to avoid wide tables with miscellaneous
attributes. In a wide table, a change in any cell can cause the whole
row to update. This amplifies small real-world changes into broad
database changes.

Narrow meaningful fragments are more stable. When each table has a clear
informational purpose, changes to that table correspond more closely to
changes in the thing it represents.

Another habit is to manage dependencies carefully. Small reference
tables can have large downstream effects. If a country reference table
fails to load, and downstream transformations depend on it early in the
pipeline, country data may disappear across many products.

These design habits are developed further in
\hyperref[docs-efficient-stable-pipeline-load-mechanics]{Load mechanics}
and \hyperref[docs-efficient-stable-pipeline-load-dependencies]{Load
dependencies}.

\section{Conclusion}\label{conclusion-3}

Uniqueness, existence, and stability are three common fault patterns.
But the possibilities of error are limitless. Parallel loads may
deadlock. Source systems may change without notice. Files may arrive
late. Users may enter unexpected values.

It is not useful to enumerate every possible fault.

The important thing is the mindset: anticipating errors.

The experienced data engineer designs for what the world might do to the
pipeline, not only what the pipeline does to the data.

This is why the data engineer builds with fault tolerance: not every
error is fatal, but none should pass silently.

\section{Key ideas}\label{key-ideas-12}

\begin{quote}
\textbf{Key ideas.}

Fault tolerance is the discipline of containing failure.

Good engineering contains errors before they spread through the
pipeline.

A fault-tolerant pipeline distinguishes between errors that should stop
everything, errors that should stop one table, and errors that should be
isolated to a small set of records.

Uniqueness faults occur when the relationship between real-world
entities and database records breaks down.

Existence faults occur when required records or values are missing, or
when records remain after their real-world counterpart has disappeared.

Stability faults occur when small real-world changes produce
disproportionate changes in the data world.

Not every error is fatal, but none should pass silently.
\end{quote}

\cleardoublepage

\thispagestyle{empty}
\phantomsection
\hypertarget{docs-efficient-stable-pipeline}{}
\addcontentsline{toc}{part}{Efficient \& stable pipeline}
\vspace*{0.34\textheight}
\begin{center}
{\Huge\bfseries Efficient \& stable pipeline}
\end{center}
\cleardoublepage

\chapter{Efficiency and
stability}\label{docs-efficient-stable-pipeline-efficiency-and-stability}

Developers often talk about fast and slow. But what do these terms mean?
Is one second fast? Is ten seconds too slow? There are situations where
one second to process ten rows is unacceptably slow, and others where
ten seconds to process ten thousand different rows is impressively fast.
These judgements depend on the infrastructure, the structure of the
incoming data, and the value of the information being processed.

A sharp intuition for whether something ``can still be faster'' is one
of the hallmarks of a veteran developer. A developer who cannot tell
that something could be faster will not attempt to improve it. On the
other hand, trying to optimise a computation that is already at its
limit is futile.

Given the number of environmental factors involved in time-to-compute,
and how slippery the idea of ``fast'' can be, this intuition is
difficult to develop. A way to make the problem more tractable is to
focus on one aspect of performance --- efficiency.

Efficiency in a data pipeline means using the minimal amount of resource
to calculate one piece of information. There are two ways to achieve
this:

\begin{enumerate}
\def\labelenumi{\arabic{enumi}.}
\item
  The pipeline processes the minimal amount of information needed to
  reflect changes in the data.
\item
  The algorithm that applies these changes uses the minimal amount of
  computational space-time resources.
\end{enumerate}

The first is informational efficiency. The second is algorithmic
efficiency.

Informational efficiency is about processing as little information as
possible. A pipeline is inefficient if it processes hundreds of millions
of rows but in the end, only a handful changed. Algorithmic efficiency
is about computing those changes as cheaply as possible on a predefined
computational environment. It is inefficient if it takes seconds to
process one row when it could be done in milliseconds. Informational
efficiency is universal across all technologies. Algorithmic efficiency
is more technology dependent and varies from site to site.

Efficiency matters because it is key to the sustainable growth of a
warehouse. It also directly affects how quickly users get their results.
It is within the data engineer's control. Decisions about server
configuration and information value may not be easy to influence.
However, the data engineer can always ensure a pipeline does just enough
work, and no more, to reflect real world changes.

Stability is a closely related concept. It is the idea that small
changes to input should lead to small changes in output. An unstable
pipeline is one where a small change leads to a large change in output.
For example, changing a country name in a reference table might ripple
through and rewrite many tables. A more dramatic instance of this was
the fear of the Y2K bug where there were widespread concerns that switch
over on the New Year's Day 2000 would cause disastrous impacts on
systems that stored years only in the final two digits.

Stability matters because it shapes the user's experience of reliability
because of the magnitude of impact makes the problem obvious. Reports
that suddenly become empty one day radically undermines trust. There is
no point explaining it was because a change in international country
names caused a ripple effect. As far as the user is concerned, the
product is unreliable.

Efficiency and stability are related but distinct. An inefficient
pipeline is when there is a lot of work, but few results change in the
table. An unstable pipeline is when there are too many changes in a
table. From a business point-of-view, instability can also be
inefficiency. For example, a large table with {[}Today's date{]} column
would be updating every day (unstable) but is inefficient because it is
much work for little business value.

That is, the number of row changes do not reflect information value
changes.

The chapters in this part explore of efficiency and stability. A
recurring theme is that meaningful fragments form the foundation of an
efficient and stable pipeline.

\cleardoublepage

\chapter{Load
mechanics}\label{docs-efficient-stable-pipeline-load-mechanics}

During a load, a table's data needs to update with the most recent set.
The simplest approach is to drop the original table and replace the data
with a new set. This is how every data engineer starts. But it violates
the key tenets of pipeline efficiency and stability.

When a table's content is replaced en masse, it is not possible to know
which rows changed. Downstream processes cannot respond with precision
because every row looks new and needs response. This breaks the goal of
efficiency to ensure database changes are minimal.

More importantly, when a table is dropped and replaced, any error that
prevents a successful load becomes catastrophic. Reports may become
empty for the user. It can also result in a domino effect of errors in
downstream tables.

In a production pipeline that transforms hundreds or thousands of
tables, the cumulative probability and effect of these problems become
untenably disruptive.

Instead of drop and replace, updating the data of a table should occur
through three steps:

\begin{enumerate}
\def\labelenumi{\arabic{enumi}.}
\item
  Stage the incoming copy of the data into a temporary location
\item
  Check the incoming data against current data to pick up changes,
  anomalies and constraint violations
\item
  Apply genuine changes to the target table, updating any timestamp
  columns.
\end{enumerate}

These are the same steps whether the incoming data is a full extract or
a subset of increments.

\section{Loading}\label{loading}

Suppose the current table is X. The goal is to update X with the latest
batch of data. X has a primary key column {[}PK col1{]}, {[}PK col2{]}.

\subsection{Step 1 -- Stage}\label{step-1-stage}

The first step is to load the latest batch of data into a temporary
table. In traditional warehousing terminology, this is the staging
table. Suppose it is called X\_staging.

X\_staging may be the full load of data to replace X. It may be less
than the full set but still cover all the changed records in a recent
batch. This is an incremental extract. An incremental extract reduces
the number of records to stage, check, and apply from hundreds of
millions to a couple of thousand. The ways to do this reliably are
covered in Tracking change and Responding to change.

\subsection{Step 2 -- Check}\label{step-2-check}

The next step is to check the content of X\_staging against X for
changes and validity before loading onto X. This is the most intensive
step and is the one usually skipped in naive approaches to loading.

The aspects to check are:

\begin{itemize}
\item
  Changes:
\item
  Which rows to insert?
\item
  Which rows to update?
\item
  Which rows to delete?
\item
  Instability:
\item
  Are there a concerningly high number of rows to insert or update?
\item
  Are there a concerningly high number of rows to delete?
\item
  Violations:
\item
  Are there rows of uniqueness violations to reject?
\item
  Are there rows of existence (not null) violations to reject?
\end{itemize}

Checking for changes is essential. Checking for instability is highly
desirable. Checking for each of the violations is not critical, but
proactively rejecting a small number of bad rows and allowing correct
rows to pass through can reduce disruption to users by allowing them to
get access to most of their data.

\subsubsection{Checking for changes}\label{checking-for-changes}

Not every row in X\_staging is a genuine change in X. In fact, during a
full load, only a few hundred out of millions may differ in the latest
batch. The first step is to detect which rows are the ones that changed.

This check is necessary even if X\_staging is incrementally extracted
with the intention of capturing only changed data. Despite the best
effort of the data engineer, it is often difficult to extract only the
changed rows, especially when there are multiple joins and
transformations. While incremental extract may narrow the number of
candidate rows from millions to thousands, some of them may still be
identical to the rows in X, requiring a check to find those that are
truly changed.

In checking for changes, there are three possibilities:

\begin{itemize}
\item
  For a row in X\_staging, is the primary key value new? If so, it is an
  insert.
\item
  For a row in X\_staging that has an existing primary key value in X,
  are any of the column values different? If so, it is an update.
\item
  For a row in X, should the primary key value still be retained? If
  not, the row should be deleted.
\end{itemize}

The first two are collectively known as an upsert. A computationally
effective way of calculating inserts and updates simultaneously in SQL
is to:

\begin{enumerate}
\def\labelenumi{\arabic{enumi}.}
\item
  Start from X\_staging and left join X on the primary key columns {[}PK
  col1{]}, {[}PK col2{]}.
\item
  Keep only the rows from X\_staging that do not identically exist in X.
  This can be done using an EXCEPT comparison in the WHERE clause.
\item
  Add a computed column {[}Is new row{]} depending on whether the X
  primary key column value is null.
\end{enumerate}

Storing the result of this query into X\_upsert gives only the rows that
changed, with a column {[}Is new row{]} to indicate whether it is an
insert or update.

\begin{quote}
\textbf{Note.}

TODO: Insert manuscript screenshot or diagram from the source draft.
Source PDF note: \texttt{{[}EXAMPLE\ SQL{]}}.
\end{quote}

Checking for deletes depends on whether X\_staging is fully extracted or
incrementally extracted.

If X\_staging is fully extracted, then it represents all the records
that should be in the final table. Rows from X that should be deleted
are those whose primary key values are not in X\_staging. An
implementation in SQL would be a left anti-join of X on X\_staging.

The primary keys to delete should be stored in a table X\_delete.

\begin{quote}
\textbf{Note.}

TODO: Insert manuscript screenshot or diagram from the source draft.
Source PDF note: \texttt{{[}EXAMPLE\ SQL{]}}.
\end{quote}

If X\_staging is incrementally extracted, the idea would not work
because it only has a fraction of the primary keys in the desired final
output. Rather, X\_delete needs to be tailored to the incremental
extract logic for X. This is explored in greater depth in Responding to
change.

\subsubsection{Checking for instability}\label{checking-for-instability}

With the tables X\_upsert and X\_delete, it is easy to check whether
there are an abnormally high number of rows to be upserted or deleted in
a batch. If there are an abnormally high number of rows, then the load
should abort.

The thresholds for abnormality can be a combination of the number of
rows in the base table and the percentage of rows. These are stability
thresholds.

For example, if the table X has a hundred thousand rows, and more than
10\% are upserted, then this can be a sign of abnormality for aborting
the load.

The threshold for deletes should be different from the threshold for
upserts. There are usually far fewer deletes than upserts. For example,
less than 5\% of a large table should be deleted on a given day.

In practice, a simple stability threshold can apply as a default across
the entire pipeline to catch most abnormalities, such as having half the
rows disappear. Specific tables may have their stability thresholds
tailored to their case. Examining a table's history of changes can be
useful in determining the threshold. This can be done if these
statistics are logged. This logging is discussed in a subsequent
subsection.

The SQL for aborting on the upsert threshold can be done by checking the
number of rows in X\_upsert versus that of X:

\begin{quote}
\textbf{Note.}

TODO: Insert manuscript screenshot or diagram from the source draft.
Source PDF note: \texttt{{[}EXAMPLE\ SQL{]}}.
\end{quote}

Checking the delete threshold is only possible on a full extract. The
SQL is as follows:

\begin{quote}
\textbf{Note.}

TODO: Insert manuscript screenshot or diagram from the source draft.
Source PDF note: \texttt{{[}EXAMPLE\ SQL{]}}.
\end{quote}

These checks are an important part of monitoring and improving a
pipeline's stability. It is very unlikely that every record of a
business process's entire history, or even more than 5\% of them,
changes their content. If the upsert rows in X\_staging do exceed this
number, it is more likely to be an unintended side-effect of a system
change rather than a change in informational content. Stability
thresholds prevent these unintended changes from propagating.

If the load is frequently aborted for a table, then the extraction logic
needs to be improved for stability, or the thresholds need to be
adjusted.

\subsubsection{Checking for violations}\label{checking-for-violations}

When there are too many rows to upsert or delete, the entire table
aborts its load. This means the user does not get any of the latest
data. This can be too drastic when only one row out of thousands is
problematic.

Instead, if the table has uniqueness constraints and existence
constraints (not null property on a column), only the rows which violate
these constraints need to be discarded, and the rest can proceed. These
are the reject rows and can be stored in X\_reject.

The algorithm for rejecting existence violations is simply to discard
rows in X\_staging which have a null value on a column with a not null
property, and send them to X\_reject.

Suppose X has a column {[}Col not null{]} with a not null property, the
SQL is as follows:

\begin{quote}
\textbf{Note.}

TODO: Insert manuscript screenshot or diagram from the source draft.
Source PDF note: \texttt{{[}Example\ SQL{]}}.
\end{quote}

The algorithm for detecting uniqueness violations has two scenarios.
Suppose X\_staging has a column {[}Col unique{]}. The possibilities are:

\begin{enumerate}
\def\labelenumi{\arabic{enumi}.}
\item
  The duplicates exist in X\_staging when two rows have the same value
  on {[}Col unique{]}.
\item
  The duplicate happens when X\_staging applies an upsert that creates a
  value in {[}Col unique{]} that already exists in X, but not if that
  same row swaps its value to another one to avoid a duplicate.
\end{enumerate}

In these cases, all but one of the rows should be sent to X\_reject.
Example SQL follows.

\begin{quote}
\textbf{Note.}

TODO: Insert manuscript screenshot or diagram from the source draft.
Source PDF note: \texttt{{[}Example\ SQL{]}}.
\end{quote}

The two criteria can also be extended to other criteria such as data
type violations. In theory, any tailored criteria can be applied to
abort a load or reject some rows on a particular table. These two cover
many of the problems that a large pipeline encounters. They can be
applied by default to every load through automation rather than manual
tailoring.

\subsection{Step 3 - Apply}\label{step-3---apply}

Once the changes are determined in X\_upsert and X\_delete, and invalid
rows are sent to X\_reject, the changes can be applied to X.

The first step is to delete the rows in X using the keys in X\_delete.
Rather than completely losing the rows, it is important to send the
deletes to a history table so that downstream processes know which rows
have been deleted. In some platforms, this can be done automatically. In
this example, the table to hold the old records is X\_history. The SQL
is as follows:

\begin{quote}
\textbf{Note.}

TODO: Insert manuscript screenshot or diagram from the source draft.
Source PDF note: \texttt{{[}Example\ SQL{]}}.
\end{quote}

Then the upserts should be applied. The rows to be inserted from
X\_upsert into X are identified by loading X\_upsert where {[}Is new
row{]} = 1. The rows in X to be updated using values from X\_upsert are
identified by {[}Is new row{]} = 0.

As with deletes, the old values in X that are to be updated should be
sent to X\_history rather than being overridden and lost. Keeping these
rows is important for incremental processing downstream.

This step of applying changes should also add architectural columns for
change datetimes. They are:

\begin{itemize}
\item
  {[}Row insert datetime{]} --- when was the row, defined by the primary
  key value, first inserted into the table.
\item
  {[}Row update datetime{]} --- when was the row, defined by the primary
  key value, most recently updated with a new value. Because of the
  Check step, these are always genuine changes.
\item
  {[}Row delete datetime{]} --- when was the row, now in X\_history,
  deleted from X.
\end{itemize}

This could be because of a full delete by its primary key or because it
was overridden by an update.

Some platforms calculate one or more of these automatically.

\begin{quote}
\textbf{Note.}

TODO: Insert manuscript screenshot or diagram from the source draft.
Source PDF note: \texttt{{[}Example\ SQL{]}}.
\end{quote}

By the end of this step, and if there were no other errors, the changes
would be applied to X. The additional artefacts -- X\_history, {[}Row
insert datetime{]}, {[}Row update datetime{]}, and {[}Row delete
datetime{]} -- provide information for downstream tables that rely on X
to incrementally extract from it.

\section{After loading}\label{after-loading}

Cleaning-up and logging are important part of the load mechanics.

\subsection{Cleaning-up}\label{cleaning-up}

If there are no errors, the tables X\_staging, X\_upsert, and X\_delete
do not need to be retained. The rows in X\_history can also be purged
after a while, once they have informed incremental processing
downstream. In Delta table context, the history rows are the tombstoned
rows, and their purge is known as vacuuming.

If there are errors, then the intermediary tables and X\_reject need to
be retained for troubleshooting. It is important to ensure they are
secured in the same way as X, and do not accidentally cause unintended
data leakage.

\subsection{Logging}\label{logging}

In addition, it is essential to log the table's bookmark, and highly
valuable to log the change statistics.

The table's refresh bookmark is the datetime of the start of the load,
if the load successfully completed. Suppose the table X draws data from
A, B, and C. The refresh bookmark of a load defines when X started
reading from its source data, and any records that appear in A, B, and C
after that datetime have come afterwards. When X is incrementally
extracted in the next batch, the updated records in the source tables
are easily identified. This is explained in greater depth in Tracking
changes.

\begin{quote}
\textbf{Note.}

TODO: Insert manuscript screenshot or diagram from the source draft.
Source PDF note: \texttt{{[}SCREENSHOT{]}}.
\end{quote}

The Check step calculates the number of rows to be inserted, updated,
and deleted. It is only a small additional effort to log these numbers.
While not essential, they form a rich history to understand a table's
evolution and can be valuable for troubleshooting.

Since they are already calculated as part of the load, it is easily
worthwhile to log them.

\begin{quote}
\textbf{Note.}

TODO: Insert manuscript screenshot or diagram from the source draft.
Source PDF note: \texttt{{[}SCREENSHOT{]}}.
\end{quote}

The success, or failure and associated failure messages, should be
logged.

\begin{quote}
\textbf{Note.}

TODO: Insert manuscript screenshot or diagram from the source draft.
Source PDF note: \texttt{{[}EXAMPLE\ SCREENSHOT{]}}.
\end{quote}

\section{Analysis}\label{analysis}

The full three steps create five extra tables in addition to X ---
X\_staging, X\_upsert, X\_delete, X\_reject, and X\_history. It also
creates three extra columns --- {[}Row insert datetime{]}, {[}Row update
datetime{]}, and {[}Row delete datetime{]}. The logic is also much more
complex than simply dropping and replacing the table X.

In a simple environment with five or even fifty tables, all the
additional work may be unjustified overhead. In a moderately sized
pipeline, the benefits offered by the extra work outweigh the cost. Past
a certain point where errors and resource become costly, they are
non-negotiable.

From the view of efficiency, the row-by-row check for genuine change
seems exorbitant. But in a table with millions of rows and only a small
fraction changing in a load, it is much faster to check for changes and
apply them than to load all rows and perform any re-indexing.

More importantly, checking each row and updating {[}Row update
datetime{]} only for rows with a genuine change informs downstream
tables how to respond. Without such a check, all downstream tables will
need to reload all rows in every batch.

Seen in a different way, the data warehouse has responsibility to keep
track of every single change that happens within it. The row check is
the foundation for this tracking.

From the view of stability, the row-by-row check plays an important role
in stabilising a pipeline of incremental extract by eliminating noisy
changes. Incremental extract typically works by selecting candidate
change rows through a change detection column. If an external factor
causes the change detection column to shift all its value but none of
the other column values changed in the row, this process of row-by-row
checking eliminates these nil changes and informs downstream tables not
to incorrectly react.

To rephrase, a pipeline of incremental extract relies on change
detection columns, but when one such column goes wrong, the effect can
be disastrous. Instead of relying solely on the change detection column,
it should be used as a first pass, while allowing the Check step to do a
second pass to eliminate any noise if things go wrong.

Finally, the Check step improves stability and fault tolerance by
limiting the impact of abnormalities or errors. To take advantage of the
fault tolerance, tables must not have unstable columns that change
frequently for no informational reason, or may suddenly change on many
rows, such as on a switch to new year. Examples of these are discussed
in the Fault tolerance chapter.

On a specific table, the additional work may seem excessive. But doing
the work consistently as a pattern is the backbone of a sophisticated
pipeline by monitoring and controlling row changes. Investment in
automation or appropriate technology ensures that the additional logic
is applied by default, with no development cost to the engineer.

\cleardoublepage

\chapter{Load stack}\label{docs-efficient-stable-pipeline-load-stack}

The chapter Load mechanics look at the details of how a single table is
loaded. This chapter looks at how load across multiple tables can be
orchestrated across the pipeline.

In a data pipeline, tables need to be loaded in the correct order. Even
with a small number of tables, such as twenty to fifty, their
dependencies can become difficult to manage. Poor implementation, such
as manually coding their load order, leads to an unmaintainable pipeline
that is hard to debug.

In addition, loading tables in parallel is an easy way to reduce the
total amount of time to process a large pipeline. Even a four-fold
reduction may mean the difference between meeting the start of the
business day or not.

Load orchestration is the task of ensuring the loads in a pipeline
happen correctly, always in order, and in parallel if computational
resources are available.

There are two main ways of orchestrating a load:

\begin{itemize}
\item
  A centralised approach where a main process assigns work to workers
\item
  A decentralised approach where workers grab work from a queue
\end{itemize}

In the centralised approach, a software orchestrates other server jobs
to load tables in the right order and manages the parallelism. It is
often difficult for data engineering teams to create such software at
scale and requires the purchase of a vendor product.

On the other hand, the decentralised approach in this chapter is easy
for any team to implement. All it takes are:

\begin{itemize}
\item
  An accurate record of pipeline lineage as metadata, that is, how one
  table draws data from other tables in the pipeline.
\item
  The ability to load one table at a time, that is, there are no
  procedures or functions that load two tables in an entangled fashion.
\end{itemize}

The decentralised approach achieves the aim by maintaining a queue of
work, making the queue visible, and allowing workers to grab from the
queue. We call this queue the load stack.

\section{Load stack and load
candidates}\label{load-stack-and-load-candidates}

The load stack is a list of tables to be loaded. Each row represents a
table to be loaded and annotated with two columns: {[}Is started{]} and
{[}Is ended{]}. These indicate whether the load for that table has begun
or finished. If {[}Is started{]} = 0, the table has not yet started
loading and needs processing.

\begin{quote}
\textbf{Note.}

TODO: Insert manuscript screenshot or diagram from the source draft.
Source PDF note: \texttt{{[}Screenshot{]}}.
\end{quote}

But not every table with {[}Is started{]} = 0 is ready to load. Loads
must follow the correct order. A table is ready only when all the tables
it depends on have finished loading, a concept known as execution
readiness. These are the load candidates.

Identifying the load candidates requires a list of each table's load
dependencies. That is, the tables which directly feed into a table
through a load.

\begin{quote}
\textbf{Note.}

TODO: Insert manuscript screenshot or diagram from the source draft.
Source PDF note: \texttt{{[}Screenshot{]}}.
\end{quote}

It is instructive to visualise the dependencies in a directional graph.
Each arrow shows a direct feed from one table into another.

\begin{quote}
\textbf{Note.}

TODO: Insert manuscript screenshot or diagram from the source draft.
Source PDF note: \texttt{{[}Screenshot{]}}.
\end{quote}

By tracing dependencies upstream, the full set of prerequisites for each
table can be found. These are the expanded dependencies of a table.

\begin{quote}
\textbf{Note.}

TODO: Insert manuscript screenshot or diagram from the source draft.
Source PDF note: \texttt{{[}Screenshot{]}}.
\end{quote}

The load stack and the expanded dependencies are all that are needed to
identify the load candidates, that is, the tables that are ready to
load.

A table is a load candidate if there are none of the tables in its
expanded dependencies that have not ended. That is, there are no
dependencies where {[}Is ended{]} = 0. This logic also works for tables
with no dependencies at all, since their expanded dependency list is
empty and therefore contains no unfinished loads. When a table has no
unfinished loads in an upstream table, it is a candidate to start
loading.

This logic can be exposed as a view. The view shows the current list of
load candidates and can be queried at any time to find which tables are
ready to load.

\begin{quote}
\textbf{Note.}

TODO: Insert manuscript screenshot or diagram from the source draft.
Source PDF note: \texttt{{[}SQL\ and\ SCREENSHOT{]}}.
\end{quote}

\section{Using the load stack}\label{using-the-load-stack}

The load stack table and the load candidate view can be used to
implement a parallel load using the following algorithm.

\subsection{Step 1}\label{step-1}

At the beginning of the load, populate the load stack and set {[}Is
started{]} and {[}Is ended{]} to 0 for all rows. This is called
refreshing the load stack and may require deleting existing rows.

\subsection{Step 2}\label{step-2}

Start any number of workers. These could be functions in sub-processes
or stored procedures in independent sessions. Each worker continuously
polls the load stack and continues polling until there are no more
tables with {[}Is started{]} = 0.

If there are still tables to be loaded, the worker queries the load
candidate view to pick up a table that is ready to load.

If the view returns a table, the worker

\begin{enumerate}
\def\labelenumi{\arabic{enumi}.}
\item
  sets {[}Is started{]} = 1 for that table in the load stack. This is
  claiming the table for loading,
\item
  invokes the function or procedure to load the table, and
\item
  sets {[}Is ended{]} = 1 when the load finishes or is aborted due to an
  error.
\end{enumerate}

If the view returns no tables, the worker pauses for two seconds before
polling again.

The worker exits when there are no more tables in the load stack with
{[}Is started{]} = 0.

Claiming from the load stack must be done in a way that prevents
multiple workers from claiming the same table. This can be achieved by
wrapping the read from the load candidate view and the update to the
load stack in a transaction. This avoids race conditions.

In summary, this is all that is necessary to implement parallel load:

\begin{itemize}
\item
  The load candidate view, informed by expanded dependencies, ensures
  that a table is only offered for loading if it has no upstream tables
  still loading. This guarantees correct load order.
\item
  Multiple workers can pick up jobs as soon as they become available.
\item
  Transactions and locks around the claim step ensure that workers do
  not clash when selecting the same table.
\end{itemize}

While these steps are sufficient for parallel loading, it is useful to
log additional columns.

\begin{itemize}
\item
  Start and end datetimes for each table in the load stack.
\item
  A single {[}Workflow ID{]} assigned to every table in the load stack
  when it is refreshed. This ID represents the batch of load.
\item
  A {[}Worker ID{]} generated by each worker, recorded against the table
  to show how jobs are distributed.
\end{itemize}

\section{Advantages}\label{advantages}

The advantages of the load stack are numerous.

\subsection{No special code for
parallelism}\label{no-special-code-for-parallelism}

There is no special code for parallelism. The logic looks identical
whether one worker or many are used. The primary purpose of the load
stack and the load candidate view is to ensure that tables load in the
correct order. The ability to load multiple tables simultaneously is a
natural side-effect of the design. The only difference between serial
and parallel load is the number of workers started at a given moment.
This does not need to be known to any central orchestrator. In fact, the
load stack does not even track how many workers are active.

\subsection{Unlimited and dynamic
workers}\label{unlimited-and-dynamic-workers}

The number of workers is unlimited and can be turned on or off at any
time. Many parallel load algorithms decide on the number of workers at
the beginning and launch them accordingly. The load stack approach
allows the platform engineer to respond to server resources in real time
by adjusting the number of workers. Only minor changes are needed to
allow workers to exit the polling loop early.

\subsection{Greedy and balanced work
assignment}\label{greedy-and-balanced-work-assignment}

The algorithm is well load-balanced by default. It uses a greedy
approach to assigning work. As soon as a load candidate is ready, a
worker picks up the job. There are no situations where a table is ready
to load and a worker is available, but the load does not happen. While
this does not guarantee optimal load time, it minimises idle workers.
For further improvements, the load candidate view can be adjusted to
return the next table based on a more fine-tuned balancing logic.

\subsection{Partial pipeline support}\label{partial-pipeline-support}

It is easy to load only a portion of the pipeline. The expanded
dependencies show exactly which tables are upstream from any given set.
Loading only part of the pipeline can be achieved by populating the load
stack with just those tables. This allows a subset of the pipeline to
run more frequently or on a different schedule than the full pipeline,
without needing changes to a central orchestration plan.

\subsection{Mid-load interference}\label{mid-load-interference}

Since load is controlled by making tables available on the load stack,
the approach supports mid-load interference by editing the load stack.
For example, tables can be put back on the load stack during a load by
setting {[}Is started{]} = 0 for those tables. This supports error
correction and retries mid-load.

\subsection{Increased load frequency}\label{increased-load-frequency}

Most warehouses run loads in a daily batch, but some business needs
require more up-to-date information that could be hourly or minutes
latency, demanding a higher load frequency.

In practice, higher load frequencies are only practical if each of the
table in the pipeline is incrementally processed. It is unlikely that
every table in a large warehouse would be implemented this way. The load
stack makes it a possible to load any portion of the pipeline. This
means that it is easy to select a path within the overall pipeline for
increased load frequency while allowing the rest to happen in usual
batches.

In the most demanding case, editing the load stack mid-load can be used
to implement continuous loads. For example, workers start to process the
pipeline, while another process periodically resets one or more tables
by setting {[}Is started{]} = 0 on the load stack. This puts the tables
back on the queue for continuous processing.

The load stack allows a data engineering team to treat low and high
frequency loads with the same orchestration. A path within a pipeline
where all the tables are efficient can be loaded more regularly,
increasing the frequency from once a day to multiples times in an hour,
and ultimately to run continuously. It only depends on the investment in
making each of the table in the path to process within the required
response time.

This graduation allows ``just-enough'' investment to meet frequency
targets. This gives teams a continuous-load option before considering
the need to adopt dedicated streaming platform.

\subsection{Cross-technology
orchestration}\label{cross-technology-orchestration}

The load stack supports orchestration across multiple technologies.
While many pipelines run on a single stack, complex problems may require
different technologies for different parts of the load. The usual
solution is to purchase orchestration software with connectors for each
technology. The load stack reverses this. The only requirement is that
the central load stack table is visible to all technologies. There is no
need for a central worker to have connectors. For example, the load
stack and load candidate view may be implemented in a SQL database, with
claiming done through a stored procedure. As long as other technologies,
such as Spark or other SQL warehouses, can connect via ODBC, they can
participate. The load stack enables cross-technology orchestration with
minimal requirements.

\section{Conclusion}\label{conclusion-4}

Load orchestration is the task of ensuring tables in a data pipeline
load correctly with respect to their dependencies, even when doing so in
parallel.

It is easy for load orchestration to go out of hand, creating pipelines
that are opaque and impossible to maintain. When the pipeline becomes
complex, some organisations may purchase orchestration software and
additional servers to host them.

The load stack approach is a simple alternative. It does not require
additional servers or software. Only minor amendments using tools
already available to a data engineer are needed. The only requirements
for the data engineer are to capture each table with its direct
dependencies in a metadata table, and to modularise code so that each
table can be loaded independently. The load candidate view computes
execution readiness by consulting the state of the load stack to know
what jobs have been complete and comparing it to the metadata table of
dependencies. This ensures the load happens in the correct order no
matter the level of complexity in dependencies, and no matter the number
of workers, even in the absence of a centralised orchestrator.

Compared to other common approaches towards orchestration, the load
stack is a queue-based coordination adjusted for lineage-awareness. At
any given point, it uses the actual state of upstream tables to
determine readiness of tables to load. This allows the load to respond
immediately to changes in circumstances, something that is not easy in
other approaches.

It is technology-agnostic, lightweight, and easy to adopt --- which
makes it ideal for any team to implement with low investment cost.

\cleardoublepage

\chapter{Load
dependencies}\label{docs-efficient-stable-pipeline-load-dependencies}

One of the biggest factors impacting efficiency and stability is the use
of dependencies in a pipeline.

Dependencies occur when a materialised table selects from other tables
for its load.

This is a load dependency, because the materialised table needs to be
loaded in the pipeline. Load dependencies propagate information from
downstream.

Dependencies also occur through views. Views also propagate information,
but they do so indirectly by propagating logic rather than persisted
data. We distinguish them as a view dependency.

Load dependency is the focus of this chapter. Whenever the word
``dependency'' is used in this chapter, it refers to load dependency.

Dependencies improve efficiency because they allow reuse rather than
recalculation of the same piece of information. In this sense, they also
improve quality by reducing the risk of logic becoming out-of-sync.

However, dependencies can harm stability because failure in one table
directly propagates downstream, sometimes leading to domino effect of
issues.

Moreover, changing the business logic of a table early in the pipeline
theoretically requires all downstream tables to reload, in case the
change affects their results. In theory, it is possible to isolate the
impact by examining the logic change. In practice, this is too
error-prone, and it is safer to reload downstream tables. This may lead
to a situation where small changes upstream cause large changes
downstream. In this scenario of logic change, dependencies can be
harmful to both efficiency and stability.

They are inefficient because large amounts of processing may occur with
little change in outcome. They are unstable because during the reload,
many tables become temporarily empty as they are reloading.

Ultimately, at every turn the data engineer must choose between
re-calculating or re-using existing information. This is a decision
based on trading efficiency and stability.

\section{Criteria for dependency}\label{criteria-for-dependency}

There are three criteria to consider when committing to a dependency.
The dependency should be:

\begin{itemize}
\item
  Valuable -- the information is worth selecting and not easily
  recomputed.
\item
  Targeted -- the source table is narrow, rather than wide with
  miscellaneous columns.
\item
  Stable -- small changes in the source should not cause massive changes
  downstream.
\end{itemize}

Valuable Using dependencies to propagate information is convenient in
the short term. However, the issues that come with a heavy chain of
dependencies often take time to surface.

This can lead new data engineers to use dependencies too casually and
not notice the problem until it is too late.

The first step in managing dependencies is to ensure that the selected
information has enough value to justify the dependency.

For example, the data engineer should avoid bringing in ``convenience''
columns from other tables. A common case is reading the country name
from the reference country table simply for display or debugging
purposes early in the pipeline. Instead, the fact table may only need
the country code, and the join to the reference table can occur later in
the presentation layer through a view.

Another method is to recalculate where it is simple. For instance,
computing the first of month from a date column is preferable to looking
it up from a calendar table. This may involve repeating simple logic,
but it avoids forming a dependency. Similarly, using a standardised
default value for unknown references can reduce unnecessary joins. For
example, if the surrogate key for unknown values is always ``1'', this
number can be hardcoded in the transaction table with little risk,
avoiding the need to select from the reference table to find the
surrogate key value by name.

Targeted Dependencies should be targeted --- selecting just enough for
their purpose. When dependencies are not well targeted, more
dependencies are formed than are actually useful, thus diluting the
value of the dependency relative to its cost.

The smallest unit of dependency is a table. That is, when a table
selects a column from a source table, the entire source table becomes a
dependency. If many columns are packed into one table, then any table
that needs just one of those columns must depend on the whole. This
leads to a concentration of dependencies on that single table. Wide
tables with miscellaneous columns not only create tangled dependencies
but also reduce clarity of information flow.

As an exaggerated example, if there were a table called Cake.Everything,
and multiple tables selected from it, then any update to Cake.Everything
would require updates to all downstream tables. It would also be unclear
what information about cake is being passed downstream.

Meaningful fragments are excellent for breaking down dependencies to
ensure they are more targeted. Continuing the example, Cake.Everything
can be split into well-defined fragments such as Cake.Sales,
Cake.Milestones, and Cake.Quality. Selecting from these fragments helps
target the dependency and makes it clearer what information is being
selected.

It is often better to start a new fragment than to add a new column to
an existing table --- especially if the new column introduces a
different concept. For example, if there is a Cake.RefCalendar table
used to describe dates, and there is a need to include sales seasons, it
may be tempting to add a column like {[}Is sales season{]}. This is
problematic because it increases the weight of dependency on
Cake.RefCalendar when the sales season is needed for analytical
computation. Instead, it is better to create Cake.RefSaleSeason with a
primary key {[}First date of week{]} to indicate that sale seasons
operate at the week grain. Cake.RefSaleSeason would be a narrow table
with just enough information to describe the sale season. Downstream
tables would select from Cake.RefSaleSeason, offloading the dependency
from Cake.RefCalendar and making it clear that the table is interested
in sale seasons. Cake.RefSaleSeason would relate to Cake.RefCalendar via
a foreign key on {[}Calendar date{]}.

Stable Stability means ensuring that small changes in the source do not
cause large, unintended changes downstream.

A common example of instability is when a transaction table selects the
country name from a reference table. If the name of a country is
updated, that single change may trigger millions of downstream updates.
This kind of ripple effect can be costly and difficult to manage.

Where possible, dependencies should be stable. Where not possible, the
risks should be mitigated. The two main sources of instability are
row-wise instability and column-wise instability.

Row-wise instability occurs when a single row drives many downstream
rows. This is typical of reference tables. A small change in a reference
table can affect large volumes of transaction data. Metadata-driven
loads are another example. If a pipeline reads a metadata table to
decide what to load, an error in that metadata can cascade into
widespread changes.

Changes tend to be amplified when a dependency selects one row and
applies it across many rows --- for example, reading a small table into
a large one. Small tables are often a source of instability for this
reason.

There are several ways to mitigate row-wise instability. In the country
name example, the name may not need to be denormalised into the
transaction table at all. It can remain in the dimension table for
presentation, while the fact table stores only the country code. More
generally, moving information from transaction tables to reference
tables, such as adding computed columns to reference tables, and
allowing the reference tables to carry the information load is best
practice. Another option is to defer lookups until later in the pipeline
or perform them only through views.

Column-wise instability occurs when the source column itself is
volatile. For example, selecting {[}Update datetime{]} into a downstream
table introduces risk. A routine system update that refreshes all
timestamps may seem trivial to the application but can cause widespread
downstream updates.

To reduce this risk, a polling table can be used. This technique is
explained in the chapter Tracking changes.

Another example is a column like {[}Is archived{]}. A digital system may
perform bulk updates using this column. From the system's perspective,
the change may only affect what appears in the UI, but for the data
pipeline, it can have a significant impact.

Column-wise instability is less common than row-wise instability, but it
is important to monitor for both.

\section{The case of surrogate keys}\label{the-case-of-surrogate-keys}

Surrogate keys act as stand-ins for the primary key of a table. Most
commonly, they allow a composite primary key made up of multiple columns
to be replaced by a single integer. Surrogate keys may be created by
source systems, but they are often introduced by the warehouse as part
of the data pipeline.

Surrogate keys are necessary in a warehouse for several reasons. First,
and most importantly, Power BI supports only single-column
relationships. If the primary key consists of multiple columns, a
surrogate key is required to implement the relationship.

Surrogate keys also simplify Power BI measures that rely on distinct
counts, making implementation more straightforward.

Second, multiple-column keys can be clumsy to carry. This is especially
true for storytelling dimensions, where the primary key may consist of
several binary columns --- sometimes more than ten. It becomes
cumbersome to record all ten columns on each transaction table and join
on all of them. The same issue arises with type II keys such as
{[}Country code{]}, {[}Start datetime{]}. A lookup to such a table
involves an inequality join on the validity period {[}Start datetime{]}
and {[}End datetime{]}, which is error-prone when repeated across the
code base.

Third, some primary keys are slow to join on. This can happen with type
II keys due to their complex predicates, or when the primary key uses
data types that are not conducive to performant joins, such as long
strings.

For these reasons, it is often desirable to retrieve a surrogate key.
This involves looking up the table using its primary key and bringing in
the surrogate key column. In this way, surrogate key retrievals can be a
valuable dependency.

However, surrogate keys can become magnets for dependency. They are
particularly unstable. If the surrogate keys are regenerated, because
the table needs to be reloaded for amending a logic, all downstream key
values reshuffle.

There are two ways to manage surrogate key dependencies.

The first is to defer the lookup to later stages of the pipeline,
perhaps only during the final load into the Power BI semantic model. If
the original business key consists of only two columns and join
performance is acceptable, this deferral may be reasonable.

However, if many downstream tables rely on this key and the data types
are not ideal for joins, it may be better to accept the cost and swap
out the keys earlier in the pipeline.

The second approach is to compute the surrogate key rather than look it
up. For example, in a storytelling dimension with five binary columns as
the primary key --- {[}Is A{]}, {[}Is B{]}, {[}Is C{]}, {[}Is D{]},
{[}Is E{]} --- the surrogate key can be calculated by treating these as
a base-two number and converting it to base-ten using a simple formula:
{[}SK{]} = 1 + 2\footnote{Is A} + 4\footnote{Is B} + 8\footnote{Is C} +
16\footnote{Is D} + 32\footnote{Is E}

Using this formula, it is possible to derive the correct surrogate key
value without performing a lookup.

Whether using this or another formula, if the computation is simple,
calculating the surrogate key on the fly can provide the benefits of a
surrogate key without forming a dependency.

\section{Views as an alternative}\label{views-as-an-alternative}

Views are also a form of dependency that can propagate information.
However, they are not a direct substitute for load dependencies. Load
dependencies propagate information itself, while views propagate logic.
For example, when two tables are joined in a view, it is the logic of
the join, not the result of the join, that is passed downstream.

This makes view dependencies and load dependencies fundamentally
different.

However, views can be effective as an alternative in a range of
scenarios, especially when the dependency is not stable. For instance,
instead of denormalising a country name into a materialised table, a
view can perform the same join. In this case, a change to the country
name in the reference table does not trigger any update operations
downstream. However, when the user queries the view, they retrieve the
updated name. This is particularly useful in parameterised reports,
where users look up sections of data at a time and it is not necessary
to have the full set materialised in advance.

In theory, a server with infinite resources could support an entire
pipeline built from views. This would eliminate the kind of instability
caused by materialised dependencies. However, views can also amplify
instability by instantly propagating any errors downstream. Materialised
tables, by contrast, offer a buffer of stability for reporting.

Another limitation of views is that they do not track changes. As
explained in the Load mechanics chapter, the Check step of a load
verifies whether a row has genuinely changed. This allows downstream
workloads, such as machine learning predictions, to respond
appropriately. Views do not support this type of row-level change
detection.

Because views propagate logic rather than information, they also lose
historical data when source systems delete older records. This can be
problematic for pipelines that require long-term persistence.

Views can be an effective alternative to load dependencies under certain
conditions.

They are suitable when query performance is adequate, when instant
propagation of changes (including potential errors) is acceptable, when
long-term persistence is not required, and when row-level tracking of
changes is not needed for downstream processes.

\section{Conclusion}\label{conclusion-5}

A good pipeline should contain a healthy level of dependencies. This is
a sign that the data engineer is adding value to raw data through
transformation, aggregation, and storytelling. A pipeline with almost no
dependencies often reflects a garbage-in, garbage-out approach, or
complex logic buried in a few scripts.

As a rule of thumb, a source system with 10 to 20 raw tables may result
in 5 or 6 layers of dependencies. In larger systems with 30 to 50
tables, especially when integrated with other systems, the
transformation pipeline may easily reach 10 or more layers. In the site
where this work is done, there are 10000+ tables, 15000+ dependencies
between them, amounting to 25+ layers of dependencies.

Dependencies are an effective tool for propagating information and
achieving efficiency. They allow the data engineer to build reusable
blocks of insight, articulate business meaning, and pass them for use.
But they must be used with care. Each dependency increases entanglement
and can reduce stability. A data engineer must be deliberate in
committing to a dependency by selecting only those that are valuable,
targeted, and stable.

\cleardoublepage

\chapter{Tracking
changes}\label{docs-efficient-stable-pipeline-tracking-changes}

The aim of information efficiency is to process only what has changed.
If there were no records changed after a table is processed, then ideal
is have spent zero processing time it. In practice, even when there are
no changes in input, a load still takes time to process.

Every unnecessary scan consumes resources without adding value.
Efficiency is achieved when downstream tables can reliably extract only
the records that have changed and ignore the rest.

Doing this well requires systematic change tracking. Without it,
pipelines fall back bespoke methods for tracking changes. This leads to
errors in incremental processing, periodic full reloads, and fragile
recovery from disruptions. A robust pattern avoids these pitfalls.

The following is a simple approach illustrate the concept of change
tracking, before unpacking the details of a more robust approach.

\section{Simple approach}\label{simple-approach}

The most basic approach is to include an update datetime column from the
source system and filter on it. This works when the source table deletes
do not occur, and the target table mirrors the source one to one.

\subsection{Example scenario}\label{example-scenario}

\begin{itemize}
\item
  Source table Raw.Event with columns {[}PK1{]}, {[}PK2{]},
  {[}Col1{]}\ldots{[}Col15{]}, {[}Update datetime{]}
\item
  Target table Curated.Event with {[}PK1{]}, {[}PK2{]}, {[}Col1{]},
  {[}Col2{]}, {[}Update datetime{]}
\end{itemize}

Incremental extract logic would be as follows:

While straightforward and intuitive, this approach suffers from a few
problems.

First, it is not scalable to complexity. Additional source tables will
break the logic because it is not practical to add the update datetime
from each source onto the target table. It is possible to create a
complex compound datetime, but it would be error prone.

Second, it is not stable. If a system update shifts all values in
{[}Update datetime{]}, the pipeline will reprocess every row, causing a
blow out in processing time for updates, and triggering similar issues
downstream.

Third, the approach is not scalable to continuous load. For example, if
Curated.Event filters Raw.Event for rare events, then the maximum update
datetime on Curated.Event will not represent latest processing time of
Raw.Event. The simple approach will force a re-scan that goes too far
back. For a continuous load scenario, this will be prohibitive.

A more scalable and reusable pattern is to separate the time-tracking
artefacts from the data content itself.

\section{Refresh bookmarks and polling
tables}\label{refresh-bookmarks-and-polling-tables}

As explained in Load mechanics, as each table is processed, the pipeline
logs the refresh bookmark of a load. That is, the starting datetime of
the load of each table is logged if the load succeeded. The refresh
bookmark for all tables in the pipeline can be stored in one table. Each
row of this table is one table that was loaded, together with the
starting datetime of that load.

This bookmark defines the cut off for what records in a source has fully
made it to the processing of a target table. Any records that have come
to the source table after this bookmark datetime would need to be
considered in the next batch of processing. If the table load aborted
because of a fault, then the bookmark should not be logged so that the
next batch can resume from last success point.

\subsection{Advanced note}\label{advanced-note}

If the source is continuously loading, then some stray records may still
come into the source table after the bookmark datetime, but before the
processing has complete. However, there are no guarantees. Consequently,
either the incoming source needs to be frozen, or the next batch should
still resume from bookmark and reprocess these late records.

The query to retrieve the latest bookmark for a table is simple:

On the other hand, keeping track of source changes depends on whether
the source table was created by the pipeline or whether it was raw data
table from external source.

In the former case, the tables created by the pipeline have reliable
update datetimes created by the process in Load mechanics -- {[}Row
insert datetime{]}, {[}Row update datetime{]} and {[}Row delete
datetime{]}. Since these are created by the pipeline itself, they are
comparable with the refresh bookmark and can be compared directly for
understanding what records have come in after the latest bookmark.

However, for tables that were not created by the pipeline, even though
they may have a column called {[}Update datetime{]}, there may be a
substantial lag between that value versus when the row made it to the
database. Hence, they cannot be used directly.

To distinguish these two cases, we would say that the update datetimes
in pipelines are in-sync while the update datetimes from source systems
that are loaded in a separate process are out-of-sync with the data
pipeline.

If the datetimes are out-of-sync with the refresh bookmark, they cannot
be directly compared. Instead, just as the target table's state of
processing can be tracked by a separate table, the source table's state
of processing can be tracked by a separate polling table.

A polling table is any table that can be consulted to know whether and
what portion of data needs to be refreshed. For a polling table to be
effective, it needs to be structured so that it can be queried rapidly
-- much more rapidly than the time to process the full data. Ideally,
this query should be close to zero time to support continuous load.

Following the example, direct comparison of Raw.Event{[}Update
datetime{]} with the latest bookmark in Curated.Event is not possible.
There can be an unknown lag between the update timestamp and its arrival
in the database. The only safe assumption is that the datetime increases
monotonically.

A polling table provides a way to map the pipeline's refresh time to the
latest update datetime observed in Raw.Event. Suppose this table is
called Raw.Bookmark, and is loaded using an append-only process. Each
appended row records the maximum update datetime from Raw.Event at the
time of the refresh.

The columns of Raw.Bookmark are:

\begin{itemize}
\item
  {[}Refresh datetime{]} --- the datetime of the check as the row is
  appended. This datetime, since it is computed by the load, is part of
  the pipeline and thus in-sync with the refresh bookmarks.
\item
  {[}Bookmark datetime{]} --- the maximum of Raw.Event{[}Update
  datetime{]} as at that time of refresh. This datetime, coming from the
  source and later loaded, is out-of-sync.
\end{itemize}

This table maps the {[}Update datetime{]} column, which is out-of-sync
with the pipeline, to a column, {[}Refresh datetime{]}, that is in-sync
with the pipeline. This mapping makes it possible to compare changes to
the incoming data tracked by the source's {[}Update datetime{]} with
processing events, which are tracked against the pipeline's record of
datetimes.

In theory, refreshing this table infinitely often would allow the
pipeline to translate its own time into the update datetime of Raw.Event
at any point. In practice, the refresh frequency aligns with the
pipeline cadence. There is no benefit in refreshing the table when
Raw.Event has not changed. In addition, Raw.Bookmark can contain the
bookmarks for multiple tables, such as Raw.Event, Raw.Event2, Raw.Event3
rather than creating one table for each.

At the start of processing Curated.Event, the polling table Raw.Bookmark
allows the pipeline to fetch the maximum bookmark of Raw.Event as of
Curated.Event's previous batch start. Any row on Raw.Event whose
{[}Update datetime{]} represents data that have come after the previous
batch start and eligible for processing.

In summary, tracking records that have changed in a source table since
the last process depends on two components:

\begin{itemize}
\item
  Tracking the target. A bookmark table for the pipeline that records
  the processing datetimes of each table whenever a load completes
  successfully.
\item
  Tracking the source. If the source table has an {[}Update datetime{]}
  column that is synchronised with the pipeline, this column can be used
  directly to retrieve updated rows. If not, a polling table is required
  to map the source table's update datetime back to the pipeline's
  datetime for comparison.
\end{itemize}

With these components, it becomes possible to identify changes in the
source relative to the target's last processed state.

\section{Change detection columns}\label{change-detection-columns}

The previous example assumes that the source table Raw.Event has a
reliable {[}Update datetime{]} column for change detection. This is not
always the case.

In the ideal scenario, the change detection column is an architectural
column from the source system. Examples include server constructs,
trigger-managed datetimes, or built-in datetimes from commercial
products that can be trusted. This is also true for pipeline tables
loaded as described in Load mechanics.

Less ideally, the change detection column is managed by the source
application through application logic. These columns are generally
reliable but can be prone to developer error or fail to capture direct
database changes.

Finally, there may be no change detection columns at all. In this case,
it is not possible to determine whether a source row has been updated.
However, it may still be possible to identify inserted records by
performing an anti-join between the source and target tables. This
approach is useful for append-only loads such as building of hub tables.

\section{Tracking deletes}\label{tracking-deletes}

Polling tables can be used for tracking updates and inserts. Tracking
deletes is equally important. However, this can be far more difficult
because, in the case of deletes, rows disappear completely from the
table itself.

In the ideal scenario, deletes are tracked in the source through
architectural artefacts.

Some database technologies implement default history tables or delta
logs. These history tables allow reliable tracking of deletes, provided
they include a change detection column to indicate when the delete
occurred.

Less ideally, some applications maintain dedicated business audit tables
that track changes, including deletes. These tables can be highly
reliable, though they are often difficult to query.

In the worst case, there is no change tracking. Rows can disappear
without a trace. The only option is to compare the source and target
using an anti-join on the primary key.

\section{The role of the Filter step}\label{the-role-of-the-filter-step}

Two themes recur, whether for tracking upserts or tracking deletes:

\begin{itemize}
\item
  Source update datetimes that are synchronised with the pipeline's
  process datetimes are trivial to query.
\item
  Architectural artefacts, rather than business artefacts, are the most
  reliable way to track inserts, updates, and deletes.
\end{itemize}

When either of these is absent, tracking changes becomes complex or even
impossible.

This is one of the reasons why pipelines begin with a Filter step that
keeps transformation minimal. The filter step is designed to be the
first interaction with the source data. On first arrival of the data,
the pipeline adds the architectural artefacts required for change
tracking. Doing this once at the start allows the rest of the pipeline
to apply transformations incrementally with reliable change tracking.

The optimal case is if the incoming data can be incrementally extracted
using polling tables and the filter step adds the architectural columns
for change-tracking. This effectively converts the source change
detection columns, which may be out-of-sync with the pipeline, to change
detection columns which are in-sync the datetimes tracked by the
pipeline.

The Filter step, which is a simple extract of the necessary rows and
columns with minimal joins or aggregations, may look like a step with
low value-add. However, this step is serving as a critical foundation
for an efficient pipeline because proper load mechanics computes
change-tracking artefacts to support efficient downstream response.

\section{Conclusion}\label{conclusion-6}

An efficient pipeline requires knowing exactly what changed in the
source so that it responds only to those records. These changes include
inserts, updates, and deletes.

A simple approach to tracking upserts is to include the source change
detection column in the target table. However, this can be problematic
for non-trivial extracts and is at risk of instability. Instead,
bookmarks to track the target and polling tables to track the source are
reliable patterns for tracking changes relative to each other.

Deletes are much harder to track without additional artefacts or audit
tables. If these artefacts are not available, a full comparison of each
batch is necessary to identify changes for upserts, deletes, or both.

Given the complexity of tracking sources and the fragility of relying on
other systems to maintain reliable architectural columns, it is
important that each pipeline begins with an uncomplicated extract from
the source data. The Load mechanics process annotates this extract with
reliable change detection artefacts for downstream processing.

\cleardoublepage

\chapter{Responding to
change}\label{docs-efficient-stable-pipeline-responding-to-change}

The previous chapter Tracking changes focused on tracking changes in the
source. The chapter builds on the previous and addresses the harder task
of responding correctly to these changes. This is often known as
incremental extract and load.

Accurately responding to change can be error-prone. This is because
inserts, updates, and deletes upstream do not translate directly into
the corresponding changes downstream. The following examples illustrate
the possible complexities:

\begin{itemize}
\item
  Inserts can trigger deletes. The Bank.CustomersToFollowUp is a
  computed table of customers who have not made recent deposits and
  require follow-up by the service team. This table takes input from the
  Bank.Transactions table. In this case, an insert into
  Bank.Transactions, such as a new deposit, would trigger a delete in
  Bank.CustomersToFollowUp.
\item
  Deletes can trigger updates. The Bank.AccountSummary is a computed
  table of account-level aggregates, such as total number of holders and
  balance per account. It takes input from Bank.AccountHolder, which
  records all individuals associated with each account. If a joint
  holder is removed, resulting in a delete from Bank.AccountHolder, this
  changes the holder count and may trigger an update in
  Bank.AccountSummary.
\item
  Updates can trigger inserts. The Bank.GoldCustomer is a computed table
  of customers whose account balance exceeds a defined threshold. It
  takes input from the Bank.AccountBalance table. In this case, updating
  a row in
\end{itemize}

Bank.AccountBalance to increase the balance may cause a customer to
newly qualify for gold status, triggering an insert into
Bank.GoldCustomer.

These examples are highly artificial. However, it is easy to imagine the
possibilities under general circumstances:

\begin{itemize}
\item
  The presence of anti-joins (e.g. not exists) means that inserts
  upstream can trigger deletes downstream.
\item
  The presence of aggregation or windowing (e.g. group by) means that
  deletes upstream can trigger updates downstream.
\item
  The presence of filters on derived expressions (e.g. having) means
  that updates upstream can trigger inserts downstream.
\end{itemize}

The following is an illustrative list of correct ways to respond to
upstream changes, and a pattern for applying this response to
incrementally loading a table.

\section{Analysing the query}\label{analysing-the-query}

Suppose there are two tables X and Y, and these are used as inputs to a
data pipeline.

The set up is as follows:

\begin{itemize}
\item
  X has columns {[}Header ID{]}, {[}Value{]}, {[}Status{]}, where
  {[}Header ID{]} is primary key.
\item
  Y has columns {[}Header ID{]}, {[}Line number{]}, {[}Value{]},
  {[}Status{]} , where {[}Header ID{]} and {[}Line number{]} is primary
  key.
\end{itemize}

And Y{[}Header ID{]} is a foreign key to X{[}Header ID{]} and joins on
this column.

There could be inserts, updates and deletes on X, as well as on Y. Thus,
there are six possible input changes each time. In practice, inserts and
updates and can often be considered as a single upsert. Thus, the inputs
are:

\begin{itemize}
\item
  upserts from X
\item
  deletes from X
\item
  upserts from Y
\item
  deletes from Y
\end{itemize}

Suppose the target table is T. The goal is to determine which rows in
the target table need to be inserted, updated, or deleted. This is done
by computing a driver set of primary keys that might be affected by
changes in source tables (X and Y). Inserts and updates can be
considered as one. Thus, the goal is to calculate:

\begin{itemize}
\item
  primary keys to upsert on T
\item
  primary keys to deletes on T
\end{itemize}

The way to calculate these depends on the nature of the query. Some
possible scenarios include:

\begin{enumerate}
\def\labelenumi{\arabic{enumi}.}
\item
  A straight select of X
\item
  A filter of X on status
\item
  X inner join Y
\item
  X inner join Y group by {[}Header ID{]} to sum over Y{[}value{]}
\item
  X left join Y coalesce on {[}Line number{]} to 0 where Y{[}Header
  ID{]} does not exist
\item
  Y left join X
\item
  X left join Y group by {[}Header ID{]} to sum over Y{[}value{]}
\item
  X union Y
\item
  X left anti-join Y
\end{enumerate}

All of these would require a different treatment. These are summarised
in the table:

{\def\LTcaptype{none} % do not increment counter
\begin{longtable}[]{@{}
  >{\raggedright\arraybackslash}p{(\linewidth - 8\tabcolsep) * \real{0.2000}}
  >{\raggedright\arraybackslash}p{(\linewidth - 8\tabcolsep) * \real{0.2000}}
  >{\raggedright\arraybackslash}p{(\linewidth - 8\tabcolsep) * \real{0.2000}}
  >{\raggedright\arraybackslash}p{(\linewidth - 8\tabcolsep) * \real{0.2000}}
  >{\raggedright\arraybackslash}p{(\linewidth - 8\tabcolsep) * \real{0.2000}}@{}}
\toprule\noalign{}
\begin{minipage}[b]{\linewidth}\raggedright
Scenario
\end{minipage} & \begin{minipage}[b]{\linewidth}\raggedright
Primary key
\end{minipage} & \begin{minipage}[b]{\linewidth}\raggedright
Approach
\end{minipage} & \begin{minipage}[b]{\linewidth}\raggedright
Primary keys to upsert in T
\end{minipage} & \begin{minipage}[b]{\linewidth}\raggedright
Primary keys to delete in T
\end{minipage} \\
\midrule\noalign{}
\endhead
\bottomrule\noalign{}
\endlastfoot
1. Straight select of X & {[}Header ID{]} & Simple case: X changes map
directly to T. & Headers from X that were inserted or updated & Headers
from X that were deleted \\
2. Filter of X on status & {[}Header ID{]} & Filter condition means
updates can add or remove rows. & Headers from X that were inserted or
updated and now pass the filter & Headers from X that were deleted, or
updated and no longer pass the filter \\
3. X inner join Y & ({[}Header ID{]}, {[}Line number{]}) & Join means
changes on either side can add or remove pairs. & Pairs for Y rows
inserted/updated, plus pairs for X headers inserted/updated with
matching Y rows & Pairs for Y rows deleted, plus all pairs for headers
deleted from X \\
4. X inner join Y with group by {[}Header ID{]} & {[}Header ID{]} &
Deleting a Y row can change the aggregate or remove the header. &
Headers from X or Y that changed and now have Y rows; include changes in
Y, deletes in Y can also change the aggregated value & Headers deleted
from X, or headers that now have zero Y rows \\
5. X left join Y with synthetic line 0 & ({[}Header ID{]}, {[}Line
number{]}) & Synthetic line 0 appears only when no Y rows exist. & Pairs
for Y rows inserted/updated; plus (Header ID, 0) for headers with no Y
rows; plus pairs for new X headers & Pairs for Y rows deleted; remove
(Header ID, 0) when Y rows exist; remove all pairs for headers deleted
from X \\
6. Y left join X & ({[}Header ID{]}, {[}Line number{]}) & T follows Y's
lines; X changes matter only if Y lines exist. & Pairs from Y
inserted/updated; plus pairs under headers in X that changed & Pairs
from Y that were deleted \\
7. X left join Y with group by {[}Header ID{]} & {[}Header ID{]} & X
controls row presence; Y only changes the aggregates. & Headers in X
inserted; headers where values changed due to X or Y changes, including
deletes from Y & Headers deleted from X \\
8. Union of {[}Header ID{]} & {[}Header ID{]} & Header exists if present
in either X or Y. & Headers inserted/updated in X or Y; headers
projected from one side but still exist on the other & Headers deleted
from both X and Y \\
9. X left anti-join Y & {[}Header ID{]} & Row appears only when X exists
and Y does not. & Headers inserted/updated in X with no Y rows; headers
where Y rows were deleted and now qualify & Headers deleted from X;
headers that gained a Y row \\
\end{longtable}
}

These show an array of considerations for two tables. If there are more
tables, even up to 10, the analysis can become complicated.

Consider a query that draws from several tables, such as:

The first step is to determine conditions for upsert in the target.

For each source table, evaluate how changes in that table may influence
the result of the query. In general:

\begin{itemize}
\tightlist
\item
  An upsert in a source table typically results in an upsert in the
  target.
\end{itemize}

Any insert or update in A, B, or C may alter projected values and
therefore requires re-computation.

\begin{itemize}
\tightlist
\item
  A delete in a source table may also result in an upsert in the target.
\end{itemize}

This occurs when the query includes:

\begin{itemize}
\tightlist
\item
  Left joins, where the absence of a row changes a value to a default
  rather than removing the target row.
\item
  Aggregations, where the removal of a contributing row modifies a
  computed measure.
\end{itemize}

Thus, for each table, the question is:

``Could an insert, update, or delete in this table change the values
produced by the query?''

If so, the corresponding key should be included in the driver set for
upserts.

The next step is to determine conditions for deletes from the target.

In most cases, this occurs only when the tables that governs row
presence loses the corresponding row. For example:

\begin{itemize}
\item
  If the query begins with A left join B left join C ..., then deleting
  a row from A removes the target row. Deleting a row from B or C
  generally does not remove the target row; it only changes values.
\item
  If the query begins with A inner join B left join C, then deleting a
  row from A or B removes the target row.
\end{itemize}

But this can also happen through updates or inserts:

\begin{itemize}
\item
  In the case of A left anti-join B, an insert to A can cause a row to
  delete.
\item
  In the case of A filtered on a row value, an update to the value can
  cause a delete.
\end{itemize}

In summary, upserts arise from any change that affects projected values.
Deletes occur only when the base row is removed. Filtering predicates
and grouping clauses may cause deletes in a source table to behave as
upserts in the target, as they alter aggregates or conditions.

The same reasoning applies to more complex constructs such as unions or
set operations, as these are ultimately composed of similar units.

Working through these scenarios can be error prone. Latest innovations
in generative AI can provide valuable assistance in this analysis. For
example, the table of solutions generated for the 9 scenarios above was
generated by GPT-5 using only column 1.

\section{Applying the change}\label{applying-the-change}

Analysing the query with respect to changes from the source tables
informs an incremental extract. Recall from Load mechanics that the
difference between incremental extract and full extract are:

\begin{itemize}
\item
  The staging table for the load has only a minimal set of records that
  is much smaller than the full set, but still covers all the records
  that would need to be upserted in the current batch. This is what
  makes the load fast.
\item
  Deletes cannot be done automatically by comparing the full set of
  primary keys between the staging and the target. Instead, it needs to
  be customised for the query by analysing the impact of changes in
  source.
\end{itemize}

The detailed implementation of an incremental extract follows a
consistent pattern.

\subsection{Write query as normal}\label{write-query-as-normal}

Begin with the full query that expresses the business logic. This is the
reference point for all subsequent steps.

\subsection{Fetch the bookmark}\label{fetch-the-bookmark}

At the start of the load, retrieve the bookmark that records the last
successful refresh.

This value defines the boundary for change detection. The bookmark for
all tables loaded by the pipeline should be logged to a central location
and easily retrieved.

\subsection{Create minimal staging
table}\label{create-minimal-staging-table}

Using the analysis of the query, create a temporary table of keys to
upsert to the target.

In these examples, the driver table has the full key. In practice, it
may be a primary key part in the case of multi-columns primary keys.

Attach the driver set as an inner join to the original query. If
performance requires, add a clustered index to the driver table.

This is now a minimal staging table for incremental load.

\subsection{Create the delete set}\label{create-the-delete-set}

Using the analysis of the query, create a temporary table of keys to
delete from the target. This can be used to delete records that should
no longer be retained by logic of the query.

Retrieving the records that are deleted would vary depending on how they
are tracked.

Consider the case where overwritten or deleted records are moved to a
history table.

Then the examples would be:

This step works even if the source rows were inserted and deleted in
between the load of the target table, for example, if the target table
failed to load in the interim.

Inexperienced engineers can make the mistake of over-deleting with the
idea of re-inserting extra records. This can be tempting because a
delete + insert can be easier to calculate than a finely targeted delete
+ upsert. However, this approach is harmful because it is a source of
potential instability. In unpredictable circumstances, this may trigger
the whole table to delete. Instead, deletes should be perfect and remove
only rows that no longer satisfy the query's presence rule.

The full script follows.

This workflow offers several advantages.

Uniformity. The same query logic underpins both full loads and
incremental loads. The difference lies only in the addition of a driver
set join.

Idempotency. The process works regardless of how many times it runs
within an interval. If a load fails and is retried, the bookmark ensures
correctness. If the load runs more frequently than usual, the outcome
remains consistent.

Graceful fallback. If upstream changes touch many rows, say due to a
system update, the driver set expands and the incremental load behaves
like a full load. The fully load staging table would be compared to the
target and only genuinely updated records are applied. There is no
explosion of deletes that would cause a blow out in operations.

\section{Best practice workflow}\label{best-practice-workflow}

Given the complexity of responding to change, this section outlines a
step-by-step recipe for developing and validating pipelines that respond
to upstream inserts, updates, and deletes. The workflow balances
technical correctness with performance assurance and safe deployment.

\subsection{Start with a unit test}\label{start-with-a-unit-test}

\begin{itemize}
\item
  Compare row counts between the target table and the expected query.
\item
  Use a datetime filter for performance, but ensure it's independent of
  the extract datetime.
\end{itemize}

\subsection{Simulate data needing
update}\label{simulate-data-needing-update}

\begin{itemize}
\tightlist
\item
  Load the table in full, then wait for source data to change to
  simulate an incremental load.
\end{itemize}

\subsection{Build the upsert driver}\label{build-the-upsert-driver}

\begin{itemize}
\item
  As each source table is added, re-run the driver query to check
  performance.
\item
  Apply indexes if performance degrades.
\end{itemize}

\subsection{Test the query filter}\label{test-the-query-filter}

\begin{itemize}
\tightlist
\item
  Ensure the inner join between the driver and the original query is
  performant.
\end{itemize}

\subsection{Calculate deletes}\label{calculate-deletes}

\begin{itemize}
\tightlist
\item
  Ensure the delete driver query runs quickly.
\end{itemize}

\subsection{Simulate deletes}\label{simulate-deletes}

\begin{itemize}
\tightlist
\item
  Rather than actually deleting, use a select to simulate deletes during
  development to avoid accidental data loss.
\end{itemize}

\subsection{Apply changes and
validate}\label{apply-changes-and-validate}

Uniqueness, existence, and stability violations

\begin{itemize}
\item
  Apply upserts and deletes.
\item
  Compare the result with a fully loaded copy of the table.
\item
  The unit test should pass.
\end{itemize}

\subsection{Zero-load benchmark}\label{zero-load-benchmark}

\begin{itemize}
\item
  After loading, re-run the upsert and delete drivers.
\item
  They should return no rows and run in near-zero time.
\item
  The time it takes to do this is the fastest the query can ever run.
\end{itemize}

\subsection{Run incrementally over
time}\label{run-incrementally-over-time}

\begin{itemize}
\item
  Continue loading incrementally over multiple days.
\item
  The unit test should continue to pass.
\end{itemize}

\section{Conclusion}\label{conclusion-7}

A transformation query can be seen as a formula that takes inputs and
produces output. Information efficiency is achieved when changes in the
input are accurately tracked and correctly translated into changes in
the output. The difficulty of doing this depends not on the volume of
data, but on the shape of the query. Anti-joins, aggregations, filters,
and window functions all introduce complexity in determining which rows
to upsert and which to delete. This determination should be made by
analysing the query through the lens of relational algebra. Modern
tools, including generative AI, can assist in this analysis by surfacing
the logical structure of the query and identifying the propagation of
change.

Once the analysis is complete, the transformation can be converted from
full to incremental using a standardised approach. This is to fetch the
refresh bookmark of the target, compute the upsert driver table,
downfilter the original query, and compute the delete driver table. Care
must be taken not to over-delete and reinsert. Deletes should perfectly
represent only the rows that should no longer be part of the target as
of query.

When coupled with the use of refresh bookmarks and the techniques
described in Tracking changes, this approach yields an efficient and
resilient pipeline. Given the intricacies involved, careful
implementation and testing is essential. A unit test is critical to
catch errors during development, and silent failures in production.

Performance evaluation should be done step by step, with each component
tested in isolation. The zero-load scenario provides a baseline
benchmark for performance and correctness. This disciplined workflow
ensures that the pipeline responds to change with both accuracy and
speed.

\cleardoublepage

\chapter{Optimising Power BI
load}\label{docs-efficient-stable-pipeline-optimising-power-bi-load}

Loading a Power BI model is often one of the longest steps. This is
particularly true for models with large fact tables. New technologies
such as DirectLake promise to remove this step entirely, yet Import mode
remains the most reliable way to deliver a fast and responsive
experience. Consequently, it is important to establish efficiency for
Power BI loads.

Loading of Power BI should focus on assembling prepared fragments. It is
not the stage to curate information. The restaurant analogy from
Meaningful fragments makes the point. Mise en place is the long process
before service time to customers. During mise en place, ingredients are
washed, cut and portioned in advance. At service the kitchen completes
the final assembly and plates dishes quickly, rather than peeling
potatoes when the orders are made. Likewise, a Power BI load should be
focusing on assembling finished fragments rather than transforming data.

There are three main ways to optimise Power BI loads:

\begin{itemize}
\item
  Using DirectQuery to avoid loading
\item
  Efficient underlying source tables
\item
  Partitioning for parallel refresh, with two further enhancements of
  rolling windows and incremental refresh
\end{itemize}

This chapter assumes the model sources data from SQL, where
purpose-built views or tables act as source tables that map one-to-one
with the Power BI model tables.

\section{Using DirectQuery to avoid
loading}\label{using-directquery-to-avoid-loading}

DirectQuery for fact tables avoids the load entirely by sending DAX
queries to the source at report time. To provide a responsive
experience, the underlying source table often needs to be materialised
as a columnstore table dedicated to this purpose. While materialising a
large fact table is costly, it can also be incrementally loaded so that
refresh remains minimal.

By using DirectQuery over am incrementally loaded columnstore table,
this achieves the highest possible efficiency in terms of information
movement. However, there are serious drawbacks for user experience:

\begin{itemize}
\item
  DAX queries become noticeably slower.
\item
  Some complex DAX expressions are not supported at all.
\item
  Certain DAX functions behave differently in Import mode versus
  DirectQuery, such as in how blanks are treated.
\item
  Power BI limits how many rows can be retrieved in DirectQuery for
  certain operations, which restricts what reports can be built.
\item
  Dimension values can appear blank in filter lists even when the table
  does not actually contain blanks.
\end{itemize}

DirectQuery works best when:

\begin{itemize}
\item
  The table is far longer than it is wide, so queries touch fewer
  attributes.
\item
  Queries against the table are simple -- limited to sums, counts,
  minimums, and maximums that work well against columnstore tables
  through their equivalent
\end{itemize}

SQL query.

\begin{itemize}
\tightlist
\item
  Users are not expected to browse unit records and do not hit Power
  BI's retrieval limits.
\end{itemize}

Where dimensions are shared by both Import and DirectQuery facts, dual
mode must be implemented. Dual mode keeps a copy of the dimension for
Import while still allowing the engine to query it in DirectQuery when
needed. This ensures relationships remain fast for Import facts and
correct for DirectQuery facts.

\section{Efficient underlying source
tables}\label{efficient-underlying-source-tables}

One of the biggest factors influencing Power BI load times is the
efficiency of the underlying source tables, which are often implemented
as views. A source view should be a straightforward join of ready-to-use
tables that have been indexed appropriately to support high-performance
execution.

The views may have simple logic such as:

\begin{itemize}
\item
  Looking up surrogate keys to replace composite primary keys
\item
  Backfilling null values with defaults
\item
  Adding on-demand columns, for example {[}Days expired since
  creation{]} using functions like GETDATE()
\end{itemize}

However, heavy transformations such as windowing or nested logic should
be avoided.

These belong upstream in curated layers where they can be tested and
reused. String aggregation for display should also be avoided, as this
can be handled with appropriate DAX measures, as explained in Designing
Measures.

If a source table still requires elaborate logic to stitch pieces
together, the issue is not SQL tuning. The solution is to reconsider
whether upstream fragments have been adequately prepared. The ideal case
is to create fragments that are both meaningful and effective for
loading into the Power BI model.

Materialisation is sometimes necessary to ensure the source table loads
rapidly. This is especially true when data needs to be rearranged to
support fast retrieval of load partitions against the partition key - a
subject we turn to now.

\section{Partitions}\label{partitions}

Defining a table into partitions is the most effective way to reduce
Power BI load times when fact tables grow large.

Partitions are akin to having a large box of books from multiple
publishing years.

Batches of books arrive at regular intervals, and the box needs to be
updated with the incoming batch. There can be new books, and old books
can be updated. Updating the whole box can take a long time if the
number of books is huge. Splitting them into boxes by publishing year
allows a divide-and-conquer approach and vastly reduces total update
time. The ability to load partitions in parallel is the first benefit of
partitioning.

If boxes of books that are too old become unnecessary, they do not need
updating and the oldest set can be dropped. Over time, this reduces the
number of boxes to update.

This ability to apply a rolling window is the second possible benefit of
partitions.

Finally, if the incoming books can be sorted by year, and in each batch
only a fraction of the years have new books, then only those boxes need
updating. It may take additional effort to track changes from the
incoming books, but this effort can dramatically reduce total update
time. This ability to incrementally refresh is the third benefit of
partitions.

Thus, partitioning a table offers increasing levels of enhancement that
can be adopted as the situation requires:

\begin{enumerate}
\def\labelenumi{\arabic{enumi}.}
\item
  Load partitions in parallel, immediately reducing time to load
\item
  Apply rolling window to drop off older partitions that are not
  necessary
\item
  Incrementally refresh only the partitions that changed
\end{enumerate}

\subsection{Parallel partitions}\label{parallel-partitions}

A partition key is a business date or datetime column that divides the
data into intervals such as daily, monthly, or yearly. Power BI can load
these partitions in parallel.

Depending on the degree of parallelism, this can lead to massive
improvements in load time with little effort.

To make partitions effective, the source must return a single partition
quickly. If the underlying view is complex, the solution is not to push
SQL harder but to move complexity upstream and materialise the result
with an index on the partition key. This ensures that Power BI can
request ``give me this month'' and receive the data immediately.
Continuing the analogy, the incoming batches of books must be readily
accessible by publishing year. If not, it may take longer to fetch each
year's books than the time gained from divide-and-conquer. The optimal
solution is to organise the incoming batches by publishing year. In SQL
terms, this means storing the source data sorted by the partition key
used to define the Power BI partition.

There is a trade-off between load time and query time. Splitting one big
box of books into multiple boxes by publishing year increases the total
space used for storage because each box has a minimum size that can lead
to wasted space. It can also take longer to fetch books because there
are now more boxes to work through. The choice of interval size between
publishing year, first of month, or publish date depends on the actual
volume of books in the processing.

When partitioning, new developers often confuse which date to use. The
partition key must be a business date rather than a change-tracking
column like {[}Row update datetime{]}. Using the same column for
partitioning and change detection will push rows toward the end of the
range and break the logic. Following the analogy, the date should be the
publishing year rather than the year the book arrived in a batch. The
latter will always be in the last batch and does not correctly update
older books.

\subsection{Rolling windows}\label{rolling-windows}

When partitions are in place, rolling windows provide the next level of
efficiency. Power BI allows a data engineer to define the number of
partitions to keep, and older partitions are dropped from the model.
This keeps the model size under control and prevents refresh times from
growing indefinitely.

Boundary points in the window are error-prone. Two details deserve
attention. First, if the SQL source table itself rolls history, its
window must be kept in sync with the Power BI table's rolling window to
avoid mismatches. Second, the boundaries of fact tables must match
related dimensions and other facts. This is particularly true of ID
dimensions. If dimensions are not kept in sync, they will contain
primary key values that no longer exist in the fact tables.

\subsection{Incremental refresh}\label{incremental-refresh}

Incremental refresh adds change detection on top of partitions so only
partitions with changes are refreshed. The native Power BI interface
allows the user to choose a datetime column in the table for change
detection. During refresh, Power BI evaluates the maximum value per
partition. This is the polling query. The maximum value is stored
against each partition as its refresh bookmark. If the next polling
query returns a different value than the bookmark, that partition is
refreshed. There is a substantial benefit to load time if only a
fraction of the polling queries differs from the bookmarks, and thus
only that fraction of partitions are refreshed.

\begin{quote}
\textbf{Note.}

TODO: Insert manuscript screenshot or diagram from the source draft.
Source PDF note:
\texttt{{[}SCREENSHOT\ of\ XMLA\ partition\ and\ polling\ query{]}}.
\end{quote}

Following the analogy, each box of books by publishing year may have a
label showing the last time it was updated. This is the bookmark. When
batches of books arrive, they can be scanned for each publishing year to
see if the update date on any book is greater than the label on the box.
This is the polling query. This is only effective if incoming batches
are organised so that it is easy to find the maximum update datetime for
each publishing year.

The native approach offered by Power BI can be limited if the source
table is a view that joins multiple tables, and thus there is no single
column representing change detection.

It is also not applicable if partitions have deletes because the change
detection column cannot reflect delete datetimes. When the native
approach is not sufficient, the user can specify a custom polling query
with arbitrary M code.

Polling tables can facilitate rapid evaluation of a polling query for
complex cases. A polling table is simply a list of partition dates and
the datetime the partition was last updated. During refresh, the polling
query looks up this update datetime for a partition date.

Following the analogy, the incoming batches of books may have a
companion spreadsheet that tracks publishing years and the last update
datetime of books for that year. During updates, this spreadsheet is
consulted to decide whether that year's box needs updating.

Power BI does not support custom polling queries in its interface.
XMLA-based configuration is required.

If Power BI can rapidly evaluate whether a partition needs refreshing
and retrieve the data for that partition quickly from the source table,
this can lead to substantial improvements in load times. Supporting
these requires the data engineer to prepare the data to facilitate these
queries through indexing, sorting, or even constructing dedicated
polling tables. This preparation itself takes time, and any base tables
or polling tables should themselves be incrementally refreshed.
Consequently, optimising Power BI loads means trading Power BI load
times with data source load times. If not done well, implementing these
artefacts can lead to even longer time than saved through defining
incremental partition refreshes.

\section{Conclusion}\label{conclusion-8}

Loading a Power BI model is time intensive. Using DirectQuery over a
columnstore table avoids the load entirely, but DirectQuery imposes
serious limits on the queries that can be made. Consequently, the most
common factor for a fast Power BI load is the efficiency of the
underlying source tables. These should be simple assemblies of
well-defined fragments that join easily and avoid complex
transformations.

When fact tables become substantial, partitions can deliver a major
speed boost by enabling parallel loads. This often requires rearranging
data or implementing dedicated indexes to support retrieval by partition
date.

When partitions are in place, rolling windows provide further benefit by
dropping older partitions and keeping load times consistent. This
introduces room for error and requires attention to detail at boundary
points.

Incremental refresh adds another layer of efficiency by refreshing only
the partitions that have changed. This relies on a polling query to
detect changes. Dedicated polling table can support instant evaluation
of the polling.

Whatever the technique, the foundation for a fast-loading Power BI model
is set far in advance in the pipeline. It depends on the availability of
meaningful fragments to avoid complex transforms, the planned placement
of indexes to support partition keys, and the careful tracking of row
changes to enable polling tables.

In the ideal case, information efficiency is maximised end to end:

\begin{enumerate}
\def\labelenumi{\arabic{enumi}.}
\item
  Raw tables enter the pipeline incrementally, using polling tables to
  track changes and annotated with change-tracking artefacts.
\item
  Downstream tables transform incrementally with in-sync datetimes and
  are designed for ease of assembly.
\item
  Source tables for Power BI, whether for DirectQuery or for partitioned
  loads, are also incrementally materialised.
\item
  Finally, Power BI tables themselves are incrementally refreshed using
  partitions, supported by incrementally refreshed polling tables for
  fast change detection.
\end{enumerate}

All this means the data engineer plans far ahead, letting Power BI's
efficiency requirements shape the pipeline even as the first table is
built.

\cleardoublepage

\thispagestyle{empty}
\phantomsection
\hypertarget{docs-beyond-techniques}{}
\addcontentsline{toc}{part}{Beyond techniques}
\vspace*{0.34\textheight}
\begin{center}
{\Huge\bfseries Beyond techniques}
\end{center}
\cleardoublepage

\chapter{The five principles of data
engineering}\label{docs-beyond-techniques-the-five-principles-of-data-engineering}

The earlier sections outlined patterns and practices for data
engineering scenarios.

They focused on the treatment of information rather than technology. The
technology-specific sections are on Power BI, and even then, we stayed
away from syntax to focus on fundamentals.

This focus on information rather than technology is deliberate. It
points data engineers to underlying principles rather than situational
techniques. This last section, Beyond techniques, is where patterns
culminate and summarised in five principles.

\section{The problem with shallow
curation}\label{the-problem-with-shallow-curation}

The most common, but incorrect, approach for new data engineers is
shallow curation.

The typical steps are:

\begin{enumerate}
\def\labelenumi{\arabic{enumi}.}
\item
  Take the source system and map to the target environment, adjusting
  for names and data types.
\item
  Receive a list of business requirements for reporting from
  stakeholders.
\item
  Using the source data, apply transformation rules to meet stated
  requirements, load them into a few big tables for Power BI. A more
  experienced data engineer may organise these as facts and dimensions.
\end{enumerate}

This mindset has many drawbacks. Firstly, it leads to garbage-in,
garbage-out. Source applications are not designed for analytical intent
and never has the information content to meet business needs.
Requirements stated upfront are rarely adequate to be a reliable guide
to address this gap. Moreover, this approach focuses on the questions at
hand rather than all the questions about the business. Finally, going
straight from source to reporting tables means there are no building
blocks for use in different scenarios. Using the information for a
different scenario, such as ML feature engineering, requires rebuilding
logic that embedded in dimensional models.

In other words, shallow curation is too easily satisfied. It does not
sufficiently push the data as found to meet the demands of business
intent. Neither does it break down complexity in building towards
advanced levels of achievement. Consequently, the difference is not
between merely ``good versus poor'' quality, but between what is
``possible versus impossible'' for a team to achieve with this mindset.

The problem with shallow curation is that there are some outputs. After
all, the data engineer has faithfully reproduced data from the source
system for users to gain access. Such an engineer may even feel
satisfied for doing an excellent job when reality is far from the
potential value. If business users are dissatisfied when business
objectives are not fully met, it is possible to blame poor data quality
at source or on the unclear business rules. The danger of shallow
curation is that, like all mindsets that starts by reducing standards,
it settles for less, and by settling for less, creates a blind spot for
the engineer who cannot see the failure to reach excellence.

\section{Five principles}\label{five-principles}

A surer basis for data engineering can be summarised in five principles:

\begin{itemize}
\item
  Instead of garbage-in-garbage-out with raw data, add value through
  expressive entities.
\item
  Instead of building giant tables, create meaningful fragments.
\item
  Instead of stopping with what works now, anticipate errors that may
  occur.
\item
  Instead of reacting to requirements, build momentum through guiding
  attention.
\item
  Instead of stopping at the symptoms, diagnose the root cause
\end{itemize}

That is, instead of finding quick fixes to what is wrong, the expert
data engineer focuses on what needs to be right for business
intent---now and into the future. These can be remembered through five
phrases: expressive entities, meaningful fragments, anticipate errors,
build momentum, root cause.

It is not a matter of ``fast workaround'' versus ``slow and proper.'' It
is the difference between what gets tangled up in a mess versus what
will not. In the long term, the five principles to fast and flexible
paths to success.

The earlier sections covered three of the principles--expressiveness
entities, meaningful fragments and anticipating errors. The remaining
section introduces the final principles--build momentum, root cause.

The section ends with the chapter Hallmarks of quality. This essay was
written years before this text and reproduced with minimal edits. The
essay guided the years of work now summarised in this book.

\cleardoublepage

\chapter{Working with
stakeholders}\label{docs-beyond-techniques-working-with-stakeholders}

The most common cause of project failure is the inability of the
delivery team to build momentum with stakeholders. When momentum stalls,
two patterns can often be found.

The first is when the delivery team asks stakeholders for a list of
``reporting requirements,'' but the stakeholders struggle to provide
one. From their perspective, the richness of their business cannot be
reduced to a simple checklist. When the team insists on detailed
requirements, the process reaches a stalemate.

The second pattern appears more promising at first. Stakeholders provide
a clear list of requirements upfront---sales aggregated by date or
region, non-compliance by product type, even precise business
definitions. The delivery team implements these, and all seems well. Yet
during testing, the product does not meet its purpose. It lacks core
features, breaks on edge cases, or proves too complex to use.

Both scenarios share the same flaw of reacting to requirements. Despite
its obvious and repeated failure, this pattern persists. The reason for
this persistence is because experts are often familiar with the
intricacies of their own discipline but underestimate the challenges of
others, expecting simplistic answers for everything else to be given in
clear-cut form so they can proceed with their part.

The world is not so straightforward. The most difficult part of a data
project is not technical complexity but in thinking deeply about the
business and its relationship with the data world. Asking stakeholders
to enumerate requirements is convenient for delivery teams because it
shifts the responsibility for this deep thinking to stakeholders,
leaving developers with the easier task of implementing clear technical
rules.

Instead, in a successful project, the delivery team guides stakeholders
to express their intent, then works together to explore the data in
light of that business intent. Through this process, the stakeholders
uncover what they need to ``see'' from the data for them to achieve
business outcomes. These discovered needs---not an initial list---are
the real requirements. This shift turns a one-way process into a two-way
dialogue where the delivery team plays a leading role by guiding
collective attention.

The fourth principle of data engineer is therefore: instead of reacting
to requirements, build momentum through guiding attention.

\section{Guided attention}\label{guided-attention}

During development, the delivery team and stakeholders explore the data
together to meet business intent. The delivery team contributes
expertise by guiding the stakeholders in how they should ``see'' the
data in light of business intent---hence guided attention. This dialogue
is an iteration between business and data, between questions and
answers, between thinking and testing, and between business and
technical expertise. The seven principles that follow provide a
practical framework for making this a success:

\begin{enumerate}
\def\labelenumi{\arabic{enumi}.}
\item
  Focus on trust
\item
  Lead by listening
\item
  Own the business intent
\item
  Anchor a vision
\item
  Gather around the solution
\item
  Design for workflows
\item
  Spot the 20\%
\end{enumerate}

\subsection{Focus on trust}\label{focus-on-trust}

Many assume that the primary mode of requirements gathering is that of
extracting information. In this view, the question becomes: ``How
effectively and accurately can I get the information from stakeholders
to build a product?'' This is the wrong focus. The first responsibility
of a delivery team is to gain the stakeholder's trust.

The reasons for focus on growing trust are:

\begin{enumerate}
\def\labelenumi{\arabic{enumi}.}
\item
  Trust is the context for information. Clearly articulating and
  refining requirements is arduous. Stakeholders will not be able to
  give the right information before they have first established a strong
  trust in the people asking for the information.
\item
  Relationship with others is the biggest factor for meaningful and
  enjoyable work.
\end{enumerate}

Strong relationships with stakeholders are more likely to lead to
enjoyable collaboration during the development phase, and subsequently
the success of the team.

\begin{enumerate}
\def\labelenumi{\arabic{enumi}.}
\setcounter{enumi}{2}
\item
  The focus on extracting information sees the stakeholders
  instrumentally as a vehicle for information. The focus on growing
  trust sees stakeholders as human beings as an end-to-themselves with
  whom we work.
\item
  Delivery teams have recurring partnerships with stakeholders that
  extend beyond a particular project. The focus on building trust takes
  the long-term view.
\end{enumerate}

Without trust, no data project can succeed. After each engagement, the
question should not be ``Do we know more about the requirements?'' but
``Do our stakeholders trust us more?''

The best ways to grow trust are transparency and active listening.

\begin{itemize}
\item
  Transparency means clear communication of the construction plan and
  regular demonstrations of progress---weekly or twice-weekly. This is
  akin to a home buyer seeing the building take shape.
\item
  Active listening means thoughtful paraphrasing and summarising
  stakeholder input. When stakeholders hear their own language
  reflected, they feel understood and affirmed. This also happens when
  they are invited to check metadata, which gives them a tangible way to
  contribute.
\end{itemize}

A common mistake is for delivery teams to assume that the purpose of
speaking is to transfer information---if they already know what
stakeholders are about to say, then they can skip to the next part. This
is a mistake. Beneath the surface, the stakeholder is:

\begin{enumerate}
\def\labelenumi{\arabic{enumi}.}
\item
  Developing a feeling of being heard.
\item
  Clarifying their own thinking through articulation.
\item
  Growing in data literacy through dialogue with technical experts.
\end{enumerate}

All these build trust. By remembering that the first responsibility of
the team is to gain trust, we resist the temptation to cut people off or
jump in to correct them. Patience is paramount.

\subsection{Lead by listening}\label{lead-by-listening}

In some types of partner dancing, there are two designated roles: Lead
and Follow. The Lead initiates all the movements and the Follow
completes the sequence with an elegant response. An analogous dynamic is
needed with stakeholders. The delivery team should formulate its own
statements as responding gracefully to a previous statement from
stakeholders. These responses typically take one of the following forms:

\begin{itemize}
\item
  Clarifying question: ``You said X, can you please clarify whether X
  means...?'' or ``You said your team needs X---can you help us
  understand how this fits into your goal?''
\item
  Playback: ``There's a lot in what you just said---can I paraphrase to
  check I'm on the right track?'' or ``I've summarised your explanation
  into workflows---can I share them for confirmation?''
\item
  Proposal: ``Based on what you said about your workflows, here are
  wireframes of our solution hypothesis---does this align?''
\item
  Amendment: ``You gave feedback on our hypothesis---here's a
  refinement. What do you think?''
\end{itemize}

By responding to stakeholder initiatives, the delivery team is, in
reality, leading. This is achieved by listening closely to what has been
said and naturally guiding the stakeholder's attention through a
technical problem. This approach ensures:

\begin{itemize}
\item
  Stakeholders are heard, feel heard, and are seen to be heard.
\item
  The conversation remains structured, logical, and focused on business
  objectives.
\end{itemize}

There will be times when the team needs to correct an error, redirect
focus, or counter unconscious bias. This principle still applies. A
common scenario is when stakeholders describe solutions before the
problem is fully explored---for example, asking for a detailed data dump
dashboard when analysis suggests a summarised view is better.

Saying ``Let's focus on requirements instead of jumping to the
solution'' sounds natural but often feels dismissive. For stakeholders,
the solution is the requirement.

Regardless of content, stakeholders are trying to communicate something
important.

This should always be affirmed. Even if the content is incorrect, the
underlying intent is valuable. Respond by continuing with their
initiative rather than breaking off. For example:

\begin{itemize}
\item
  Explore the intent: ``You said your goal is to detect potential
  non-compliance--- how does the data dump help?''
\item
  Explore the utility: ``You want a detailed data dump---can you give an
  example of a workflow where this is used?''
\item
  Re-orient: ``Thanks for raising the data dump idea. You previously
  said your goal is detecting non-compliance. Can we explore scenarios
  where that happens, then revisit the data dump in that context?''
\end{itemize}

These approaches respect intent and apply the first principle: focus on
trust. It also applies the fifth principle: gather around the solution.

This principle should never be used as a shield for blame (``I just did
what you told me to do'').

\subsection{Own the business intent}\label{own-the-business-intent}

The task of data engineering is to align data to business intent. It is
an experimental process that demands the data engineer see the
stakeholder's perspective ``first-hand.'' Without this, the team will
never truly see what the stakeholder sees, and the potential of the
solution will remain unrealised. This act of viewing the data through
the stakeholder's lens is owning the business intent.

The data engineer who owns the business intent must be committed to
in-depth business analysis.

This document is not a guide to business analysis, except to note that
business processes often conform to recurring patterns. Mastering these
patterns helps the team achieve high-quality analysis quickly, even for
complex scenarios. The team can do this by consistently asking
stakeholders five key questions:

\begin{itemize}
\item
  Intent: What is the business intent?
\item
  Measure: How does the business measure the achievement of this intent,
  either directly or indirectly?
\item
  Sensor: What instruments does the business have for knowing whether
  things are going well or not?
\item
  Controls: What levers does the business have to influence this
  measure?
\item
  Drivers: What external events outside of does the business's control
  may influence this measure?
\end{itemize}

Business analysis needs to permeate the project from start to finish.
For example, concepts such as ``good vs bad'' entities or milestones for
measuring processes are rarely defined in source systems. It is
unrealistic to expect stakeholders to define these upfront before seeing
the data. Instead, these definitions evolve through exploration and
require input by technical expertise and experience with similar
problems. This, once again, demands that the team see the business
problem first-hand.

Owning the business intent is an extension of the first principle of
trust. Nothing gains trust with stakeholders as quickly as a team who
can fluently speak the details of a business language. From this
perspective, business analysis is a delivery team's way of active
listening.

Whether a team owns the business intent will define its passion, drive,
business knowledge, and creativity. These qualities will determine
whether its output is a mediocre reflection of the current state or a
solution that pushes the business forward to excellence.

\subsection{Anchor the vision}\label{anchor-the-vision}

The natural extension of owning the business intent is a vision for the
business. A vision is a view of the value the team aspires to achieve.
It is both ambitious and concrete.

Whilst having a vision is important for any project, it is particularly
critical in exploratory data projects for the following reasons:

\begin{itemize}
\item
  Improves the probability of arrival. Complex data projects are often
  ambiguous and uncertain. Teams that articulate and revisit a vision
  are more likely to reach the intended outcome.
\item
  Sustains engagement. Long projects involve hard work, such as
  resolving edge cases. A clear vision inspires a sense of meaning to
  maintain engagement.
\item
  Keeps debates in perspective. Under constraints, debates are
  inevitable and can become flashpoints of tension. A vision keeps these
  debates in perspective.
\end{itemize}

Teams often relax around specific problems when they see those problems
as only one part of the whole.

\begin{itemize}
\tightlist
\item
  Structures discussion. Data projects are open-ended and can easily
  drift.
\end{itemize}

Anchoring the project in a vision, and deriving each step from that
vision, keeps the team on track.

Using the vision to structure discussions begins with agreement on the
overall intent of the business. The project objective is situated within
the business objective, which is itself a part of the broader
organisation's objectives. The discussion on each delivery feature
should be anchored to this hierarchy, with constant reference back to
the higher-level intent. When the conversation becomes lost, the team
moves up one level and reorients.

Example dialogue:

Team \& stakeholder: Introductions Stakeholder: Jumps into details about
needing a dashboard Team: ``Can you start by telling us the purpose of
your business, your role, and what success means to you?''

Stakeholder: Provides some context, still detailed Team: ``So your
organisation's goal is to improve customer experience, and your team
supports this by monitoring service quality?''

Stakeholder: ``Yes, but also\ldots''

Team: ``Okay, so you track service quality, but you also need to
identify recurring issues to prevent them. Is that right?''

Stakeholder: ``Yes\ldots''

Team: ``At the start you mentioned a dashboard. What is its purpose in
this context? Is it for your team's internal use, or to share insights
with other areas?''

This pattern---bringing stakeholders up to the highest level and then
unpacking details step by step---requires fast thinking and familiarity
with the organisation's goals. It is harder than it appears, but
essential for clarity.

When priorities conflict, reference to the higher-level intent provides
common ground.

Agreement is easier at higher levels and can be used to resolve issues.
This approach works only when the conversation has been structured from
the vision downward.

Unlike a building project, which relies on well-defined specifications
of the target building and strict timelines to track progress, data
projects are exploratory and open-ended. In this setting, a clear
vision, rather than detailed specifications, plays the role of guiding
the team. It serves as a north star and compass that sustains quality
and momentum across a long project. For this reason, team leaders should
insist on vision, quality, and momentum rather than rigid deadlines.

\subsection{Gather around the
solution}\label{gather-around-the-solution}

Data projects often waste time through too much talking and not enough
doing.

Instead, the best way to collaborate is to gather around a
solution---start with a pen-and-paper solution hypothesis and refine it
through discussion and experimentation.

The project should develop hypotheses as early as possible, even by the
second engagement. Early development matters because:

\begin{itemize}
\item
  A convincing solution is the only proof the team has understood the
  requirements.
\item
  Hypotheses provide a concrete point for clarifying stakeholder
  thinking.
\item
  Refinement creates excitement as all parties see an early vision come
  to life.
\end{itemize}

A solution hypothesis is the most important part of requirements
gathering because it is the only real proof the team has done the work.
In an analogy of commissioning a building, it is futile if the builder
only has a list of requirements (four rooms, lots of sunlight). The
builder wins the contract by providing an architectural sketch---nothing
less will do. No buyer would pay a mortgage deposit for a building
without a sketch.

Why expect stakeholders to work with a team without showing a solution
hypothesis?

When projects meander, it is often because direction is unclear. The
best advice for such situations: ``When the project is stuck, draw a
picture of the solution.''

The best way to refine the solution hypothesis is through an open
workbench format. In an open workbench, the delivery team and
stakeholders meet twice a week to explore the data model, experiment
with new features, test definitions, and provide feedback for the next
iteration. Stakeholders may come from different business areas, making
the data model a centre for converging perspectives.

This format requires developers to be confident hosting discussions and
responding to impromptu questions from stakeholders. The entire team
should support them in doing so. When done well, the open workbench has
a transformative effect: it builds trust through transparency and
enables genuine dialogue.

Teams sometimes say, ``Let's not jump into solution mode.'' This
statement is not acceptable with stakeholders. For them, the solution is
the requirement. When asked to describe requirements, most stakeholders
are painting a mental picture of the solution. Even if it is a bad
solution, it still describes what they imagine they are working toward.
The team should playback the intent or seek clarification.

The error is to draw a hard line between requirements and solution. This
is a distinction developers make, not stakeholders. The real distinction
is between the why (intent) and the what (solution). Stakeholders will
often mix both. It is the team's job to discern one from the other.

This cannot be stressed enough: to stakeholders, the solution is the
requirement---not something separate from it. In a data project, it is
never too early to gather around the solution.

Gathering around the solution is the ultimate way the delivery team
guides stakeholders toward clarity. It helps them see how reshaped data
interacts with business intent, turning abstract requirements into
something tangible. Through iterative refinement---testing, visualising,
and adjusting---the solution evolves until it displays fidelity with
business reality. In doing so, attention is guided deliberately,
ensuring that every discussion converges on what matters most: a
solution that stakeholders can recognise, trust, and act upon.

\subsection{Design for workflows}\label{design-for-workflows}

It is common for data projects to deliver reports that are rarely used.
Stakeholders may specify many requirements--- ``I want to see count of X
per country''---and show excitement during development and testing. Yet
after deployment, usage often drops quickly.

This happens when solutions are not anchored in real workflows.
Anchoring a solution with workflows ensure that it will have usage when
deployed.

Every workflow has at least two elements: an intent and a trigger. As an
example of a business working at the helpdesk, see Figure 1.

Figure 1. Example of a workflow.

A trigger can be a specific event (e.g., receiving an email from a
customer) or a schedule (e.g., publishing a quarterly report).

Contextualising requirements in workflows offers key advantages:

\begin{enumerate}
\def\labelenumi{\arabic{enumi}.}
\item
  Ensures the product delivers value because it meets the intent of an
  established workflow.
\item
  Ensures the product will be used because real-world triggers exist for
  its use.
\item
  Integrates the product seamlessly into the user's daily role.
\item
  Helps stakeholders recognise when a requirement is not important and
  can be dropped.
\item
  Provides a natural way to rationalise requirements---for example, one
  report page per workflow.
\end{enumerate}

Designing for workflows is a simple tool to ensure that project teams
are grounded in real world problems.

\subsection{Spot the 20\%}\label{spot-the-20}

The 80/20 rule says 80\% of results come from 20\% of effort. Its flip
side is that the remaining 20\% consumes 80\% of the effort. This
creates an illusion that most work is done when the hardest part
remains---the hidden, complex 20\%.

The 20/80 trap is dangerous because:

\begin{enumerate}
\def\labelenumi{\arabic{enumi}.}
\item
  Foundational elements are missed until production.
\item
  Expectations are mismanaged when stakeholders think the work is
  complete.
\item
  Technical debt piles up when essentials are discovered late.
\end{enumerate}

Teams avoid the trap by:

\begin{enumerate}
\def\labelenumi{\arabic{enumi}.}
\item
  Grounding requirements in thorough business analysis to uncover hidden
  assumptions.
\item
  Focusing equally on what works and what doesn't---the parts that work
  are often forgotten.
\item
  Asking not only ``What should this do?'' but ``What could go wrong?''
\item
  Balancing common cases and edge cases without bias.
\item
  Developing in sound engineering order, not simply in the order
  stakeholders want to see.
\end{enumerate}

As an analogy, home buyers focus on visible features---``How many
rooms?'' or ``Is there good ventilation?''---based on obvious needs and
past pain. A good builder starts with the foundation. Likewise, a good
data team begins with the start of the business process and
systematically work through the business processes in order, rather than
jumping to flashy features.

Because of this, the team will often run counter to stakeholder
instincts. This is healthy.

In effect, the team acts as:

\begin{itemize}
\item
  A counterweight to natural biases that overlook fundamentals.
\item
  A safeguard against missed edge cases that matter.
\end{itemize}

Foreseeing hidden elements is hard. It takes strong business knowledge,
technical experience, and high engagement skills to counter bias. Only
the best teams do it well.

\section{Conclusion}\label{conclusion-9}

The hardest part of any data project is thinking deeply about the
business and how it relates to data. This can only be achieved when
stakeholders and the delivery team explore the data together---thinking
together through a dialogue focused on the solution. In this dialogue,
the relationship is not symmetrical. Stakeholders, often new to data, do
not know what to look for. This is why the experienced delivery team
plays an important role in guiding the dialogue's attention to relevant
observations.

The seven principles of engagement capture hard-earned wisdom for
facilitating these dialogues. They run counter to ingrained habits that
focus on visible details while ignoring complexity. They are treasures
for any delivery team.

The principles change the team's view of engagement from receiving
requirements and getting product sign-off to that of building trust with
stakeholders. Trust is the guiding principle. Fluently speaking the
business language establishes trust early in the project, gathering the
solution and continuing to demonstrate progress on the business vision
carries trust to the end.

These principles move the team beyond reacting to requirements. Applied
well, they enable the team to guide stakeholders through complexity and
build momentum toward a solution that truly meets business intent.

\cleardoublepage

\chapter{Construction
planning}\label{docs-beyond-techniques-construction-planning}

The fourth principle of data engineering is to build momentum through
guiding attention rather than reacting to requirements. This can be
challenging in projects that are lengthy and overwhelmingly complex. In
such cases, the engineer often feels at a loss for where to begin.
Because the data engineer sits on the critical path of data projects,
this uncertainty directly impacts project momentum. Conversely, a data
engineer who delivers outputs in an orderly and predictable way propels
the project forward--- participants align on the goal and instinctively
sense that the team is taking concrete steps toward achieving it.

In complex projects, a data engineer stays organised and guides the
delivery team by formulating a plan. This plan is expected from an
experienced engineer because:

\begin{enumerate}
\def\labelenumi{\arabic{enumi}.}
\item
  No project waits in silence. The engineer cannot disappear for a month
  to apply ideal patterns and return later. Unless the engineer provides
  the team and sponsors with confidence about the weeks or months ahead,
  there will be no engineering work at all.
\item
  Logical sequencing matters in a build. Doing tasks in a sensible order
  helps the engineer move forward with confidence, avoids code
  entanglement, and reduces mistakes.
\item
  Only the engineer knows the effort required. The data engineer alone
  understands how long quality work takes and has the sole
  responsibility to advocate for that time.
\end{enumerate}

If the engineer does not actively formulate a plan, there will be no
plan---leading the project to meander. Or someone else will create the
plan for the engineer---leading the project to rush. In either case, the
project suffers.

Projects can succeed or fail depending on the data engineer's plan. It
is one of the data engineer's leadership roles when guiding
stakeholders. Formulating a plan for the team is one of the most
advanced skills to master and a mark of a confident engineer.

\section{An effective plan}\label{an-effective-plan}

A construction plan lays out, in sequence, the components the data
engineer will build and provides a forecast for how long each part will
take. This plan gives users a clear expectation of what features will
arrive and when. In addition, a construction plan includes non-feature
components such as unit tests, technical debt clean-up, and performance
tuning. The plan defines the project's momentum.

A plan is effective if it provides the project with confidence and
momentum. For this purpose, the focus is on sequence rather than precise
timings. While it is important not to blow out project timelines, the
sequence is more critical than exact estimates because delivering one
feature after another is strongly associated with user satisfaction and
project momentum.

An effective plan is akin to a travel itinerary for exploring a new
country. An itinerary is not effective merely by meeting a predefined
schedule; it is effective when it moves forward logically, providing a
quality experience without wasting time.

A data engineer can craft an effective plan by aiming for three
characteristics:

\begin{itemize}
\item
  Ambitious yet grounded
\item
  Orderly yet flexible
\item
  Rotating new features with backend consolidation
\end{itemize}

Being ambitious means having a view to tackle the full business
processes and designing a model that addresses all questions, rather
than limiting the scope to specific requirements raised by current
stakeholders. Being grounded means supporting this vision with
well-informed business analysis so that the plan is neither vague nor
impossible, and ensuring that necessary information is available for
reliable time estimates.

An orderly yet flexible plan is akin to a travel itinerary that is
physically sensible (not looping back across half the country) but still
organised into reasonable segments that can be shifted. For a data
engineer, being orderly means isolating releases into promotes code
clarity, avoid code entanglement, and minimises rework. Flexibility
means structuring the plan so that key business priorities can shift
forward or backward as a block without massively disrupting the
schedule. This requires grouping deliverables into logical blocks
centred on core business processes, rather than maintaining a long
miscellaneous list of business questions.

New features excite stakeholders and propel project momentum. However,
the data engineer must balance these visible features with invisible but
equally important consolidation work. An effective plan achieves both by
alternating between them. This rotation not only supports quality work
by building consolidation into the schedule but also gives stakeholders
time to absorb one feature before moving on to the next. This immersion
is necessary because data insights are exploratory in nature. Testing a
new feature requires users to interact with the output, grasp its
implications, and check edge cases thoroughly.

\section{Formulating an effective
plan}\label{formulating-an-effective-plan}

A data engineer can craft an effective plan in four stages. Each stage
involves increasing commitment to specific outcomes. A simple project
may require only the first stage, while a complex project may demand a
detailed plan at the fourth stage. The four stages are:

\begin{enumerate}
\def\labelenumi{\arabic{enumi}.}
\item
  Discovery
\item
  Vision
\item
  Scope
\item
  Build
\end{enumerate}

\subsection{Stage 1 -- Discovery}\label{stage-1-discovery}

Formulating a plan begins with discovering facts about the business. The
purpose of discovery is to answer key questions. At the highest level,
these include:

\begin{itemize}
\item
  What are the major business processes in the end-to-end business
  lifecycle?
\item
  What information is captured by each process, and where is it stored?
\item
  What is the stakeholders' interest in each process?
\item
  What reporting artefacts already exist, and how are they currently
  used in workflows?
\end{itemize}

The team answers these questions by reviewing artefacts relevant to the
business processes and interviewing stakeholders.

To sharpen the focus on business objectives, the team should also
answer:

\begin{itemize}
\item
  Intent: What is the business intent?
\item
  Measure: How does the business measure the achievement of this intent,
  either directly or indirectly?
\item
  Sensor: What instruments does the business have for knowing whether
  things are going well or not?
\item
  Controls: What levers does the business have to influence this
  measure?
\item
  Drivers: What external events outside of does the business's control
  may influence this measure?
\end{itemize}

The outputs of discovery should be expressed as linear process diagrams
and cumulative information diagrams, as explained in the chapter
Anticipating Questions.

Recall that a linear process diagram is a simple diagram that lays out
the major business processes in chronological order, without loops or
branches. Its purpose is to give a clear, high-level view of the
end-to-end flow, making it easy to see where each process fits in the
overall lifecycle. A cumulative information diagram is a table that
lists business information as rows and business processes as columns,
showing which process captures or inherits each piece of information.
This diagram helps identify what data is available at each stage and
ensures completeness for answering business questions.

\subsection{Stage 2 -- Vision}\label{stage-2-vision}

The vision for all data engineering projects is the same: advance
business intent by obtaining insights across all underlying business
processes. This vision is adapted to each project using findings from
the discovery phase, especially the linear process diagrams.

For example, if the core business processes for an organisation are
manufacturing, quality control, sales, shipping, and customer feedback,
the vision statement could be:

``Understand the factors that drive sales profit and customer
satisfaction through an integrated data source of complete business
processes---from manufacturing, quality control, sales, shipping, to
customer feedback. Obtaining timely and accurate insights into measures
such as manufacturing turnaround, early detection of quality issues,
changes in sales trend, efficiencies in shipping and latest customer
sentiment.''

This vision statement is ambitious yet grounded. It resonates with
business stakeholders, inspires confidence, and establishes trust by
demonstrating that the project team understands the business. Being
neither too vague nor too specific, the vision frames the project at the
right level for prioritisation and dialogue.

Such vision statements are easy to craft and extend naturally from the
discovery questions. Yet many teams produce poor alternatives that lead
projects to failure. For example, it is common to see vision statements
such as:

``Reduce pain points of manual processes and create near real-time
dashboard of the operation.''

These statements are not a deep analysis of business intent but a
reaction to user complaints. As such, they lack the substance to serve
as a foundation for the project.

As part of vision setting, it is important to sketch pen-and-paper
wireframes of sample reports that user can build off the data. The
wireframe is a play-back what the team has heard, and the beginning of a
solution hypothesis. Its purpose is to build trust that the team has
understood the needs and inspire the project to strive for the finish
line. It is akin to an artistic model of a public building prior to
build---a good model generates excitement and builds public confidence.
To achieve this purpose, the wireframe itself needs to be ambitious yet
grounded.

When the vision is framed as above, defining the project scope becomes
straightforward.

\subsection{Stage 3 -- Scope}\label{stage-3-scope}

The vision is an ambition to cover the business processes end to end. It
is a visionary statement of the linear process diagram.

The full vision may be too large to tackle in one project. The scope
defines a subset of business processes to focus on. For example, rather
than covering manufacturing through to customer feedback, the project
scope may concentrate on the first two--- manufacturing and quality
control---or on priorities such as sales and shipping. Dividing the
vision this way allows the creation of project phases: the first phase
focuses on two business processes, the second on the next two, and so
on.

The selection of business processes should consider three factors:

\begin{enumerate}
\def\labelenumi{\arabic{enumi}.}
\item
  Availability of source data. If the information captured by a business
  process resides in systems that are difficult to access, it may need
  to be deferred to later phases.
\item
  Order of the business lifecycle. Where possible, work in the natural
  sequence---manufacturing, then quality control, then sales, and so
  forth.
\end{enumerate}

Information accumulates along the lifecycle, and working in the same
order creates building blocks for the next stage.

\begin{enumerate}
\def\labelenumi{\arabic{enumi}.}
\setcounter{enumi}{2}
\tightlist
\item
  Business priorities. These often pull toward the tail end of the
  lifecycle, where most activity occurs. The data engineer must judge
  based on dependencies. For example, starting with sales without
  manufacturing may be acceptable, but starting with shipping without
  sales orders is not---no matter how much users want insights on
  shipping delays.
\end{enumerate}

The cumulative information diagram, which tracks the flow of information
and its storage, is useful for informing scope decisions.

The scope must include a time estimate for project completion. This
requires input from an experienced engineer who can make a judgment
based on discovery findings.

Like the vision statement, the scope is a straightforward application of
the linear process diagram.

A common error is to scope by system. For example, rather than working
through the business process one by one, the team works through
integrating systems one by one.

This is the wrong perspective. One reason, amongst many, is that it
easily leads to a situation where the project creates an unusable
product because the system does not have the minimal set of information
to describe one business process.

\subsection{Stage 4 -- Build}\label{stage-4-build}

The final stage is the construction plan. This is the practical roadmap
for delivery that translates scope into actionable steps and sets out
how the project will progress.

A construction plan is akin to a six-month travel itinerary in a large
country. The country is the organisation, and the towns are the major
business processes. Features are the specific sites within a town.
Building blocks are times of being en route, such as taking a long bus,
while technical consolidation is the time of rest in the journey.

This travel itinerary is exploratory in nature, so it needs flexibility,
but it also needs order so that the physical route makes sense and there
is a continuous flow of experiences alternating with rest. This cannot
be achieved if the itinerary is too vague: ``Let's see the country.''
Nor can it be too detailed: ``Here's a long, ad hoc list of things we
want to do all across the country in the next six months.'' Instead, an
orderly yet flexible itinerary sketches out the major towns and
important sites, while the details of visiting a town only need to be
fleshed out when getting close.

If the itinerary is well organised, it can be changed mid-way, such as
swapping the order of towns or sites, with minimal disruption to the
overall journey.

An engineer guiding stakeholders through a data project is akin to a
tour guide crafting a travel itinerary for a traveller experiencing a
country. The analogous concepts apply.

The two primary considerations are: how the delivery is divided into
releases, and the sequencing of those releases.

At the highest level, build should proceed in the sequence of business
processes. For example, if the scope includes manufacturing, quality
control, sales, and shipping, then the build releases should follow this
order, process by process. This contrasts with tackling a mix of
manufacturing and sales, or sales and shipping, in one release. If the
data engineer finds this cannot be done, then the abstraction of
business processes may need to be revisited. For example, if the data
for sales and shipping are so entangled that they cannot be separated
for build, then it is arguable that sales and shipping in the
organisation are better seen as one process.

If possible, proceed in the same order as the business process. For
example, avoid jumping to shipping and then doubling back to sales.
However, if reverse order is necessary due to business priorities, then
organising the build along the boundary of business processes can still
avoid creating a mess.

The releases should be further broken down within a business process.
For example, a complex project may expect to build 50 tables in total;
this may be divided into releases of 3 to 8 tables at a time. The more
challenging the project, the smaller the release size.

Each release should be centred on a group of business information to
expose. For example, if the scope is manufacturing, one release may
focus on basic attributes such as product types, production status, and
production date; another release on manufacturing details such as
material inputs and outputs; and another release on aggregated
information such as time to manufacture. In grouping releases,
priorities should be given to being as logical as possible from a
business point of view.

Once defined, the construction plan must build in order of the pipeline
by following table dependencies. This means a gradual increase in
computation complexity. For example, the first release focuses on basic
attributes with little transformation, the second release adds more
complex computation, and a third release aggregates information from
prior tables. This is the natural progression for a data engineer
following the filter -\textgreater{} compute -\textgreater{} reduce
steps explained in Creating information. This gradual layering of
computation keeps complexity manageable and ensures each part is tested
systematically along the way.

In a complex project with many attributes, it is bad practice to work on
all the underlying tables and then all the Power BI tables and measures
as one release.

Instead, Power BI releases should be interleaved as information is
curated in the pipeline. This means that features are surfaced to users
continuously, so they get exposure to business attributes as soon as
they are added to the pipeline.

Releases should be tightly coupled with unit tests, fault tolerance, and
metadata.

Whenever a feature is released, unit tests and metadata should come with
the release, rather than as a bulk afterthought.

Releases should be kept short in time---approximately two weeks, or
three weeks at most. They need to be small enough for thorough code
review.

In a construction plan, not all releases will be user features. There is
also important backend work such as building blocks or performance
tuning. It is important to earmark the feature releases because they are
the key driver of project momentum. Users should know when features are
coming and receive them continuously in small chunks.

User features are typically exposed to users in the self-service Power
BI model.

User features should rotate with backend work. Not only does this keep
up project momentum, it allows data engineers to spend time on quality
consolidation and gives users space to fully immerse in testing features
between releases.

By anchoring the construction plan on the user features, and keeping in
mind how releases build up into a user feature, it becomes easier to
reshuffle development as priorities change without ending up in a
tangled mess.

Finally, it is not necessary to forecast the plan in full detail for the
entire project. This rigidity is not fit for the exploratory nature of
data analytics. At any given moment, only the next four or so releases
need to be planned in detail, while features further out do not need
detailed planning. For example, when working on manufacturing, the data
engineer should know the plan for the next few releases to complete the
work on manufacturing, while the further processes such as sales and
shipping only need to be noted.

A construction plan designed in this manner is, like a good travel
itinerary, both orderly and flexible.

The following is an example of a construction plan for manufacture and
quality control.

{\def\LTcaptype{none} % do not increment counter
\begin{longtable}[]{@{}
  >{\raggedright\arraybackslash}p{(\linewidth - 6\tabcolsep) * \real{0.2500}}
  >{\raggedright\arraybackslash}p{(\linewidth - 6\tabcolsep) * \real{0.2500}}
  >{\raggedright\arraybackslash}p{(\linewidth - 6\tabcolsep) * \real{0.2500}}
  >{\raggedright\arraybackslash}p{(\linewidth - 6\tabcolsep) * \real{0.2500}}@{}}
\toprule\noalign{}
\begin{minipage}[b]{\linewidth}\raggedright
Release
\end{minipage} & \begin{minipage}[b]{\linewidth}\raggedright
Information focus
\end{minipage} & \begin{minipage}[b]{\linewidth}\raggedright
Description
\end{minipage} & \begin{minipage}[b]{\linewidth}\raggedright
Artefacts
\end{minipage} \\
\midrule\noalign{}
\endhead
\bottomrule\noalign{}
\endlastfoot
1 & Manufacture core tables & Basic tables from source system to
describe the manufacture process, with metadata and minor clean-up &
Cake.Manufacture; Cake.RefManufactureSite; Cake.RefManufactureStatus;
Cake.ManufactureStep; Cake.RefProductType \\
2 & Manufacture SLA tables & Define milestone points in the manufacture
process, and compute time between milestones and total
time-to-manufacture by aggregating over Cake.ManufactureStep, compare
the results against SLA for different product types &
Cake.ManufactureBatchMilestone; Cake.RefManufactureProductSla \\
3 & Quality control core tables & Basic tables from source system to
describe the quality control, with metadata and minor clean-up &
Cake.SampleTesting; Cake.RefQualityControlCriteria;
Cake.RefQualityControlResult \\
4 & Quality control summary & Aggregate results from different quality
control criteria testing, to compute an overall result for a batch, with
the concepts of failure severity & Cake.ManufactureBatchOutcome;
Cake.RefBatchOutcome \\
5 & Power BI model for exposure & Create the views that define the fact
and dimension tables for a Power BI report, and expose them in a Power
BI model with corresponding measures, including metadata & Driver views
for facts \& dimensions: PBI\_VIEW.Manufacture;
PBI\_VIEW.QualityControl; PBI\_VIEW.ReportingCalendar;
PBI\_VIEW.BatchOutcome; PBI\_VIEW.ProductType etc. Corresponding Power
BI tables: Manufacture, Quality control, Reporting calendar, Product
type etc. Measures. \\
\end{longtable}
}

\section{Stakeholder requirements}\label{stakeholder-requirements}

A typical IT project is driven by stakeholder requirements. But
stakeholder requirements cannot determine the whole plan, any more than
a town planner can plan a town-build simply by collecting a list of
requirements from citizens. On the contrary, the town planner is
expected to contribute expertise in dialogue with community.

In practice, stakeholder requirements swing between two extremes---too
vague to guide development, or too specific so that the forest is lost
for the trees. Consequently, the delivery team should not react to the
requirements and build the first thing that is asked. Such an approach
is destined to lead to a tangled build if the requirements are too
specific, or meandering if they are too vague.

Instead, the data engineer takes stakeholder input as part of a broader
consideration for a data model that anticipates questions across
end-to-end business processes.

This can be done by reframing stakeholder requirements as objectives
about core business processes. If the requirements are too specific,
such as a long list of details like ``I want to see count of sales by
product type,'' they can be brought up a level by categorising them as
objectives on sales. If the requirements are too vague, such as ``I just
want to know about my business,'' they can be made more concrete by
articulating objectives on key processes such as sales, manufacturing,
or shipping.

In short, stakeholder requirements should be reframed to the level
explained in the chapter Anticipating questions. The delivery team has
responsibility in guiding the stakeholders comfortably on this journey
as described in Working with stakeholders.

\section{Conclusion}\label{conclusion-10}

In traditional IT projects, data engineers are often seen as ``doers''
who implement specified rules. This is far from their full potential.
Being closest to the data, they actively shape how the organisation
understands its business by reorganising information at a fundamental
level. Data engineers are true influencers.

This chapter continues the theme of a data engineer's influence. In a
high-intensity project, the data engineer---working with other team
members---formulates a plan that not only addresses business
requirements but reshapes them, using this plan to build project
momentum. The data engineer is doing the build, and only the data
engineer can craft the plan.

The way to create an effective plan is to use the tools explained in
Anticipating Questions: linear process diagrams and cumulative
information diagrams. These tools break down organisational complexity
into a level that is useful for all parties and especially conducive to
the data engineer shaping data under these frames. In this sense, the
linear process diagram becomes a common ground between business
stakeholders and the delivery team for planning and prioritisation.

An effective plan lays the groundwork for success. In a complex project
that requires a defined scope, a rough estimate is that:

\begin{itemize}
\item
  50\% of success is determined by the scope
\item
  35\% by the construction planning
\item
  15\% by the actual development work
\end{itemize}

This is no surprise. As developers mature, development becomes more
mechanical and predictable. Familiarity with patterns makes
implementation routine. The more experience builds, the more developers
can foresee outcomes and issues in advance.

Consequently, success hinges increasingly on planning. This is why the
mark of an experienced data engineer is an effective plan.

\cleardoublepage

\chapter{Sound judgement}\label{docs-beyond-techniques-sound-judgement}

Sound judgement is the defining capability of a mature data engineer.

It is tempting to think of data engineering as a sequence of technical
tasks: building pipelines, fixing errors, optimising performance, and
integrating systems. In practice, these tasks are secondary. What
distinguishes effective engineers is not how quickly they act, but where
they act.

Sound judgement is the ability to recognise at what level a problem
should be addressed --- whether a situation calls for surface
correction, structural repair, or a reconsideration of what ``good''
even means. It is exercised continuously, not only when something is
visibly broken.

For this reason, troubleshooting is not a special activity reserved for
failures. It is the most explicit expression of a posture that pervades
the entire discipline. The same judgement used to diagnose a broken
pipeline is used when deciding how to model an entity, where to
aggregate information, or whether a system should be conformed at all.

This chapter presents a simple diagnostic framework that makes sound
judgement explicit. This framework defines the final principle of data
engineering---instead of treating symptoms, diagnose root cause.

\section{A diagnostic framework for sound
judgement}\label{a-diagnostic-framework-for-sound-judgement}

Sound judgement can be practised deliberately by distinguishing four
stages that are often confused:

{\def\LTcaptype{none} % do not increment counter
\begin{longtable}[]{@{}
  >{\raggedright\arraybackslash}p{(\linewidth - 2\tabcolsep) * \real{0.5000}}
  >{\raggedright\arraybackslash}p{(\linewidth - 2\tabcolsep) * \real{0.5000}}@{}}
\toprule\noalign{}
\begin{minipage}[b]{\linewidth}\raggedright
Stage
\end{minipage} & \begin{minipage}[b]{\linewidth}\raggedright
Diagnostic question
\end{minipage} \\
\midrule\noalign{}
\endhead
\bottomrule\noalign{}
\endlastfoot
Symptom & What is observably wrong? \\
Trigger & What event caused the symptom to surface? \\
Root cause & What structural violation allows this to recur? \\
Final effect & What outcome must hold for the system to be genuinely
fixed? \\
\end{longtable}
}

This framework does not prescribe solutions. Its purpose is to
discipline attention. It prevents engineers from reacting to what is
loud instead of addressing what is consequential.

\subsection{A simple example}\label{a-simple-example}

Consider someone learning piano.

\begin{itemize}
\item
  Symptom: the playing sounds offbeat.
\item
  Trigger: the finger is lifted too late before the next note.
\item
  Root cause: the learner relies on playing by ear rather than learning
  the mechanics of timing and articulation.
\item
  Final effect: the piece is played according to the intended rhythm,
  not merely
\end{itemize}

``close enough'' to how it should sound.

Two observations matter.

First, the trigger is not the root cause. The trigger explains how the
symptom occurs, not why it persists.

Second, the final effect is not the absence of mistakes. It is the
presence of the intended outcome.

Sound judgement lies in acting at the level that prevents recurrence,
not merely silences error.

\subsection{Step 1: Symptoms --- seeing clearly before
acting}\label{step-1-symptoms-seeing-clearly-before-acting}

The first act of sound judgement is to systematically canvas the
symptoms, rather than reacting to the first issue reported.

Stakeholders describe what they see. They do not describe the full
boundary of failure.

Beginning work from a single reported symptom allows the reporting
channel to define the problem --- and that definition is almost always
incomplete.

At this stage, interpretation is a liability. The task is mechanical
observation.

In medicine, uncertainty is met with vital checks. In data systems, the
equivalent is to examine surface health before forming explanations.

Typical symptom canvassing includes:

\begin{itemize}
\item
  If one report is broken, are others broken?
\item
  If one metric is wrong, are related metrics wrong?
\item
  If a domain disappears, did the underlying rows disappear or only
  their classification?
\item
  Did row counts change? Did distributions shift? Did null rates spike?
\item
  Did upstream pipelines succeed? Did downstream consumers fail?
\end{itemize}

Jumping ahead at this stage is a failure of judgement. Engineers tend to
do so for two reasons:

\begin{itemize}
\item
  pattern-matching based on past experience
\item
  taking the reported symptom as a complete description
\end{itemize}

Both lead to work that is precise but misdirected.

Sound judgement at this stage consists in restraint. Seeing clearly
comes before acting decisively.

\subsection{Step 2: Triggers --- locating the point of
exposure}\label{step-2-triggers-locating-the-point-of-exposure}

The second act of sound judgement is to identify the trigger --- the
event that caused the symptom to surface.

Triggers are time-bound and concrete. They narrow the search space and
provide orientation.

Common triggers in data engineering include:

\begin{itemize}
\item
  schema changes in source systems
\item
  new or invalid records entering a load
\item
  code changes or refactors
\item
  infrastructure or platform migrations
\item
  permission or security changes
\item
  changes in load ordering or scheduling
\end{itemize}

Triggers are useful because they are often discoverable, and because
they sometimes allow temporary containment: rollbacks, bypasses, or
delayed execution.

However, sound judgement requires a strict distinction:

A trigger explains when a symptom appeared.

A root cause explains why it can reappear.

Treating the trigger as the cause leads to fragile systems that fail
again under different conditions.

\subsection{Step 3: Root causes --- acting at the level that
matters}\label{step-3-root-causes-acting-at-the-level-that-matters}

The third act of sound judgement is to diagnose the root cause.

Root causes explain recurrence. They are structural violations that
allow different triggers to produce the same class of symptom.

Root cause diagnosis may be:

\begin{itemize}
\item
  Theoretical: reading the system and identifying what must be true for
  the failure to occur
\item
  Experimental: isolating variables through controlled tests
\end{itemize}

Both require skill. Theoretical diagnosis depends on structural
understanding.

Experimental diagnosis depends on technical competence and disciplined
inquiry.

A useful rule is:

Root causes violate core principles that define the ideal system.

Triggers merely activate those violations.

In data engineering, recurring root causes include:

\begin{itemize}
\item
  incorrect or ambiguous grain
\item
  missing, unstable, or misused keys
\item
  semantic mismatches hidden behind identical column names
\item
  wide tables that collapse unrelated concepts
\item
  implicit dependencies on ordering, timing, or existence
\item
  weak business processes that produce unreliable identity
\item
  compounded logic that mixes extraction, transformation, and
  aggregation
\end{itemize}

Two engineers may observe the same symptom and trigger. One applies a
patch. The other reshapes the system. The difference is not effort, but
judgement.

Sound judgement is the willingness to act where the system is wrong, not
merely where it is noisy.

\subsection{Step 4: Final effect --- knowing when you are
done}\label{step-4-final-effect-knowing-when-you-are-done}

Between diagnosing root cause and completing work lies the act of
applying a fix. Fixes vary: containment, refactoring, redesign, or
process change.

Regardless of the fix, the final act of sound judgement is to check the
final effect.

This includes:

\begin{itemize}
\item
  visual validation: not ``does the code run'', but ``are the outputs
  correct''
\item
  environment validation: behaviour across development, pre-production,
  and production
\item
  end-to-end validation: whether the information now supports the
  intended decision
\end{itemize}

A critical rule applies:

The final effect is not the absence of the symptom.

Error suppression creates silence, not correctness. Silence is often
mistaken for success, until the failure reappears elsewhere.

Knowing the final effect is difficult because it requires seeing what
``good'' should be.

Engineers often fail to recognise slow systems because they do not know
what fast systems look like. They fail to recognise weak models because
they do not know what expressive models look like.

In data engineering, the final effect is always the same:

The business receives the information it needs to make the decision.

Pipeline success, test success, and error absence are necessary --- but
they are not the end.

\section{Sound judgement throughout data
engineering}\label{sound-judgement-throughout-data-engineering}

Sound judgement is not exercised only when something breaks. It is
present throughout the discipline of a data engineer:

\begin{itemize}
\item
  Expressiveness: o Symptom-level response: rename columns or add
  documentation when reports confuse users. o Root-cause response:
  reshape the model so business meaning is explicit in its structure.
\item
  Fragment modelling: o Symptom-level response: patch special cases into
  existing tables as complexity grows. o Root-cause response: extract
  independent fragments and isolate responsibilities.
\item
  Reference data and conformance: o Symptom-level response: reconcile
  mismatched numbers in reports. o Root-cause response: introduce shared
  reference data and apply it consistently.
\item
  Aggregation and grain: o Symptom-level response: simplify visuals or
  add filters when reports are slow or inconsistent. o Root-cause
  response: aggregate correctly in the pipeline at a grain aligned with
  decisions.
\end{itemize}

In each case, nothing is ``broken'' in an obvious sense. Yet judgement
is still required.

Mature engineers recognise these situations early and act at the level
that prevents future failure.

\section{Conclusion}\label{conclusion-11}

Sound judgement is the ability to act at the right level of abstraction,
at the right time, for the right reason.

The diagnostic framework introduced here makes that judgement explicit:

\begin{enumerate}
\def\labelenumi{\arabic{enumi}.}
\item
  Canvas the symptoms
\item
  Identify the trigger
\item
  Diagnose the root cause
\item
  Check the final effect
\end{enumerate}

Common failures arise from:

\begin{itemize}
\item
  acting on symptoms without understanding scope
\item
  mistaking triggers for causes
\item
  equating silence with success
\end{itemize}

The goal is not to avoid problems. It is to address them at the level
where they stop recurring. Thus, instead of treating the symptom,
diagnose the root cause.

Sound judgement is not a talent. It is a discipline --- and it is the
thread that runs through all effective data engineering work.

\cleardoublepage

\chapter{Getting started as a data
engineer}\label{docs-beyond-techniques-getting-started-as-a-data-engineer}

The full set of concepts and patterns can feel overwhelming for a new
data engineer. It is not necessary to master everything upfront. In
fact, a good data engineer is defined by much more than technical
proficiency.

Starting out as a data engineer is less about mastering every technique
and more about developing the habits and perspective that sustain
growth. The priority is to internalise the engineering principles until
they become second nature.

\section{Developing good habits}\label{developing-good-habits}

New engineers often focus on reaching great technical feats. This is the
wrong approach. Instead, all new engineers should aim for two
goals---exposure and technical discipline.

The greatest and most common danger for engineers is becoming
narrow---focusing purely on engineering techniques rather than the
business or wider architectural perspectives. Early exposure to diverse
aspects of a data project and team is a healthy antidote.

The second most common danger for engineers is focusing on flashy
techniques. On the contrary, the best engineer is built on a foundation
of consistent discipline in technical hygiene. That discipline is formed
early through small problems.

Exposure and technical discipline can be developed through the following
habits:

\subsection{Aim to be helpful, especially on small
tasks}\label{aim-to-be-helpful-especially-on-small-tasks}

A wide range of small tasks gives exposure to the scenarios an engineer
deals with. The best tasks involve addressing unit-test failures or
responding to low-urgency production errors. These build the mindset of
anticipating errors through exposure to fixing them.

\subsection{Talk to the team}\label{talk-to-the-team}

Constant and open conversation with team members contributes to
exposure. Engineers on the frontline are often the first to see an
issue. Consequently, the ability to raise risks and articulate
uncertainties develops through working with a team. Learning from
experienced team members is a core part of gaining exposure, though not
all team members are good models. It is important to be aware of their
track record. In any case, no engineer should work silently in
isolation. A new engineer is able to develop good habits of
communication before the habit of solo-working takes hold.

\subsection{Talk to the business}\label{talk-to-the-business}

A new engineer benefits from learning the business problem
``first-hand,'' as if part of the business team. This is owning the
business intent. Instead of focusing solely on code, early habits should
include writing quality metadata and descriptions. Confidence in live
demonstrations and hands-on experimentation with the business grows over
time and requires comfort in acknowledging ignorance or mistakes.

\subsection{Understand the
architecture}\label{understand-the-architecture}

Developers often focus only on the development environment rather than
the complete architecture. Yet multiple environments, CI/CD practices,
and orchestration of loads all influence the engineer's work. Thinking a
few steps ahead---not just about development but also about deployment
and monitoring--- becomes an important habit.

\subsection{Technical hygiene}\label{technical-hygiene}

Technical hygiene begins with code management. This starts with
elementary habits such as using Git version control as second nature.
The most common error for new engineers is creating massive branches
that release many features at once. Smaller, logically isolated releases
build the capability to break down complex problems into simplified
steps.

\subsection{Check, don't assume}\label{check-dont-assume}

New engineers often make assumptions---about data, people, business,
communication, or architecture. A habit of checking every step carefully
is a prerequisite for tackling complex problems.

\subsection{Estimate and refine}\label{estimate-and-refine}

Engineers often shy away from estimating effort and time-to-complete,
developing an unhealthy habit of making excuses for why estimation is
too difficult. Yet the ability to estimate accurately is one of the best
indicators of technical competency. This will not come on day one.
Developers commonly underestimate before work starts and overestimate
once underway. A useful habit is making a rough guess of
time-to-complete before starting any task and checking the estimate
afterward.

\subsection{Focus on the principles}\label{focus-on-the-principles}

Data engineers must develop mastery of techniques necessary for quality
engineering. Most new engineers learn better by doing rather than
reading a book. However, understanding techniques through a theoretical
framework from first principles remains important. Instead of being
satisfied with ticking off a delivery feature, a new engineer benefits
from reflecting on how the task links to the patterns in this book,
especially the principles of data engineering.

\section{Conclusion}\label{conclusion-12}

A new engineer benefits from balancing the mastery of techniques with
confidence in the span of responsibility. A data engineer is more than a
developer applying techniques to build a data model. The role involves
engaging with the business, understanding its perspective, and alerting
the team to issues that may not be visible.

Being helpful remains the best attitude for a new engineer.

Mastery of techniques is essential, though it develops over time.
Progress comes from reflecting on these techniques consciously from
first principles. This theoretical understanding forms the foundation
for tackling more complex tasks.

Progress for a new engineer can be measured by the ability to estimate
effort. The ultimate sign of success is the capacity to formulate an
effective construction plan.

This does not require to-the-day accuracy in scheduling, but rather an
orderly and flexible plan that provides clarity and builds momentum for
a project team. Reasonable precision in estimates is sufficient for this
purpose.

The ability to formulate an effective construction plan is built on
consistent technical discipline, and that discipline begins with the
earliest practices of a new engineer.

These habits are not optional---they serve as the foundation for
everything else in this book.

\cleardoublepage

\chapter{Closing essay: Hallmarks of
quality}\label{docs-beyond-techniques-closing-essay-hallmarks-of-quality}

As developers, and broadly as a team with the goal of delivering
products, we are under the constant temptation to build products that
answer to the demands of the day.

These demands are often driven by emotions we feel towards other people.
They may be happiness of appreciations from a client or the fear of
reprimand if we do not deliver on time.

However, the instinctive reactions of personalities, and especially of
personalities who barely understand our craft, are not reliable guides
to quality. While these reactions may be justified---for example, if a
project has been dragging on too long, then the people around us are
right to be angry, they cannot be relied upon as the quality of work.

Our work is of quality if it enables the organisation to make sound,
positive decisions.

Nevertheless, this is too abstract a goal for day-to-day work. Even
other measures such as product usages or sentiment analysis of feedback
are too distant to help developers make decisions in their daily work.
We must seek other guides.

Our purpose is to lay down hallmarks of quality that help a team to
recognise and celebrate quality. They are nothing new but are
time-honoured practices.

As a developer, we must deliver the features, but also complete the
following:

\begin{enumerate}
\def\labelenumi{\arabic{enumi}.}
\item
  Expressive entities
\item
  Well-written explanation
\item
  Thoughtful unit-tests
\item
  Monitor assumptions
\item
  Anticipate errors
\item
  Adhere to patterns
\item
  Optimised code
\end{enumerate}

These are the hallmarks of quality. When done well, we have not only
delivered a feature but delivered quality work. The work will be high
value, long lasting, will win trust and importantly, give us pride in
the work we do. Not everyone will carry out all aspects in every
instance, but as a team, these are non-negotiables to quality.

\section{Expressive entities}\label{expressive-entities}

When looking at data, it is easy to forget we are dealing with a real
world. Yet the real world is what decision makers deal with. All models,
whether they are simulation models or data models or machine learning
models, are models of the world.

Correspondence to the world by the agent trying to influence it is
therefore the final arbitrator of whether a data model is successful.

Expressive entities refer to the idea that the tables and relationships
we create in the data warehouse should not merely reflect the data as we
found it but corresponds strongly to real world business processes. It
is expressive because the consumer or reviewer of the model can clearly
recognise the world the model is trying to approximate. If a competent
layperson cannot easily relate the model to the real world, then the
model is not expressive.

The process of creating expressive entities happen at a local level when
we use table and column names that are business meaningful. At the level
of a single model, we often say that we are ``crafting'' entities or
tables. At the warehouse level across several models, we often say that
we are conforming entities.

A data model is not a quality data model unless it is a useful
representation of the world. A data model is, therefore, a
worldview---more than a worldview, but not less. A developer without a
strong worldview cannot create a sound data model.

In a rapid development environment with the pressure to deliver,
crafting and conforming entities are difficult tasks to achieve. They
are hard because it is hard to think clearly about the world in a
systematic way and then articulate it through data.

There is a temptation to keep moving on to the next feature. Moreover,
the benefits of crafting entities are long term and not easily
understood by stakeholders. People understand testing and monitoring
assumptions, but what does a ``conceptual model'' mean? However, the
implication of inexpressive entities is that no one had spent time
thinking about the world and the processes with which we try to
influence it.

A quality data model is only achievable by a developer who has thought
seriously about the world, how data relates to it, and how to organise
the world in a way that makes sense for decisions. The ability to create
data entities that are accurate and expressive of the world is an
infallible mark of quality.

\section{Well-written explanation}\label{well-written-explanation}

Documentation is particularly important for data analysis (as opposed to
say, a game programmer) because data is a projection of real-world
processes and a data code is acting as an "interpreter" between data and
the world. This interpretation has a plain English equivalent.
Documentation is not simply "code comments" from one developer to
another but is an articulation to oneself and to others how the output
is reflecting reality. The writer C.S Lewis said, ``If you cannot say
something in common language, you either do not believe in it or you do
not understand it.'' If we cannot explain the logic behind our code in
plain English, why should anyone believe we know what we are doing?

The primary purpose of a well-written explanation of our code is to
articulate the logic back to ourselves. What are the entities we
created? What do they really mean? What assumptions have we made?
Writing a clear explanation is to step back from the code and do the
work again from a different angle. Often, in articulating an
explanation, one finds new ways of approaching the problem or discover
assumptions that we missed.

The secondary purpose of documentation is to win trust. To
non-developers, our work is a black box. A well-written explanation is
part of the transparency that are integral to trust. Business areas also
have a clearer understanding of the quality work and effort put in
development. These builds trust.

A developer does not need to writer documentation alone. It is the one
task where other people can be intimately involved. Sharing and
collaborating on documentation with subject matter experts are avenues
to clarify concepts and listen to feedback.

It is interesting to see a shift of habit as developers mature. A new
developer often write code with no documentation. Over time, the
developer gains a habit to write documentation after completing code,
then a habit to write documentation together with code. When mature, it
is common for a developer to start with code documentation for complex
scenarios (usually for the purpose for circulation). There is nothing
surprising about this. The transition merely reflects the fact that the
developer has mastered patterns in the work environment and is shifting
focus from the mechanics of code-writing to the reality behind what the
code is about.

\section{Thoughtful unit-tests}\label{thoughtful-unit-tests}

At the minimum, writing proper unit-tests means that developer has
solved the problem more than once and from different angles. This is
enough to justify the writing of the test. Often in the writing of a
test, the developer discovers a logical error or edge case that escaped
attention during development.

In the long term, tests exist because the world is not static. What is
true at the time of writing the code may not be true the day after. The
data itself can change, or the code itself may change. Sometimes these
changes can lead to significant errors that lead to incorrect decisions
or harm trust.

In isolated instances, one code element may have low probability of
going wrong due to external changes. When running several hundred
thousand lines of codes against several thousand entities, and a dozen
developers are changing several hundred lines per week, the
probabilities of error add up.

Tests are about quality and not quantity nor coverage. Poorly written
tests can lead to false confidence.

Resist the temptation to delay a test. If instinct says a code should be
tested, do not delay to the next week. It is astonishing how often the
test of which we think ``I probably should write this test, but I will
do this a bit later'' is the one that would have caught an error on the
week's product release.

\section{Assumptions are monitored}\label{assumptions-are-monitored}

All reasoning involves assumptions. Dealing with incomplete data
especially requires assumptions. Assumptions is the X in the ``If X,
then Y.'' Assumptions is the short-cut to Y. They allow one to defer
insignificant problems to a future date.

The world is not static. Assumptions valid at the time of writing the
code may be invalid the next day. When assumptions fail, the consequence
is always negative. This is because if the failure of the assumption has
no impact, then the assumption is not necessary in the first place.

The failure of assumptions is often incremental rather than sudden. The
world and its data usually drift from where we thought it would be
rather than making a drastic break.

This is particularly relevant with the rise of AI---models that work
today will behave differently as the world changes.

Whether they are sudden or gradual, we must monitor assumptions to
mitigate the impact for the inevitable day when they fail. If we
interpret fault-tolerance as the ability to go ahead despite errors,
then assumptions monitoring is a form of macro fault-tolerance.

Developers often assume uniqueness and existence of data. Without such
assumptions, there can be no analysis. But what if the product table is
not unique by {[}Product ID{]} tomorrow because the application had to
add a {[}Product sub-type{]}? This is quite possible if the source table
came from a non-relational system. If this happens, our logic will fall
apart. On the other hand, we become paralysed by risks if we constantly
assume source tables can change drastically. Therefore, automated
treatment of common types of assumptions should be built into any mature
processes.

The key step to monitoring assumptions is to know when we are making
one. More often, the issue is not that we are not monitoring the
assumption, but that we are not conscious of making them. We assume the
source tables will load in the daily ELT process. What if one of them
does not load, or if they load but not in the order we expect? Another
common type of assumption is to use hard-coded values in code.

Once assumptions have been identified, then it is a matter of being
disciplined and monitoring them through an automated framework.

\section{Anticipate errors}\label{anticipate-errors}

Tests and assumptions fall under the general category of error handling.
The question which concerns a seasoned developer is not whether the code
runs today, but how it may fail tomorrow. Many can write code against
the actual, but writing code against the possible is a whole new level.
One can often identify a superior developer by an instinctive,
second-nature, obsession about how the code may fail.

Even the most seemingly inconsequential actions can trigger an avalanche
of unintended consequences. By learning from these errors and improving
resilience of the process, these errors decrease. However, any large,
structural change, are still likely to cause unintended effects. The
first task of a developer then, is to expect errors to happen and to
anticipate specific ones. A developer must assume that things will
always go wrong.

The primary purpose of error handling is to limit damage and prevent
catastrophe of ripple effects. If we make it a rule to expect errors,
the easiest response is to limit the damage through fault-tolerant code.
This is the task of preventing or silencing errors.

The second purpose of error handling is to reveal silent errors. Not all
errors cause panic. In fact, the most dangerous errors are silent ones.
They can be silent by intention or by accident. In either case, a human
should be alerted to apply a fix, or the code should be designed to
recover automatically before damage propagates. This is the task of
detecting or recovering gracefully from errors.

The third purpose of error handling is for the team to grow
continuously. Each error reveals a blind spot. When errors are handled
properly, the team will be able to learn from the error and improve on
processes. This may be implementing new automated checks or rethinking
the algorithm from scratch. With each error the system becomes more
mature and more resilient. This is the task of learning from errors.
Learning from errors requires an open acknowledgement of failure and a
willingness to share the lesson with others. This is why humility is
important for developers.

Error handling is hard because the developer needs to look beyond the
now and anticipate any changes in the world. It is also challenging
because error handling requires an end-to-end understanding of the data
architecture. The ability to anticipate, prevent, silence, detect,
recover gracefully, and learn from error is a hallmark of genius. This
should apply both to technical situation (what happens to the join if
the column is null?) and to business rules (does the logic given by the
business make sense in all circumstances?).

A delivery team that can continuously learn from errors is unstoppable.
While error handling is challenging, a new developer can start by
adopting the habit of expecting errors.

\section{Adherence to patterns}\label{adherence-to-patterns}

There are slavish adherences to patterns that are counterproductive.
However, patterns, whether they be coding styles, naming conventions,
architecture, or the use of Git, are an indispensable part of a strong
team. There are bad teams that follow patterns, but there are no good
teams that do not have them.

The primary purpose of patterns is to pass on hard-earned wisdom
acquired by the team. In a complex environment, there are solutions that
are difficult to arrive at, or are counter-intuitive or are invisible to
the untrained eye. This wisdom has been acquired by predecessors at
great cost. It is not proper for developers to remain ignorant of their
value due to laziness in observation or a naive confidence in the
developer\textquotesingle s own judgement.

The second purpose of patterns is to create a consistent experience for
end users. This consistent experience ranges from design language such
as table names to more abstract experience such as joining a set of
unfamiliar tables. This consistency of experience over a body of work is
a felt quality that cannot be overstated. This hallmark of quality is
achievable only by a team where every developer plays their part
adhering to the whole.

The third purpose of patterns is to support thinking by filtering out
mental noise. Yes, CamelCase is no better than snake\_case for table
names. We could have used {[}Datetime created{]} as well as {[}Creation
datetime{]} or indeed {[}Created datetime{]}. In isolation, whether text
files have hard tabs, or 4-space, or 3-space soft tabs matter little. On
the other hand, when the same standard decision is applied across the
codebase, the pattern creates an important stability. In fact, when
patterns are consistently applied, anomalies stand out---even anomalies
which have nothing to do with the pattern. During review of code that is
not written in accordance with a usual pattern, the reviewer is forced
to carefully parse and interpret every line and every word.

These mental noises reduce the probability of detecting an error. It is
astonishing how, after the developer fixes some code indentation,
previously unseen logical errors stand out.

Judicious use of patterns in a team will advance quality by promoting
solutions that survived the test of time. Patterns also accelerate
delivery by removing developer uncertainty and simplifying decision
making. On the other hand, an unthinking adherence to patterns through a
``copy-and-paste'' mindset can lead to overengineering or stifle
innovation. To strike a balance between individual innovation and
accumulated wisdom, team members should see themselves as participants
in a conversation of innovation in which each team member can contribute
with an attitude of ``inquire, not debate.''

\section{Optimised code}\label{optimised-code}

Performant code carries out a task with as little server resource as
possible. Elegant code promotes readability and simplicity. Taken
together, we say the code has been optimised.

The primary purpose of optimising code is to support long term
maintainability. In a rapid development environment, a multitude of
non-performant code can easily overwhelm the server. If code is not
elegant, we increase the burden of changing the code. In isolated
instances the problem may be insignificant. Over several hundred
thousand lines of code, the problem compounds.

A direct way to write performant code in a data warehouse is to use
incremental extract logic. This is by no means easy. It can be difficult
to do accurately when integrating several datasets. Incremental extract
can influence overall design. For example, one may be selective about
what tables to use in a single script to avoid a dependency that would
ruin the potential to incrementally extract.

The secondary purpose of optimising code is to think about the problem
again. It is harder, and takes more thought, to write simple code than a
complicated one. Often in making code elegant, the developer sees the
concept in a different light. The process of a re-think can lead to
conceptual breakthroughs.

Not all good developers can write optimised code, but bad developers can
never write elegant, optimised code. Optimised code is one of the harder
hallmarks to achieve. We can all improve with practice.

\section{Three aspects of quality}\label{three-aspects-of-quality}

The hallmarks describe three aspects of data-based code: expressiveness,
error handling and elegant code.

\subsection{Expressiveness}\label{expressiveness-1}

Creating expressive entities, well-written documentation.

Expressiveness is grounded in the fact that data is not a real-world
entity in and of itself but is a projection of real-world processes onto
computer databases. Such projections are often fuzzy, imperfect, and
contingent. When we create expressive entities, our goal is to
reconstruct a clear and useful picture of the world from whence the data
came.

Well-written documentation serves the same purpose by articulating in
plain language how a data model relates to the real world. Analysts who
continually take data at face value (each row in the Task table is a
task right?), or mechanically applying statistical algorithms from a
textbook is divorced from the reality behind the data.

\subsection{Error handling}\label{error-handling}

Anticipating errors, thoughtful unit-tests, monitoring assumptions.

Error handling is to code what seat belts and safety breaks are to cars
or monitoring and kill-switches are to a power plant. They are not "nice
to haves" but are indispensable components to a complex computational
environment. Unit tests and monitoring assumptions all help to mitigate
against the inevitable day when an error occurs. A developer creating
code without considering what may go wrong is akin to a builder who adds
a floor to a building without strengthening the foundations---it will
work, for now.

\subsection{Elegant code}\label{elegant-code}

Adherence to patterns, optimised code.

Elegant code is key to sustainable growth. In a complex and fast-moving
environment, we cannot afford spaghetti code nor free-for-all approaches
to common problems.

Performance-tuned code is obvious for the simple reason that the
computational resources are finite.

\section{Final words}\label{final-words}

These hallmarks of quality are not revolutionary. They have been around
since code development began. The issue is not that developers do not
know about them but that they do not do them because of lack of
discipline or pressures to deliver.

There is the idea that doing work properly cost time and delay delivery.
The retort is that doing work improperly cost more time. Writing good
tests help capture errors earlier and reduces the time it takes to fix
things up. Crafted entities vastly reduce the time to deliver complex
features because the model becomes more powerful. Writing a clear
explanation helps all get on-board, establishes trust and builds
momentum. On the other hand, technical debt quickly comes back with a
vengeance. At the level of a single development cycle, doing the work
properly may increase development time by 20 to 40\%. Over the course of
several months, the dividends of quality work pay off and enable faster
delivery than hasty development could ever achieve.

The developer is constantly under the pressures of deadlines. However,
if one does not have time to do the work well, the solution is not to do
the work poorly. The best way to achieve the hallmarks is simply
consider these as part of, and not something in addition to, the work
itself. Once we make the mental leap, it becomes more natural to build
these into the daily work and to the pace of projects.

\end{document}
